% =========================================================================
% CONSOLIDATED STRATEGIC PRIORITY CLUSTERS
% File: consolidated_priorities_v1.tex
% Integration: % =========================================================================
% CONSOLIDATED STRATEGIC PRIORITY CLUSTERS
% File: consolidated_priorities_v1.tex
% Integration: % =========================================================================
% CONSOLIDATED STRATEGIC PRIORITY CLUSTERS
% File: consolidated_priorities_v1.tex
% Integration: % =========================================================================
% CONSOLIDATED STRATEGIC PRIORITY CLUSTERS
% File: consolidated_priorities_v1.tex
% Integration: \input{consolidated_priorities_v1} in master document
%
% Required packages (add to master preamble if not present):
%   \usepackage[table]{xcolor}
%   \usepackage{multirow}
%   \usepackage{array}
%
% New bibliography keys required (entries provided at end of file):
%   fips203, fips204, fips205, nistir8413, nistsp800207,
%   etsi103744, etsizsmoo9, tgcpm20, isoiec11889,
%   opentitan, keystone, threeGPP23288, threeGPP33501, dora, gdpr
% =========================================================================

% ---- Shared colors/macro helpers when included from a master document ----
\providecolor{gaterow}{RGB}{240,240,238}
\providecolor{phase1row}{RGB}{237,245,253}
\providecolor{phase2row}{RGB}{237,249,244}
\providecolor{phase3row}{RGB}{245,244,254}
\providecolor{colheader}{RGB}{245,245,243}
\providecolor{clusterheader}{RGB}{220,232,248}

\providecolor{sslredbg}{RGB}{252,235,235}
\providecolor{sslredtx}{RGB}{163,45,45}
\providecolor{sslamberbg}{RGB}{250,238,218}
\providecolor{sslambtx}{RGB}{133,79,11}
\providecolor{ssltransbg}{RGB}{230,241,251}
\providecolor{ssltrantx}{RGB}{24,95,165}
\providecolor{ssltealbg}{RGB}{225,245,238}
\providecolor{ssltealtx}{RGB}{15,110,86}

\providecommand{\sslbadge}[3]{%
      \parbox{\linewidth}{\raggedright
            \colorbox{#1}{\scriptsize\parbox{0.92\linewidth}{\raggedright
                  \textbf{\textcolor{#2}{#3}}}}}}

\providecommand{\kpistrip}[1]{%
      \par\vspace{2pt}\noindent\rule{\linewidth}{0.3pt}%
      \par\noindent\textit{\scriptsize #1}}

\providecommand{\cellhead}[1]{\noindent\textbf{#1}\par\smallskip}
\providecommand{\gatelbl}[1]{\scriptsize\textit{Gate~#1:}}
\providecommand{\subp}[1]{\textit{\scriptsize(#1)}}

\section{Strategic Priority Clusters}
\label{ch:strategic_priority_clusters}

% \subsection{Consolidation Rationale and Structural Logic}
% \label{sec:consolidation_rationale}

% The ten executive priorities identified through the 4-Filter Prioritisation Framework represent
% distinct technology domains, but they share deep structural dependencies and converge on three
% strategic imperatives that define Europe's 15-year cybersecurity trajectory for the Cloud-Edge-IoT
% continuum. Maintaining ten parallel profiles in a dual-audience deliverable risks fragmenting
% the narrative for European Commission technical evaluators while obscuring the systemic
% interdependencies that determine whether individual priorities succeed or fail.

% The consolidation presented in this chapter reorganises the ten priorities into three strategic
% clusters, each with a unifying architectural thesis. To preserve the analytical precision required
% by IEEE Access peer review, each cluster retains its constituent sub-priorities as explicitly
% labelled components within a nested progression matrix, allowing individual technology maturation
% trajectories to be tracked independently while the cluster-level framing communicates strategic
% coherence to policy audiences.

% The mapping is as follows:
% \begin{itemize}
%   \item \textbf{Priority Cluster 1 --- Sovereign Trusted Infrastructure:}
%         P1 (Hardware Roots of Trust), P2 (PQC Agility), P3 (Lightweight
%         Cryptography and Physical-Layer Security), P8 (Confidential
%         Computing), P9 (Resilient-by-Design / Island Mode).
%   \item \textbf{Priority Cluster 2 --- Cognitive Security and Autonomous
%         Cyber Defence:} P6 (Cognitive SecOps / LLMSecOps), P10 (Proactive
%         APT Management), and the AI-as-target dimension drawn from
%         Cross-Layers CL1.2, CL1.3, and CL4.7.
%   \item \textbf{Priority Cluster 3 --- Federated Operational Sovereignty:}
%         P4 (Security Interoperability via Simpl), P5 (Supply Chain
%         Verification), P7 (Pan-European NaaS Federation).
% \end{itemize}

% One critical cross-priority dependency must be made explicit before the cluster profiles are
% presented. Hardware supply chain attestation (P5) and Runtime SBOM verification constitute
% a shared gate condition for both Cluster 1 (TEE and secure element integrity) and Cluster 3
% (federated supply chain governance). Progression in P8 (Confidential Computing) is contingent
% on the availability of trustworthy RISC-V secure elements from P1; P3 (Lightweight PhySec)
% and P2 (PQC) both depend on EU Chips Act fabrication capacity as a common Phase 2 enabler.
% These shared bottlenecks are identified explicitly in the Cross-Priority Bottleneck assessment
% in Section~\ref{sec:bottlenecks} and traced to specific progression matrix gate conditions.

% =========================================================================
% PRIORITY CLUSTER 1: SOVEREIGN TRUSTED INFRASTRUCTURE
% =========================================================================

% \newpage
\section{Priority Cluster 1: Sovereign Trusted Infrastructure}
\label{sec:cluster1}

\noindent\textbf{Constituent Sub-Priorities:} P1 (Hardware Roots of Trust \& EUCC `High'),
P2 (PQC Agility for Edge/IoT), P3 (Lightweight Cryptography \& Physical-Layer Security),
P8 (Confidential Computing \& Privacy-Preserving AI), P9 (Resilient-by-Design \& Island Mode).

\medskip
\noindent\textbf{Cluster Thesis:} Before any federated or AI-driven security capability can be
trusted, the substrate on which it executes must provide cryptographically verifiable integrity,
quantum-resistant confidentiality, and autonomous survivability independent of external
infrastructure. This cluster defines the verifiable hardware and cryptographic foundation that
all higher-order capabilities in Clusters 2 and 3 depend upon as an architectural prerequisite.

\subsection*{7-Dimensional Analysis: Priority Cluster 1}

\begin{enumerate}

% ---- DIMENSION 1 ----
\item \textbf{Baseline \& Target Objective}

\textit{Baseline (2025--2026):}
European critical infrastructure and edge deployments operate under high foreign dependency
across all five constituent sub-priorities. Hardware roots of trust are dominated by non-EU
Trusted Platform Module (TPM) manufacturers --- Infineon SLB9670, NXP SE050, and
STMicroelectronics ST33 series collectively account for the majority of the EU market under
ISO/IEC~11889~\cite{isoiec11889} and TCG TPM~2.0 specifications~\cite{tgcpm20}. The Trusted
Execution Environment (TEE) market is dominated by Intel TDX, AMD SEV-SNP, and ARM TrustZone,
none of which are EU-designed or manufactured. Post-Quantum Cryptography (PQC) standardisation
reached a definitive milestone with NIST's publication of FIPS~203 (ML-KEM)~\cite{fips203},
FIPS~204 (ML-DSA)~\cite{fips204}, and FIPS~205 (SLH-DSA)~\cite{fips205} in August~2024;
however, hardware-optimised implementations for 8/16-bit microcontrollers remain at the
Applied~R\&D stage, as confirmed by NIST IR~8413~\cite{nistir8413}. Physical-layer security
for extreme-edge mMTC devices lacks standardised specifications. Edge nodes universally operate
with hard cloud dependencies and no autonomous fallback architectures meeting the operational
resilience mandates of NIS2~\cite{nis2} and DORA~\cite{dora}.

Current cluster ODP: Applied R\&D (driven by PQC, Lightweight PhySec, and Confidential
Computing sovereign alternatives). Current cluster SSL: High Foreign Dependency.

\textit{Target (2040):}
At least 60\% of EU critical infrastructure edge nodes deploy EU-designed RISC-V open-hardware
secure elements certified at EUCC `High' assurance level~\cite{eucc}. 100\% of new EU critical
IoT deployments support hybrid or full FIPS~203/204/205-compliant PQC. Zero-energy nodes
execute quantum-resistant device authentication within a sub-100\,$\mu$J energy budget per
operation. Trusted Execution Environments built on EU open-hardware architectures (Keystone~\cite{keystone},
OpenTitan~\cite{opentitan}) cover at least 50\% of new EU cloud-edge deployments. All NIS2-essential
and DORA-critical edge nodes demonstrate validated 72-hour autonomous Island Mode survivability
independent of WAN connectivity.

% ---- DIMENSION 2 ----
\item \textbf{Primary Drivers}

\textit{Threat Drivers:}
\begin{itemize}
  \item State-sponsored APT exploitation of JTAG debug interfaces and electromagnetic
        side-channel analysis (EMSCA) on field-deployed edge nodes; the Spectre and Meltdown
        class of microarchitectural side-channel vulnerabilities establishes that even
        production-certified hardware can be compromised at scale, and far-edge nodes
        operating without physical security controls present a materially higher exposure
        surface than hyperscaler data centres.
  \item ``Harvest Now, Decrypt Later'' (HNDL) campaigns: state-level adversaries are
        actively archiving encrypted EU government and industrial communications under
        the expectation of Cryptographically Relevant Quantum Computer (CRQC)
        availability. NIST IR~8413 places the Q-Day horizon in the 10--20 year range
        with a highly uncertain lower bound~\cite{nistir8413}; HNDL renders this a
        present-tense threat against data with multi-decade confidentiality
        requirements.
  \item Physical tampering and environmental attacks on unattended far-edge nodes
        (smart grid substations, O-RAN radio units, industrial gateways) where the
        physical security controls available in hyperscaler facilities are absent.
  \item Cascading WAN disconnection failures propagating through centrally-dependent edge
        architectures, exploited by adversaries to degrade critical infrastructure
        services during hybrid conflict scenarios.
\end{itemize}

\textit{Regulatory Drivers:}
\begin{itemize}
  \item CRA Articles~10 and~11: mandatory vulnerability handling and security update
        obligations for products with digital elements; enforcement begins in 2027 for
        most product categories~\cite{cra}.
  \item NIS2 Article~21(2)(e): cryptographic policies and key management for essential
        and important entities; full enforcement began October 2024~\cite{nis2}.
  \item DORA Article~9: ICT resilience testing including hardware-level resilience
        requirements for financial sector critical infrastructure~\cite{dora}.
  \item AI Act Article~9: risk management system requirements for high-risk AI
        systems; critical infrastructure AI components require hardware integrity
        guarantees that only EUCC `High' certified secure elements can
        provide~\cite{aiact}.
\end{itemize}

% ---- DIMENSION 3 ----
\item \textbf{High Impact Activities}

\begin{enumerate}[label=(\alph*)]
  \item \textbf{[Phase~1]} Develop sovereign European Secure Elements based on
        RISC\-V open-hardware (building on OpenTitan~\cite{opentitan} and
        Keystone~\cite{keystone} frameworks), targeting EUCC `High' assurance level
        evaluation; implement automated hardware-level remote attestation protocols
        interoperable with TPM~2.0~\cite{tgcpm20} and ISO/IEC~11889~\cite{isoiec11889}
        for heterogeneous large-scale edge fleets.

  \item \textbf{[Phase~1--2]} Optimise NIST FIPS~203 (ML-KEM)~\cite{fips203} and
        FIPS~204 (ML-DSA)~\cite{fips204} for 8/16-bit microcontrollers under severe
        energy and memory constraints (target: ML-KEM-512 within 32\,KB RAM on
        Cortex-M0+ class devices); develop crypto-agility OTA update middleware enabling
        algorithm rotation without hardware replacement, compatible with ETSI TS~103~744
        hybrid key exchange profiles~\cite{etsi103744}.

  \item \textbf{[Phase~1--2]} Implement ultra-lightweight PQC and keyless
        Physical-Layer Security (RF fingerprinting exploiting Channel State Information
        and Carrier Frequency Offset, CSI/CFO-based device authentication) for mMTC
        devices operating under zero-energy or near-zero-energy constraints; target
        integration with 6G ISAC (Integrated Sensing and Communication) interfaces per
        SNS~JU~SRIA~\cite{snsju}.

  \item \textbf{[Phase~2]} Deploy multi-party confidential computing pipelines
        using TEE abstraction layers (Keystone enclaves on RISC-V~\cite{keystone})
        combined with Fully Homomorphic Encryption (FHE) schemes for collaborative AI
        training across multi-provider nodes without raw data exposure; align with
        EUCS `High+' certification requirements~\cite{eucs}.

  \item \textbf{[Phase~2--3]} Engineer autonomous Island Mode architecture: soft-state
        edge orchestration using cached credentials with 72-hour validity windows,
        LoRaWAN LPWAN fallback micro-slices, eBPF-based kernel-level policy enforcement
        for autonomous node isolation, and bio-inspired self-healing swarm protocols
        enabling edge nodes to sustain critical service delivery without cloud
        connectivity or human intervention.
\end{enumerate}

% ---- DIMENSION 4 ----
\item \textbf{Catalysts \& Enablers}

\textit{Funding Alignment:}
Horizon Europe Cluster~4 (Digital, Industry and Space) for sovereign secure silicon and
PQC hardware implementation; Horizon Europe Cluster~3 (Civil Security for Society) for
edge resilience and Island Mode architectures; EU Chips Act (Regulation (EU) 2023/1781)
pilot fabrication lines for security-critical silicon via the Chips Joint Undertaking
(KDT~JU); ECCC Strategic Agenda priority areas ``Secure Components and Systems'' and
``Post-Quantum Cryptography transition infrastructure''~\cite{eccc}; Digital Europe
Programme for PQC migration tooling and crypto-agility middleware.

\textit{Standardisation:}
NIST NCCoE ``Migration to Post-Quantum Cryptography'' project (reference architecture
and migration playbooks); CEN/CENELEC JTC~13/WG~8 (hardware security standards aligned
with EUCC Common Criteria); ETSI TC CYBER (ETSI TS~103~744 quantum-safe hybrid key
exchanges~\cite{etsi103744}); ENISA PQC transition guidelines and cryptographic inventory
tooling; 3GPP SA3 for physical-layer security integration with 5G/6G
specifications~\cite{threeGPP33501}.

\textit{Ecosystem Enablers:}
RISC-V International open-hardware community; OpenTitan coalition pilot deployments
for EU sovereign RoT~\cite{opentitan}; SNS~JU pilot networks for 6G-integrated PhySec
validation~\cite{snsju}; federated Cyber Range infrastructure (CL3.5) for Island Mode
and TEE attestation stress-testing.

% ---- DIMENSION 5 ----
\item \textbf{Disruptors \& Systemic Risks}

\begin{itemize}
  \item \textbf{Q-Day cliff:} CRQC becoming operational before EU PQC migration
        completes; NIST IR~8413 places this horizon at 10--20 years with a highly
        uncertain lower bound~\cite{nistir8413}. The EU PQC migration timeline targeting
        critical infrastructure completion by 2030 has limited buffer against the
        pessimistic lower bound, and HNDL campaigns make this a present-tense threat
        against long-lived secrets.
  \item \textbf{EU Chips Act EUV lithography shortfall:} RISC-V secure element
        production at the node geometries required for EUCC `High' certification depends
        on advanced lithography; EU 2030 Chips Act targets (20\% global market share)
        do not specify security-critical silicon capacity independently, and ASML EUV
        tooling remains subject to US--Netherlands export control dynamics that could
        constrain sovereign fabrication timelines beyond 2028.
  \item \textbf{Side-channel analysis breakthroughs against FIPS~203 software
        implementations:} lattice-based algorithms have demonstrated timing and
        cache-timing attack surfaces in software implementations on constrained devices;
        hardware countermeasures are not yet standardised at the time of writing
        [CITATION REQUIRES VERIFICATION --- recent SCA against ML-KEM implementations].
  \item \textbf{TEE architectural monoculture:} if EU open-hardware TEE adoption
        concentrates on a single architecture (e.g., Keystone only~\cite{keystone}), a
        single discovered vulnerability propagates across the entire EU critical
        infrastructure TEE estate without architectural diversity as a containment layer.
  \item \textbf{LPWAN regulatory fragmentation:} Island Mode fallback using LoRaWAN
        operates on unlicensed 868\,MHz spectrum in the EU; spectrum policy
        fragmentation across Member States could prevent coordinated emergency
        micro-slice activation under unified crisis response scenarios.
\end{itemize}

% ---- DIMENSION 6 ----
\item \textbf{Failsafe Mechanisms \& Resilience}

\begin{itemize}
  \item \textbf{Hybrid cryptographic transition:} deploy ML-KEM (FIPS~203)~\cite{fips203}
        combined with ECDH-P384 in a hybrid key encapsulation scheme per
        ETSI TS~103~744~\cite{etsi103744}; if ML-KEM is found vulnerable via cryptanalytic
        breakthrough, ECDH-P384 maintains classical confidentiality while SLH-DSA
        (FIPS~205)~\cite{fips205} provides backup hash-based signature integrity with
        distinct mathematical foundations.
  \item \textbf{TEE implementation diversity:} deploy a minimum of two architecturally
        distinct TEE technologies across critical EU infrastructure (ARM CCA on
        server-class hardware and RISC-V Keystone~\cite{keystone} on edge nodes) to
        prevent single-architecture compromise from cascading across the entire
        confidential computing estate.
  \item \textbf{Autonomous node quarantine via eBPF:} if TEE remote attestation fails,
        eBPF-based network policy enforcement at the Linux kernel level (e.g., via
        Cilium or Falco eBPF rulesets) isolates the affected node from the orchestration
        plane within sub-second latency, blocking lateral movement before human
        intervention is available.
  \item \textbf{Island Mode failsafe:} edge nodes maintain cryptographically signed
        cached credentials with defined 72-hour validity windows; on WAN disconnection,
        a LoRaWAN LPWAN micro-slice activates a minimal viable network carrying
        life-safety telemetry; if LPWAN also fails, nodes revert to local-only operation
        with integrity-protected audit logging for post-incident forensics.
  \item \textbf{RF fingerprinting fallback:} for zero-energy IoT nodes where
        cryptographic processing exceeds the available energy budget, Physical-Layer
        authentication via RF fingerprinting (CSI/CFO-based device identity) provides
        keyless device verification using intrinsic wireless channel characteristics,
        requiring no cryptographic key storage or computation.
\end{itemize}

% ---- DIMENSION 7 ----
\item \textbf{Transversal Meta-Layer}

\textit{Sovereignty \& Autonomy:}
This cluster eliminates approximately 100\% foreign dependency on non-EU TPMs (Infineon,
NXP, STMicroelectronics dominate the EU market under TCG TPM~2.0~\cite{tgcpm20}) and
non-EU TEE architectures (Intel TDX, AMD SEV-SNP, ARM TrustZone, all non-EU designed and
manufactured). Crypto-agility frameworks additionally reduce dependency on US-controlled
NIST algorithm supply chains by enabling algorithm substitution if standardisation processes
are constrained by geopolitical factors.

\textit{Trust \& Ethics:}
GDPR Article~25 (data protection by design and by default) is operationalised through
TEE-enforced data minimisation and in-use encryption in multi-tenant environments,
ensuring that sensitive data processing cannot be accessed by platform operators~\cite{gdpr}.
AI Act Article~9 risk management system requirements for high-risk AI components require
hardware integrity guarantees --- only EUCC `High' certified secure elements provide the
mathematically verifiable trust anchor this obligation demands~\cite{aiact}.

\textit{Sustainable Security (Green IT):}
CRYSTALS-Kyber (ML-KEM) demonstrates approximately 1.5$\times$ computational overhead
relative to ECDH at equivalent security levels on ARM Cortex-M4, per NIST IR~8413~\cite{nistir8413},
representing a manageable energy penalty. Lightweight PQC research targets sub-100\,$\mu$J
per authentication operation for zero-energy IoT nodes. EU-sovereign RISC-V secure
elements designed for edge power constraints enable more energy-efficient cryptographic
operations than general-purpose processor TPM emulation, reducing per-device carbon
footprint at continental deployment scale.

\end{enumerate}



% =========================================================================
% CLUSTER 1: SOVEREIGN TRUSTED INFRASTRUCTURE
% =========================================================================
\newpage

\subsection{Progression Matrix: Priority Cluster 1 --- Sovereign Trusted Infrastructure}

\begin{landscape}
\setlength{\textwidth}{26.5cm}\setlength{\linewidth}{26.5cm}%
\setlength{\hsize}{26.5cm}\setlength{\columnwidth}{26.5cm}
\setlength{\tabcolsep}{5pt}
\renewcommand{\arraystretch}{1.35}

\begin{longtable}{%
|>{\footnotesize\raggedright\arraybackslash}p{2.8cm}%
|>{\footnotesize\raggedright\arraybackslash}p{5.4cm}%
|>{\footnotesize\raggedright\arraybackslash}p{5.4cm}%
|>{\footnotesize\raggedright\arraybackslash}p{5.4cm}%
|>{\footnotesize\raggedright\arraybackslash}p{5.4cm}|}

\caption{\textbf{Cluster~1: Sovereign Trusted Infrastructure} (P1, P2, P3, P8, P9).
Transition matrix across the four methodology dimensions. Sub-priority labels in
parentheses indicate where individual priorities diverge from the cluster-level
descriptor within a cell. ODP scale: Applied R\&D $\to$ Fragmented Piloting
$\to$ Standardised Integration $\to$ Ubiquitous Operation. SSL scale:
High Foreign Dependency $\to$ Emerging EU Ecosystem $\to$ Full Sovereign Autonomy.}
\label{tab:cluster1}\\

\hline
\rowcolor{colheader}
\textbf{Phase / SSL} &
\textbf{R\&I Maturity}{\scriptsize ODP indicator per sub-priority} &
\textbf{Policy, Regulatory \& Standards}{\scriptsize enforcement $\cdot$
  certification $\cdot$ standardisation} &
\textbf{Operational Deployment \& Adoption}{\scriptsize actor types $\cdot$
  scale $\cdot$ sector coverage} &
\textbf{Market Uptake \& Strategic Sovereignty}{\scriptsize SSL indicator
  $\cdot$ dependency reduction $\cdot$ industrial capacity} \\
\hline
\endfirsthead

\hline
\rowcolor{colheader}
\textbf{Phase / SSL} &
\textbf{R\&I Maturity} &
\textbf{Policy, Regulatory \& Standards} &
\textbf{Operational Deployment \& Adoption} &
\textbf{Market Uptake \& Strategic Sovereignty} \\
\hline
\endhead
\hline\multicolumn{5}{r}{\footnotesize\textit{Continued\ldots}}\\ \endfoot
\hline \endlastfoot

% ---- PHASE 1 ----
\rowcolor{phase1row}
\textbf{Phase~1}
\newline\textit{Foundation \&}
\newline\textit{Compliance}
\newline\textit{2025--2030}
\newline\sslbadge{sslredbg}{sslredtx}{SSL: High Foreign Dependency}
&
\cellhead{Cluster ODP: Applied R\&D / Fragmented Piloting}
The cluster is split at baseline. \subp{P1, P9} reach Fragmented
Piloting: RISC-V secure element projects targeting EUCC `High' are
funded under Horizon Europe and EU Chips Act; Island Mode soft-state
protocols are validated in 5G research testbeds.
\subp{P2, P3, P8} remain at Applied R\&D: NIST FIPS~203/204/205
algorithms are being optimised for 8/16-bit microcontrollers under
severe energy and memory constraints; Keystone and OpenTitan
prototypes are under EUCC pre-evaluation; RF fingerprinting
Physical-Layer Security is validated in controlled 5G lab environments;
Fully Homomorphic Encryption approaches practical latency bounds in
research settings only.
\kpistrip{The cluster ODP is Applied R\&D, driven by the majority of
sub-priorities. Deployment claims for P2, P3, P8 must be grounded in
named Horizon Europe or ECCC project outcomes, not technology
previews.}
&
\cellhead{Regulatory environment: enforcing, certification nascent}
NIS2 Article~21(2)(e) cryptographic policies for essential entities
entered full enforcement in October~2024. CRA Articles~10 and~11
(vulnerability handling and security update obligations) begin
enforcement in 2027, creating immediate procurement-level demand for
hardware security certification. AI~Act Articles~9 and~15 establish
hardware integrity requirements for high-risk AI systems and enter
applicability in 2025--2026. ENISA PQC migration guidelines are
published; EUCC evaluation processes are initiated for the first RISC-V
prototypes, but no EU-specific EUCC~`High' certified sovereign secure
element exists yet. CEN/CENELEC JTC~13 standards for hardware security
MIMs are in draft phase.
\kpistrip{The regulatory framework is enforcing existing obligations
but lacks EU-specific hardware certification outputs; the critical gap
is the absence of a ratified EUCC~`High' EU-origin product.}
&
\cellhead{Deployment: research consortia and national agencies only}
RISC-V secure element pilots are operated by 3--5 Horizon Europe
consortia and national cybersecurity agencies in controlled
environments. Intel~TDX, AMD~SEV-SNP, and ARM~TrustZone dominate
all operational TEE deployments across EU cloud and critical
infrastructure. PQC is confined to cloud-side software implementations
in pilot contexts; far-edge IoT devices run classical cryptography
exclusively. Edge nodes operate with hard cloud dependencies and no
validated autonomous fallback architecture meeting DORA or NIS2
business continuity thresholds. SMEs have no access to EU-sovereign
hardware alternatives at viable cost.
\kpistrip{Operational deployment at Phase~1 is institutional and
research-concentrated. The gap between pilot existence and operational
availability for regulated-sector procurement must be explicitly
acknowledged.}
&
\cellhead{SSL: High Foreign Dependency}
Non-EU TPMs (Infineon~SLB9670, NXP~SE050, STMicroelectronics~ST33)
hold $\geq$90\% of EU market share in hardware security modules.
Intel, AMD, and ARM control the entire operational TEE market. No
EU-designed and EU-manufactured secure element is commercially
available. EU Chips Act fabrication investment for security-critical
silicon is initiated but pre-production. US NIST algorithm supply chain
is the sole PQC standardisation reference. Structural vulnerability is
direct: geopolitical restriction of access to any of the three major
TPM suppliers would immediately impair EU critical infrastructure
attestation capability.
\kpistrip{Quantified dependency at Phase~1 baseline: $\geq$90\% foreign
market share, zero EU-certified sovereign hardware alternatives.
Sovereignty progress is measured from this specific figure.}
\\
\hline

% ---- GATE 1->2 ----
\rowcolor{gaterow}
\gatelbl{1 $\rightarrow$ 2} &
\multicolumn{4}{>{\footnotesize\raggedright\arraybackslash}p{22.0cm}|}{%
\cellcolor{gaterow}%
\textbf{R\&I trigger:} At least one EU-designed RISC-V secure element prototype
has been independently validated against EUCC~`High' evaluation criteria and has
entered the formal ENISA/national scheme certification process.
\textbf{Policy trigger:} CRA Articles~10 and~11 are in active enforcement with
documented non-compliance decisions, confirming that market incentives for
certified EU hardware alternatives are financially material.
\textbf{Operational trigger:} Documented operational deployment of hybrid
ML-KEM (FIPS~203) plus ECDH-P384 PQC schemes in at least one EU critical
infrastructure sector (energy or financial) at production scale, not pilot scale.
\textbf{Market trigger:} EU Chips Act security-critical silicon fabrication
pilot line has produced and delivered its first commercially viable batch of
EU-origin secure elements, confirmed by independent yield and cost assessment.}
\\
\hline

% ---- PHASE 2 ----
\rowcolor{phase2row}
\textbf{Phase~2}
\newline\textit{Applied \& Industrial}
\newline\textit{R\&D alongside}
\newline\textit{Federation \&}
\newline\textit{Standardisation}
\newline\textit{2030--2035}
\newline\sslbadge{ssltransbg}{ssltrantx}{SSL: Emerging EU Ecosystem}
&
\cellhead{Cluster ODP: Standardised Integration / continued Applied R\&D}
\subp{P1, P9} reach Standardised Integration: the first EUCC~`High' certified
EU-manufactured RISC-V secure elements are in production and being integrated
into critical infrastructure; automated remote attestation at fleet scale is
validated; 72-hour Island Mode achieves standardised certification under
3GPP/ETSI fallback API profiles.
\subp{P2, P3} transition from Applied R\&D to Fragmented Piloting: hybrid
ML-KEM plus ECDH-P384 schemes are validated at Industrial IoT scale; lightweight
PQC for sub-100~$\mu$J zero-energy nodes is demonstrated under 6G ISAC testbeds.
\subp{P8} reaches Fragmented Piloting: RISC-V Keystone TEEs are operational
in multi-provider Data Spaces; FHE performance approaches the threshold for
selective production use in high-value financial analytics.
\kpistrip{Phase~2 R\&I is not purely consolidation: three of five sub-priorities
are still in active applied and industrial R\&D. Describing this phase as
``standardisation only'' misrepresents the maturation arc.}
&
\cellhead{EUCC certification and PQC standards formalised}
The EUCC~`High' certification of the first EU-manufactured RISC-V secure
element is the definitive Phase~2 policy gate: it transforms applied research
into a market-accessible, legally certifiable product and unlocks the procurement
pathway for EU critical infrastructure operators. NIST FIPS~203/204/205 are
adopted into ETSI~TS profiles for EU deployment contexts. ENISA publishes a
mandatory PQC migration timeline for critical infrastructure operators, referenced
in NIS2 sector-specific implementing acts. EU Chips Act fabrication investments
are active and measurable. CEN/CENELEC hardware security standards for EU-origin
secure elements enter ratification.
\kpistrip{The single Phase~2 policy gate is EUCC~`High' certification of a
EU-origin product. All other policy milestones are enabling conditions, not
gate conditions.}
&
\cellhead{Critical infrastructure and regulated sectors deploying EU alternatives}
EU EUCC~`High' certified secure elements are deployed in energy sector
substations and telecommunications critical infrastructure in at least five
Member States. Hybrid PQC schemes are standard in new Industrial IoT and
financial sector (DORA-compliant) deployments. RISC-V Keystone TEEs operate
in federated multi-provider Data Spaces alongside existing Intel TDX
infrastructure. Validated 72-hour Island Mode is deployed in at least three
EU 5G cross-border corridor pilots under SNS~JU governance. SME-accessible
PQC tooling and SBOM compliance support are available commercially.
\kpistrip{Phase~2 deployment evidence requires named sector and Member State
scope; anonymous ``multi-sector deployment'' claims do not constitute
verifiable operational maturity.}
&
\cellhead{SSL: Emerging EU Ecosystem}
EU-designed secure elements achieve 15--30\% market share in new critical
infrastructure procurement, rising from a near-zero baseline. EU Chips Act
fabrication investments have created a domestic production capacity that is
commercially credible, if not yet dominant. Foreign TPMs remain the standard
for legacy deployments but are no longer the only option for new regulated-sector
procurement. EU open-hardware TEE ecosystem (Keystone, OpenTitan) holds
commercially certifiable products. The conditions for Full Sovereign Autonomy
are structurally in place --- fabrication capacity, certification infrastructure,
standards leadership --- but have not yet reached continental operational scale.
\kpistrip{SSL transition requires industrial substance, not research demonstration.
15--30\% new critical infra market share is the quantified Phase~2 SSL threshold.}
\\
\hline

% ---- GATE 2->3 ----
\rowcolor{gaterow}
\gatelbl{2 $\rightarrow$ 3} &
\multicolumn{4}{>{\footnotesize\raggedright\arraybackslash}p{22.0cm}|}{%
\cellcolor{gaterow}%
\textbf{R\&I trigger:} Production-ready PQC-native RISC-V silicon is available
from EU-domiciled fabrication at commercially viable yields, confirmed by
independent cost-per-unit assessment against equivalent non-EU alternatives.
\textbf{Policy trigger:} EUCC~`High' certification is mandated as a procurement
requirement for all new EU critical infrastructure edge node deployments, with
no waiver provisions for new procurements (legacy replacement programmes are
time-bound).
\textbf{Operational trigger:} Validated 72-hour Island Mode is deployed as a
standard operational baseline --- not a pilot --- in $\geq$50\% of new EU
NIS2-essential entity edge node procurements in the energy and
telecommunications sectors.
\textbf{Market trigger:} EU Chips Act sovereign fabrication capacity achieves
$\geq$50\% self-sufficiency for security-critical silicon in new EU critical
infrastructure procurement, measured by share of new orders fulfilled by
EU-origin fabrication.}
\\
\hline

% ---- PHASE 3 ----
\rowcolor{phase3row}
\textbf{Phase~3}
\newline\textit{Autonomy \&}
\newline\textit{Leadership}
\newline\textit{2035--2040}
\newline\sslbadge{ssltealbg}{ssltealtx}{SSL: Full Sovereign Autonomy}
&
\cellhead{Cluster ODP: Ubiquitous Operation}
All five sub-priorities operate at Ubiquitous Operation. EU-designed
RISC-V secure elements are the hardware default across all new EU ICT
products. Fully Homomorphic Encryption is operationalised for
end-to-end confidential processing from IoT to HPC, not confined to
high-value analytics. Zero-energy PhySec authentication via RF
fingerprinting operates as a standard device identity mechanism for
battery-less mMTC nodes. Bio-inspired swarm protocols enable autonomous
cyber-physical recovery without cloud intervention. Phase~3 R\&I is
oriented toward the next-generation frontier: neuromorphic secure
compute, post-second-generation quantum hardware, and post-CMOS
security architectures --- EU research groups define the agenda rather
than adopting it.
\kpistrip{Ubiquitous Operation requires deployment statistics, not
deployment plans. Phase~3 evidence must cite procurement share data
and operational audit findings, not programme milestone completions.}
&
\cellhead{EU as international hardware security standards reference}
EUCC~`High' certification is mandatory for all new EU ICT products
without waiver provisions. PQC-native silicon is the EU procurement
default. EU hardware security standards are referenced in ISO/IEC
international standardisation and in bilateral digital trade
agreements. Continuous automated EUCC lifecycle maintenance replaces
point-in-time audit models completely. The EU sets the terms of
hardware security compliance for market access, transforming its
internal market standard into a de facto international reference.
\kpistrip{Phase~3 policy maturity requires evidence of international
standards influence: at least one EU hardware security standard adopted
or formally referenced in an ISO/IEC or ITU-T international standard.}
&
\cellhead{Security-by-Design the unconditional hardware and crypto default}
$\geq$60\% of EU critical infrastructure edge nodes deploy EUCC~`High'
certified EU-designed secure elements. FHE is operational end-to-end
from IoT sensor to HPC compute node. Autonomous Island Mode provides
validated 72-hour survivability as a mandatory operational baseline
across all NIS2-essential and DORA-critical edge deployments. No new
EU ICT product enters the market without a hardware root of trust.
The question ``does this deployment have hardware-level integrity?''
has no non-trivial answer for any new EU deployment.
\kpistrip{Phase~3 deployment normalisation: $\geq$60\% market penetration
for EU secure elements in critical infra; 100\% of new EU ICT products
with hardware RoT as a default.}
&
\cellhead{SSL: Full Sovereign Autonomy}
EU-designed and EU-manufactured secure elements hold $\geq$60\% of new
critical infrastructure procurement market share. Foreign legacy
dependency operates under binding replacement programmes with defined
timelines. EU sovereign RISC-V silicon and open-hardware TEE
architectures are exported to and adopted by partner jurisdictions.
EU PQC and hardware TEE standards are the international reference.
Structural exposure to geopolitical disruption of any single non-EU
hardware supplier has been eliminated for new deployments. The EU
has moved from defensive sovereignty (reducing dependency) to offensive
strategic posture (defining the terms of global hardware security
competition).
\kpistrip{Full Sovereign Autonomy is verified by $\geq$60\% market share
in new critical infra procurement plus the existence of binding
replacement programmes for legacy foreign hardware --- not by 100\%
market share.}
\\
\hline

\end{longtable}
\end{landscape}


% ---- PROGRESSION MATRIX: CLUSTER 1 ----

% \noindent
% \resizebox{\textwidth}{!}{%
% \begin{tabular}{|p{4.2cm}|p{5.6cm}|p{5.6cm}|p{5.6cm}|}
% \hline
% \rowcolor{gray!30}
% \multicolumn{4}{|c|}{\textbf{Progression Matrix: Priority Cluster 1 ---
%   Sovereign Trusted Infrastructure (P1, P2, P3, P8, P9)}} \\
% \hline
% \rowcolor{gray!15}
% \textbf{Metric} &
% \textbf{Phase 1: Foundation \& Compliance (2025--2030)} &
% \textbf{Phase 2: Federation \& Standardisation (2030--2035)} &
% \textbf{Phase 3: Autonomy \& Leadership (2035--2040)} \\
% \hline
% \rowcolor{gray!10}
% \multicolumn{4}{|l|}{\textit{\textbf{Cluster-Level Progression Indicators}}} \\
% \hline
% \textbf{Cluster ODP} &
% Applied R\&D &
% Standardised Integration &
% Ubiquitous Operation \\
% \hline
% \textbf{Cluster SSL} &
% High Foreign Dependency &
% Emerging EU Ecosystem &
% Full Sovereign Autonomy \\
% \hline
% \rowcolor{gray!10}
% \multicolumn{4}{|l|}{\textit{\textbf{Sub-Priority Decomposition (Individual ODP / SSL per Phase)}}} \\
% \hline
% \textit{P1: Hardware RoT \newline \& EUCC `High'} &
% ODP: Fragmented Piloting \newline
% SSL: High Foreign Dependency &
% ODP: Standardised Integration \newline
% SSL: Emerging EU Ecosystem &
% ODP: Ubiquitous Operation \newline
% SSL: Full Sovereign Autonomy \\
% \hline
% \textit{P2: PQC Agility \newline for Edge/IoT} &
% ODP: Applied R\&D \newline
% SSL: High Foreign Dependency &
% ODP: Fragmented Piloting \newline
% SSL: Emerging EU Ecosystem &
% ODP: Ubiquitous Operation \newline
% SSL: Full Sovereign Autonomy \\
% \hline
% \textit{P3: Lightweight Crypto \newline \& Physical-Layer Sec.} &
% ODP: Applied R\&D \newline
% SSL: High Foreign Dependency &
% ODP: Standardised Integration \newline
% SSL: Emerging EU Ecosystem &
% ODP: Ubiquitous Operation \newline
% SSL: Full Sovereign Autonomy \\
% \hline
% \textit{P8: Confidential Computing \newline \& Privacy-Pres. AI} &
% ODP: Applied R\&D \newline
% SSL: High Foreign Dependency &
% ODP: Fragmented Piloting \newline
% SSL: Emerging EU Ecosystem &
% ODP: Ubiquitous Operation \newline
% SSL: Full Sovereign Autonomy \\
% \hline
% \textit{P9: Resilient-by-Design \newline \& Island Mode} &
% ODP: Fragmented Piloting \newline
% SSL: High Foreign Dependency &
% ODP: Standardised Integration \newline
% SSL: Contained Local Survivability &
% ODP: Ubiquitous Operation \newline
% SSL: Systemic EU Resilience \\
% \hline
% \rowcolor{blue!7}
% \textbf{Operational State} &
% EU-funded RISC-V secure element pilot projects commence EUCC evaluation;
% NIST FIPS~203/204/205 algorithms are adopted by research consortia and
% optimised for constrained devices in laboratory settings; Intel TDX
% and AMD SEV-SNP dominate all operational TEE deployments while
% EU open-hardware alternatives remain at prototype stage;
% edge nodes operate with hard cloud dependencies and lack
% autonomous fallback architectures validated at operational scale. &
% EU-manufactured EUCC `High' certified secure elements enter
% critical infrastructure deployments across energy and
% telecommunications sectors; hybrid ML-KEM plus ECDH-P384 schemes
% are standardised and deployed in Industrial IoT corridors;
% RISC-V Keystone-based TEEs are operational in federated
% multi-provider Data Spaces; validated 72-hour Island Mode
% survivability demonstrated in designated EU 5G cross-border
% corridor pilots under SNS~JU governance. &
% EU Chips Act sovereign fabrication supplies $\geq$60\% of EU demand
% for security-critical silicon; all new EU ICT products ship with
% EUCC `High' certified hardware roots of trust by default;
% fully homomorphic encryption is operationalised for
% end-to-end confidential processing from IoT to HPC;
% bio-inspired swarm protocols coordinate autonomous
% cyber-physical recovery without any cloud intervention,
% achieving DORA-mandated recovery time objectives
% autonomously at continental scale. \\
% \hline
% \rowcolor{orange!12}
% \textbf{Critical Enabler} &
% CRA Articles~10 and~11 enforcement (2024--2027) creating
% mandatory market demand for hardware security certification,
% establishing the commercial conditions for EU secure element
% pilot investment &
% EUCC `High' certification of the first EU-manufactured
% RISC-V secure element, constituting the gate condition
% for the hardware sovereignty transition and unlocking
% EU Chips Act fabrication investment justification &
% EU Chips Act sovereign EUV-capable fabrication capacity
% achieving $\geq$60\% market self-sufficiency in
% security-critical silicon, removing the last structural
% foreign dependency in the infrastructure trust chain \\
% \hline
% \end{tabular}%
% }

% =========================================================================
% PRIORITY CLUSTER 2: COGNITIVE SECURITY AND AUTONOMOUS CYBER DEFENCE
% =========================================================================

\newpage
\section{Priority Cluster 2: Cognitive Security and Autonomous Cyber Defence}
\label{sec:cluster2}

\noindent\textbf{Constituent Sub-Priorities:} P6 (Cognitive SecOps / LLMSecOps \& Automated
Threat Intelligence), P10 (Proactive APT Management), and the AI-as-target dimension
drawn from CL1.2 (Model Integrity / DataBOMs), CL1.3 (Dynamic Supply Chain Verification /
Runtime SBOMs), and CL4.7 (AI Trustworthiness \& Algorithmic Accountability).

\medskip
\noindent\textbf{Cluster Thesis:} Artificial intelligence is simultaneously the most powerful
instrument available for autonomous cyber defence and the most consequential new attack surface
introduced into the security architecture. A roadmap that deploys AI agents for autonomous
threat containment without hardening those agents against adversarial manipulation creates a
critical second-order vulnerability. This cluster therefore addresses three inseparable
sub-dimensions: AI as security instrument (LLMSecOps and APT hunting), AI as adversary
amplifier (AI-generated APT tooling), and AI as target (model poisoning, adversarial evasion,
and algorithmic accountability for autonomous containment decisions).

\subsection*{7-Dimensional Analysis: Priority Cluster 2}

\begin{enumerate}

% ---- DIMENSION 1 ----
\item \textbf{Baseline \& Target Objective}

\textit{Baseline (2025--2026):}
European Security Operations Centres (SOCs) rely predominantly on rule-based SIEM platforms ---
Splunk (US-domiciled, acquired by Cisco in 2023) and IBM QRadar dominate enterprise
deployments. AI-based tools are deployed as analyst ``copilots'' with human analysts
retaining all containment authority; no EU-sovereign foundation models for security
operations exist at production scale. APT dwell times in European enterprise environments
average 150--200 days, with state-sponsored threat actors demonstrating extended persistence
in EU critical infrastructure. The adversarial robustness of AI security agents --- their
resistance to prompt injection, model evasion, and training data poisoning --- is addressed
only in research contexts without operational certification frameworks. NWDAF integration
per 3GPP TS~23.288~\cite{threeGPP23288} exists in 5G standalone deployments but remains
fragmented across operators, and no EU-mandated framework for the trustworthiness of
autonomous security AI exists prior to AI Act enforcement.

Current cluster ODP: Applied R\&D (AI-as-target and sovereign LLM development); Siloed AI
Copilots (P6); Multi-Layered Alert Correlation (P10). Cluster ODP: Applied R\&D.
Cluster SSL: High Reliance on US Foundation Models / High Dependency on Global Threat Feeds.

\textit{Target (2040):}
\begin{sloppypar}
Swarm-based agentic AI defends the EU Cloud-Edge Continuum autonomously, with APT dwell
time reduced to under 24 hours in certified critical infrastructure deployments. EU-sovereign
edge Small Language Models (SLMs) and Large Context Models (LCMs) for security operations
are certified under AI Act Article~9 and cover at least 50\% of EU critical infrastructure
SOC deployments. All autonomous security agents pass mandatory adversarial robustness
validation before operational deployment. Moving Target Defense reduces APT reconnaissance
success rates below 5\% in validated simulation environments. EU AI security frameworks
are adopted as global reference implementations by partner jurisdictions.
\end{sloppypar}

% ---- DIMENSION 2 ----
\item \textbf{Primary Drivers}

\textit{Threat Drivers:}
\begin{itemize}
  \item Nation-state APT campaigns (attributed by EU agencies to APT28/GRU,
        APT29/SVR, and Sandworm) increasingly exploit LLM hallucination and
        prompt injection vulnerabilities to manipulate autonomous security agents
        into executing incorrect containment actions --- weaponising the defence
        architecture against itself.
  \item AI-generated spear-phishing and poly\-morphic malware
        enabling automated, highly-targeted APT campaigns that adapt
        faster than signature-based
        detection cycles, invalidating the static detection rules that rule-based
        SIEM platforms depend upon.
  \item Supply chain attacks targeting AI model training pipelines: data poisoning
        and backdoor insertion into security AI models deployed at continental scale
        can silently degrade detection capability without triggering observable
        anomalies in standard monitoring.
  \item Systematic abuse of MITRE ATT\&CK-catalogued techniques (T1027 Obfuscated
        Files or Information; T1055 Process Injection) within AI-assisted
        reconnaissance tools, enabling more efficient bypass of detection rules
        derived from those same frameworks.
\end{itemize}

\textit{Regulatory Drivers:}
\begin{itemize}
  \item AI Act Article~9: risk management system requirements for high-risk AI
        systems --- autonomous security containment agents making access revocation
        or network isolation decisions qualify as high-risk systems requiring
        mandatory conformity assessment before operational deployment~\cite{aiact}.
  \item AI Act Article~15: accuracy, robustness, and cybersecurity requirements
        for high-risk AI, including specific resilience requirements against
        adversarial manipulation of inputs and training data~\cite{aiact}.
  \item AI Act Article~13: transparency and provision of information for high-risk
        AI --- autonomous security decisions must be explainable to affected
        entities and competent national authorities~\cite{aiact}.
  \item NIS2 Article~21(2)(b): incident handling procedures, including
        AI-augmented detection and response, for essential entities~\cite{nis2}.
\end{itemize}

% ---- DIMENSION 3 ----
\item \textbf{High Impact Activities}

\begin{enumerate}[label=(\alph*)]
  \item \textbf{[Phase~1]} Deploy NWDAF-integrated (3GPP TS~23.288~\cite{threeGPP23288})
        telemetry pipelines feeding ML-based UEBA (User and Entity Behaviour Analytics)
        for APT indicator correlation; establish behavioural baseline profiles across
        federated enterprise environments targeting reduction of APT dwell time from
        the current 150--200 day average to under 72 hours in pilot deployments via
        ML-assisted Indicator of Compromise (IoC) pattern detection.

  \item \textbf{[Phase~1--2]} Develop EU-sovereign Small Language Models (SLMs,
        under 10B parameters) and Large Context Models (LCMs) for edge-native SOC
        operations (LLMSecOps), trained on MITRE ATT\&CK-aligned threat intelligence
        corpora under GDPR-compliant data governance~\cite{gdpr}; establish mandatory
        AI Act Article~9 conformity assessment pipelines for all autonomous security
        agent deployments prior to production authorisation~\cite{aiact}.

  \item \textbf{[Phase~2]} Deploy edge-native ``Guardian'' nodes executing autonomous
        Zero-Trust reconfigurations per NIST SP~800-207 ZTA principles~\cite{nistsp800207}
        and Moving Target Defense (MTD) --- dynamically shuffling IP addressing, network
        topology, and software diversity to disrupt APT reconnaissance --- governed by
        ETSI GS~ZSM~009 closed-loop automation standards~\cite{etsizsm009}.

  \item \textbf{[Phase~2]} Establish mandatory adversarial robustness validation
        pipelines for all AI security agents prior to operational deployment: automated
        red-team adversarial testing (prompt injection, model evasion, backdoor
        detection via CL1.4 purple-team automation), with continuous certification
        under a CEN/CENELEC-standardised AI security agent assurance framework
        aligned with CL4.7 AI Trustworthiness requirements.

  \item \textbf{[Phase~3]} Deploy swarm-based agentic AI coordination protocols
        enabling decentralised, self-evolving threat intelligence synthesis across
        the EU Cloud-Edge Continuum; integrate with Deception Technology
        infrastructure (CL2.6: honeypots, honeynets, credential canaries, fake API
        endpoints) to generate high-fidelity, low-noise adversarial ground truth for
        continuous autonomous agent retraining.
\end{enumerate}

% ---- DIMENSION 4 ----
\item \textbf{Catalysts \& Enablers}

\textit{Funding Alignment:}
EuroHPC Joint Undertaking AI Factories for training EU-sovereign security foundation
models at the compute scale required; ECCC Strategic Agenda priority areas ``Cybersecurity
of AI'' and ``AI-enabled cybersecurity tools and services''~\cite{eccc}; Horizon Europe
Cluster~3 (Civil Security for Society) calls targeting AI-driven threat intelligence and
autonomous incident response; EU AI Office for conformity assessment infrastructure
supporting high-risk security AI certification under the AI Act.

\textit{Standardisation:}
ETSI GS~ZSM~009 (Zero-Touch Network and Service Management, Closed-Loop
Automation~\cite{etsizsm009}); 3GPP TS~23.288 (NWDAF architecture for AI-native 5G
network analytics~\cite{threeGPP23288}); MITRE ATT\&CK framework (behavioural threat
intelligence taxonomy for agent training corpora); NIST SP~800-207 (Zero Trust
Architecture~\cite{nistsp800207}); CEN/CENELEC JTC~13 working groups on AI security
agent assurance standardisation.

\textit{Ecosystem Enablers:}
ENISA AI cybersecurity threat landscape guidelines; federated Cyber Range infrastructure
(CL3.5) for adversarial agent testing under realistic attack scenarios; CONCORDIA
successor programme for EU-wide SOC federation and cross-border CTI sharing~\cite{concordia};
EU AI Act notified conformity assessment bodies for high-risk AI product certification.

% ---- DIMENSION 5 ----
\item \textbf{Disruptors \& Systemic Risks}

\begin{itemize}
  \item \textbf{Adversarial AI arms race:} threat actors weaponising LLMs for
        automated zero-day discovery and polymorphic exploit generation may outpace
        the defensive AI maturation timeline --- specifically, AI-generated
        vulnerability discovery could exceed human patching bandwidth within the
        Phase~1--2 window, before EU-sovereign autonomous response agents are
        certified for operational deployment.
  \item \textbf{LLM hallucination in containment agents:} false positive isolation
        of critical infrastructure nodes (e.g., a hospital network segment incorrectly
        quarantined during a medical emergency) constitutes an operational safety risk
        that may impose worse harm than the threat the agent was designed to
        contain --- this is the central argument for mandatory HITL escalation in Phase~1.
  \item \textbf{EU sovereign LLM compute deficit:} training security-relevant
        foundation models requires EuroHPC-scale compute; as of 2025--2026, EU frontier
        model training capacity lags US and China by approximately one order of
        magnitude, constraining Phase~1 EU-sovereign models to SLM/LCM scale
        ($<$10B parameters) and limiting capability relative to US foundation models.
  \item \textbf{AI model poisoning via CTI feed corruption:} if adversaries inject
        false indicators into NWDAF analytics pipelines or Cyber Threat Intelligence
        (CTI) sharing platforms (CL2.1), behavioural baselines trained on poisoned
        data will silently misclassify legitimate traffic and fail to detect real APTs
        --- a harder failure mode to detect than outright system unavailability.
  \item \textbf{AI Act certification bottleneck:} mandatory conformity assessment
        for high-risk autonomous security agents may create a 24--36 month regulatory
        gap during which EU organisations operate without certified autonomous
        containment, relying on slower human-in-the-loop processes that cannot match
        AI-assisted adversary speed.
\end{itemize}

% ---- DIMENSION 6 ----
\item \textbf{Failsafe Mechanisms \& Resilience}

\begin{itemize}
  \item \textbf{Shadow mode operation:} all autonomous AI security agents run in
        parallel with human-supervised rule-based SIEM during Phases~1--2;
        containment actions affecting Tier-1 critical infrastructure (energy,
        healthcare, transport per NIS2 Annex~I) require mandatory Human-In-The-Loop
        (HITL) escalation regardless of AI confidence score, with reversible quarantine
        actions bounded by a defined 60-minute rollback window~\cite{nis2}.
  \item \textbf{Model diversity and behavioural circuit breakers:} deploy a minimum
        of two architecturally distinct AI detection models (e.g., transformer-based
        LLM plus graph neural network UEBA) to prevent a single-model adversarial
        exploit from disabling the entire detection capability; behavioural circuit
        breakers automatically suspend any agent exhibiting an anomalous action rate
        or pattern deviating beyond $3\sigma$ from its established operational
        baseline, reverting to rule-based SIEM until re-certification.
  \item \textbf{RLHF continuous correction:} incorporate analyst override decisions
        into continuous Reinforcement Learning from Human Feedback (RLHF) retraining
        to prevent model drift from training distribution; quarterly adversarial
        red-team audits (CL1.4) detect capability degradation before it reaches
        operational significance.
  \item \textbf{Air-gapped CTI validation:} all external CTI feeds ingested by
        NWDAF analytics pipelines pass through an isolated validation environment
        where indicator quality is assessed against independent behavioural ground
        truth before poisoned data can influence operational AI model weights.
  \item \begin{sloppypar}\textbf{Deception-backed calibration:} honeynet and credential canary
        infrastructure (CL2.6) provides high-fidelity, low-noise ground truth against
        which AI behavioural models are continuously calibrated, reducing false
        positive rates and providing an independent signal for detecting model
        poisoning via discrepancy analysis between canary alerts and AI model
        predictions.\end{sloppypar}
\end{itemize}

% ---- DIMENSION 7 ----
\item \textbf{Transversal Meta-Layer}

\textit{Sovereignty \& Autonomy:}
This cluster eliminates reliance on US-domiciled SIEM platforms (Splunk/Cisco, IBM QRadar)
and US foundation models for operational security intelligence. EU-sovereign SLMs trained
under GDPR-compliant data governance~\cite{gdpr} ensure that sensitive threat intelligence
and behavioural telemetry does not transit non-EU cloud infrastructure. EU leadership in
trustworthy AI security tools creates an exportable technology capability aligned with
the AI Act's ambition for global regulatory standard-setting.

\textit{Trust \& Ethics:}
AI Act Article~13 (transparency requirements) is mandated for all autonomous security
decisions --- affected entities must receive meaningful explanations for AI-driven access
revocations or network isolation actions~\cite{aiact}. GDPR Article~22 (restrictions on
automated individual decision-making) applies where AI containment decisions affect
individual data subjects; human review mechanisms must be structurally embedded, not
optional add-ons~\cite{gdpr}. AI Act Article~15 adversarial robustness requirements
ensure security AI systems cannot be weaponised against the infrastructure they are
designed to protect~\cite{aiact}.

\textit{Sustainable Security (Green IT):}
Edge-native SLMs under 10B parameters reduce inference energy consumption by two to
three orders of magnitude compared to centralised GPT-class model inference, enabling
real-time threat classification at the point of detection without data centre energy
overhead. Federated training paradigms reduce repeated data transmission overhead
compared to centralised training architectures and are compatible with the privacy
preservation requirements of GDPR Article~25~\cite{gdpr}.

\end{enumerate}

% =========================================================================
% CLUSTER 2: COGNITIVE SECURITY AND AUTONOMOUS CYBER DEFENCE
% =========================================================================
\newpage
\subsection{Progression Matrix: Priority Cluster 2 --- Cognitive Security and Autonomous
Cyber Defence}

\begin{landscape}
\setlength{\textwidth}{26.5cm}\setlength{\linewidth}{26.5cm}%
\setlength{\hsize}{26.5cm}\setlength{\columnwidth}{26.5cm}
\setlength{\tabcolsep}{5pt}
\renewcommand{\arraystretch}{1.35}

\begin{longtable}{%
|>{\footnotesize\raggedright\arraybackslash}p{2.8cm}%
|>{\footnotesize\raggedright\arraybackslash}p{5.4cm}%
|>{\footnotesize\raggedright\arraybackslash}p{5.4cm}%
|>{\footnotesize\raggedright\arraybackslash}p{5.4cm}%
|>{\footnotesize\raggedright\arraybackslash}p{5.4cm}|}

\caption{\textbf{Cluster~2: Cognitive Security \& Autonomous Cyber Defence}
(P6, P10, AI-as-target dimension from CL1.2, CL1.3, CL4.7).
Transition matrix across the four methodology dimensions. The AI-as-target
dimension (model integrity, adversarial robustness, algorithmic accountability)
is labelled \textit{(AI-T)} where it diverges from P6 or P10 descriptors.}
\label{tab:cluster2}\\

\hline
\rowcolor{colheader}
\textbf{Phase / SSL} &
\textbf{R\&I Maturity}{\scriptsize ODP indicator per sub-priority} &
\textbf{Policy, Regulatory \& Standards}{\scriptsize enforcement $\cdot$
  certification $\cdot$ standardisation} &
\textbf{Operational Deployment \& Adoption}{\scriptsize actor types $\cdot$
  scale $\cdot$ sector coverage} &
\textbf{Market Uptake \& Strategic Sovereignty}{\scriptsize SSL indicator
  $\cdot$ dependency reduction $\cdot$ industrial capacity} \\
\hline
\endfirsthead

\hline
\rowcolor{colheader}
\textbf{Phase / SSL} &
\textbf{R\&I Maturity} &
\textbf{Policy, Regulatory \& Standards} &
\textbf{Operational Deployment \& Adoption} &
\textbf{Market Uptake \& Strategic Sovereignty} \\
\hline
\endhead
\hline\multicolumn{5}{r}{\footnotesize\textit{Continued\ldots}}\\ \endfoot
\hline \endlastfoot

% ---- PHASE 1 ----
\rowcolor{phase1row}
\textbf{Phase~1}
\newline\textit{Foundation \&}
\newline\textit{Compliance}
\newline\textit{2025--2030}
\newline\sslbadge{sslredbg}{sslredtx}{SSL: High Reliance on}
\newline\sslbadge{sslredbg}{sslredtx}{US Foundation Models}
&
\cellhead{Cluster ODP: Applied R\&D / Siloed AI Copilots}
\subp{P6} reaches Siloed AI Copilots: commercial LLMs (GPT-4-class)
assist human SOC analysts in alert triage, threat report generation, and
MITRE ATT\&CK technique mapping; no EU-sovereign LLM for security
operations exists at production scale; NWDAF (3GPP TS~23.288)
telemetry pipelines are integrated in 5G SA pilot deployments for
anomaly detection. \subp{P10} is at Multi-Layered Alert Correlation:
ML-based UEBA is applied to APT IoC detection in research and pilot
contexts; APT dwell time reduction from the 150--200 day average is
demonstrated in controlled settings. \subp{AI-T} is at Applied R\&D:
adversarial robustness testing (prompt injection, model evasion,
DataBOM verification) is an academic research function without
operational certification frameworks.
\kpistrip{Phase~1 R\&I distinguishes three sub-dimensions with
different maturity levels. Conflating them into a single ODP cell
obscures that AI-as-target capability is the least mature and the
most consequential gap.}
&
\cellhead{AI Act enforcing; conformity assessment infrastructure absent}
AI~Act Articles~9 (risk management), 13 (transparency), and 15
(adversarial robustness) enter applicability in 2025--2026, establishing
mandatory requirements for autonomous security AI as high-risk systems.
However, no EU conformity assessment bodies with capacity to certify
autonomous security agents are operational at the start of Phase~1.
ETSI GS~ZSM~009 (closed-loop automation) is published but not mandated.
NIS2 Article~21(2)(b) incident handling procedures are being
operationalised across essential entities. ENISA AI cybersecurity threat
landscape guidelines are published. The critical regulatory gap is the
absence of a standardised adversarial robustness test suite applicable
to security-specific AI agents.
\kpistrip{The AI Act creates obligations before the conformity
assessment infrastructure exists to fulfil them. This 24--36 month gap
is the dominant Phase~1 policy risk for this cluster.}
&
\cellhead{Deployment: SOC copilot use by large enterprises and national CSIRTs}
US-platform LLMs (GPT-4-class, under data processor agreements) operate
as analyst copilots in large enterprise and national CSIRT SOCs.
Rule-based SIEMs (Splunk/Cisco, IBM~QRadar) remain the primary
detection infrastructure across EU regulated sectors. NWDAF-integrated
AI analytics operate in 3--5 5G SA deployments run by major EU
telecom operators. No certified autonomous security containment agent
is in operational deployment anywhere in the EU. APT dwell times of
150--200 days in EU enterprise environments reflect the gap between
AI-assisted detection capability and operational deployment at scale.
SMEs lack access to AI-powered SOC tooling at viable cost.
\kpistrip{Phase~1 deployment is concentrated in the largest institutional
actors (national CSIRTs, Tier-1 operators, major financial entities).
The gap between these deployments and SME access defines the adoption
challenge for Phase~2.}
&
\cellhead{SSL: High Reliance on US Foundation Models}
The EU SOC tooling market is dominated by US-domiciled platforms:
Splunk (Cisco), IBM QRadar, and Microsoft Sentinel collectively control
the EU enterprise SIEM market. The foundation models used for SOC
copilot functions are US-origin (OpenAI, Anthropic, Google) and are
accessed under data processor agreements, creating persistent
extra-territorial data exposure. No EU-sovereign foundation model at
production scale for security operations exists. EuroHPC compute
capacity is insufficient for frontier-class security model training.
The adversarial robustness validation market is pre-commercial.
\kpistrip{Quantified dependency: US platforms hold $\geq$80\% of EU
enterprise SIEM market; US-origin models are the sole source for
LLM-based SOC tooling. This specific dependency structure is what
Phases~2 and~3 must reduce.}
\\
\hline

% ---- GATE 1->2 ----
\rowcolor{gaterow}
\gatelbl{1 $\rightarrow$ 2} &
\multicolumn{4}{>{\footnotesize\raggedright\arraybackslash}p{22.0cm}|}{%
\cellcolor{gaterow}%
\textbf{R\&I trigger:} At least one EU-sovereign Small Language Model
($\leq$10B parameters) trained on GDPR-compliant security corpora has
been validated against an independently benchmarked EU critical infrastructure
SOC use case at production-equivalent scale.
\textbf{Policy trigger:} AI~Act conformity assessment bodies with capacity
to certify autonomous security AI as high-risk systems under Articles~9 and~15
are operational and have issued at least one certificate to a security AI
agent product.
\textbf{Operational trigger:} At least one EU-sovereign LLMSecOps agent is
in certified operational deployment in a named NIS2-essential entity,
demonstrating APT dwell time below 72 hours in a documented production
environment.
\textbf{Market trigger:} A standardised adversarial robustness test suite
for security AI agents, developed under CEN/CENELEC or ETSI TC~CYBER
governance, has been published and is referenced in at least one EU
procurement framework for critical infrastructure AI tools.}
\\
\hline

% ---- PHASE 2 ----
\rowcolor{phase2row}
\textbf{Phase~2}
\newline\textit{Applied \& Industrial}
\newline\textit{R\&D alongside}
\newline\textit{Federation \&}
\newline\textit{Standardisation}
\newline\textit{2030--2035}
\newline\sslbadge{ssltransbg}{ssltrantx}{SSL: Sovereign EU Edge}
\newline\sslbadge{ssltransbg}{ssltrantx}{LLMs Emerging}
&
\cellhead{Cluster ODP: Federated Edge Agents / Automated Triage}
\subp{P6} reaches Federated Edge Agents: EU-sovereign SLMs ($\leq$10B
parameters) trained on MITRE ATT\&CK-aligned corpora are operational
at the edge for real-time SOC analytics without cloud inference
dependency; Guardian nodes execute autonomous Zero-Trust
reconfigurations per NIST~SP~800-207. \subp{P10} reaches Automated
Triage and LLMSecOps Containment: LLMSecOps agents ingest CTI feeds,
track APT lateral movement against behavioural frameworks, and
orchestrate incident triage autonomously under ETSI GS~ZSM~009
governance; APT dwell time $\leq$72 hours in certified pilot deployments.
\subp{AI-T} reaches Standardised Integration: adversarial robustness
validation pipelines (prompt injection, model evasion, backdoor
detection) are standardised and mandatory under AI~Act~Article~9
conformity assessment; continuous RLHF retraining and circuit-breaker
mechanisms are in production.
\kpistrip{Phase~2 R\&I includes active applied and industrial R\&D for
all three sub-dimensions --- swarm AI protocols, post-quantum AI
security, and next-generation adversarial testing frameworks --- not
merely consolidation of Phase~1 results.}
&
\cellhead{AI Act conformity assessment operational: the Phase~2 gate}
The single Phase~2 policy gate is the operationalisation of a
standardised AI~Act conformity assessment framework for autonomous
security agents, enabling certified deployment without per-case
regulatory approval. ETSI GS~ZSM~009 is mandated within EU NaaS
federation governance frameworks. A CEN/CENELEC-standardised adversarial
robustness test suite for security AI is published. MITRE ATT\&CK
behavioural framework is formally referenced in NIS2 sector-specific
incident handling guidance. EuroHPC AI Factory compute contracts for
EU-sovereign security model training are active.
\kpistrip{Without operational AI~Act conformity assessment
infrastructure, autonomous security agents cannot be certified for
regulated-sector deployment regardless of technical readiness.
This is the dominant Phase~2 policy dependency.}
&
\cellhead{Certified EU sovereign AI agents in critical infrastructure SOCs}
EU SLM-based LLMSecOps agents are in certified operational deployment
in energy, financial (DORA), and telecommunications critical
infrastructure SOCs in $\geq$3 Member States. Guardian nodes execute
autonomous Zero-Trust reconfigurations under human-in-the-loop
escalation protocols for Tier-1 critical actions. NWDAF-integrated
analytics are mandatory in 5G SA deployments. Deception infrastructure
(honeynets, credential canaries, fake API endpoints per CL2.6) is
deployed alongside AI detection systems, providing high-fidelity
ground truth for model calibration. APT dwell time $\leq$72 hours
demonstrated in pilot critical infrastructure deployments.
\kpistrip{Phase~2 deployment evidence requires named critical
infrastructure sector, Member State scope, and a documented APT dwell
time metric from an operational --- not simulated --- environment.}
&
\cellhead{SSL: Sovereign EU Edge LLMs Emerging}
EU-sovereign SLMs cover $\geq$30\% of critical infrastructure SOC
deployments in regulated sectors, rising from zero at baseline.
EuroHPC AI Factories supply training infrastructure for EU-origin
security models, reducing dependence on US hyperscaler compute.
US-platform SIEMs retain market share in legacy deployments but
face procurement competition from certified EU alternatives.
Adversarial robustness validation is a commercially available service
with EU-domiciled providers. EU AI security tool exports to partner
jurisdictions have begun, indicating the transition from security
consumer to security producer posture.
\kpistrip{SSL evidence at Phase~2: $\geq$30\% critical infra SOC market
share for EU-sovereign AI tools; at least one documented EU AI security
tool export event to a non-EU partner jurisdiction.}
\\
\hline

% ---- GATE 2->3 ----
\rowcolor{gaterow}
\gatelbl{2 $\rightarrow$ 3} &
\multicolumn{4}{>{\footnotesize\raggedright\arraybackslash}p{22.0cm}|}{%
\cellcolor{gaterow}%
\textbf{R\&I trigger:} EU-certified swarm AI coordination protocols
enabling decentralised autonomous threat intelligence synthesis have passed
independent security evaluation and are deployed in at least one cross-border
EU critical infrastructure corridor under ETSI GS~ZSM~009 governance.
\textbf{Policy trigger:} EU AI security standards (adversarial robustness
profiles, autonomous agent accountability frameworks) are adopted or formally
referenced in at least one ISO/IEC or ITU-T international standard, confirming
EU regulatory leadership beyond domestic jurisdiction.
\textbf{Operational trigger:} APT dwell time $\leq$24 hours is documented in
a named NIS2-essential entity operational deployment (not a simulation),
attributable to autonomous AI-driven detection and containment.
\textbf{Market trigger:} EU-sovereign AI security tools hold $\geq$50\% market
share in new critical infrastructure SOC deployments across at least five
Member States, displacing US-platform SIEMs as the procurement default in
regulated sectors.}
\\
\hline

% ---- PHASE 3 ----
\rowcolor{phase3row}
\textbf{Phase~3}
\newline\textit{Autonomy \&}
\newline\textit{Leadership}
\newline\textit{2035--2040}
\newline\sslbadge{ssltealbg}{ssltealtx}{SSL: Global Leadership in}
\newline\sslbadge{ssltealbg}{ssltealtx}{Trustworthy AI Security}
&
\cellhead{Cluster ODP: Ubiquitous AI Immune System}
All three sub-dimensions reach Ubiquitous Operation. Swarm-based
agentic AI defends the EU Cloud-Edge Continuum via self-evolving,
decentralised threat intelligence without any central coordination
dependency. Moving Target Defense continuously shuffles IP addressing,
network topology, and software diversity; validated APT reconnaissance
success rates fall below 5\% in certified operational environments.
Adversarial robustness validation is continuous and fully automated:
no security AI agent operates in EU critical infrastructure without
an active conformity certification. Phase~3 R\&I addresses
post-swarm frontiers: quantum-native AI security architectures,
brain-inspired neuromorphic threat detection, and AI-to-AI
adversarial dynamics beyond current threat models.
\kpistrip{Phase~3 ODP evidence requires operational statistics from
production deployments: APT dwell time figures, MTD effectiveness
metrics, and autonomous containment accuracy rates from certified
environments --- not from cyber range simulations.}
&
\cellhead{EU AI security governance as international reference}
EU AI Act frameworks for trustworthy autonomous security agents are
adopted internationally and referenced in bilateral digital
partnership agreements. Autonomous agent certification operates at
continental scale without per-case regulatory review. Zero-touch
network management under ETSI GS~ZSM~009 governs all EU NaaS
operations. The EU AI Office and ENISA jointly define the international
standard for security AI accountability. EU AI security governance
has transitioned from reactive adaptation to agenda-setting: other
jurisdictions adopt EU frameworks to access EU market and partnership
programmes.
\kpistrip{Phase~3 policy leadership is verified by international
adoption events --- not by the number of EU instruments in force.
At least one non-EU jurisdiction formally adopting EU security AI
governance standards constitutes the Phase~3 policy evidence threshold.}
&
\cellhead{Autonomous continental-scale AI cyber defence as operational default}
APT dwell time $\leq$24 hours in all certified EU NIS2-essential
critical infrastructure deployments. MTD is deployed as a mandatory
operational baseline. Swarm AI coordination operates without central
cloud dependency: each EU Cloud-Edge node participates in a
self-organising threat intelligence mesh. Deception infrastructure
is fully integrated into the autonomous response loop, providing
continuous model calibration ground truth. Human intervention is
reserved for exceptional legal escalation scenarios; routine security
operations are fully autonomous within the governance framework.
\kpistrip{The Phase~3 deployment test: is autonomous AI defence
the unconditional operational baseline, or is it still a premium
programme? 100\% of new critical infrastructure deployments
incorporating autonomous AI security constitutes the evidence threshold.}
&
\cellhead{SSL: Global Leadership in Trustworthy AI Security}
EU-sovereign AI security tools hold $\geq$50\% of new critical
infrastructure SOC deployments across the EU. US-platform SIEMs
are legacy infrastructure under active replacement programmes with
defined timelines. EU AI security models are exported to and
commercially adopted by partner jurisdictions. EU sets the commercial
terms for security AI market access: certification against EU AI~Act
Article~9/15 profiles is a de facto global procurement requirement.
The EU has transitioned from defensive sovereignty to offensive
strategic posture: it defines the technical terms on which the
global security AI industry must compete.
\kpistrip{Full sovereign leadership in this cluster requires both
domestic dominance ($\geq$50\% critical infra SOC market share)
and international influence (at least two partner jurisdictions
formally adopting EU security AI governance). Both conditions must
be met.}
\\
\hline

\end{longtable}
\end{landscape}

% ---- PROGRESSION MATRIX: CLUSTER 2 ----

% \noindent
% \resizebox{\textwidth}{!}{%
% \begin{tabular}{|p{4.2cm}|p{5.6cm}|p{5.6cm}|p{5.6cm}|}
% \hline
% \rowcolor{gray!30}
% \multicolumn{4}{|c|}{\textbf{Progression Matrix: Priority Cluster 2 ---
%   Cognitive Security and Autonomous Cyber Defence (P6, P10, AI-as-Target)}} \\
% \hline
% \rowcolor{gray!15}
% \textbf{Metric} &
% \textbf{Phase 1: Foundation \& Compliance (2025--2030)} &
% \textbf{Phase 2: Federation \& Standardisation (2030--2035)} &
% \textbf{Phase 3: Autonomy \& Leadership (2035--2040)} \\
% \hline
% \rowcolor{gray!10}
% \multicolumn{4}{|l|}{\textit{\textbf{Cluster-Level Progression Indicators}}} \\
% \hline
% \textbf{Cluster ODP} &
% Applied R\&D / Siloed AI Copilots &
% Federated Edge Agents &
% Ubiquitous AI Immune System \\
% \hline
% \textbf{Cluster SSL} &
% High Reliance on US Foundation Models &
% Sovereign European Edge LLMs / SLMs &
% Global Leadership in Trustworthy AI Security \\
% \hline
% \rowcolor{gray!10}
% \multicolumn{4}{|l|}{\textit{\textbf{Sub-Priority Decomposition (Individual ODP / SSL per Phase)}}} \\
% \hline
% \textit{P6: Cognitive SecOps \newline (LLMSecOps)} &
% ODP: Siloed AI Copilots \newline
% SSL: High Reliance on US Foundation Models &
% ODP: Federated Edge Agents \newline
% SSL: Sovereign EU Edge LLMs Emerging &
% ODP: Ubiquitous AI Immune System \newline
% SSL: Global Leadership in Trustworthy AI \\
% \hline
% \textit{P10: Proactive APT \newline Management} &
% ODP: Multi-Layered Alert \newline Correlation \& Hunting \newline
% SSL: High Dependency on Global Threat Feeds &
% ODP: Automated Triage \& \newline LLMSecOps Containment \newline
% SSL: Sovereign Pan-EU CTI \& Behavioural Analytics &
% ODP: Ubiquitous Moving \newline Target Defense (MTD) \newline
% SSL: Full Autonomous Counter-Adversarial Resilience \\
% \hline
% \textit{AI-as-Target \newline (CL1.2, CL1.3, CL4.7)} &
% ODP: Applied R\&D \newline
% SSL: High Foreign Dependency &
% ODP: Standardised Integration \newline
% SSL: Emerging EU Ecosystem &
% ODP: Ubiquitous Operation \newline
% SSL: Full Sovereign Autonomy \\
% \hline
% \rowcolor{blue!7}
% \textbf{Operational State} &
% LLMs assist human SOC analysts in alert triage and threat report
% synthesis; NWDAF telemetry is routed to centralised AI for anomaly
% detection in 5G standalone deployments; APT dwell times are
% reduced via ML-assisted IoC correlation in pilot enterprise environments;
% adversarial robustness validation of AI security agents is conducted
% in academic and research settings without operational certification
% frameworks in force. &
% Edge-native EU-sovereign SLMs execute autonomous Zero-Trust
% reconfigurations and traffic throttling in certified critical
% infrastructure deployments; LLMSecOps agents ingest CTI and
% track lateral movement against MITRE ATT\&CK behavioural
% frameworks in real time; adversarial robustness validation
% is standardised under AI Act Article~9 conformity assessment;
% APT dwell time reduced to under 72 hours in EU critical
% infrastructure pilot deployments. &
% Swarm-based agentic AI defends the EU Cloud-Edge Continuum
% via self-evolving, decentralised threat intelligence without
% central coordination dependency; Moving Target Defense
% continuously shuffles network topology, IP addressing, and
% software diversity, achieving validated APT reconnaissance
% success rates below 5\% in operational simulation environments;
% EU AI security models are formally adopted as reference
% implementations by partner jurisdictions. \\
% \hline
% \rowcolor{orange!12}
% \textbf{Critical Enabler} &
% AI Act Articles~9 and~15 enforcement (2025--2026) establishing
% mandatory risk management and adversarial robustness
% requirements for autonomous security AI, creating the
% regulatory framework for certified Phase~2 deployment &
% Standardised AI Act conformity assessment framework for
% autonomous security agents, enabling certified operational
% deployment without per-case regulatory approval and
% unlocking Phase~3 swarm-level autonomy &
% EU-certified self-evolving AI swarm protocols achieving
% full operational deployment under ETSI GS~ZSM~009
% zero-touch automation governance with Pan-EU
% CTI-sharing infrastructure \\
% \hline
% \end{tabular}%
% }

% =========================================================================
% PRIORITY CLUSTER 3: FEDERATED OPERATIONAL SOVEREIGNTY
% =========================================================================

\newpage
\section{Priority Cluster 3: Federated Operational Sovereignty}
\label{sec:cluster3}

\noindent\textbf{Constituent Sub-Priorities:} P4 (Security Interoperability via Simpl
Middleware), P5 (Supply Chain Verification and Secured Infrastructure), P7 (Pan-European
NaaS Federation --- ``Schengen for Cloud-Edge/AI'').

\medskip
\noindent\textbf{Cluster Thesis:} Technical excellence in Cluster~1 and autonomous
defence capability in Cluster~2 deliver no strategic value if they cannot operate
verifiably across borders, across providers, and across supply chains. This cluster
defines the governance fabric, interoperability standards, and supply chain integrity
mechanisms that transform individually sovereign components into a collectively sovereign
European computing continuum. Critically, it addresses the legal and semantic as well as
the technical interoperability problem --- without a ``Schengen for Cloud'' legal framework,
cross-border compute federation remains commercially fragile regardless of technical
standardisation.

\subsection*{7-Dimensional Analysis: Priority Cluster 3}

\begin{enumerate}

% ---- DIMENSION 1 ----
\item \textbf{Baseline \& Target Objective}

\textit{Baseline (2025--2026):}
European cloud and edge infrastructure remains highly fragmented across proprietary provider
silos with incompatible security policies, authentication mechanisms, and data exchange
protocols. Common European Data Spaces (Manufacturing, Health, Energy, Agriculture) are in
early piloting stages, with security interoperability relying on bilateral agreements rather
than standardised automated mechanisms aligned with the Simpl open-source middleware
stack~\cite{cloud_edge_roadmap}. Supply chain verification is largely static: CRA-mandated
Software Bills of Materials (SBOMs)~\cite{cra} are point-in-time declarations without runtime
verification; Digital Product Passport (DPP) frameworks remain in pilot stage under the EU
Ecodesign Regulation. NaaS federation across EU telecom operators does not exist at
operational scale --- cross-border edge compute sharing requires bespoke commercial agreements.
CAMARA Network APIs are being standardised by GSMA but remain fragmentedly implemented
across European operators.

Current cluster ODP: Fragmented Piloting (P4, P7); early supply chain disclosure phase
(P5). Cluster SSL: Emerging EU Ecosystem (P4); Hyperscaler Platform Dominance (P7);
High reliance on fragmented non-EU supply chain auditing (P5).

\textit{Target (2040):}
A unified European NaaS platform enables seamless cross-border compute roaming under a
``Schengen for Cloud'' legal framework, legally operable as a single jurisdiction for
data workloads. All Common European Data Spaces operate on Simpl-based security
interoperability with automated MIM compliance --- provider-agnostic, cryptographically
verifiable, and free of vendor lock-in. Runtime SBOMs with Zero-Knowledge Proof (ZKP)
verification eliminate unverified components from EU critical infrastructure supply chains.
EU achieves full sovereign control over supply chain transparency with pan-European Cyber
Threat Intelligence sharing for supply chain compromise indicators.

% ---- DIMENSION 2 ----
\item \textbf{Primary Drivers}

\textit{Threat Drivers:}
\begin{itemize}
  \item Supply chain attacks exploiting the disaggregation complexity of Open RAN
        deployments, where a single compromised firmware component sourced from one
        of 50-plus hardware and software vendors per macro cell deployment can affect
        thousands of base stations across multiple operators before detection.
  \item ``Weakest link'' cascading failures in federated NaaS environments: a
        single provider with a compromised Minimal Interoperability Mechanism (MIM)
        implementation can propagate malicious policies across the entire federated
        fabric before semantic validation detects the enforcement inconsistency.
  \item US CLOUD Act extra-territorial reach over non-EU cloud providers processing
        EU data creates a structural sovereignty risk that persists regardless of
        technical security measures if the legal framework enabling EU-sovereign
        federated compute is not established.
  \item Silent policy translation failures: automated security policy translation
        between heterogeneous provider MIM implementations may produce semantically
        equivalent but enforcement-different policies, creating exploitable gaps
        invisible to standard compliance auditing procedures.
\end{itemize}

\textit{Regulatory Drivers:}
\begin{itemize}
  \item Data Act interoperability obligations (Articles~23--31): mandatory
        requirements for cloud service providers covering security interface
        standardisation and data portability~\cite{nis2}.
        [CITATION REQUIRES VERIFICATION --- confirm final Data Act article numbering
        for interoperability provisions.]
  \item NIS2 Article~21(2)(h): supply chain security measures mandatory for essential
        entities, requiring documented risk assessment of all ICT suppliers~\cite{nis2}.
  \item CRA Annex~I Part~II: supply chain vulnerability handling requirements for
        products with digital elements, including SBOM disclosure
        obligations~\cite{cra}.
  \item DORA Article~28: ICT third-party risk management requiring continuous
        monitoring and documented exit strategies for critical ICT
        suppliers~\cite{dora}.
\end{itemize}

% ---- DIMENSION 3 ----
\item \textbf{High Impact Activities}

\begin{enumerate}[label=(\alph*)]
  \item \textbf{[Phase~1]} Define and standardise security-specific Minimal
        Interoperability Mechanisms (MIMs) within the Simpl open-source middleware
        stack~\cite{cloud_edge_roadmap}; deploy pilots in three Common European Data
        Spaces (Manufacturing, Health, Energy) with automated cross-provider policy
        enforcement, federated identity attestation, and cryptographically verifiable
        compliance reporting.

  \item \textbf{[Phase~1--2]} Implement Runtime SBOMs with cryptographic proof of
        execution for all containerised microservices and AI components migrating across
        the multi-provider continuum; integrate Digital Product Passports (DPP) for
        hardware component traceability from silicon fabrication to operational
        deployment, aligned with CRA SBOM obligations~\cite{cra}.

  \item \textbf{[Phase~1--2]} Deploy harmonised CAMARA Network APIs across a minimum
        of five EU telecom operators establishing standardised interfaces for cross-border
        edge compute roaming; test Automated Compliance Negotiators in designated EU
        5G cross-border corridor deployments under IPCEI-CIS governance.

  \item \textbf{[Phase~2]} Establish ``Data Visa'' mechanisms enabling dynamic
        cross-border AI model and data migration without manual legal friction, governed
        by a common semantic security ontology (CL3.4) aligned with STIX~2.1/TAXII
        threat intelligence standards and automated EUCS compliance verification~\cite{eucs}.

  \item \textbf{[Phase~2--3]} Deploy Zero-Knowledge Proofs (ZKPs) for IP-protected
        supply chain component verification: hardware vendors prove FIPS~203/204/205
        algorithm compliance and absence of known vulnerabilities without revealing
        proprietary design details~\cite{fips203,fips204,fips205}; operationalise
        pan-European supply chain compromise indicator sharing via standardised
        CTI lifecycle platforms (CL2.5).
\end{enumerate}

% ---- DIMENSION 4 ----
\item \textbf{Catalysts \& Enablers}

\textit{Funding Alignment:}
IPCEI-CIS (Important Project of Common European Interest on Next Generation Cloud
Infrastructure and Services) as the primary funding vehicle for cross-border NaaS federation
pilots; Digital Europe Programme for Simpl development, Common Data Space pilots, and
MIM standardisation infrastructure; ECCC Strategic Agenda priority areas ``Supply Chain
Security'' and ``Secure Interoperability of Systems and Networks''~\cite{eccc}; European
Alliance for Industrial Data, Edge and Cloud for governance and pilot ecosystem
coordination~\cite{cloud_edge_roadmap}.

\textit{Standardisation:}
CEN/CENELEC JTC~13 for security MIM standardisation within the Simpl ecosystem;
GSMA CAMARA API Alliance for harmonised Network API standardisation across operators;
ETSI TC CYBER for supply chain security standards; OASIS STIX~2.1/TAXII for threat
intelligence interoperability; W3C Verifiable Credentials Data Model for SSI-based
supply chain attestation chains.

\textit{Ecosystem Enablers:}
Gaia-X Federation Services trust framework alignment with Simpl MIM security policies;
Digital Product Passport pilot programmes under EU Ecodesign Regulation~2024;
ENISA supply chain risk management guidelines and threat landscape reports for ICT
sectors; dedicated SME compliance support programmes for CRA-mandated SBOM tooling,
recognising that sub-50 person software vendors represent the most vulnerable link in
the supply chain verification architecture~\cite{cra}.

% ---- DIMENSION 5 ----
\item \textbf{Disruptors \& Systemic Risks}

\begin{itemize}
  \item \textbf{Hyperscaler resistance to Simpl MIM adoption:} AWS, Azure, and
        GCP collectively control approximately 65\% of the EU cloud market as of
        2024 and have strong commercial incentives to maintain proprietary APIs;
        without mandatory EU Cloud Rulebook provisions requiring MIM compliance,
        voluntary adoption rates are unlikely to achieve genuine interoperability
        within Phase~1 timelines.
  \item \textbf{IPCEI-CIS implementation delays:} complex EU state aid notification
        procedures have historically delayed IPCEI projects by 18--36 months
        (IPCEI-H2 delayed by state aid clearance processes); P7 NaaS federation
        Phase~1 milestones carry material schedule risk if IPCEI-CIS governance
        does not streamline implementation mechanisms.
  \item \textbf{Semantic interoperability gap:} automated security policy translation
        between heterogeneous MIM implementations is architecturally harder than
        syntactic interoperability --- subtle semantic differences in ``deny'' versus
        ``drop'' versus ``quarantine'' semantics across provider implementations
        can create silent enforcement gaps undetectable through standard compliance
        auditing.
  \item \textbf{ZKP computational overhead:} current ZKP schemes for supply chain
        verification carry 10--100$\times$ computational overhead relative to
        standard cryptographic verification, limiting practical deployment to
        batch verification rather than real-time per-transaction attestation
        in Phases~1--2.
  \item \textbf{EU semiconductor supply chain concentration:} advanced node
        fabrication (sub-7\,nm) required for cutting-edge supply chain attestation
        hardware remains concentrated in Taiwan (TSMC) and South Korea (Samsung);
        a Taiwan Strait contingency would immediately constrain EU supply chain
        verification capability at the silicon level, exposing a gap the EU Chips Act
        does not address within the 2030 timeframe.
\end{itemize}

% ---- DIMENSION 6 ----
\item \textbf{Failsafe Mechanisms \& Resilience}

\begin{itemize}
  \item \textbf{Autonomous unfederation:} if a provider fails Runtime SBOM
        verification or MIM compliance checks, eBPF-based zero-trust policy
        enforcement severs that provider from the federated NaaS fabric and redirects
        workloads to verified compliant nodes within a sub-60-second SLA; provider
        re-admission requires fresh ZKP-verified attestation, not manual
        recertification.
  \item \textbf{Graceful national fallback:} cross-border SLA failure or ``Data
        Visa'' mechanism collapse triggers workload redirect to national infrastructure
        maintained at sufficient sovereign capacity for essential services, fulfilling
        NIS2 Article~21(2)(c) business continuity obligations and DORA Article~9 ICT
        resilience requirements~\cite{nis2,dora}.
  \item \textbf{Cryptographic attestation of MIM enforcement:} all Simpl MIM security
        controls include verifiable proofs that policies are enforced as declared --- not
        merely configured --- using TEE-backed attestation integrated with Cluster~1
        secure element infrastructure~\cite{cloud_edge_roadmap}, eliminating
        compliance-theatre scenarios where declared and actual enforcement diverge.
  \item \textbf{Supply chain quarantine:} unverified or high-risk components are
        isolated into quarantined network slices (VLAN/VxLAN-isolated segments with
        eBPF-enforced east-west traffic blocking) pending ZKP verification;
        quarantined components cannot interact with production workloads or access
        sensitive data paths.
  \item \textbf{Multi-jurisdictional residency fallback:} if cross-border ``Data
        Visa'' mechanisms fail legally or technically, automated GDPR Article~44
        compliance enforcement pins workloads to the originating jurisdiction,
        preventing unauthorised cross-border data transfers under federation failure
        conditions~\cite{gdpr}.
\end{itemize}

% ---- DIMENSION 7 ----
\item \textbf{Transversal Meta-Layer}

\textit{Sovereignty \& Autonomy:}
This cluster breaks non-EU hyperscaler API lock-in (AWS, Azure, GCP controlling
approximately 65\% of EU cloud market as of 2024) by enabling European providers to
compete on standardised Simpl-based infrastructure. A ``Schengen for Cloud'' legal
framework eliminates US CLOUD Act extra-territorial jurisdiction exposure for EU data
processed on federated EU-sovereign NaaS infrastructure. Runtime SBOM verification and
ZKP-based supply chain attestation reduce EU dependency on non-EU vendor self-certification
declarations --- the dominant model under current CRA-baseline SBOM obligations.

\textit{Trust \& Ethics:}
GDPR Article~25 (privacy by design and by default) is operationalised through
MIM-enforced data minimisation policies that are cryptographically verifiable, not
declarative~\cite{gdpr}. Data Act interoperability obligations are fulfilled through
Simpl MIM compliance, providing auditable trust chains for GDPR Article~5(2) data
processing accountability. Transparent Runtime SBOM verification provides the
supply chain integrity that GDPR Article~25 implicitly requires for processors
handling EU personal data at scale~\cite{gdpr}.

\textit{Sustainable Security (Green IT):}
Federated NaaS resource pooling across EU operators eliminates redundant edge compute
over-provisioning by individual operators, reducing total EU edge infrastructure carbon
footprint through shared utilisation efficiency. DPP-based supply chain traceability
aligns with EU Ecodesign Regulation circular economy requirements for ICT hardware,
enabling component reuse lifecycle tracking that reduces electronic waste from
premature infrastructure replacement driven by opacity in component provenance.

\end{enumerate}


% =========================================================================
% CLUSTER 3: FEDERATED OPERATIONAL SOVEREIGNTY
% =========================================================================
\newpage
\subsection{Progression Matrix: Priority Cluster 3 --- Federated Operational Sovereignty}

\begin{landscape}
\setlength{\textwidth}{26.5cm}\setlength{\linewidth}{26.5cm}%
\setlength{\hsize}{26.5cm}\setlength{\columnwidth}{26.5cm}
\setlength{\tabcolsep}{5pt}
\renewcommand{\arraystretch}{1.35}

\begin{longtable}{%
|>{\footnotesize\raggedright\arraybackslash}p{2.8cm}%
|>{\footnotesize\raggedright\arraybackslash}p{5.4cm}%
|>{\footnotesize\raggedright\arraybackslash}p{5.4cm}%
|>{\footnotesize\raggedright\arraybackslash}p{5.4cm}%
|>{\footnotesize\raggedright\arraybackslash}p{5.4cm}|}

\caption{\textbf{Cluster~3: Federated Operational Sovereignty} (P4, P5, P7).
Transition matrix across the four methodology dimensions. P4 (Simpl
Interoperability), P5 (Supply Chain Verification), and P7 (NaaS Federation)
are labelled where their maturity trajectories diverge within a cell.}
\label{tab:cluster3}\\

\hline
\rowcolor{colheader}
\textbf{Phase / SSL} &
\textbf{R\&I Maturity}{\scriptsize ODP indicator per sub-priority} &
\textbf{Policy, Regulatory \& Standards}{\scriptsize enforcement $\cdot$
  certification $\cdot$ standardisation} &
\textbf{Operational Deployment \& Adoption}{\scriptsize actor types $\cdot$
  scale $\cdot$ sector coverage} &
\textbf{Market Uptake \& Strategic Sovereignty}{\scriptsize SSL indicator
  $\cdot$ dependency reduction $\cdot$ industrial capacity} \\
\hline
\endfirsthead

\hline
\rowcolor{colheader}
\textbf{Phase / SSL} &
\textbf{R\&I Maturity} &
\textbf{Policy, Regulatory \& Standards} &
\textbf{Operational Deployment \& Adoption} &
\textbf{Market Uptake \& Strategic Sovereignty} \\
\hline
\endhead
\hline\multicolumn{5}{r}{\footnotesize\textit{Continued\ldots}}\\ \endfoot
\hline \endlastfoot

% ---- PHASE 1 ----
\rowcolor{phase1row}
\textbf{Phase~1}
\newline\textit{Foundation \&}
\newline\textit{Compliance}
\newline\textit{2025--2030}
\newline\sslbadge{sslredbg}{sslredtx}{SSL: Fragmented; Hyperscaler}
\newline\sslbadge{sslredbg}{sslredtx}{Dominance (P7); Emerging}
\newline\sslbadge{sslredbg}{sslredtx}{EU Ecosystem (P4)}
&
\cellhead{Cluster ODP: Fragmented Piloting (P4, P7) / Disclosure Phase (P5)}
\subp{P4} is at Fragmented Piloting: Simpl security MIM pilots are
launched in three Common European Data Spaces (Manufacturing, Health,
Energy) with manual policy coordination and proof-of-concept
implementations. \subp{P7} is at Bilateral Interconnection: harmonised
CAMARA Network APIs are deployed in 2--3 bilateral 5G corridor agreements
between individual EU telecom operators; cross-border edge compute
roaming does not exist at operational scale. \subp{P5} is at Supply
Chain Disclosure: CRA-mandated static SBOM declarations are being
submitted as point-in-time compliance artefacts without runtime
verification; Digital Product Passport pilot programmes are initiated
under the EU Ecodesign Regulation; Zero-Knowledge Proof verification
for supply chain attestation is at Applied R\&D stage.
\kpistrip{The three sub-priorities have meaningfully different Phase~1
maturity levels. P5 is the least mature and the most consequential
shared dependency: Runtime SBOM verification is a gate condition for
both P4 (MIM compliance trust) and Cluster~1 (TEE integrity claims).}
&
\cellhead{Multiple legislative obligations activating; standards nascent}
CRA~SBOM obligations (Annex~I Part~II) enter enforcement in 2027,
mandating static SBOM disclosure for all products with digital elements
entering the EU market. NIS2~Article~21(2)(h) makes supply chain risk
assessment mandatory for essential entities. DORA~Article~28 requires
continuous ICT third-party risk monitoring for financial sector entities.
Data Act interoperability obligations activate, creating procurement-level
pressure for cross-provider security interface standardisation. The
Simpl-Open official launch with initial MIM definitions constitutes the
Phase~1 technical policy trigger for P4. IPCEI-CIS enters its
deployment phase. GSMA CAMARA API Alliance standardisation work is
active but not yet mandated in EU regulatory instruments.
\kpistrip{Phase~1 policy activity is high-volume but fragmented:
multiple legislative instruments activate simultaneously without
a unified interoperability mandate that would create coherent market
pull for Simpl-based standardisation.}
&
\cellhead{Deployment: bilateral pilots and Data Space proof-of-concepts}
Three Common European Data Space security pilots operate Simpl-based
MIM proof-of-concepts with manual policy coordination; cross-provider
security policy enforcement is not automated. CAMARA API deployments
are confined to 2--3 bilateral 5G corridor agreements between single
operator pairs. Static SBOM declarations are submitted as CRA
compliance artefacts but are not verified at runtime; supply chain
risk assessment is vendor-declared rather than independently attested.
Hyperscalers (AWS, Azure, GCP) retain $\approx$65\% of EU cloud market
through proprietary API ecosystems. Cross-border AI model migration
requires bespoke bilateral legal review per transfer event. SME-accessible
SBOM tooling and MIM compliance support are not yet standardised or
commercially affordable.
\kpistrip{Phase~1 deployment is bilateral and pilot-scale. The gap
between proof-of-concept MIM implementations and automated enforcement
across the majority of Data Spaces defines the core operational
challenge for Phase~2.}
&
\cellhead{SSL: Mixed --- Emerging EU Ecosystem (P4) / Hyperscaler Dominance (P7)}
\subp{P4} has reached Emerging EU Ecosystem: the Simpl open-source
stack exists and is being piloted with EU institutional backing.
\subp{P7} remains at Hyperscaler Platform Dominance: AWS, Azure, and GCP
collectively control $\approx$65\% of EU cloud market; no EU NaaS
federation operates at commercially relevant scale; cross-border edge
compute is fragmented into national silos. \subp{P5} is at High
Reliance on Fragmented Non-EU Auditing: supply chain verification
is vendor-declared under US-dominated certification ecosystems (Common
Criteria national schemes, SOC2). The cluster's dominant SSL profile
is foreign dependency, driven by the weight of P7 market structure.
\kpistrip{P4 and P7 require separate SSL tracking in Phase~1: conflating
them into a single cluster SSL figure obscures that Simpl has achieved
Emerging Ecosystem while NaaS federation has not.}
\\
\hline

% ---- GATE 1->2 ----
\rowcolor{gaterow}
\gatelbl{1 $\rightarrow$ 2} &
\multicolumn{4}{>{\footnotesize\raggedright\arraybackslash}p{22.0cm}|}{%
\cellcolor{gaterow}%
\textbf{R\&I trigger:} At least one Common European Data Space has deployed
automated (not manually coordinated) Simpl MIM-based cross-provider security
policy enforcement, with documented evidence that enforcement is cryptographically
verifiable rather than declaratively asserted.
\textbf{Policy trigger:} Simpl security MIM definitions are published in
candidate CEN/CENELEC standard form and referenced as a compliance pathway in
at least one EU sector-specific NIS2 or Data Act implementing act.
\textbf{Operational trigger:} CAMARA-based cross-border edge compute roaming is
operational between at least three distinct EU telecom operator pairs (not a single
bilateral agreement), with documented SLA performance and compliance evidence.
\textbf{Market trigger:} Runtime SBOM tooling with cryptographic execution proofs
is available as a commercially supported product or service from at least two
EU-domiciled vendors, establishing the market infrastructure for Phase~2
mandatory adoption.}
\\
\hline

% ---- PHASE 2 ----
\rowcolor{phase2row}
\textbf{Phase~2}
\newline\textit{Applied \& Industrial}
\newline\textit{R\&D alongside}
\newline\textit{Federation \&}
\newline\textit{Standardisation}
\newline\textit{2030--2035}
\newline\sslbadge{ssltransbg}{ssltrantx}{SSL: Growing Pan-EU}
\newline\sslbadge{ssltransbg}{ssltrantx}{Sovereignty}
&
\cellhead{Cluster ODP: Standardised Integration (P4) / Federated Edge Roaming (P7) / Dynamic Runtime Verification (P5)}
\subp{P4} reaches Standardised Integration: automated policy translation
via standardised Simpl MIMs is deployed across the majority of active
Common European Data Spaces; federated identity attestation and
cross-provider compliance reporting are automated. \subp{P7} reaches
Federated Edge Roaming: CAMARA-based cross-border roaming operates
across $\geq$5 EU operator pairs; ``Data Visa'' mechanisms enable
dynamic AI model migration without manual legal review in certified
corridors. \subp{P5} transitions from disclosure to dynamic runtime
verification: Runtime SBOMs with cryptographic proof-of-execution
verify containerised components as they migrate across multi-provider
nodes; ZKP batch verification of supply chain attestation is operationally
validated; Digital Product Passports are mandatory under the EU
Ecodesign Regulation.
\kpistrip{Phase~2 R\&I activity includes industrial R\&D for ZKP
real-time verification (currently batch-only due to 10--100$\times$
overhead), semantic security ontology standardisation, and legal
engineering of the ``Schengen for Cloud'' framework --- not solely
standardisation of existing prototypes.}
&
\cellhead{Security MIMs as CEN/CENELEC standards: the Phase~2 gate}
The single Phase~2 policy gate is the adoption of Simpl security MIMs
as ratified CEN/CENELEC standards, referenced as mandatory compliance
requirements in the EU Cloud Rulebook. This transforms MIM adoption
from voluntary best-practice to market-access obligation for providers
serving EU-regulated sectors. DPP becomes mandatory under the EU
Ecodesign Regulation~2024 for ICT hardware categories. Data Act full
enforcement creates binding procurement-level interoperability
obligations. DORA continuous ICT third-party monitoring is fully
operationalised in the financial sector. IPCEI-CIS cross-border NaaS
governance framework is operational.
\kpistrip{The Phase~2 policy gate is the CEN/CENELEC ratification of
security MIMs. Voluntary Simpl adoption without this standards event
does not constitute Phase~2 policy maturity.}
&
\cellhead{Cross-sector rollout; SME participation; Runtime SBOM in production}
Simpl MIMs are deployed across the majority of active Common European
Data Spaces with automated cross-provider enforcement. Runtime SBOMs
with cryptographic execution proofs verify supply chain components in
critical infrastructure deployments in at least five Member States.
CAMARA cross-border roaming operates across $\geq$5 EU operator pairs
with documented SLA performance. ``Data Visa'' mechanisms are
operational in certified cross-border AI model migration corridors.
SME-accessible Simpl MIM compliance tooling and Runtime SBOM services
are commercially available from EU-domiciled providers. Unverified
supply chain components are automatically quarantined in isolated
network slices.
\kpistrip{Phase~2 deployment must show geographic breadth ($\geq$5 Member
States), cross-sector presence ($\geq$3 Data Spaces under automated MIM
enforcement), and SME participation in at least one deployment context.
Single-sector rollout does not constitute Phase~2 operational maturity.}
&
\cellhead{SSL: Growing Pan-EU Sovereignty}
EU Data Space infrastructure operating under Simpl achieves $\geq$30\%
non-hyperscaler share in new regulated-sector cloud procurement. Runtime
SBOM ecosystem is commercially mature with multiple competing EU-domiciled
vendors. EU NaaS consortium (under IPCEI-CIS) operates cross-border
compute at pilot commercial scale, demonstrating that EU providers
can federate to compete with hyperscaler geographic coverage. DPP
traceability reduces procurement risk for EU critical infrastructure
buyers by making supply chain provenance verifiable. Foreign hyperscaler
market share in EU regulated sectors declines for the first time from
its $\approx$65\% Phase~1 baseline.
\kpistrip{Phase~2 SSL evidence: $\geq$30\% non-hyperscaler share in new
regulated procurement; first documented market share decline for
non-EU hyperscalers in EU regulated cloud procurement.}
\\
\hline

% ---- GATE 2->3 ----
\rowcolor{gaterow}
\gatelbl{2 $\rightarrow$ 3} &
\multicolumn{4}{>{\footnotesize\raggedright\arraybackslash}p{22.0cm}|}{%
\cellcolor{gaterow}%
\textbf{R\&I trigger:} ZKP verification for supply chain attestation has
achieved real-time per-transaction performance (eliminating the current
batch-only constraint) in at least one production critical infrastructure
deployment, confirmed by independent performance benchmarking.
\textbf{Policy trigger:} A ``Schengen for Cloud'' legal framework --- or
equivalent EU legal instrument resolving US CLOUD Act extra-territorial
jurisdiction conflict for federated EU-sovereign NaaS workloads --- has
been enacted, providing the jurisdictional basis for single-market operation
of federated compute.
\textbf{Operational trigger:} The EU NaaS federation (under IPCEI-CIS or its
successor) operates at commercial scale across $\geq$10 EU telecom operator
members, offering cross-border edge compute services as a standard commercial
product with documented customer deployments outside its founding consortium.
\textbf{Market trigger:} EU supply chain governance frameworks (Runtime SBOM
plus ZKP attestation) are adopted or formally referenced by at least two
non-EU partner jurisdictions as part of digital trade or partnership agreements,
confirming that EU governance leadership is internationally exportable.}
\\
\hline

% ---- PHASE 3 ----
\rowcolor{phase3row}
\textbf{Phase~3}
\newline\textit{Autonomy \&}
\newline\textit{Leadership}
\newline\textit{2035--2040}
\newline\sslbadge{ssltealbg}{ssltealtx}{SSL: Full Federated}
\newline\sslbadge{ssltealbg}{ssltealtx}{Sovereign Autonomy}
&
\cellhead{Cluster ODP: Ubiquitous Operation across all sub-priorities}
All three sub-priorities reach Ubiquitous Operation. Simpl-based
security is the universal standard for cross-provider Data Space
operation in the EU; no alternative interoperability model holds
market relevance in regulated sectors. ZKP-verified supply chain
attestation operates at real-time transaction scale for all critical
infrastructure component verification. The EU NaaS federation operates
as a single commercial entity legally equivalent to a unified
jurisdiction for data workloads. Phase~3 R\&I pursues next-generation
frontiers: post-ZKP cryptographic verification, semantic AI-to-AI
policy negotiation, and legal engineering of the next-generation
``Digital Single Market'' framework beyond cloud-edge compute.
\kpistrip{Phase~3 ODP evidence: Simpl MIM compliance rate in new regulated
EU cloud procurement; ZKP verification throughput in production critical
infrastructure deployments; EU NaaS federation revenue relative to
hyperscaler EU revenue in regulated sectors.}
&
\cellhead{EU federated governance as international model}
The ``Schengen for Cloud'' legal framework is enacted and operational,
resolving US CLOUD Act jurisdiction conflicts for federated EU-sovereign
workloads and enabling single-market legal operation of cross-border
compute. EU Cloud Rulebook is mandatory for all new cloud deployments
in EU regulated sectors without exception. EU supply chain governance
standards (Runtime SBOM plus ZKP) are referenced in international
supply chain security standards (ISO/IEC) and in digital trade
agreements with partner jurisdictions. The EU has transitioned from
rule-taker to rule-setter in international cloud and supply chain
governance.
\kpistrip{Phase~3 governance leadership is verified by two events:
enactment of the ``Schengen for Cloud'' framework (domestic) and
formal adoption of EU supply chain governance standards by $\geq$2
non-EU partner jurisdictions (international).}
&
\cellhead{Provider-agnostic EU trust fabric as continental operational default}
A single provider-agnostic EU trust fabric operates continuously across
all active Common European Data Spaces with real-time automated security
orchestration and zero manual policy coordination. ZKP-verified supply
chain attestation eliminates all unverified components from EU critical
infrastructure deployments in regulated sectors. The EU NaaS federation
is a standard commercial service operating at hyperscaler-comparable
scale, with EU operators collectively competitive on cost and coverage.
Security-by-Design is not an add-on: every new cloud-edge deployment
in the EU incorporates Simpl MIM compliance, Runtime SBOM verification,
and cross-border sovereignty assurance as architectural defaults.
\kpistrip{Phase~3 deployment normalisation: Simpl MIM compliance in
$\geq$90\% of new regulated EU cloud deployments; ZKP supply chain
verification in 100\% of new critical infrastructure component
procurement; EU NaaS federation operating as a single commercial entity
with documented non-founding-member customers.}
&
\cellhead{SSL: Full Federated Sovereign Autonomy}
Non-EU hyperscaler market share in EU regulated cloud procurement
falls below 40\% for new deployments (from $\approx$65\% at baseline),
with a binding trajectory toward continued reduction. EU NaaS
federation achieves hyperscaler-comparable geographic coverage and
commercial viability across the EU. Supply chain governance
internationally adopted: at least two partner jurisdictions have
formally adopted EU Runtime SBOM plus ZKP attestation standards as
domestic procurement requirements. The EU has removed the structural
dependency that allowed US CLOUD Act jurisdiction to apply to
EU-regulated data: federated sovereignty is legally guaranteed, not
only technically implemented.
\kpistrip{Full Federated Sovereign Autonomy requires three
concurrent conditions: $<$40\% hyperscaler new-procurement share;
``Schengen for Cloud'' legally enacted; $\geq$2 partner jurisdictions
formally adopting EU supply chain governance standards. Achieving
any two without the third does not constitute this SSL level.}
\\
\hline

\end{longtable}
\end{landscape}

% ---- PROGRESSION MATRIX: CLUSTER 3 ----


% \noindent
% \resizebox{\textwidth}{!}{%
% \begin{tabular}{|p{4.2cm}|p{5.6cm}|p{5.6cm}|p{5.6cm}|}
% \hline
% \rowcolor{gray!30}
% \multicolumn{4}{|c|}{\textbf{Progression Matrix: Priority Cluster 3 ---
%   Federated Operational Sovereignty (P4, P5, P7)}} \\
% \hline
% \rowcolor{gray!15}
% \textbf{Metric} &
% \textbf{Phase 1: Foundation \& Compliance (2025--2030)} &
% \textbf{Phase 2: Federation \& Standardisation (2030--2035)} &
% \textbf{Phase 3: Autonomy \& Leadership (2035--2040)} \\
% \hline
% \rowcolor{gray!10}
% \multicolumn{4}{|l|}{\textit{\textbf{Cluster-Level Progression Indicators}}} \\
% \hline
% \textbf{Cluster ODP} &
% Fragmented Piloting &
% Standardised Integration &
% Ubiquitous Operation \\
% \hline
% \textbf{Cluster SSL} &
% Emerging EU Ecosystem &
% Growing Pan-EU Sovereignty &
% Full Federated Sovereign Autonomy \\
% \hline
% \rowcolor{gray!10}
% \multicolumn{4}{|l|}{\textit{\textbf{Sub-Priority Decomposition (Individual ODP / SSL per Phase)}}} \\
% \hline
% \textit{P4: Security Interoperability \newline via Simpl Middleware} &
% ODP: Fragmented Piloting \newline
% SSL: Emerging EU Ecosystem &
% ODP: Standardised Integration \newline
% SSL: Full Sovereign Autonomy &
% ODP: Ubiquitous Operation \newline
% SSL: Full Sovereign Autonomy \\
% \hline
% \textit{P5: Supply Chain \newline Verification \& Infra.} &
% ODP: Full Supply Chain Disclosure \newline
% SSL: High Reliance on \newline Fragmented Non-EU Auditing &
% ODP: Partial Restriction \& \newline Dynamic Runtime Verification \newline
% SSL: Emerging EU Supply Chain \newline Tracking Frameworks &
% ODP: Full Restriction Process \& \newline ZKP-Verified Attestation \newline
% SSL: Full Sovereign Supply \newline Chain Control \& Transparency \\
% \hline
% \textit{P7: Pan-European NaaS \newline Federation} &
% ODP: Bilateral Interconnection \newline
% SSL: Hyperscaler Platform Dominance &
% ODP: Federated Edge Roaming \newline
% SSL: Pan-EU Consortium Competes &
% ODP: True Digital Single Market \newline
% SSL: European Scale Parity Achieved \\
% \hline
% \rowcolor{blue!7}
% \textbf{Operational State} &
% Initial Simpl security pilots are deployed in three Common European
% Data Spaces with manual policy coordination and proof-of-concept
% MIM implementations; CRA-baseline SBOM declarations are submitted
% by hardware and software vendors as static point-in-time documents
% without runtime verification; harmonised CAMARA APIs are tested in
% bilateral 5G corridor agreements between individual EU telecom
% operators; hyperscalers retain dominant market share through
% proprietary API ecosystems. &
% Automated policy translation via standardised Simpl MIMs is
% deployed across the majority of active Common European Data Spaces;
% Runtime SBOMs with cryptographic execution proofs verify
% containerised components dynamically as they migrate across
% multi-provider continuum nodes; harmonised CAMARA-based cross-border
% edge roaming operates across five or more EU operator pairs;
% ``Data Visa'' mechanisms enable AI model migration without
% manual legal review in certified corridor deployments. &
% A single provider-agnostic European trust fabric operates
% continuously across all EU Common Data Spaces with real-time
% automated security orchestration; ZKP-verified supply chain
% attestation eliminates unverified components from all
% EU critical infrastructure deployments; a unified NaaS
% platform legally operable as a single data jurisdiction
% enables EU operators to achieve hyperscaler-comparable
% scale; EU supply chain governance frameworks are adopted
% as reference implementations internationally. \\
% \hline
% \rowcolor{orange!12}
% \textbf{Critical Enabler} &
% Simpl-Open official launch with initial security MIM definitions
% published as open standards, establishing the technical
% baseline for cross-provider security policy automation &
% Security MIMs adopted as CEN/CENELEC standards with
% mandatory reference in EU Cloud Rulebook, creating the
% regulatory basis for non-voluntary hyperscaler compliance &
% ``Schengen for Cloud'' legal framework enacted, resolving
% US CLOUD Act extra-territorial jurisdiction conflict and
% enabling single-jurisdiction data workload operation at
% continental scale \\
% \hline
% \end{tabular}%
% }

% =========================================================================
% CROSS-PRIORITY BOTTLENECK ASSESSMENT
% =========================================================================

\newpage
\section{Cross-Priority Bottleneck Assessment}
\label{sec:bottlenecks}

Three shared structural dependencies determine whether individual sub-priority
progressions can advance on their stated timelines, regardless of cluster-level investment.
These bottlenecks are identified here because they constitute cross-cutting gate conditions
that cannot be resolved by any single cluster in isolation.

\textbf{Bottleneck 1 --- EU Chips Act Sovereign Fabrication Capacity.}
The Phase~2 gate conditions for P1 (EUCC `High' certified RISC-V secure elements),
P2 (PQC-native RISC-V silicon), P3 (zero-energy secure hardware), and P8 (RISC-V
Keystone TEE scaling) all converge on the same physical dependency: sovereign EU
fabrication capacity for security-critical silicon at commercially viable yields.
The Phase~2 critical enabler for P1 (``EUCC `High' certification of first
EU-manufactured RISC-V secure element'') and the Phase~3 enabler (``EU Chips Act
sovereign EUV-capable fabrication achieving $\geq$60\% self-sufficiency'') constitute
a single two-stage gate that controls the SSL trajectory of four of the five Cluster~1
sub-priorities simultaneously.

\textbf{Bottleneck 2 --- AI Act Conformity Assessment Infrastructure.}
The Phase~2 gate condition for Cluster~2 (``standardised AI Act conformity assessment
framework for autonomous security agents'') is a prerequisite not only for P6 and P10
but also for any AI-driven capability within Cluster~1 (autonomous Island Mode swarm
protocols) and Cluster~3 (Automated Compliance Negotiators in NaaS federation) that
incorporates autonomous decision-making. The absence of EU AI Act conformity assessment
bodies with sufficient capacity to certify high-risk security AI at scale represents a
systemic bottleneck that, if unresolved by 2028--2030, will delay the Phase~2 SSL
transition across all three clusters.

\textbf{Bottleneck 3 --- Runtime SBOM and Supply Chain Attestation Infrastructure.}
Runtime SBOM verification (P5) is a prerequisite for both the Cluster~1 TEE integrity
claim (P8 confidential computing pipelines cannot be trusted without verified component
provenance) and the Cluster~3 MIM compliance model (Simpl-based interoperability
cannot provide cryptographic enforcement guarantees if component integrity is
unverifiable). The Phase~1 critical enabler for Cluster~3 (Simpl-Open MIM definitions)
and the Phase~1 ODP for P5 (full supply chain disclosure) must therefore proceed in
coordination, not sequentially, with the Phase~1 applied research activities in
Cluster~1. This interdependency must be explicitly reflected in Horizon Europe project
call design to ensure that SBOM infrastructure investment is not siloed within supply
chain security calls but is structurally integrated into hardware security and
interoperability funding instruments.

% =========================================================================
% NEW BIBLIOGRAPHY ENTRIES
% =========================================================================

% Add the following entries to the master \begin{thebibliography}{99} block:
 in master document
%
% Required packages (add to master preamble if not present):
%   \usepackage[table]{xcolor}
%   \usepackage{multirow}
%   \usepackage{array}
%
% New bibliography keys required (entries provided at end of file):
%   fips203, fips204, fips205, nistir8413, nistsp800207,
%   etsi103744, etsizsmoo9, tgcpm20, isoiec11889,
%   opentitan, keystone, threeGPP23288, threeGPP33501, dora, gdpr
% =========================================================================

% ---- Shared colors/macro helpers when included from a master document ----
\providecolor{gaterow}{RGB}{240,240,238}
\providecolor{phase1row}{RGB}{237,245,253}
\providecolor{phase2row}{RGB}{237,249,244}
\providecolor{phase3row}{RGB}{245,244,254}
\providecolor{colheader}{RGB}{245,245,243}
\providecolor{clusterheader}{RGB}{220,232,248}

\providecolor{sslredbg}{RGB}{252,235,235}
\providecolor{sslredtx}{RGB}{163,45,45}
\providecolor{sslamberbg}{RGB}{250,238,218}
\providecolor{sslambtx}{RGB}{133,79,11}
\providecolor{ssltransbg}{RGB}{230,241,251}
\providecolor{ssltrantx}{RGB}{24,95,165}
\providecolor{ssltealbg}{RGB}{225,245,238}
\providecolor{ssltealtx}{RGB}{15,110,86}

\providecommand{\sslbadge}[3]{%
      \parbox{\linewidth}{\raggedright
            \colorbox{#1}{\scriptsize\parbox{0.92\linewidth}{\raggedright
                  \textbf{\textcolor{#2}{#3}}}}}}

\providecommand{\kpistrip}[1]{%
      \par\vspace{2pt}\noindent\rule{\linewidth}{0.3pt}%
      \par\noindent\textit{\scriptsize #1}}

\providecommand{\cellhead}[1]{\noindent\textbf{#1}\par\smallskip}
\providecommand{\gatelbl}[1]{\scriptsize\textit{Gate~#1:}}
\providecommand{\subp}[1]{\textit{\scriptsize(#1)}}

\section{Strategic Priority Clusters}
\label{ch:strategic_priority_clusters}

% \subsection{Consolidation Rationale and Structural Logic}
% \label{sec:consolidation_rationale}

% The ten executive priorities identified through the 4-Filter Prioritisation Framework represent
% distinct technology domains, but they share deep structural dependencies and converge on three
% strategic imperatives that define Europe's 15-year cybersecurity trajectory for the Cloud-Edge-IoT
% continuum. Maintaining ten parallel profiles in a dual-audience deliverable risks fragmenting
% the narrative for European Commission technical evaluators while obscuring the systemic
% interdependencies that determine whether individual priorities succeed or fail.

% The consolidation presented in this chapter reorganises the ten priorities into three strategic
% clusters, each with a unifying architectural thesis. To preserve the analytical precision required
% by IEEE Access peer review, each cluster retains its constituent sub-priorities as explicitly
% labelled components within a nested progression matrix, allowing individual technology maturation
% trajectories to be tracked independently while the cluster-level framing communicates strategic
% coherence to policy audiences.

% The mapping is as follows:
% \begin{itemize}
%   \item \textbf{Priority Cluster 1 --- Sovereign Trusted Infrastructure:}
%         P1 (Hardware Roots of Trust), P2 (PQC Agility), P3 (Lightweight
%         Cryptography and Physical-Layer Security), P8 (Confidential
%         Computing), P9 (Resilient-by-Design / Island Mode).
%   \item \textbf{Priority Cluster 2 --- Cognitive Security and Autonomous
%         Cyber Defence:} P6 (Cognitive SecOps / LLMSecOps), P10 (Proactive
%         APT Management), and the AI-as-target dimension drawn from
%         Cross-Layers CL1.2, CL1.3, and CL4.7.
%   \item \textbf{Priority Cluster 3 --- Federated Operational Sovereignty:}
%         P4 (Security Interoperability via Simpl), P5 (Supply Chain
%         Verification), P7 (Pan-European NaaS Federation).
% \end{itemize}

% One critical cross-priority dependency must be made explicit before the cluster profiles are
% presented. Hardware supply chain attestation (P5) and Runtime SBOM verification constitute
% a shared gate condition for both Cluster 1 (TEE and secure element integrity) and Cluster 3
% (federated supply chain governance). Progression in P8 (Confidential Computing) is contingent
% on the availability of trustworthy RISC-V secure elements from P1; P3 (Lightweight PhySec)
% and P2 (PQC) both depend on EU Chips Act fabrication capacity as a common Phase 2 enabler.
% These shared bottlenecks are identified explicitly in the Cross-Priority Bottleneck assessment
% in Section~\ref{sec:bottlenecks} and traced to specific progression matrix gate conditions.

% =========================================================================
% PRIORITY CLUSTER 1: SOVEREIGN TRUSTED INFRASTRUCTURE
% =========================================================================

% \newpage
\section{Priority Cluster 1: Sovereign Trusted Infrastructure}
\label{sec:cluster1}

\noindent\textbf{Constituent Sub-Priorities:} P1 (Hardware Roots of Trust \& EUCC `High'),
P2 (PQC Agility for Edge/IoT), P3 (Lightweight Cryptography \& Physical-Layer Security),
P8 (Confidential Computing \& Privacy-Preserving AI), P9 (Resilient-by-Design \& Island Mode).

\medskip
\noindent\textbf{Cluster Thesis:} Before any federated or AI-driven security capability can be
trusted, the substrate on which it executes must provide cryptographically verifiable integrity,
quantum-resistant confidentiality, and autonomous survivability independent of external
infrastructure. This cluster defines the verifiable hardware and cryptographic foundation that
all higher-order capabilities in Clusters 2 and 3 depend upon as an architectural prerequisite.

\subsection*{7-Dimensional Analysis: Priority Cluster 1}

\begin{enumerate}

% ---- DIMENSION 1 ----
\item \textbf{Baseline \& Target Objective}

\textit{Baseline (2025--2026):}
European critical infrastructure and edge deployments operate under high foreign dependency
across all five constituent sub-priorities. Hardware roots of trust are dominated by non-EU
Trusted Platform Module (TPM) manufacturers --- Infineon SLB9670, NXP SE050, and
STMicroelectronics ST33 series collectively account for the majority of the EU market under
ISO/IEC~11889~\cite{isoiec11889} and TCG TPM~2.0 specifications~\cite{tgcpm20}. The Trusted
Execution Environment (TEE) market is dominated by Intel TDX, AMD SEV-SNP, and ARM TrustZone,
none of which are EU-designed or manufactured. Post-Quantum Cryptography (PQC) standardisation
reached a definitive milestone with NIST's publication of FIPS~203 (ML-KEM)~\cite{fips203},
FIPS~204 (ML-DSA)~\cite{fips204}, and FIPS~205 (SLH-DSA)~\cite{fips205} in August~2024;
however, hardware-optimised implementations for 8/16-bit microcontrollers remain at the
Applied~R\&D stage, as confirmed by NIST IR~8413~\cite{nistir8413}. Physical-layer security
for extreme-edge mMTC devices lacks standardised specifications. Edge nodes universally operate
with hard cloud dependencies and no autonomous fallback architectures meeting the operational
resilience mandates of NIS2~\cite{nis2} and DORA~\cite{dora}.

Current cluster ODP: Applied R\&D (driven by PQC, Lightweight PhySec, and Confidential
Computing sovereign alternatives). Current cluster SSL: High Foreign Dependency.

\textit{Target (2040):}
At least 60\% of EU critical infrastructure edge nodes deploy EU-designed RISC-V open-hardware
secure elements certified at EUCC `High' assurance level~\cite{eucc}. 100\% of new EU critical
IoT deployments support hybrid or full FIPS~203/204/205-compliant PQC. Zero-energy nodes
execute quantum-resistant device authentication within a sub-100\,$\mu$J energy budget per
operation. Trusted Execution Environments built on EU open-hardware architectures (Keystone~\cite{keystone},
OpenTitan~\cite{opentitan}) cover at least 50\% of new EU cloud-edge deployments. All NIS2-essential
and DORA-critical edge nodes demonstrate validated 72-hour autonomous Island Mode survivability
independent of WAN connectivity.

% ---- DIMENSION 2 ----
\item \textbf{Primary Drivers}

\textit{Threat Drivers:}
\begin{itemize}
  \item State-sponsored APT exploitation of JTAG debug interfaces and electromagnetic
        side-channel analysis (EMSCA) on field-deployed edge nodes; the Spectre and Meltdown
        class of microarchitectural side-channel vulnerabilities establishes that even
        production-certified hardware can be compromised at scale, and far-edge nodes
        operating without physical security controls present a materially higher exposure
        surface than hyperscaler data centres.
  \item ``Harvest Now, Decrypt Later'' (HNDL) campaigns: state-level adversaries are
        actively archiving encrypted EU government and industrial communications under
        the expectation of Cryptographically Relevant Quantum Computer (CRQC)
        availability. NIST IR~8413 places the Q-Day horizon in the 10--20 year range
        with a highly uncertain lower bound~\cite{nistir8413}; HNDL renders this a
        present-tense threat against data with multi-decade confidentiality
        requirements.
  \item Physical tampering and environmental attacks on unattended far-edge nodes
        (smart grid substations, O-RAN radio units, industrial gateways) where the
        physical security controls available in hyperscaler facilities are absent.
  \item Cascading WAN disconnection failures propagating through centrally-dependent edge
        architectures, exploited by adversaries to degrade critical infrastructure
        services during hybrid conflict scenarios.
\end{itemize}

\textit{Regulatory Drivers:}
\begin{itemize}
  \item CRA Articles~10 and~11: mandatory vulnerability handling and security update
        obligations for products with digital elements; enforcement begins in 2027 for
        most product categories~\cite{cra}.
  \item NIS2 Article~21(2)(e): cryptographic policies and key management for essential
        and important entities; full enforcement began October 2024~\cite{nis2}.
  \item DORA Article~9: ICT resilience testing including hardware-level resilience
        requirements for financial sector critical infrastructure~\cite{dora}.
  \item AI Act Article~9: risk management system requirements for high-risk AI
        systems; critical infrastructure AI components require hardware integrity
        guarantees that only EUCC `High' certified secure elements can
        provide~\cite{aiact}.
\end{itemize}

% ---- DIMENSION 3 ----
\item \textbf{High Impact Activities}

\begin{enumerate}[label=(\alph*)]
  \item \textbf{[Phase~1]} Develop sovereign European Secure Elements based on
        RISC\-V open-hardware (building on OpenTitan~\cite{opentitan} and
        Keystone~\cite{keystone} frameworks), targeting EUCC `High' assurance level
        evaluation; implement automated hardware-level remote attestation protocols
        interoperable with TPM~2.0~\cite{tgcpm20} and ISO/IEC~11889~\cite{isoiec11889}
        for heterogeneous large-scale edge fleets.

  \item \textbf{[Phase~1--2]} Optimise NIST FIPS~203 (ML-KEM)~\cite{fips203} and
        FIPS~204 (ML-DSA)~\cite{fips204} for 8/16-bit microcontrollers under severe
        energy and memory constraints (target: ML-KEM-512 within 32\,KB RAM on
        Cortex-M0+ class devices); develop crypto-agility OTA update middleware enabling
        algorithm rotation without hardware replacement, compatible with ETSI TS~103~744
        hybrid key exchange profiles~\cite{etsi103744}.

  \item \textbf{[Phase~1--2]} Implement ultra-lightweight PQC and keyless
        Physical-Layer Security (RF fingerprinting exploiting Channel State Information
        and Carrier Frequency Offset, CSI/CFO-based device authentication) for mMTC
        devices operating under zero-energy or near-zero-energy constraints; target
        integration with 6G ISAC (Integrated Sensing and Communication) interfaces per
        SNS~JU~SRIA~\cite{snsju}.

  \item \textbf{[Phase~2]} Deploy multi-party confidential computing pipelines
        using TEE abstraction layers (Keystone enclaves on RISC-V~\cite{keystone})
        combined with Fully Homomorphic Encryption (FHE) schemes for collaborative AI
        training across multi-provider nodes without raw data exposure; align with
        EUCS `High+' certification requirements~\cite{eucs}.

  \item \textbf{[Phase~2--3]} Engineer autonomous Island Mode architecture: soft-state
        edge orchestration using cached credentials with 72-hour validity windows,
        LoRaWAN LPWAN fallback micro-slices, eBPF-based kernel-level policy enforcement
        for autonomous node isolation, and bio-inspired self-healing swarm protocols
        enabling edge nodes to sustain critical service delivery without cloud
        connectivity or human intervention.
\end{enumerate}

% ---- DIMENSION 4 ----
\item \textbf{Catalysts \& Enablers}

\textit{Funding Alignment:}
Horizon Europe Cluster~4 (Digital, Industry and Space) for sovereign secure silicon and
PQC hardware implementation; Horizon Europe Cluster~3 (Civil Security for Society) for
edge resilience and Island Mode architectures; EU Chips Act (Regulation (EU) 2023/1781)
pilot fabrication lines for security-critical silicon via the Chips Joint Undertaking
(KDT~JU); ECCC Strategic Agenda priority areas ``Secure Components and Systems'' and
``Post-Quantum Cryptography transition infrastructure''~\cite{eccc}; Digital Europe
Programme for PQC migration tooling and crypto-agility middleware.

\textit{Standardisation:}
NIST NCCoE ``Migration to Post-Quantum Cryptography'' project (reference architecture
and migration playbooks); CEN/CENELEC JTC~13/WG~8 (hardware security standards aligned
with EUCC Common Criteria); ETSI TC CYBER (ETSI TS~103~744 quantum-safe hybrid key
exchanges~\cite{etsi103744}); ENISA PQC transition guidelines and cryptographic inventory
tooling; 3GPP SA3 for physical-layer security integration with 5G/6G
specifications~\cite{threeGPP33501}.

\textit{Ecosystem Enablers:}
RISC-V International open-hardware community; OpenTitan coalition pilot deployments
for EU sovereign RoT~\cite{opentitan}; SNS~JU pilot networks for 6G-integrated PhySec
validation~\cite{snsju}; federated Cyber Range infrastructure (CL3.5) for Island Mode
and TEE attestation stress-testing.

% ---- DIMENSION 5 ----
\item \textbf{Disruptors \& Systemic Risks}

\begin{itemize}
  \item \textbf{Q-Day cliff:} CRQC becoming operational before EU PQC migration
        completes; NIST IR~8413 places this horizon at 10--20 years with a highly
        uncertain lower bound~\cite{nistir8413}. The EU PQC migration timeline targeting
        critical infrastructure completion by 2030 has limited buffer against the
        pessimistic lower bound, and HNDL campaigns make this a present-tense threat
        against long-lived secrets.
  \item \textbf{EU Chips Act EUV lithography shortfall:} RISC-V secure element
        production at the node geometries required for EUCC `High' certification depends
        on advanced lithography; EU 2030 Chips Act targets (20\% global market share)
        do not specify security-critical silicon capacity independently, and ASML EUV
        tooling remains subject to US--Netherlands export control dynamics that could
        constrain sovereign fabrication timelines beyond 2028.
  \item \textbf{Side-channel analysis breakthroughs against FIPS~203 software
        implementations:} lattice-based algorithms have demonstrated timing and
        cache-timing attack surfaces in software implementations on constrained devices;
        hardware countermeasures are not yet standardised at the time of writing
        [CITATION REQUIRES VERIFICATION --- recent SCA against ML-KEM implementations].
  \item \textbf{TEE architectural monoculture:} if EU open-hardware TEE adoption
        concentrates on a single architecture (e.g., Keystone only~\cite{keystone}), a
        single discovered vulnerability propagates across the entire EU critical
        infrastructure TEE estate without architectural diversity as a containment layer.
  \item \textbf{LPWAN regulatory fragmentation:} Island Mode fallback using LoRaWAN
        operates on unlicensed 868\,MHz spectrum in the EU; spectrum policy
        fragmentation across Member States could prevent coordinated emergency
        micro-slice activation under unified crisis response scenarios.
\end{itemize}

% ---- DIMENSION 6 ----
\item \textbf{Failsafe Mechanisms \& Resilience}

\begin{itemize}
  \item \textbf{Hybrid cryptographic transition:} deploy ML-KEM (FIPS~203)~\cite{fips203}
        combined with ECDH-P384 in a hybrid key encapsulation scheme per
        ETSI TS~103~744~\cite{etsi103744}; if ML-KEM is found vulnerable via cryptanalytic
        breakthrough, ECDH-P384 maintains classical confidentiality while SLH-DSA
        (FIPS~205)~\cite{fips205} provides backup hash-based signature integrity with
        distinct mathematical foundations.
  \item \textbf{TEE implementation diversity:} deploy a minimum of two architecturally
        distinct TEE technologies across critical EU infrastructure (ARM CCA on
        server-class hardware and RISC-V Keystone~\cite{keystone} on edge nodes) to
        prevent single-architecture compromise from cascading across the entire
        confidential computing estate.
  \item \textbf{Autonomous node quarantine via eBPF:} if TEE remote attestation fails,
        eBPF-based network policy enforcement at the Linux kernel level (e.g., via
        Cilium or Falco eBPF rulesets) isolates the affected node from the orchestration
        plane within sub-second latency, blocking lateral movement before human
        intervention is available.
  \item \textbf{Island Mode failsafe:} edge nodes maintain cryptographically signed
        cached credentials with defined 72-hour validity windows; on WAN disconnection,
        a LoRaWAN LPWAN micro-slice activates a minimal viable network carrying
        life-safety telemetry; if LPWAN also fails, nodes revert to local-only operation
        with integrity-protected audit logging for post-incident forensics.
  \item \textbf{RF fingerprinting fallback:} for zero-energy IoT nodes where
        cryptographic processing exceeds the available energy budget, Physical-Layer
        authentication via RF fingerprinting (CSI/CFO-based device identity) provides
        keyless device verification using intrinsic wireless channel characteristics,
        requiring no cryptographic key storage or computation.
\end{itemize}

% ---- DIMENSION 7 ----
\item \textbf{Transversal Meta-Layer}

\textit{Sovereignty \& Autonomy:}
This cluster eliminates approximately 100\% foreign dependency on non-EU TPMs (Infineon,
NXP, STMicroelectronics dominate the EU market under TCG TPM~2.0~\cite{tgcpm20}) and
non-EU TEE architectures (Intel TDX, AMD SEV-SNP, ARM TrustZone, all non-EU designed and
manufactured). Crypto-agility frameworks additionally reduce dependency on US-controlled
NIST algorithm supply chains by enabling algorithm substitution if standardisation processes
are constrained by geopolitical factors.

\textit{Trust \& Ethics:}
GDPR Article~25 (data protection by design and by default) is operationalised through
TEE-enforced data minimisation and in-use encryption in multi-tenant environments,
ensuring that sensitive data processing cannot be accessed by platform operators~\cite{gdpr}.
AI Act Article~9 risk management system requirements for high-risk AI components require
hardware integrity guarantees --- only EUCC `High' certified secure elements provide the
mathematically verifiable trust anchor this obligation demands~\cite{aiact}.

\textit{Sustainable Security (Green IT):}
CRYSTALS-Kyber (ML-KEM) demonstrates approximately 1.5$\times$ computational overhead
relative to ECDH at equivalent security levels on ARM Cortex-M4, per NIST IR~8413~\cite{nistir8413},
representing a manageable energy penalty. Lightweight PQC research targets sub-100\,$\mu$J
per authentication operation for zero-energy IoT nodes. EU-sovereign RISC-V secure
elements designed for edge power constraints enable more energy-efficient cryptographic
operations than general-purpose processor TPM emulation, reducing per-device carbon
footprint at continental deployment scale.

\end{enumerate}



% =========================================================================
% CLUSTER 1: SOVEREIGN TRUSTED INFRASTRUCTURE
% =========================================================================
\newpage

\subsection{Progression Matrix: Priority Cluster 1 --- Sovereign Trusted Infrastructure}

\begin{landscape}
\setlength{\textwidth}{26.5cm}\setlength{\linewidth}{26.5cm}%
\setlength{\hsize}{26.5cm}\setlength{\columnwidth}{26.5cm}
\setlength{\tabcolsep}{5pt}
\renewcommand{\arraystretch}{1.35}

\begin{longtable}{%
|>{\footnotesize\raggedright\arraybackslash}p{2.8cm}%
|>{\footnotesize\raggedright\arraybackslash}p{5.4cm}%
|>{\footnotesize\raggedright\arraybackslash}p{5.4cm}%
|>{\footnotesize\raggedright\arraybackslash}p{5.4cm}%
|>{\footnotesize\raggedright\arraybackslash}p{5.4cm}|}

\caption{\textbf{Cluster~1: Sovereign Trusted Infrastructure} (P1, P2, P3, P8, P9).
Transition matrix across the four methodology dimensions. Sub-priority labels in
parentheses indicate where individual priorities diverge from the cluster-level
descriptor within a cell. ODP scale: Applied R\&D $\to$ Fragmented Piloting
$\to$ Standardised Integration $\to$ Ubiquitous Operation. SSL scale:
High Foreign Dependency $\to$ Emerging EU Ecosystem $\to$ Full Sovereign Autonomy.}
\label{tab:cluster1}\\

\hline
\rowcolor{colheader}
\textbf{Phase / SSL} &
\textbf{R\&I Maturity}{\scriptsize ODP indicator per sub-priority} &
\textbf{Policy, Regulatory \& Standards}{\scriptsize enforcement $\cdot$
  certification $\cdot$ standardisation} &
\textbf{Operational Deployment \& Adoption}{\scriptsize actor types $\cdot$
  scale $\cdot$ sector coverage} &
\textbf{Market Uptake \& Strategic Sovereignty}{\scriptsize SSL indicator
  $\cdot$ dependency reduction $\cdot$ industrial capacity} \\
\hline
\endfirsthead

\hline
\rowcolor{colheader}
\textbf{Phase / SSL} &
\textbf{R\&I Maturity} &
\textbf{Policy, Regulatory \& Standards} &
\textbf{Operational Deployment \& Adoption} &
\textbf{Market Uptake \& Strategic Sovereignty} \\
\hline
\endhead
\hline\multicolumn{5}{r}{\footnotesize\textit{Continued\ldots}}\\ \endfoot
\hline \endlastfoot

% ---- PHASE 1 ----
\rowcolor{phase1row}
\textbf{Phase~1}
\newline\textit{Foundation \&}
\newline\textit{Compliance}
\newline\textit{2025--2030}
\newline\sslbadge{sslredbg}{sslredtx}{SSL: High Foreign Dependency}
&
\cellhead{Cluster ODP: Applied R\&D / Fragmented Piloting}
The cluster is split at baseline. \subp{P1, P9} reach Fragmented
Piloting: RISC-V secure element projects targeting EUCC `High' are
funded under Horizon Europe and EU Chips Act; Island Mode soft-state
protocols are validated in 5G research testbeds.
\subp{P2, P3, P8} remain at Applied R\&D: NIST FIPS~203/204/205
algorithms are being optimised for 8/16-bit microcontrollers under
severe energy and memory constraints; Keystone and OpenTitan
prototypes are under EUCC pre-evaluation; RF fingerprinting
Physical-Layer Security is validated in controlled 5G lab environments;
Fully Homomorphic Encryption approaches practical latency bounds in
research settings only.
\kpistrip{The cluster ODP is Applied R\&D, driven by the majority of
sub-priorities. Deployment claims for P2, P3, P8 must be grounded in
named Horizon Europe or ECCC project outcomes, not technology
previews.}
&
\cellhead{Regulatory environment: enforcing, certification nascent}
NIS2 Article~21(2)(e) cryptographic policies for essential entities
entered full enforcement in October~2024. CRA Articles~10 and~11
(vulnerability handling and security update obligations) begin
enforcement in 2027, creating immediate procurement-level demand for
hardware security certification. AI~Act Articles~9 and~15 establish
hardware integrity requirements for high-risk AI systems and enter
applicability in 2025--2026. ENISA PQC migration guidelines are
published; EUCC evaluation processes are initiated for the first RISC-V
prototypes, but no EU-specific EUCC~`High' certified sovereign secure
element exists yet. CEN/CENELEC JTC~13 standards for hardware security
MIMs are in draft phase.
\kpistrip{The regulatory framework is enforcing existing obligations
but lacks EU-specific hardware certification outputs; the critical gap
is the absence of a ratified EUCC~`High' EU-origin product.}
&
\cellhead{Deployment: research consortia and national agencies only}
RISC-V secure element pilots are operated by 3--5 Horizon Europe
consortia and national cybersecurity agencies in controlled
environments. Intel~TDX, AMD~SEV-SNP, and ARM~TrustZone dominate
all operational TEE deployments across EU cloud and critical
infrastructure. PQC is confined to cloud-side software implementations
in pilot contexts; far-edge IoT devices run classical cryptography
exclusively. Edge nodes operate with hard cloud dependencies and no
validated autonomous fallback architecture meeting DORA or NIS2
business continuity thresholds. SMEs have no access to EU-sovereign
hardware alternatives at viable cost.
\kpistrip{Operational deployment at Phase~1 is institutional and
research-concentrated. The gap between pilot existence and operational
availability for regulated-sector procurement must be explicitly
acknowledged.}
&
\cellhead{SSL: High Foreign Dependency}
Non-EU TPMs (Infineon~SLB9670, NXP~SE050, STMicroelectronics~ST33)
hold $\geq$90\% of EU market share in hardware security modules.
Intel, AMD, and ARM control the entire operational TEE market. No
EU-designed and EU-manufactured secure element is commercially
available. EU Chips Act fabrication investment for security-critical
silicon is initiated but pre-production. US NIST algorithm supply chain
is the sole PQC standardisation reference. Structural vulnerability is
direct: geopolitical restriction of access to any of the three major
TPM suppliers would immediately impair EU critical infrastructure
attestation capability.
\kpistrip{Quantified dependency at Phase~1 baseline: $\geq$90\% foreign
market share, zero EU-certified sovereign hardware alternatives.
Sovereignty progress is measured from this specific figure.}
\\
\hline

% ---- GATE 1->2 ----
\rowcolor{gaterow}
\gatelbl{1 $\rightarrow$ 2} &
\multicolumn{4}{>{\footnotesize\raggedright\arraybackslash}p{22.0cm}|}{%
\cellcolor{gaterow}%
\textbf{R\&I trigger:} At least one EU-designed RISC-V secure element prototype
has been independently validated against EUCC~`High' evaluation criteria and has
entered the formal ENISA/national scheme certification process.
\textbf{Policy trigger:} CRA Articles~10 and~11 are in active enforcement with
documented non-compliance decisions, confirming that market incentives for
certified EU hardware alternatives are financially material.
\textbf{Operational trigger:} Documented operational deployment of hybrid
ML-KEM (FIPS~203) plus ECDH-P384 PQC schemes in at least one EU critical
infrastructure sector (energy or financial) at production scale, not pilot scale.
\textbf{Market trigger:} EU Chips Act security-critical silicon fabrication
pilot line has produced and delivered its first commercially viable batch of
EU-origin secure elements, confirmed by independent yield and cost assessment.}
\\
\hline

% ---- PHASE 2 ----
\rowcolor{phase2row}
\textbf{Phase~2}
\newline\textit{Applied \& Industrial}
\newline\textit{R\&D alongside}
\newline\textit{Federation \&}
\newline\textit{Standardisation}
\newline\textit{2030--2035}
\newline\sslbadge{ssltransbg}{ssltrantx}{SSL: Emerging EU Ecosystem}
&
\cellhead{Cluster ODP: Standardised Integration / continued Applied R\&D}
\subp{P1, P9} reach Standardised Integration: the first EUCC~`High' certified
EU-manufactured RISC-V secure elements are in production and being integrated
into critical infrastructure; automated remote attestation at fleet scale is
validated; 72-hour Island Mode achieves standardised certification under
3GPP/ETSI fallback API profiles.
\subp{P2, P3} transition from Applied R\&D to Fragmented Piloting: hybrid
ML-KEM plus ECDH-P384 schemes are validated at Industrial IoT scale; lightweight
PQC for sub-100~$\mu$J zero-energy nodes is demonstrated under 6G ISAC testbeds.
\subp{P8} reaches Fragmented Piloting: RISC-V Keystone TEEs are operational
in multi-provider Data Spaces; FHE performance approaches the threshold for
selective production use in high-value financial analytics.
\kpistrip{Phase~2 R\&I is not purely consolidation: three of five sub-priorities
are still in active applied and industrial R\&D. Describing this phase as
``standardisation only'' misrepresents the maturation arc.}
&
\cellhead{EUCC certification and PQC standards formalised}
The EUCC~`High' certification of the first EU-manufactured RISC-V secure
element is the definitive Phase~2 policy gate: it transforms applied research
into a market-accessible, legally certifiable product and unlocks the procurement
pathway for EU critical infrastructure operators. NIST FIPS~203/204/205 are
adopted into ETSI~TS profiles for EU deployment contexts. ENISA publishes a
mandatory PQC migration timeline for critical infrastructure operators, referenced
in NIS2 sector-specific implementing acts. EU Chips Act fabrication investments
are active and measurable. CEN/CENELEC hardware security standards for EU-origin
secure elements enter ratification.
\kpistrip{The single Phase~2 policy gate is EUCC~`High' certification of a
EU-origin product. All other policy milestones are enabling conditions, not
gate conditions.}
&
\cellhead{Critical infrastructure and regulated sectors deploying EU alternatives}
EU EUCC~`High' certified secure elements are deployed in energy sector
substations and telecommunications critical infrastructure in at least five
Member States. Hybrid PQC schemes are standard in new Industrial IoT and
financial sector (DORA-compliant) deployments. RISC-V Keystone TEEs operate
in federated multi-provider Data Spaces alongside existing Intel TDX
infrastructure. Validated 72-hour Island Mode is deployed in at least three
EU 5G cross-border corridor pilots under SNS~JU governance. SME-accessible
PQC tooling and SBOM compliance support are available commercially.
\kpistrip{Phase~2 deployment evidence requires named sector and Member State
scope; anonymous ``multi-sector deployment'' claims do not constitute
verifiable operational maturity.}
&
\cellhead{SSL: Emerging EU Ecosystem}
EU-designed secure elements achieve 15--30\% market share in new critical
infrastructure procurement, rising from a near-zero baseline. EU Chips Act
fabrication investments have created a domestic production capacity that is
commercially credible, if not yet dominant. Foreign TPMs remain the standard
for legacy deployments but are no longer the only option for new regulated-sector
procurement. EU open-hardware TEE ecosystem (Keystone, OpenTitan) holds
commercially certifiable products. The conditions for Full Sovereign Autonomy
are structurally in place --- fabrication capacity, certification infrastructure,
standards leadership --- but have not yet reached continental operational scale.
\kpistrip{SSL transition requires industrial substance, not research demonstration.
15--30\% new critical infra market share is the quantified Phase~2 SSL threshold.}
\\
\hline

% ---- GATE 2->3 ----
\rowcolor{gaterow}
\gatelbl{2 $\rightarrow$ 3} &
\multicolumn{4}{>{\footnotesize\raggedright\arraybackslash}p{22.0cm}|}{%
\cellcolor{gaterow}%
\textbf{R\&I trigger:} Production-ready PQC-native RISC-V silicon is available
from EU-domiciled fabrication at commercially viable yields, confirmed by
independent cost-per-unit assessment against equivalent non-EU alternatives.
\textbf{Policy trigger:} EUCC~`High' certification is mandated as a procurement
requirement for all new EU critical infrastructure edge node deployments, with
no waiver provisions for new procurements (legacy replacement programmes are
time-bound).
\textbf{Operational trigger:} Validated 72-hour Island Mode is deployed as a
standard operational baseline --- not a pilot --- in $\geq$50\% of new EU
NIS2-essential entity edge node procurements in the energy and
telecommunications sectors.
\textbf{Market trigger:} EU Chips Act sovereign fabrication capacity achieves
$\geq$50\% self-sufficiency for security-critical silicon in new EU critical
infrastructure procurement, measured by share of new orders fulfilled by
EU-origin fabrication.}
\\
\hline

% ---- PHASE 3 ----
\rowcolor{phase3row}
\textbf{Phase~3}
\newline\textit{Autonomy \&}
\newline\textit{Leadership}
\newline\textit{2035--2040}
\newline\sslbadge{ssltealbg}{ssltealtx}{SSL: Full Sovereign Autonomy}
&
\cellhead{Cluster ODP: Ubiquitous Operation}
All five sub-priorities operate at Ubiquitous Operation. EU-designed
RISC-V secure elements are the hardware default across all new EU ICT
products. Fully Homomorphic Encryption is operationalised for
end-to-end confidential processing from IoT to HPC, not confined to
high-value analytics. Zero-energy PhySec authentication via RF
fingerprinting operates as a standard device identity mechanism for
battery-less mMTC nodes. Bio-inspired swarm protocols enable autonomous
cyber-physical recovery without cloud intervention. Phase~3 R\&I is
oriented toward the next-generation frontier: neuromorphic secure
compute, post-second-generation quantum hardware, and post-CMOS
security architectures --- EU research groups define the agenda rather
than adopting it.
\kpistrip{Ubiquitous Operation requires deployment statistics, not
deployment plans. Phase~3 evidence must cite procurement share data
and operational audit findings, not programme milestone completions.}
&
\cellhead{EU as international hardware security standards reference}
EUCC~`High' certification is mandatory for all new EU ICT products
without waiver provisions. PQC-native silicon is the EU procurement
default. EU hardware security standards are referenced in ISO/IEC
international standardisation and in bilateral digital trade
agreements. Continuous automated EUCC lifecycle maintenance replaces
point-in-time audit models completely. The EU sets the terms of
hardware security compliance for market access, transforming its
internal market standard into a de facto international reference.
\kpistrip{Phase~3 policy maturity requires evidence of international
standards influence: at least one EU hardware security standard adopted
or formally referenced in an ISO/IEC or ITU-T international standard.}
&
\cellhead{Security-by-Design the unconditional hardware and crypto default}
$\geq$60\% of EU critical infrastructure edge nodes deploy EUCC~`High'
certified EU-designed secure elements. FHE is operational end-to-end
from IoT sensor to HPC compute node. Autonomous Island Mode provides
validated 72-hour survivability as a mandatory operational baseline
across all NIS2-essential and DORA-critical edge deployments. No new
EU ICT product enters the market without a hardware root of trust.
The question ``does this deployment have hardware-level integrity?''
has no non-trivial answer for any new EU deployment.
\kpistrip{Phase~3 deployment normalisation: $\geq$60\% market penetration
for EU secure elements in critical infra; 100\% of new EU ICT products
with hardware RoT as a default.}
&
\cellhead{SSL: Full Sovereign Autonomy}
EU-designed and EU-manufactured secure elements hold $\geq$60\% of new
critical infrastructure procurement market share. Foreign legacy
dependency operates under binding replacement programmes with defined
timelines. EU sovereign RISC-V silicon and open-hardware TEE
architectures are exported to and adopted by partner jurisdictions.
EU PQC and hardware TEE standards are the international reference.
Structural exposure to geopolitical disruption of any single non-EU
hardware supplier has been eliminated for new deployments. The EU
has moved from defensive sovereignty (reducing dependency) to offensive
strategic posture (defining the terms of global hardware security
competition).
\kpistrip{Full Sovereign Autonomy is verified by $\geq$60\% market share
in new critical infra procurement plus the existence of binding
replacement programmes for legacy foreign hardware --- not by 100\%
market share.}
\\
\hline

\end{longtable}
\end{landscape}


% ---- PROGRESSION MATRIX: CLUSTER 1 ----

% \noindent
% \resizebox{\textwidth}{!}{%
% \begin{tabular}{|p{4.2cm}|p{5.6cm}|p{5.6cm}|p{5.6cm}|}
% \hline
% \rowcolor{gray!30}
% \multicolumn{4}{|c|}{\textbf{Progression Matrix: Priority Cluster 1 ---
%   Sovereign Trusted Infrastructure (P1, P2, P3, P8, P9)}} \\
% \hline
% \rowcolor{gray!15}
% \textbf{Metric} &
% \textbf{Phase 1: Foundation \& Compliance (2025--2030)} &
% \textbf{Phase 2: Federation \& Standardisation (2030--2035)} &
% \textbf{Phase 3: Autonomy \& Leadership (2035--2040)} \\
% \hline
% \rowcolor{gray!10}
% \multicolumn{4}{|l|}{\textit{\textbf{Cluster-Level Progression Indicators}}} \\
% \hline
% \textbf{Cluster ODP} &
% Applied R\&D &
% Standardised Integration &
% Ubiquitous Operation \\
% \hline
% \textbf{Cluster SSL} &
% High Foreign Dependency &
% Emerging EU Ecosystem &
% Full Sovereign Autonomy \\
% \hline
% \rowcolor{gray!10}
% \multicolumn{4}{|l|}{\textit{\textbf{Sub-Priority Decomposition (Individual ODP / SSL per Phase)}}} \\
% \hline
% \textit{P1: Hardware RoT \newline \& EUCC `High'} &
% ODP: Fragmented Piloting \newline
% SSL: High Foreign Dependency &
% ODP: Standardised Integration \newline
% SSL: Emerging EU Ecosystem &
% ODP: Ubiquitous Operation \newline
% SSL: Full Sovereign Autonomy \\
% \hline
% \textit{P2: PQC Agility \newline for Edge/IoT} &
% ODP: Applied R\&D \newline
% SSL: High Foreign Dependency &
% ODP: Fragmented Piloting \newline
% SSL: Emerging EU Ecosystem &
% ODP: Ubiquitous Operation \newline
% SSL: Full Sovereign Autonomy \\
% \hline
% \textit{P3: Lightweight Crypto \newline \& Physical-Layer Sec.} &
% ODP: Applied R\&D \newline
% SSL: High Foreign Dependency &
% ODP: Standardised Integration \newline
% SSL: Emerging EU Ecosystem &
% ODP: Ubiquitous Operation \newline
% SSL: Full Sovereign Autonomy \\
% \hline
% \textit{P8: Confidential Computing \newline \& Privacy-Pres. AI} &
% ODP: Applied R\&D \newline
% SSL: High Foreign Dependency &
% ODP: Fragmented Piloting \newline
% SSL: Emerging EU Ecosystem &
% ODP: Ubiquitous Operation \newline
% SSL: Full Sovereign Autonomy \\
% \hline
% \textit{P9: Resilient-by-Design \newline \& Island Mode} &
% ODP: Fragmented Piloting \newline
% SSL: High Foreign Dependency &
% ODP: Standardised Integration \newline
% SSL: Contained Local Survivability &
% ODP: Ubiquitous Operation \newline
% SSL: Systemic EU Resilience \\
% \hline
% \rowcolor{blue!7}
% \textbf{Operational State} &
% EU-funded RISC-V secure element pilot projects commence EUCC evaluation;
% NIST FIPS~203/204/205 algorithms are adopted by research consortia and
% optimised for constrained devices in laboratory settings; Intel TDX
% and AMD SEV-SNP dominate all operational TEE deployments while
% EU open-hardware alternatives remain at prototype stage;
% edge nodes operate with hard cloud dependencies and lack
% autonomous fallback architectures validated at operational scale. &
% EU-manufactured EUCC `High' certified secure elements enter
% critical infrastructure deployments across energy and
% telecommunications sectors; hybrid ML-KEM plus ECDH-P384 schemes
% are standardised and deployed in Industrial IoT corridors;
% RISC-V Keystone-based TEEs are operational in federated
% multi-provider Data Spaces; validated 72-hour Island Mode
% survivability demonstrated in designated EU 5G cross-border
% corridor pilots under SNS~JU governance. &
% EU Chips Act sovereign fabrication supplies $\geq$60\% of EU demand
% for security-critical silicon; all new EU ICT products ship with
% EUCC `High' certified hardware roots of trust by default;
% fully homomorphic encryption is operationalised for
% end-to-end confidential processing from IoT to HPC;
% bio-inspired swarm protocols coordinate autonomous
% cyber-physical recovery without any cloud intervention,
% achieving DORA-mandated recovery time objectives
% autonomously at continental scale. \\
% \hline
% \rowcolor{orange!12}
% \textbf{Critical Enabler} &
% CRA Articles~10 and~11 enforcement (2024--2027) creating
% mandatory market demand for hardware security certification,
% establishing the commercial conditions for EU secure element
% pilot investment &
% EUCC `High' certification of the first EU-manufactured
% RISC-V secure element, constituting the gate condition
% for the hardware sovereignty transition and unlocking
% EU Chips Act fabrication investment justification &
% EU Chips Act sovereign EUV-capable fabrication capacity
% achieving $\geq$60\% market self-sufficiency in
% security-critical silicon, removing the last structural
% foreign dependency in the infrastructure trust chain \\
% \hline
% \end{tabular}%
% }

% =========================================================================
% PRIORITY CLUSTER 2: COGNITIVE SECURITY AND AUTONOMOUS CYBER DEFENCE
% =========================================================================

\newpage
\section{Priority Cluster 2: Cognitive Security and Autonomous Cyber Defence}
\label{sec:cluster2}

\noindent\textbf{Constituent Sub-Priorities:} P6 (Cognitive SecOps / LLMSecOps \& Automated
Threat Intelligence), P10 (Proactive APT Management), and the AI-as-target dimension
drawn from CL1.2 (Model Integrity / DataBOMs), CL1.3 (Dynamic Supply Chain Verification /
Runtime SBOMs), and CL4.7 (AI Trustworthiness \& Algorithmic Accountability).

\medskip
\noindent\textbf{Cluster Thesis:} Artificial intelligence is simultaneously the most powerful
instrument available for autonomous cyber defence and the most consequential new attack surface
introduced into the security architecture. A roadmap that deploys AI agents for autonomous
threat containment without hardening those agents against adversarial manipulation creates a
critical second-order vulnerability. This cluster therefore addresses three inseparable
sub-dimensions: AI as security instrument (LLMSecOps and APT hunting), AI as adversary
amplifier (AI-generated APT tooling), and AI as target (model poisoning, adversarial evasion,
and algorithmic accountability for autonomous containment decisions).

\subsection*{7-Dimensional Analysis: Priority Cluster 2}

\begin{enumerate}

% ---- DIMENSION 1 ----
\item \textbf{Baseline \& Target Objective}

\textit{Baseline (2025--2026):}
European Security Operations Centres (SOCs) rely predominantly on rule-based SIEM platforms ---
Splunk (US-domiciled, acquired by Cisco in 2023) and IBM QRadar dominate enterprise
deployments. AI-based tools are deployed as analyst ``copilots'' with human analysts
retaining all containment authority; no EU-sovereign foundation models for security
operations exist at production scale. APT dwell times in European enterprise environments
average 150--200 days, with state-sponsored threat actors demonstrating extended persistence
in EU critical infrastructure. The adversarial robustness of AI security agents --- their
resistance to prompt injection, model evasion, and training data poisoning --- is addressed
only in research contexts without operational certification frameworks. NWDAF integration
per 3GPP TS~23.288~\cite{threeGPP23288} exists in 5G standalone deployments but remains
fragmented across operators, and no EU-mandated framework for the trustworthiness of
autonomous security AI exists prior to AI Act enforcement.

Current cluster ODP: Applied R\&D (AI-as-target and sovereign LLM development); Siloed AI
Copilots (P6); Multi-Layered Alert Correlation (P10). Cluster ODP: Applied R\&D.
Cluster SSL: High Reliance on US Foundation Models / High Dependency on Global Threat Feeds.

\textit{Target (2040):}
\begin{sloppypar}
Swarm-based agentic AI defends the EU Cloud-Edge Continuum autonomously, with APT dwell
time reduced to under 24 hours in certified critical infrastructure deployments. EU-sovereign
edge Small Language Models (SLMs) and Large Context Models (LCMs) for security operations
are certified under AI Act Article~9 and cover at least 50\% of EU critical infrastructure
SOC deployments. All autonomous security agents pass mandatory adversarial robustness
validation before operational deployment. Moving Target Defense reduces APT reconnaissance
success rates below 5\% in validated simulation environments. EU AI security frameworks
are adopted as global reference implementations by partner jurisdictions.
\end{sloppypar}

% ---- DIMENSION 2 ----
\item \textbf{Primary Drivers}

\textit{Threat Drivers:}
\begin{itemize}
  \item Nation-state APT campaigns (attributed by EU agencies to APT28/GRU,
        APT29/SVR, and Sandworm) increasingly exploit LLM hallucination and
        prompt injection vulnerabilities to manipulate autonomous security agents
        into executing incorrect containment actions --- weaponising the defence
        architecture against itself.
  \item AI-generated spear-phishing and poly\-morphic malware
        enabling automated, highly-targeted APT campaigns that adapt
        faster than signature-based
        detection cycles, invalidating the static detection rules that rule-based
        SIEM platforms depend upon.
  \item Supply chain attacks targeting AI model training pipelines: data poisoning
        and backdoor insertion into security AI models deployed at continental scale
        can silently degrade detection capability without triggering observable
        anomalies in standard monitoring.
  \item Systematic abuse of MITRE ATT\&CK-catalogued techniques (T1027 Obfuscated
        Files or Information; T1055 Process Injection) within AI-assisted
        reconnaissance tools, enabling more efficient bypass of detection rules
        derived from those same frameworks.
\end{itemize}

\textit{Regulatory Drivers:}
\begin{itemize}
  \item AI Act Article~9: risk management system requirements for high-risk AI
        systems --- autonomous security containment agents making access revocation
        or network isolation decisions qualify as high-risk systems requiring
        mandatory conformity assessment before operational deployment~\cite{aiact}.
  \item AI Act Article~15: accuracy, robustness, and cybersecurity requirements
        for high-risk AI, including specific resilience requirements against
        adversarial manipulation of inputs and training data~\cite{aiact}.
  \item AI Act Article~13: transparency and provision of information for high-risk
        AI --- autonomous security decisions must be explainable to affected
        entities and competent national authorities~\cite{aiact}.
  \item NIS2 Article~21(2)(b): incident handling procedures, including
        AI-augmented detection and response, for essential entities~\cite{nis2}.
\end{itemize}

% ---- DIMENSION 3 ----
\item \textbf{High Impact Activities}

\begin{enumerate}[label=(\alph*)]
  \item \textbf{[Phase~1]} Deploy NWDAF-integrated (3GPP TS~23.288~\cite{threeGPP23288})
        telemetry pipelines feeding ML-based UEBA (User and Entity Behaviour Analytics)
        for APT indicator correlation; establish behavioural baseline profiles across
        federated enterprise environments targeting reduction of APT dwell time from
        the current 150--200 day average to under 72 hours in pilot deployments via
        ML-assisted Indicator of Compromise (IoC) pattern detection.

  \item \textbf{[Phase~1--2]} Develop EU-sovereign Small Language Models (SLMs,
        under 10B parameters) and Large Context Models (LCMs) for edge-native SOC
        operations (LLMSecOps), trained on MITRE ATT\&CK-aligned threat intelligence
        corpora under GDPR-compliant data governance~\cite{gdpr}; establish mandatory
        AI Act Article~9 conformity assessment pipelines for all autonomous security
        agent deployments prior to production authorisation~\cite{aiact}.

  \item \textbf{[Phase~2]} Deploy edge-native ``Guardian'' nodes executing autonomous
        Zero-Trust reconfigurations per NIST SP~800-207 ZTA principles~\cite{nistsp800207}
        and Moving Target Defense (MTD) --- dynamically shuffling IP addressing, network
        topology, and software diversity to disrupt APT reconnaissance --- governed by
        ETSI GS~ZSM~009 closed-loop automation standards~\cite{etsizsm009}.

  \item \textbf{[Phase~2]} Establish mandatory adversarial robustness validation
        pipelines for all AI security agents prior to operational deployment: automated
        red-team adversarial testing (prompt injection, model evasion, backdoor
        detection via CL1.4 purple-team automation), with continuous certification
        under a CEN/CENELEC-standardised AI security agent assurance framework
        aligned with CL4.7 AI Trustworthiness requirements.

  \item \textbf{[Phase~3]} Deploy swarm-based agentic AI coordination protocols
        enabling decentralised, self-evolving threat intelligence synthesis across
        the EU Cloud-Edge Continuum; integrate with Deception Technology
        infrastructure (CL2.6: honeypots, honeynets, credential canaries, fake API
        endpoints) to generate high-fidelity, low-noise adversarial ground truth for
        continuous autonomous agent retraining.
\end{enumerate}

% ---- DIMENSION 4 ----
\item \textbf{Catalysts \& Enablers}

\textit{Funding Alignment:}
EuroHPC Joint Undertaking AI Factories for training EU-sovereign security foundation
models at the compute scale required; ECCC Strategic Agenda priority areas ``Cybersecurity
of AI'' and ``AI-enabled cybersecurity tools and services''~\cite{eccc}; Horizon Europe
Cluster~3 (Civil Security for Society) calls targeting AI-driven threat intelligence and
autonomous incident response; EU AI Office for conformity assessment infrastructure
supporting high-risk security AI certification under the AI Act.

\textit{Standardisation:}
ETSI GS~ZSM~009 (Zero-Touch Network and Service Management, Closed-Loop
Automation~\cite{etsizsm009}); 3GPP TS~23.288 (NWDAF architecture for AI-native 5G
network analytics~\cite{threeGPP23288}); MITRE ATT\&CK framework (behavioural threat
intelligence taxonomy for agent training corpora); NIST SP~800-207 (Zero Trust
Architecture~\cite{nistsp800207}); CEN/CENELEC JTC~13 working groups on AI security
agent assurance standardisation.

\textit{Ecosystem Enablers:}
ENISA AI cybersecurity threat landscape guidelines; federated Cyber Range infrastructure
(CL3.5) for adversarial agent testing under realistic attack scenarios; CONCORDIA
successor programme for EU-wide SOC federation and cross-border CTI sharing~\cite{concordia};
EU AI Act notified conformity assessment bodies for high-risk AI product certification.

% ---- DIMENSION 5 ----
\item \textbf{Disruptors \& Systemic Risks}

\begin{itemize}
  \item \textbf{Adversarial AI arms race:} threat actors weaponising LLMs for
        automated zero-day discovery and polymorphic exploit generation may outpace
        the defensive AI maturation timeline --- specifically, AI-generated
        vulnerability discovery could exceed human patching bandwidth within the
        Phase~1--2 window, before EU-sovereign autonomous response agents are
        certified for operational deployment.
  \item \textbf{LLM hallucination in containment agents:} false positive isolation
        of critical infrastructure nodes (e.g., a hospital network segment incorrectly
        quarantined during a medical emergency) constitutes an operational safety risk
        that may impose worse harm than the threat the agent was designed to
        contain --- this is the central argument for mandatory HITL escalation in Phase~1.
  \item \textbf{EU sovereign LLM compute deficit:} training security-relevant
        foundation models requires EuroHPC-scale compute; as of 2025--2026, EU frontier
        model training capacity lags US and China by approximately one order of
        magnitude, constraining Phase~1 EU-sovereign models to SLM/LCM scale
        ($<$10B parameters) and limiting capability relative to US foundation models.
  \item \textbf{AI model poisoning via CTI feed corruption:} if adversaries inject
        false indicators into NWDAF analytics pipelines or Cyber Threat Intelligence
        (CTI) sharing platforms (CL2.1), behavioural baselines trained on poisoned
        data will silently misclassify legitimate traffic and fail to detect real APTs
        --- a harder failure mode to detect than outright system unavailability.
  \item \textbf{AI Act certification bottleneck:} mandatory conformity assessment
        for high-risk autonomous security agents may create a 24--36 month regulatory
        gap during which EU organisations operate without certified autonomous
        containment, relying on slower human-in-the-loop processes that cannot match
        AI-assisted adversary speed.
\end{itemize}

% ---- DIMENSION 6 ----
\item \textbf{Failsafe Mechanisms \& Resilience}

\begin{itemize}
  \item \textbf{Shadow mode operation:} all autonomous AI security agents run in
        parallel with human-supervised rule-based SIEM during Phases~1--2;
        containment actions affecting Tier-1 critical infrastructure (energy,
        healthcare, transport per NIS2 Annex~I) require mandatory Human-In-The-Loop
        (HITL) escalation regardless of AI confidence score, with reversible quarantine
        actions bounded by a defined 60-minute rollback window~\cite{nis2}.
  \item \textbf{Model diversity and behavioural circuit breakers:} deploy a minimum
        of two architecturally distinct AI detection models (e.g., transformer-based
        LLM plus graph neural network UEBA) to prevent a single-model adversarial
        exploit from disabling the entire detection capability; behavioural circuit
        breakers automatically suspend any agent exhibiting an anomalous action rate
        or pattern deviating beyond $3\sigma$ from its established operational
        baseline, reverting to rule-based SIEM until re-certification.
  \item \textbf{RLHF continuous correction:} incorporate analyst override decisions
        into continuous Reinforcement Learning from Human Feedback (RLHF) retraining
        to prevent model drift from training distribution; quarterly adversarial
        red-team audits (CL1.4) detect capability degradation before it reaches
        operational significance.
  \item \textbf{Air-gapped CTI validation:} all external CTI feeds ingested by
        NWDAF analytics pipelines pass through an isolated validation environment
        where indicator quality is assessed against independent behavioural ground
        truth before poisoned data can influence operational AI model weights.
  \item \begin{sloppypar}\textbf{Deception-backed calibration:} honeynet and credential canary
        infrastructure (CL2.6) provides high-fidelity, low-noise ground truth against
        which AI behavioural models are continuously calibrated, reducing false
        positive rates and providing an independent signal for detecting model
        poisoning via discrepancy analysis between canary alerts and AI model
        predictions.\end{sloppypar}
\end{itemize}

% ---- DIMENSION 7 ----
\item \textbf{Transversal Meta-Layer}

\textit{Sovereignty \& Autonomy:}
This cluster eliminates reliance on US-domiciled SIEM platforms (Splunk/Cisco, IBM QRadar)
and US foundation models for operational security intelligence. EU-sovereign SLMs trained
under GDPR-compliant data governance~\cite{gdpr} ensure that sensitive threat intelligence
and behavioural telemetry does not transit non-EU cloud infrastructure. EU leadership in
trustworthy AI security tools creates an exportable technology capability aligned with
the AI Act's ambition for global regulatory standard-setting.

\textit{Trust \& Ethics:}
AI Act Article~13 (transparency requirements) is mandated for all autonomous security
decisions --- affected entities must receive meaningful explanations for AI-driven access
revocations or network isolation actions~\cite{aiact}. GDPR Article~22 (restrictions on
automated individual decision-making) applies where AI containment decisions affect
individual data subjects; human review mechanisms must be structurally embedded, not
optional add-ons~\cite{gdpr}. AI Act Article~15 adversarial robustness requirements
ensure security AI systems cannot be weaponised against the infrastructure they are
designed to protect~\cite{aiact}.

\textit{Sustainable Security (Green IT):}
Edge-native SLMs under 10B parameters reduce inference energy consumption by two to
three orders of magnitude compared to centralised GPT-class model inference, enabling
real-time threat classification at the point of detection without data centre energy
overhead. Federated training paradigms reduce repeated data transmission overhead
compared to centralised training architectures and are compatible with the privacy
preservation requirements of GDPR Article~25~\cite{gdpr}.

\end{enumerate}

% =========================================================================
% CLUSTER 2: COGNITIVE SECURITY AND AUTONOMOUS CYBER DEFENCE
% =========================================================================
\newpage
\subsection{Progression Matrix: Priority Cluster 2 --- Cognitive Security and Autonomous
Cyber Defence}

\begin{landscape}
\setlength{\textwidth}{26.5cm}\setlength{\linewidth}{26.5cm}%
\setlength{\hsize}{26.5cm}\setlength{\columnwidth}{26.5cm}
\setlength{\tabcolsep}{5pt}
\renewcommand{\arraystretch}{1.35}

\begin{longtable}{%
|>{\footnotesize\raggedright\arraybackslash}p{2.8cm}%
|>{\footnotesize\raggedright\arraybackslash}p{5.4cm}%
|>{\footnotesize\raggedright\arraybackslash}p{5.4cm}%
|>{\footnotesize\raggedright\arraybackslash}p{5.4cm}%
|>{\footnotesize\raggedright\arraybackslash}p{5.4cm}|}

\caption{\textbf{Cluster~2: Cognitive Security \& Autonomous Cyber Defence}
(P6, P10, AI-as-target dimension from CL1.2, CL1.3, CL4.7).
Transition matrix across the four methodology dimensions. The AI-as-target
dimension (model integrity, adversarial robustness, algorithmic accountability)
is labelled \textit{(AI-T)} where it diverges from P6 or P10 descriptors.}
\label{tab:cluster2}\\

\hline
\rowcolor{colheader}
\textbf{Phase / SSL} &
\textbf{R\&I Maturity}{\scriptsize ODP indicator per sub-priority} &
\textbf{Policy, Regulatory \& Standards}{\scriptsize enforcement $\cdot$
  certification $\cdot$ standardisation} &
\textbf{Operational Deployment \& Adoption}{\scriptsize actor types $\cdot$
  scale $\cdot$ sector coverage} &
\textbf{Market Uptake \& Strategic Sovereignty}{\scriptsize SSL indicator
  $\cdot$ dependency reduction $\cdot$ industrial capacity} \\
\hline
\endfirsthead

\hline
\rowcolor{colheader}
\textbf{Phase / SSL} &
\textbf{R\&I Maturity} &
\textbf{Policy, Regulatory \& Standards} &
\textbf{Operational Deployment \& Adoption} &
\textbf{Market Uptake \& Strategic Sovereignty} \\
\hline
\endhead
\hline\multicolumn{5}{r}{\footnotesize\textit{Continued\ldots}}\\ \endfoot
\hline \endlastfoot

% ---- PHASE 1 ----
\rowcolor{phase1row}
\textbf{Phase~1}
\newline\textit{Foundation \&}
\newline\textit{Compliance}
\newline\textit{2025--2030}
\newline\sslbadge{sslredbg}{sslredtx}{SSL: High Reliance on}
\newline\sslbadge{sslredbg}{sslredtx}{US Foundation Models}
&
\cellhead{Cluster ODP: Applied R\&D / Siloed AI Copilots}
\subp{P6} reaches Siloed AI Copilots: commercial LLMs (GPT-4-class)
assist human SOC analysts in alert triage, threat report generation, and
MITRE ATT\&CK technique mapping; no EU-sovereign LLM for security
operations exists at production scale; NWDAF (3GPP TS~23.288)
telemetry pipelines are integrated in 5G SA pilot deployments for
anomaly detection. \subp{P10} is at Multi-Layered Alert Correlation:
ML-based UEBA is applied to APT IoC detection in research and pilot
contexts; APT dwell time reduction from the 150--200 day average is
demonstrated in controlled settings. \subp{AI-T} is at Applied R\&D:
adversarial robustness testing (prompt injection, model evasion,
DataBOM verification) is an academic research function without
operational certification frameworks.
\kpistrip{Phase~1 R\&I distinguishes three sub-dimensions with
different maturity levels. Conflating them into a single ODP cell
obscures that AI-as-target capability is the least mature and the
most consequential gap.}
&
\cellhead{AI Act enforcing; conformity assessment infrastructure absent}
AI~Act Articles~9 (risk management), 13 (transparency), and 15
(adversarial robustness) enter applicability in 2025--2026, establishing
mandatory requirements for autonomous security AI as high-risk systems.
However, no EU conformity assessment bodies with capacity to certify
autonomous security agents are operational at the start of Phase~1.
ETSI GS~ZSM~009 (closed-loop automation) is published but not mandated.
NIS2 Article~21(2)(b) incident handling procedures are being
operationalised across essential entities. ENISA AI cybersecurity threat
landscape guidelines are published. The critical regulatory gap is the
absence of a standardised adversarial robustness test suite applicable
to security-specific AI agents.
\kpistrip{The AI Act creates obligations before the conformity
assessment infrastructure exists to fulfil them. This 24--36 month gap
is the dominant Phase~1 policy risk for this cluster.}
&
\cellhead{Deployment: SOC copilot use by large enterprises and national CSIRTs}
US-platform LLMs (GPT-4-class, under data processor agreements) operate
as analyst copilots in large enterprise and national CSIRT SOCs.
Rule-based SIEMs (Splunk/Cisco, IBM~QRadar) remain the primary
detection infrastructure across EU regulated sectors. NWDAF-integrated
AI analytics operate in 3--5 5G SA deployments run by major EU
telecom operators. No certified autonomous security containment agent
is in operational deployment anywhere in the EU. APT dwell times of
150--200 days in EU enterprise environments reflect the gap between
AI-assisted detection capability and operational deployment at scale.
SMEs lack access to AI-powered SOC tooling at viable cost.
\kpistrip{Phase~1 deployment is concentrated in the largest institutional
actors (national CSIRTs, Tier-1 operators, major financial entities).
The gap between these deployments and SME access defines the adoption
challenge for Phase~2.}
&
\cellhead{SSL: High Reliance on US Foundation Models}
The EU SOC tooling market is dominated by US-domiciled platforms:
Splunk (Cisco), IBM QRadar, and Microsoft Sentinel collectively control
the EU enterprise SIEM market. The foundation models used for SOC
copilot functions are US-origin (OpenAI, Anthropic, Google) and are
accessed under data processor agreements, creating persistent
extra-territorial data exposure. No EU-sovereign foundation model at
production scale for security operations exists. EuroHPC compute
capacity is insufficient for frontier-class security model training.
The adversarial robustness validation market is pre-commercial.
\kpistrip{Quantified dependency: US platforms hold $\geq$80\% of EU
enterprise SIEM market; US-origin models are the sole source for
LLM-based SOC tooling. This specific dependency structure is what
Phases~2 and~3 must reduce.}
\\
\hline

% ---- GATE 1->2 ----
\rowcolor{gaterow}
\gatelbl{1 $\rightarrow$ 2} &
\multicolumn{4}{>{\footnotesize\raggedright\arraybackslash}p{22.0cm}|}{%
\cellcolor{gaterow}%
\textbf{R\&I trigger:} At least one EU-sovereign Small Language Model
($\leq$10B parameters) trained on GDPR-compliant security corpora has
been validated against an independently benchmarked EU critical infrastructure
SOC use case at production-equivalent scale.
\textbf{Policy trigger:} AI~Act conformity assessment bodies with capacity
to certify autonomous security AI as high-risk systems under Articles~9 and~15
are operational and have issued at least one certificate to a security AI
agent product.
\textbf{Operational trigger:} At least one EU-sovereign LLMSecOps agent is
in certified operational deployment in a named NIS2-essential entity,
demonstrating APT dwell time below 72 hours in a documented production
environment.
\textbf{Market trigger:} A standardised adversarial robustness test suite
for security AI agents, developed under CEN/CENELEC or ETSI TC~CYBER
governance, has been published and is referenced in at least one EU
procurement framework for critical infrastructure AI tools.}
\\
\hline

% ---- PHASE 2 ----
\rowcolor{phase2row}
\textbf{Phase~2}
\newline\textit{Applied \& Industrial}
\newline\textit{R\&D alongside}
\newline\textit{Federation \&}
\newline\textit{Standardisation}
\newline\textit{2030--2035}
\newline\sslbadge{ssltransbg}{ssltrantx}{SSL: Sovereign EU Edge}
\newline\sslbadge{ssltransbg}{ssltrantx}{LLMs Emerging}
&
\cellhead{Cluster ODP: Federated Edge Agents / Automated Triage}
\subp{P6} reaches Federated Edge Agents: EU-sovereign SLMs ($\leq$10B
parameters) trained on MITRE ATT\&CK-aligned corpora are operational
at the edge for real-time SOC analytics without cloud inference
dependency; Guardian nodes execute autonomous Zero-Trust
reconfigurations per NIST~SP~800-207. \subp{P10} reaches Automated
Triage and LLMSecOps Containment: LLMSecOps agents ingest CTI feeds,
track APT lateral movement against behavioural frameworks, and
orchestrate incident triage autonomously under ETSI GS~ZSM~009
governance; APT dwell time $\leq$72 hours in certified pilot deployments.
\subp{AI-T} reaches Standardised Integration: adversarial robustness
validation pipelines (prompt injection, model evasion, backdoor
detection) are standardised and mandatory under AI~Act~Article~9
conformity assessment; continuous RLHF retraining and circuit-breaker
mechanisms are in production.
\kpistrip{Phase~2 R\&I includes active applied and industrial R\&D for
all three sub-dimensions --- swarm AI protocols, post-quantum AI
security, and next-generation adversarial testing frameworks --- not
merely consolidation of Phase~1 results.}
&
\cellhead{AI Act conformity assessment operational: the Phase~2 gate}
The single Phase~2 policy gate is the operationalisation of a
standardised AI~Act conformity assessment framework for autonomous
security agents, enabling certified deployment without per-case
regulatory approval. ETSI GS~ZSM~009 is mandated within EU NaaS
federation governance frameworks. A CEN/CENELEC-standardised adversarial
robustness test suite for security AI is published. MITRE ATT\&CK
behavioural framework is formally referenced in NIS2 sector-specific
incident handling guidance. EuroHPC AI Factory compute contracts for
EU-sovereign security model training are active.
\kpistrip{Without operational AI~Act conformity assessment
infrastructure, autonomous security agents cannot be certified for
regulated-sector deployment regardless of technical readiness.
This is the dominant Phase~2 policy dependency.}
&
\cellhead{Certified EU sovereign AI agents in critical infrastructure SOCs}
EU SLM-based LLMSecOps agents are in certified operational deployment
in energy, financial (DORA), and telecommunications critical
infrastructure SOCs in $\geq$3 Member States. Guardian nodes execute
autonomous Zero-Trust reconfigurations under human-in-the-loop
escalation protocols for Tier-1 critical actions. NWDAF-integrated
analytics are mandatory in 5G SA deployments. Deception infrastructure
(honeynets, credential canaries, fake API endpoints per CL2.6) is
deployed alongside AI detection systems, providing high-fidelity
ground truth for model calibration. APT dwell time $\leq$72 hours
demonstrated in pilot critical infrastructure deployments.
\kpistrip{Phase~2 deployment evidence requires named critical
infrastructure sector, Member State scope, and a documented APT dwell
time metric from an operational --- not simulated --- environment.}
&
\cellhead{SSL: Sovereign EU Edge LLMs Emerging}
EU-sovereign SLMs cover $\geq$30\% of critical infrastructure SOC
deployments in regulated sectors, rising from zero at baseline.
EuroHPC AI Factories supply training infrastructure for EU-origin
security models, reducing dependence on US hyperscaler compute.
US-platform SIEMs retain market share in legacy deployments but
face procurement competition from certified EU alternatives.
Adversarial robustness validation is a commercially available service
with EU-domiciled providers. EU AI security tool exports to partner
jurisdictions have begun, indicating the transition from security
consumer to security producer posture.
\kpistrip{SSL evidence at Phase~2: $\geq$30\% critical infra SOC market
share for EU-sovereign AI tools; at least one documented EU AI security
tool export event to a non-EU partner jurisdiction.}
\\
\hline

% ---- GATE 2->3 ----
\rowcolor{gaterow}
\gatelbl{2 $\rightarrow$ 3} &
\multicolumn{4}{>{\footnotesize\raggedright\arraybackslash}p{22.0cm}|}{%
\cellcolor{gaterow}%
\textbf{R\&I trigger:} EU-certified swarm AI coordination protocols
enabling decentralised autonomous threat intelligence synthesis have passed
independent security evaluation and are deployed in at least one cross-border
EU critical infrastructure corridor under ETSI GS~ZSM~009 governance.
\textbf{Policy trigger:} EU AI security standards (adversarial robustness
profiles, autonomous agent accountability frameworks) are adopted or formally
referenced in at least one ISO/IEC or ITU-T international standard, confirming
EU regulatory leadership beyond domestic jurisdiction.
\textbf{Operational trigger:} APT dwell time $\leq$24 hours is documented in
a named NIS2-essential entity operational deployment (not a simulation),
attributable to autonomous AI-driven detection and containment.
\textbf{Market trigger:} EU-sovereign AI security tools hold $\geq$50\% market
share in new critical infrastructure SOC deployments across at least five
Member States, displacing US-platform SIEMs as the procurement default in
regulated sectors.}
\\
\hline

% ---- PHASE 3 ----
\rowcolor{phase3row}
\textbf{Phase~3}
\newline\textit{Autonomy \&}
\newline\textit{Leadership}
\newline\textit{2035--2040}
\newline\sslbadge{ssltealbg}{ssltealtx}{SSL: Global Leadership in}
\newline\sslbadge{ssltealbg}{ssltealtx}{Trustworthy AI Security}
&
\cellhead{Cluster ODP: Ubiquitous AI Immune System}
All three sub-dimensions reach Ubiquitous Operation. Swarm-based
agentic AI defends the EU Cloud-Edge Continuum via self-evolving,
decentralised threat intelligence without any central coordination
dependency. Moving Target Defense continuously shuffles IP addressing,
network topology, and software diversity; validated APT reconnaissance
success rates fall below 5\% in certified operational environments.
Adversarial robustness validation is continuous and fully automated:
no security AI agent operates in EU critical infrastructure without
an active conformity certification. Phase~3 R\&I addresses
post-swarm frontiers: quantum-native AI security architectures,
brain-inspired neuromorphic threat detection, and AI-to-AI
adversarial dynamics beyond current threat models.
\kpistrip{Phase~3 ODP evidence requires operational statistics from
production deployments: APT dwell time figures, MTD effectiveness
metrics, and autonomous containment accuracy rates from certified
environments --- not from cyber range simulations.}
&
\cellhead{EU AI security governance as international reference}
EU AI Act frameworks for trustworthy autonomous security agents are
adopted internationally and referenced in bilateral digital
partnership agreements. Autonomous agent certification operates at
continental scale without per-case regulatory review. Zero-touch
network management under ETSI GS~ZSM~009 governs all EU NaaS
operations. The EU AI Office and ENISA jointly define the international
standard for security AI accountability. EU AI security governance
has transitioned from reactive adaptation to agenda-setting: other
jurisdictions adopt EU frameworks to access EU market and partnership
programmes.
\kpistrip{Phase~3 policy leadership is verified by international
adoption events --- not by the number of EU instruments in force.
At least one non-EU jurisdiction formally adopting EU security AI
governance standards constitutes the Phase~3 policy evidence threshold.}
&
\cellhead{Autonomous continental-scale AI cyber defence as operational default}
APT dwell time $\leq$24 hours in all certified EU NIS2-essential
critical infrastructure deployments. MTD is deployed as a mandatory
operational baseline. Swarm AI coordination operates without central
cloud dependency: each EU Cloud-Edge node participates in a
self-organising threat intelligence mesh. Deception infrastructure
is fully integrated into the autonomous response loop, providing
continuous model calibration ground truth. Human intervention is
reserved for exceptional legal escalation scenarios; routine security
operations are fully autonomous within the governance framework.
\kpistrip{The Phase~3 deployment test: is autonomous AI defence
the unconditional operational baseline, or is it still a premium
programme? 100\% of new critical infrastructure deployments
incorporating autonomous AI security constitutes the evidence threshold.}
&
\cellhead{SSL: Global Leadership in Trustworthy AI Security}
EU-sovereign AI security tools hold $\geq$50\% of new critical
infrastructure SOC deployments across the EU. US-platform SIEMs
are legacy infrastructure under active replacement programmes with
defined timelines. EU AI security models are exported to and
commercially adopted by partner jurisdictions. EU sets the commercial
terms for security AI market access: certification against EU AI~Act
Article~9/15 profiles is a de facto global procurement requirement.
The EU has transitioned from defensive sovereignty to offensive
strategic posture: it defines the technical terms on which the
global security AI industry must compete.
\kpistrip{Full sovereign leadership in this cluster requires both
domestic dominance ($\geq$50\% critical infra SOC market share)
and international influence (at least two partner jurisdictions
formally adopting EU security AI governance). Both conditions must
be met.}
\\
\hline

\end{longtable}
\end{landscape}

% ---- PROGRESSION MATRIX: CLUSTER 2 ----

% \noindent
% \resizebox{\textwidth}{!}{%
% \begin{tabular}{|p{4.2cm}|p{5.6cm}|p{5.6cm}|p{5.6cm}|}
% \hline
% \rowcolor{gray!30}
% \multicolumn{4}{|c|}{\textbf{Progression Matrix: Priority Cluster 2 ---
%   Cognitive Security and Autonomous Cyber Defence (P6, P10, AI-as-Target)}} \\
% \hline
% \rowcolor{gray!15}
% \textbf{Metric} &
% \textbf{Phase 1: Foundation \& Compliance (2025--2030)} &
% \textbf{Phase 2: Federation \& Standardisation (2030--2035)} &
% \textbf{Phase 3: Autonomy \& Leadership (2035--2040)} \\
% \hline
% \rowcolor{gray!10}
% \multicolumn{4}{|l|}{\textit{\textbf{Cluster-Level Progression Indicators}}} \\
% \hline
% \textbf{Cluster ODP} &
% Applied R\&D / Siloed AI Copilots &
% Federated Edge Agents &
% Ubiquitous AI Immune System \\
% \hline
% \textbf{Cluster SSL} &
% High Reliance on US Foundation Models &
% Sovereign European Edge LLMs / SLMs &
% Global Leadership in Trustworthy AI Security \\
% \hline
% \rowcolor{gray!10}
% \multicolumn{4}{|l|}{\textit{\textbf{Sub-Priority Decomposition (Individual ODP / SSL per Phase)}}} \\
% \hline
% \textit{P6: Cognitive SecOps \newline (LLMSecOps)} &
% ODP: Siloed AI Copilots \newline
% SSL: High Reliance on US Foundation Models &
% ODP: Federated Edge Agents \newline
% SSL: Sovereign EU Edge LLMs Emerging &
% ODP: Ubiquitous AI Immune System \newline
% SSL: Global Leadership in Trustworthy AI \\
% \hline
% \textit{P10: Proactive APT \newline Management} &
% ODP: Multi-Layered Alert \newline Correlation \& Hunting \newline
% SSL: High Dependency on Global Threat Feeds &
% ODP: Automated Triage \& \newline LLMSecOps Containment \newline
% SSL: Sovereign Pan-EU CTI \& Behavioural Analytics &
% ODP: Ubiquitous Moving \newline Target Defense (MTD) \newline
% SSL: Full Autonomous Counter-Adversarial Resilience \\
% \hline
% \textit{AI-as-Target \newline (CL1.2, CL1.3, CL4.7)} &
% ODP: Applied R\&D \newline
% SSL: High Foreign Dependency &
% ODP: Standardised Integration \newline
% SSL: Emerging EU Ecosystem &
% ODP: Ubiquitous Operation \newline
% SSL: Full Sovereign Autonomy \\
% \hline
% \rowcolor{blue!7}
% \textbf{Operational State} &
% LLMs assist human SOC analysts in alert triage and threat report
% synthesis; NWDAF telemetry is routed to centralised AI for anomaly
% detection in 5G standalone deployments; APT dwell times are
% reduced via ML-assisted IoC correlation in pilot enterprise environments;
% adversarial robustness validation of AI security agents is conducted
% in academic and research settings without operational certification
% frameworks in force. &
% Edge-native EU-sovereign SLMs execute autonomous Zero-Trust
% reconfigurations and traffic throttling in certified critical
% infrastructure deployments; LLMSecOps agents ingest CTI and
% track lateral movement against MITRE ATT\&CK behavioural
% frameworks in real time; adversarial robustness validation
% is standardised under AI Act Article~9 conformity assessment;
% APT dwell time reduced to under 72 hours in EU critical
% infrastructure pilot deployments. &
% Swarm-based agentic AI defends the EU Cloud-Edge Continuum
% via self-evolving, decentralised threat intelligence without
% central coordination dependency; Moving Target Defense
% continuously shuffles network topology, IP addressing, and
% software diversity, achieving validated APT reconnaissance
% success rates below 5\% in operational simulation environments;
% EU AI security models are formally adopted as reference
% implementations by partner jurisdictions. \\
% \hline
% \rowcolor{orange!12}
% \textbf{Critical Enabler} &
% AI Act Articles~9 and~15 enforcement (2025--2026) establishing
% mandatory risk management and adversarial robustness
% requirements for autonomous security AI, creating the
% regulatory framework for certified Phase~2 deployment &
% Standardised AI Act conformity assessment framework for
% autonomous security agents, enabling certified operational
% deployment without per-case regulatory approval and
% unlocking Phase~3 swarm-level autonomy &
% EU-certified self-evolving AI swarm protocols achieving
% full operational deployment under ETSI GS~ZSM~009
% zero-touch automation governance with Pan-EU
% CTI-sharing infrastructure \\
% \hline
% \end{tabular}%
% }

% =========================================================================
% PRIORITY CLUSTER 3: FEDERATED OPERATIONAL SOVEREIGNTY
% =========================================================================

\newpage
\section{Priority Cluster 3: Federated Operational Sovereignty}
\label{sec:cluster3}

\noindent\textbf{Constituent Sub-Priorities:} P4 (Security Interoperability via Simpl
Middleware), P5 (Supply Chain Verification and Secured Infrastructure), P7 (Pan-European
NaaS Federation --- ``Schengen for Cloud-Edge/AI'').

\medskip
\noindent\textbf{Cluster Thesis:} Technical excellence in Cluster~1 and autonomous
defence capability in Cluster~2 deliver no strategic value if they cannot operate
verifiably across borders, across providers, and across supply chains. This cluster
defines the governance fabric, interoperability standards, and supply chain integrity
mechanisms that transform individually sovereign components into a collectively sovereign
European computing continuum. Critically, it addresses the legal and semantic as well as
the technical interoperability problem --- without a ``Schengen for Cloud'' legal framework,
cross-border compute federation remains commercially fragile regardless of technical
standardisation.

\subsection*{7-Dimensional Analysis: Priority Cluster 3}

\begin{enumerate}

% ---- DIMENSION 1 ----
\item \textbf{Baseline \& Target Objective}

\textit{Baseline (2025--2026):}
European cloud and edge infrastructure remains highly fragmented across proprietary provider
silos with incompatible security policies, authentication mechanisms, and data exchange
protocols. Common European Data Spaces (Manufacturing, Health, Energy, Agriculture) are in
early piloting stages, with security interoperability relying on bilateral agreements rather
than standardised automated mechanisms aligned with the Simpl open-source middleware
stack~\cite{cloud_edge_roadmap}. Supply chain verification is largely static: CRA-mandated
Software Bills of Materials (SBOMs)~\cite{cra} are point-in-time declarations without runtime
verification; Digital Product Passport (DPP) frameworks remain in pilot stage under the EU
Ecodesign Regulation. NaaS federation across EU telecom operators does not exist at
operational scale --- cross-border edge compute sharing requires bespoke commercial agreements.
CAMARA Network APIs are being standardised by GSMA but remain fragmentedly implemented
across European operators.

Current cluster ODP: Fragmented Piloting (P4, P7); early supply chain disclosure phase
(P5). Cluster SSL: Emerging EU Ecosystem (P4); Hyperscaler Platform Dominance (P7);
High reliance on fragmented non-EU supply chain auditing (P5).

\textit{Target (2040):}
A unified European NaaS platform enables seamless cross-border compute roaming under a
``Schengen for Cloud'' legal framework, legally operable as a single jurisdiction for
data workloads. All Common European Data Spaces operate on Simpl-based security
interoperability with automated MIM compliance --- provider-agnostic, cryptographically
verifiable, and free of vendor lock-in. Runtime SBOMs with Zero-Knowledge Proof (ZKP)
verification eliminate unverified components from EU critical infrastructure supply chains.
EU achieves full sovereign control over supply chain transparency with pan-European Cyber
Threat Intelligence sharing for supply chain compromise indicators.

% ---- DIMENSION 2 ----
\item \textbf{Primary Drivers}

\textit{Threat Drivers:}
\begin{itemize}
  \item Supply chain attacks exploiting the disaggregation complexity of Open RAN
        deployments, where a single compromised firmware component sourced from one
        of 50-plus hardware and software vendors per macro cell deployment can affect
        thousands of base stations across multiple operators before detection.
  \item ``Weakest link'' cascading failures in federated NaaS environments: a
        single provider with a compromised Minimal Interoperability Mechanism (MIM)
        implementation can propagate malicious policies across the entire federated
        fabric before semantic validation detects the enforcement inconsistency.
  \item US CLOUD Act extra-territorial reach over non-EU cloud providers processing
        EU data creates a structural sovereignty risk that persists regardless of
        technical security measures if the legal framework enabling EU-sovereign
        federated compute is not established.
  \item Silent policy translation failures: automated security policy translation
        between heterogeneous provider MIM implementations may produce semantically
        equivalent but enforcement-different policies, creating exploitable gaps
        invisible to standard compliance auditing procedures.
\end{itemize}

\textit{Regulatory Drivers:}
\begin{itemize}
  \item Data Act interoperability obligations (Articles~23--31): mandatory
        requirements for cloud service providers covering security interface
        standardisation and data portability~\cite{nis2}.
        [CITATION REQUIRES VERIFICATION --- confirm final Data Act article numbering
        for interoperability provisions.]
  \item NIS2 Article~21(2)(h): supply chain security measures mandatory for essential
        entities, requiring documented risk assessment of all ICT suppliers~\cite{nis2}.
  \item CRA Annex~I Part~II: supply chain vulnerability handling requirements for
        products with digital elements, including SBOM disclosure
        obligations~\cite{cra}.
  \item DORA Article~28: ICT third-party risk management requiring continuous
        monitoring and documented exit strategies for critical ICT
        suppliers~\cite{dora}.
\end{itemize}

% ---- DIMENSION 3 ----
\item \textbf{High Impact Activities}

\begin{enumerate}[label=(\alph*)]
  \item \textbf{[Phase~1]} Define and standardise security-specific Minimal
        Interoperability Mechanisms (MIMs) within the Simpl open-source middleware
        stack~\cite{cloud_edge_roadmap}; deploy pilots in three Common European Data
        Spaces (Manufacturing, Health, Energy) with automated cross-provider policy
        enforcement, federated identity attestation, and cryptographically verifiable
        compliance reporting.

  \item \textbf{[Phase~1--2]} Implement Runtime SBOMs with cryptographic proof of
        execution for all containerised microservices and AI components migrating across
        the multi-provider continuum; integrate Digital Product Passports (DPP) for
        hardware component traceability from silicon fabrication to operational
        deployment, aligned with CRA SBOM obligations~\cite{cra}.

  \item \textbf{[Phase~1--2]} Deploy harmonised CAMARA Network APIs across a minimum
        of five EU telecom operators establishing standardised interfaces for cross-border
        edge compute roaming; test Automated Compliance Negotiators in designated EU
        5G cross-border corridor deployments under IPCEI-CIS governance.

  \item \textbf{[Phase~2]} Establish ``Data Visa'' mechanisms enabling dynamic
        cross-border AI model and data migration without manual legal friction, governed
        by a common semantic security ontology (CL3.4) aligned with STIX~2.1/TAXII
        threat intelligence standards and automated EUCS compliance verification~\cite{eucs}.

  \item \textbf{[Phase~2--3]} Deploy Zero-Knowledge Proofs (ZKPs) for IP-protected
        supply chain component verification: hardware vendors prove FIPS~203/204/205
        algorithm compliance and absence of known vulnerabilities without revealing
        proprietary design details~\cite{fips203,fips204,fips205}; operationalise
        pan-European supply chain compromise indicator sharing via standardised
        CTI lifecycle platforms (CL2.5).
\end{enumerate}

% ---- DIMENSION 4 ----
\item \textbf{Catalysts \& Enablers}

\textit{Funding Alignment:}
IPCEI-CIS (Important Project of Common European Interest on Next Generation Cloud
Infrastructure and Services) as the primary funding vehicle for cross-border NaaS federation
pilots; Digital Europe Programme for Simpl development, Common Data Space pilots, and
MIM standardisation infrastructure; ECCC Strategic Agenda priority areas ``Supply Chain
Security'' and ``Secure Interoperability of Systems and Networks''~\cite{eccc}; European
Alliance for Industrial Data, Edge and Cloud for governance and pilot ecosystem
coordination~\cite{cloud_edge_roadmap}.

\textit{Standardisation:}
CEN/CENELEC JTC~13 for security MIM standardisation within the Simpl ecosystem;
GSMA CAMARA API Alliance for harmonised Network API standardisation across operators;
ETSI TC CYBER for supply chain security standards; OASIS STIX~2.1/TAXII for threat
intelligence interoperability; W3C Verifiable Credentials Data Model for SSI-based
supply chain attestation chains.

\textit{Ecosystem Enablers:}
Gaia-X Federation Services trust framework alignment with Simpl MIM security policies;
Digital Product Passport pilot programmes under EU Ecodesign Regulation~2024;
ENISA supply chain risk management guidelines and threat landscape reports for ICT
sectors; dedicated SME compliance support programmes for CRA-mandated SBOM tooling,
recognising that sub-50 person software vendors represent the most vulnerable link in
the supply chain verification architecture~\cite{cra}.

% ---- DIMENSION 5 ----
\item \textbf{Disruptors \& Systemic Risks}

\begin{itemize}
  \item \textbf{Hyperscaler resistance to Simpl MIM adoption:} AWS, Azure, and
        GCP collectively control approximately 65\% of the EU cloud market as of
        2024 and have strong commercial incentives to maintain proprietary APIs;
        without mandatory EU Cloud Rulebook provisions requiring MIM compliance,
        voluntary adoption rates are unlikely to achieve genuine interoperability
        within Phase~1 timelines.
  \item \textbf{IPCEI-CIS implementation delays:} complex EU state aid notification
        procedures have historically delayed IPCEI projects by 18--36 months
        (IPCEI-H2 delayed by state aid clearance processes); P7 NaaS federation
        Phase~1 milestones carry material schedule risk if IPCEI-CIS governance
        does not streamline implementation mechanisms.
  \item \textbf{Semantic interoperability gap:} automated security policy translation
        between heterogeneous MIM implementations is architecturally harder than
        syntactic interoperability --- subtle semantic differences in ``deny'' versus
        ``drop'' versus ``quarantine'' semantics across provider implementations
        can create silent enforcement gaps undetectable through standard compliance
        auditing.
  \item \textbf{ZKP computational overhead:} current ZKP schemes for supply chain
        verification carry 10--100$\times$ computational overhead relative to
        standard cryptographic verification, limiting practical deployment to
        batch verification rather than real-time per-transaction attestation
        in Phases~1--2.
  \item \textbf{EU semiconductor supply chain concentration:} advanced node
        fabrication (sub-7\,nm) required for cutting-edge supply chain attestation
        hardware remains concentrated in Taiwan (TSMC) and South Korea (Samsung);
        a Taiwan Strait contingency would immediately constrain EU supply chain
        verification capability at the silicon level, exposing a gap the EU Chips Act
        does not address within the 2030 timeframe.
\end{itemize}

% ---- DIMENSION 6 ----
\item \textbf{Failsafe Mechanisms \& Resilience}

\begin{itemize}
  \item \textbf{Autonomous unfederation:} if a provider fails Runtime SBOM
        verification or MIM compliance checks, eBPF-based zero-trust policy
        enforcement severs that provider from the federated NaaS fabric and redirects
        workloads to verified compliant nodes within a sub-60-second SLA; provider
        re-admission requires fresh ZKP-verified attestation, not manual
        recertification.
  \item \textbf{Graceful national fallback:} cross-border SLA failure or ``Data
        Visa'' mechanism collapse triggers workload redirect to national infrastructure
        maintained at sufficient sovereign capacity for essential services, fulfilling
        NIS2 Article~21(2)(c) business continuity obligations and DORA Article~9 ICT
        resilience requirements~\cite{nis2,dora}.
  \item \textbf{Cryptographic attestation of MIM enforcement:} all Simpl MIM security
        controls include verifiable proofs that policies are enforced as declared --- not
        merely configured --- using TEE-backed attestation integrated with Cluster~1
        secure element infrastructure~\cite{cloud_edge_roadmap}, eliminating
        compliance-theatre scenarios where declared and actual enforcement diverge.
  \item \textbf{Supply chain quarantine:} unverified or high-risk components are
        isolated into quarantined network slices (VLAN/VxLAN-isolated segments with
        eBPF-enforced east-west traffic blocking) pending ZKP verification;
        quarantined components cannot interact with production workloads or access
        sensitive data paths.
  \item \textbf{Multi-jurisdictional residency fallback:} if cross-border ``Data
        Visa'' mechanisms fail legally or technically, automated GDPR Article~44
        compliance enforcement pins workloads to the originating jurisdiction,
        preventing unauthorised cross-border data transfers under federation failure
        conditions~\cite{gdpr}.
\end{itemize}

% ---- DIMENSION 7 ----
\item \textbf{Transversal Meta-Layer}

\textit{Sovereignty \& Autonomy:}
This cluster breaks non-EU hyperscaler API lock-in (AWS, Azure, GCP controlling
approximately 65\% of EU cloud market as of 2024) by enabling European providers to
compete on standardised Simpl-based infrastructure. A ``Schengen for Cloud'' legal
framework eliminates US CLOUD Act extra-territorial jurisdiction exposure for EU data
processed on federated EU-sovereign NaaS infrastructure. Runtime SBOM verification and
ZKP-based supply chain attestation reduce EU dependency on non-EU vendor self-certification
declarations --- the dominant model under current CRA-baseline SBOM obligations.

\textit{Trust \& Ethics:}
GDPR Article~25 (privacy by design and by default) is operationalised through
MIM-enforced data minimisation policies that are cryptographically verifiable, not
declarative~\cite{gdpr}. Data Act interoperability obligations are fulfilled through
Simpl MIM compliance, providing auditable trust chains for GDPR Article~5(2) data
processing accountability. Transparent Runtime SBOM verification provides the
supply chain integrity that GDPR Article~25 implicitly requires for processors
handling EU personal data at scale~\cite{gdpr}.

\textit{Sustainable Security (Green IT):}
Federated NaaS resource pooling across EU operators eliminates redundant edge compute
over-provisioning by individual operators, reducing total EU edge infrastructure carbon
footprint through shared utilisation efficiency. DPP-based supply chain traceability
aligns with EU Ecodesign Regulation circular economy requirements for ICT hardware,
enabling component reuse lifecycle tracking that reduces electronic waste from
premature infrastructure replacement driven by opacity in component provenance.

\end{enumerate}


% =========================================================================
% CLUSTER 3: FEDERATED OPERATIONAL SOVEREIGNTY
% =========================================================================
\newpage
\subsection{Progression Matrix: Priority Cluster 3 --- Federated Operational Sovereignty}

\begin{landscape}
\setlength{\textwidth}{26.5cm}\setlength{\linewidth}{26.5cm}%
\setlength{\hsize}{26.5cm}\setlength{\columnwidth}{26.5cm}
\setlength{\tabcolsep}{5pt}
\renewcommand{\arraystretch}{1.35}

\begin{longtable}{%
|>{\footnotesize\raggedright\arraybackslash}p{2.8cm}%
|>{\footnotesize\raggedright\arraybackslash}p{5.4cm}%
|>{\footnotesize\raggedright\arraybackslash}p{5.4cm}%
|>{\footnotesize\raggedright\arraybackslash}p{5.4cm}%
|>{\footnotesize\raggedright\arraybackslash}p{5.4cm}|}

\caption{\textbf{Cluster~3: Federated Operational Sovereignty} (P4, P5, P7).
Transition matrix across the four methodology dimensions. P4 (Simpl
Interoperability), P5 (Supply Chain Verification), and P7 (NaaS Federation)
are labelled where their maturity trajectories diverge within a cell.}
\label{tab:cluster3}\\

\hline
\rowcolor{colheader}
\textbf{Phase / SSL} &
\textbf{R\&I Maturity}{\scriptsize ODP indicator per sub-priority} &
\textbf{Policy, Regulatory \& Standards}{\scriptsize enforcement $\cdot$
  certification $\cdot$ standardisation} &
\textbf{Operational Deployment \& Adoption}{\scriptsize actor types $\cdot$
  scale $\cdot$ sector coverage} &
\textbf{Market Uptake \& Strategic Sovereignty}{\scriptsize SSL indicator
  $\cdot$ dependency reduction $\cdot$ industrial capacity} \\
\hline
\endfirsthead

\hline
\rowcolor{colheader}
\textbf{Phase / SSL} &
\textbf{R\&I Maturity} &
\textbf{Policy, Regulatory \& Standards} &
\textbf{Operational Deployment \& Adoption} &
\textbf{Market Uptake \& Strategic Sovereignty} \\
\hline
\endhead
\hline\multicolumn{5}{r}{\footnotesize\textit{Continued\ldots}}\\ \endfoot
\hline \endlastfoot

% ---- PHASE 1 ----
\rowcolor{phase1row}
\textbf{Phase~1}
\newline\textit{Foundation \&}
\newline\textit{Compliance}
\newline\textit{2025--2030}
\newline\sslbadge{sslredbg}{sslredtx}{SSL: Fragmented; Hyperscaler}
\newline\sslbadge{sslredbg}{sslredtx}{Dominance (P7); Emerging}
\newline\sslbadge{sslredbg}{sslredtx}{EU Ecosystem (P4)}
&
\cellhead{Cluster ODP: Fragmented Piloting (P4, P7) / Disclosure Phase (P5)}
\subp{P4} is at Fragmented Piloting: Simpl security MIM pilots are
launched in three Common European Data Spaces (Manufacturing, Health,
Energy) with manual policy coordination and proof-of-concept
implementations. \subp{P7} is at Bilateral Interconnection: harmonised
CAMARA Network APIs are deployed in 2--3 bilateral 5G corridor agreements
between individual EU telecom operators; cross-border edge compute
roaming does not exist at operational scale. \subp{P5} is at Supply
Chain Disclosure: CRA-mandated static SBOM declarations are being
submitted as point-in-time compliance artefacts without runtime
verification; Digital Product Passport pilot programmes are initiated
under the EU Ecodesign Regulation; Zero-Knowledge Proof verification
for supply chain attestation is at Applied R\&D stage.
\kpistrip{The three sub-priorities have meaningfully different Phase~1
maturity levels. P5 is the least mature and the most consequential
shared dependency: Runtime SBOM verification is a gate condition for
both P4 (MIM compliance trust) and Cluster~1 (TEE integrity claims).}
&
\cellhead{Multiple legislative obligations activating; standards nascent}
CRA~SBOM obligations (Annex~I Part~II) enter enforcement in 2027,
mandating static SBOM disclosure for all products with digital elements
entering the EU market. NIS2~Article~21(2)(h) makes supply chain risk
assessment mandatory for essential entities. DORA~Article~28 requires
continuous ICT third-party risk monitoring for financial sector entities.
Data Act interoperability obligations activate, creating procurement-level
pressure for cross-provider security interface standardisation. The
Simpl-Open official launch with initial MIM definitions constitutes the
Phase~1 technical policy trigger for P4. IPCEI-CIS enters its
deployment phase. GSMA CAMARA API Alliance standardisation work is
active but not yet mandated in EU regulatory instruments.
\kpistrip{Phase~1 policy activity is high-volume but fragmented:
multiple legislative instruments activate simultaneously without
a unified interoperability mandate that would create coherent market
pull for Simpl-based standardisation.}
&
\cellhead{Deployment: bilateral pilots and Data Space proof-of-concepts}
Three Common European Data Space security pilots operate Simpl-based
MIM proof-of-concepts with manual policy coordination; cross-provider
security policy enforcement is not automated. CAMARA API deployments
are confined to 2--3 bilateral 5G corridor agreements between single
operator pairs. Static SBOM declarations are submitted as CRA
compliance artefacts but are not verified at runtime; supply chain
risk assessment is vendor-declared rather than independently attested.
Hyperscalers (AWS, Azure, GCP) retain $\approx$65\% of EU cloud market
through proprietary API ecosystems. Cross-border AI model migration
requires bespoke bilateral legal review per transfer event. SME-accessible
SBOM tooling and MIM compliance support are not yet standardised or
commercially affordable.
\kpistrip{Phase~1 deployment is bilateral and pilot-scale. The gap
between proof-of-concept MIM implementations and automated enforcement
across the majority of Data Spaces defines the core operational
challenge for Phase~2.}
&
\cellhead{SSL: Mixed --- Emerging EU Ecosystem (P4) / Hyperscaler Dominance (P7)}
\subp{P4} has reached Emerging EU Ecosystem: the Simpl open-source
stack exists and is being piloted with EU institutional backing.
\subp{P7} remains at Hyperscaler Platform Dominance: AWS, Azure, and GCP
collectively control $\approx$65\% of EU cloud market; no EU NaaS
federation operates at commercially relevant scale; cross-border edge
compute is fragmented into national silos. \subp{P5} is at High
Reliance on Fragmented Non-EU Auditing: supply chain verification
is vendor-declared under US-dominated certification ecosystems (Common
Criteria national schemes, SOC2). The cluster's dominant SSL profile
is foreign dependency, driven by the weight of P7 market structure.
\kpistrip{P4 and P7 require separate SSL tracking in Phase~1: conflating
them into a single cluster SSL figure obscures that Simpl has achieved
Emerging Ecosystem while NaaS federation has not.}
\\
\hline

% ---- GATE 1->2 ----
\rowcolor{gaterow}
\gatelbl{1 $\rightarrow$ 2} &
\multicolumn{4}{>{\footnotesize\raggedright\arraybackslash}p{22.0cm}|}{%
\cellcolor{gaterow}%
\textbf{R\&I trigger:} At least one Common European Data Space has deployed
automated (not manually coordinated) Simpl MIM-based cross-provider security
policy enforcement, with documented evidence that enforcement is cryptographically
verifiable rather than declaratively asserted.
\textbf{Policy trigger:} Simpl security MIM definitions are published in
candidate CEN/CENELEC standard form and referenced as a compliance pathway in
at least one EU sector-specific NIS2 or Data Act implementing act.
\textbf{Operational trigger:} CAMARA-based cross-border edge compute roaming is
operational between at least three distinct EU telecom operator pairs (not a single
bilateral agreement), with documented SLA performance and compliance evidence.
\textbf{Market trigger:} Runtime SBOM tooling with cryptographic execution proofs
is available as a commercially supported product or service from at least two
EU-domiciled vendors, establishing the market infrastructure for Phase~2
mandatory adoption.}
\\
\hline

% ---- PHASE 2 ----
\rowcolor{phase2row}
\textbf{Phase~2}
\newline\textit{Applied \& Industrial}
\newline\textit{R\&D alongside}
\newline\textit{Federation \&}
\newline\textit{Standardisation}
\newline\textit{2030--2035}
\newline\sslbadge{ssltransbg}{ssltrantx}{SSL: Growing Pan-EU}
\newline\sslbadge{ssltransbg}{ssltrantx}{Sovereignty}
&
\cellhead{Cluster ODP: Standardised Integration (P4) / Federated Edge Roaming (P7) / Dynamic Runtime Verification (P5)}
\subp{P4} reaches Standardised Integration: automated policy translation
via standardised Simpl MIMs is deployed across the majority of active
Common European Data Spaces; federated identity attestation and
cross-provider compliance reporting are automated. \subp{P7} reaches
Federated Edge Roaming: CAMARA-based cross-border roaming operates
across $\geq$5 EU operator pairs; ``Data Visa'' mechanisms enable
dynamic AI model migration without manual legal review in certified
corridors. \subp{P5} transitions from disclosure to dynamic runtime
verification: Runtime SBOMs with cryptographic proof-of-execution
verify containerised components as they migrate across multi-provider
nodes; ZKP batch verification of supply chain attestation is operationally
validated; Digital Product Passports are mandatory under the EU
Ecodesign Regulation.
\kpistrip{Phase~2 R\&I activity includes industrial R\&D for ZKP
real-time verification (currently batch-only due to 10--100$\times$
overhead), semantic security ontology standardisation, and legal
engineering of the ``Schengen for Cloud'' framework --- not solely
standardisation of existing prototypes.}
&
\cellhead{Security MIMs as CEN/CENELEC standards: the Phase~2 gate}
The single Phase~2 policy gate is the adoption of Simpl security MIMs
as ratified CEN/CENELEC standards, referenced as mandatory compliance
requirements in the EU Cloud Rulebook. This transforms MIM adoption
from voluntary best-practice to market-access obligation for providers
serving EU-regulated sectors. DPP becomes mandatory under the EU
Ecodesign Regulation~2024 for ICT hardware categories. Data Act full
enforcement creates binding procurement-level interoperability
obligations. DORA continuous ICT third-party monitoring is fully
operationalised in the financial sector. IPCEI-CIS cross-border NaaS
governance framework is operational.
\kpistrip{The Phase~2 policy gate is the CEN/CENELEC ratification of
security MIMs. Voluntary Simpl adoption without this standards event
does not constitute Phase~2 policy maturity.}
&
\cellhead{Cross-sector rollout; SME participation; Runtime SBOM in production}
Simpl MIMs are deployed across the majority of active Common European
Data Spaces with automated cross-provider enforcement. Runtime SBOMs
with cryptographic execution proofs verify supply chain components in
critical infrastructure deployments in at least five Member States.
CAMARA cross-border roaming operates across $\geq$5 EU operator pairs
with documented SLA performance. ``Data Visa'' mechanisms are
operational in certified cross-border AI model migration corridors.
SME-accessible Simpl MIM compliance tooling and Runtime SBOM services
are commercially available from EU-domiciled providers. Unverified
supply chain components are automatically quarantined in isolated
network slices.
\kpistrip{Phase~2 deployment must show geographic breadth ($\geq$5 Member
States), cross-sector presence ($\geq$3 Data Spaces under automated MIM
enforcement), and SME participation in at least one deployment context.
Single-sector rollout does not constitute Phase~2 operational maturity.}
&
\cellhead{SSL: Growing Pan-EU Sovereignty}
EU Data Space infrastructure operating under Simpl achieves $\geq$30\%
non-hyperscaler share in new regulated-sector cloud procurement. Runtime
SBOM ecosystem is commercially mature with multiple competing EU-domiciled
vendors. EU NaaS consortium (under IPCEI-CIS) operates cross-border
compute at pilot commercial scale, demonstrating that EU providers
can federate to compete with hyperscaler geographic coverage. DPP
traceability reduces procurement risk for EU critical infrastructure
buyers by making supply chain provenance verifiable. Foreign hyperscaler
market share in EU regulated sectors declines for the first time from
its $\approx$65\% Phase~1 baseline.
\kpistrip{Phase~2 SSL evidence: $\geq$30\% non-hyperscaler share in new
regulated procurement; first documented market share decline for
non-EU hyperscalers in EU regulated cloud procurement.}
\\
\hline

% ---- GATE 2->3 ----
\rowcolor{gaterow}
\gatelbl{2 $\rightarrow$ 3} &
\multicolumn{4}{>{\footnotesize\raggedright\arraybackslash}p{22.0cm}|}{%
\cellcolor{gaterow}%
\textbf{R\&I trigger:} ZKP verification for supply chain attestation has
achieved real-time per-transaction performance (eliminating the current
batch-only constraint) in at least one production critical infrastructure
deployment, confirmed by independent performance benchmarking.
\textbf{Policy trigger:} A ``Schengen for Cloud'' legal framework --- or
equivalent EU legal instrument resolving US CLOUD Act extra-territorial
jurisdiction conflict for federated EU-sovereign NaaS workloads --- has
been enacted, providing the jurisdictional basis for single-market operation
of federated compute.
\textbf{Operational trigger:} The EU NaaS federation (under IPCEI-CIS or its
successor) operates at commercial scale across $\geq$10 EU telecom operator
members, offering cross-border edge compute services as a standard commercial
product with documented customer deployments outside its founding consortium.
\textbf{Market trigger:} EU supply chain governance frameworks (Runtime SBOM
plus ZKP attestation) are adopted or formally referenced by at least two
non-EU partner jurisdictions as part of digital trade or partnership agreements,
confirming that EU governance leadership is internationally exportable.}
\\
\hline

% ---- PHASE 3 ----
\rowcolor{phase3row}
\textbf{Phase~3}
\newline\textit{Autonomy \&}
\newline\textit{Leadership}
\newline\textit{2035--2040}
\newline\sslbadge{ssltealbg}{ssltealtx}{SSL: Full Federated}
\newline\sslbadge{ssltealbg}{ssltealtx}{Sovereign Autonomy}
&
\cellhead{Cluster ODP: Ubiquitous Operation across all sub-priorities}
All three sub-priorities reach Ubiquitous Operation. Simpl-based
security is the universal standard for cross-provider Data Space
operation in the EU; no alternative interoperability model holds
market relevance in regulated sectors. ZKP-verified supply chain
attestation operates at real-time transaction scale for all critical
infrastructure component verification. The EU NaaS federation operates
as a single commercial entity legally equivalent to a unified
jurisdiction for data workloads. Phase~3 R\&I pursues next-generation
frontiers: post-ZKP cryptographic verification, semantic AI-to-AI
policy negotiation, and legal engineering of the next-generation
``Digital Single Market'' framework beyond cloud-edge compute.
\kpistrip{Phase~3 ODP evidence: Simpl MIM compliance rate in new regulated
EU cloud procurement; ZKP verification throughput in production critical
infrastructure deployments; EU NaaS federation revenue relative to
hyperscaler EU revenue in regulated sectors.}
&
\cellhead{EU federated governance as international model}
The ``Schengen for Cloud'' legal framework is enacted and operational,
resolving US CLOUD Act jurisdiction conflicts for federated EU-sovereign
workloads and enabling single-market legal operation of cross-border
compute. EU Cloud Rulebook is mandatory for all new cloud deployments
in EU regulated sectors without exception. EU supply chain governance
standards (Runtime SBOM plus ZKP) are referenced in international
supply chain security standards (ISO/IEC) and in digital trade
agreements with partner jurisdictions. The EU has transitioned from
rule-taker to rule-setter in international cloud and supply chain
governance.
\kpistrip{Phase~3 governance leadership is verified by two events:
enactment of the ``Schengen for Cloud'' framework (domestic) and
formal adoption of EU supply chain governance standards by $\geq$2
non-EU partner jurisdictions (international).}
&
\cellhead{Provider-agnostic EU trust fabric as continental operational default}
A single provider-agnostic EU trust fabric operates continuously across
all active Common European Data Spaces with real-time automated security
orchestration and zero manual policy coordination. ZKP-verified supply
chain attestation eliminates all unverified components from EU critical
infrastructure deployments in regulated sectors. The EU NaaS federation
is a standard commercial service operating at hyperscaler-comparable
scale, with EU operators collectively competitive on cost and coverage.
Security-by-Design is not an add-on: every new cloud-edge deployment
in the EU incorporates Simpl MIM compliance, Runtime SBOM verification,
and cross-border sovereignty assurance as architectural defaults.
\kpistrip{Phase~3 deployment normalisation: Simpl MIM compliance in
$\geq$90\% of new regulated EU cloud deployments; ZKP supply chain
verification in 100\% of new critical infrastructure component
procurement; EU NaaS federation operating as a single commercial entity
with documented non-founding-member customers.}
&
\cellhead{SSL: Full Federated Sovereign Autonomy}
Non-EU hyperscaler market share in EU regulated cloud procurement
falls below 40\% for new deployments (from $\approx$65\% at baseline),
with a binding trajectory toward continued reduction. EU NaaS
federation achieves hyperscaler-comparable geographic coverage and
commercial viability across the EU. Supply chain governance
internationally adopted: at least two partner jurisdictions have
formally adopted EU Runtime SBOM plus ZKP attestation standards as
domestic procurement requirements. The EU has removed the structural
dependency that allowed US CLOUD Act jurisdiction to apply to
EU-regulated data: federated sovereignty is legally guaranteed, not
only technically implemented.
\kpistrip{Full Federated Sovereign Autonomy requires three
concurrent conditions: $<$40\% hyperscaler new-procurement share;
``Schengen for Cloud'' legally enacted; $\geq$2 partner jurisdictions
formally adopting EU supply chain governance standards. Achieving
any two without the third does not constitute this SSL level.}
\\
\hline

\end{longtable}
\end{landscape}

% ---- PROGRESSION MATRIX: CLUSTER 3 ----


% \noindent
% \resizebox{\textwidth}{!}{%
% \begin{tabular}{|p{4.2cm}|p{5.6cm}|p{5.6cm}|p{5.6cm}|}
% \hline
% \rowcolor{gray!30}
% \multicolumn{4}{|c|}{\textbf{Progression Matrix: Priority Cluster 3 ---
%   Federated Operational Sovereignty (P4, P5, P7)}} \\
% \hline
% \rowcolor{gray!15}
% \textbf{Metric} &
% \textbf{Phase 1: Foundation \& Compliance (2025--2030)} &
% \textbf{Phase 2: Federation \& Standardisation (2030--2035)} &
% \textbf{Phase 3: Autonomy \& Leadership (2035--2040)} \\
% \hline
% \rowcolor{gray!10}
% \multicolumn{4}{|l|}{\textit{\textbf{Cluster-Level Progression Indicators}}} \\
% \hline
% \textbf{Cluster ODP} &
% Fragmented Piloting &
% Standardised Integration &
% Ubiquitous Operation \\
% \hline
% \textbf{Cluster SSL} &
% Emerging EU Ecosystem &
% Growing Pan-EU Sovereignty &
% Full Federated Sovereign Autonomy \\
% \hline
% \rowcolor{gray!10}
% \multicolumn{4}{|l|}{\textit{\textbf{Sub-Priority Decomposition (Individual ODP / SSL per Phase)}}} \\
% \hline
% \textit{P4: Security Interoperability \newline via Simpl Middleware} &
% ODP: Fragmented Piloting \newline
% SSL: Emerging EU Ecosystem &
% ODP: Standardised Integration \newline
% SSL: Full Sovereign Autonomy &
% ODP: Ubiquitous Operation \newline
% SSL: Full Sovereign Autonomy \\
% \hline
% \textit{P5: Supply Chain \newline Verification \& Infra.} &
% ODP: Full Supply Chain Disclosure \newline
% SSL: High Reliance on \newline Fragmented Non-EU Auditing &
% ODP: Partial Restriction \& \newline Dynamic Runtime Verification \newline
% SSL: Emerging EU Supply Chain \newline Tracking Frameworks &
% ODP: Full Restriction Process \& \newline ZKP-Verified Attestation \newline
% SSL: Full Sovereign Supply \newline Chain Control \& Transparency \\
% \hline
% \textit{P7: Pan-European NaaS \newline Federation} &
% ODP: Bilateral Interconnection \newline
% SSL: Hyperscaler Platform Dominance &
% ODP: Federated Edge Roaming \newline
% SSL: Pan-EU Consortium Competes &
% ODP: True Digital Single Market \newline
% SSL: European Scale Parity Achieved \\
% \hline
% \rowcolor{blue!7}
% \textbf{Operational State} &
% Initial Simpl security pilots are deployed in three Common European
% Data Spaces with manual policy coordination and proof-of-concept
% MIM implementations; CRA-baseline SBOM declarations are submitted
% by hardware and software vendors as static point-in-time documents
% without runtime verification; harmonised CAMARA APIs are tested in
% bilateral 5G corridor agreements between individual EU telecom
% operators; hyperscalers retain dominant market share through
% proprietary API ecosystems. &
% Automated policy translation via standardised Simpl MIMs is
% deployed across the majority of active Common European Data Spaces;
% Runtime SBOMs with cryptographic execution proofs verify
% containerised components dynamically as they migrate across
% multi-provider continuum nodes; harmonised CAMARA-based cross-border
% edge roaming operates across five or more EU operator pairs;
% ``Data Visa'' mechanisms enable AI model migration without
% manual legal review in certified corridor deployments. &
% A single provider-agnostic European trust fabric operates
% continuously across all EU Common Data Spaces with real-time
% automated security orchestration; ZKP-verified supply chain
% attestation eliminates unverified components from all
% EU critical infrastructure deployments; a unified NaaS
% platform legally operable as a single data jurisdiction
% enables EU operators to achieve hyperscaler-comparable
% scale; EU supply chain governance frameworks are adopted
% as reference implementations internationally. \\
% \hline
% \rowcolor{orange!12}
% \textbf{Critical Enabler} &
% Simpl-Open official launch with initial security MIM definitions
% published as open standards, establishing the technical
% baseline for cross-provider security policy automation &
% Security MIMs adopted as CEN/CENELEC standards with
% mandatory reference in EU Cloud Rulebook, creating the
% regulatory basis for non-voluntary hyperscaler compliance &
% ``Schengen for Cloud'' legal framework enacted, resolving
% US CLOUD Act extra-territorial jurisdiction conflict and
% enabling single-jurisdiction data workload operation at
% continental scale \\
% \hline
% \end{tabular}%
% }

% =========================================================================
% CROSS-PRIORITY BOTTLENECK ASSESSMENT
% =========================================================================

\newpage
\section{Cross-Priority Bottleneck Assessment}
\label{sec:bottlenecks}

Three shared structural dependencies determine whether individual sub-priority
progressions can advance on their stated timelines, regardless of cluster-level investment.
These bottlenecks are identified here because they constitute cross-cutting gate conditions
that cannot be resolved by any single cluster in isolation.

\textbf{Bottleneck 1 --- EU Chips Act Sovereign Fabrication Capacity.}
The Phase~2 gate conditions for P1 (EUCC `High' certified RISC-V secure elements),
P2 (PQC-native RISC-V silicon), P3 (zero-energy secure hardware), and P8 (RISC-V
Keystone TEE scaling) all converge on the same physical dependency: sovereign EU
fabrication capacity for security-critical silicon at commercially viable yields.
The Phase~2 critical enabler for P1 (``EUCC `High' certification of first
EU-manufactured RISC-V secure element'') and the Phase~3 enabler (``EU Chips Act
sovereign EUV-capable fabrication achieving $\geq$60\% self-sufficiency'') constitute
a single two-stage gate that controls the SSL trajectory of four of the five Cluster~1
sub-priorities simultaneously.

\textbf{Bottleneck 2 --- AI Act Conformity Assessment Infrastructure.}
The Phase~2 gate condition for Cluster~2 (``standardised AI Act conformity assessment
framework for autonomous security agents'') is a prerequisite not only for P6 and P10
but also for any AI-driven capability within Cluster~1 (autonomous Island Mode swarm
protocols) and Cluster~3 (Automated Compliance Negotiators in NaaS federation) that
incorporates autonomous decision-making. The absence of EU AI Act conformity assessment
bodies with sufficient capacity to certify high-risk security AI at scale represents a
systemic bottleneck that, if unresolved by 2028--2030, will delay the Phase~2 SSL
transition across all three clusters.

\textbf{Bottleneck 3 --- Runtime SBOM and Supply Chain Attestation Infrastructure.}
Runtime SBOM verification (P5) is a prerequisite for both the Cluster~1 TEE integrity
claim (P8 confidential computing pipelines cannot be trusted without verified component
provenance) and the Cluster~3 MIM compliance model (Simpl-based interoperability
cannot provide cryptographic enforcement guarantees if component integrity is
unverifiable). The Phase~1 critical enabler for Cluster~3 (Simpl-Open MIM definitions)
and the Phase~1 ODP for P5 (full supply chain disclosure) must therefore proceed in
coordination, not sequentially, with the Phase~1 applied research activities in
Cluster~1. This interdependency must be explicitly reflected in Horizon Europe project
call design to ensure that SBOM infrastructure investment is not siloed within supply
chain security calls but is structurally integrated into hardware security and
interoperability funding instruments.

% =========================================================================
% NEW BIBLIOGRAPHY ENTRIES
% =========================================================================

% Add the following entries to the master \begin{thebibliography}{99} block:
 in master document
%
% Required packages (add to master preamble if not present):
%   \usepackage[table]{xcolor}
%   \usepackage{multirow}
%   \usepackage{array}
%
% New bibliography keys required (entries provided at end of file):
%   fips203, fips204, fips205, nistir8413, nistsp800207,
%   etsi103744, etsizsmoo9, tgcpm20, isoiec11889,
%   opentitan, keystone, threeGPP23288, threeGPP33501, dora, gdpr
% =========================================================================

% ---- Shared colors/macro helpers when included from a master document ----
\providecolor{gaterow}{RGB}{240,240,238}
\providecolor{phase1row}{RGB}{237,245,253}
\providecolor{phase2row}{RGB}{237,249,244}
\providecolor{phase3row}{RGB}{245,244,254}
\providecolor{colheader}{RGB}{245,245,243}
\providecolor{clusterheader}{RGB}{220,232,248}

\providecolor{sslredbg}{RGB}{252,235,235}
\providecolor{sslredtx}{RGB}{163,45,45}
\providecolor{sslamberbg}{RGB}{250,238,218}
\providecolor{sslambtx}{RGB}{133,79,11}
\providecolor{ssltransbg}{RGB}{230,241,251}
\providecolor{ssltrantx}{RGB}{24,95,165}
\providecolor{ssltealbg}{RGB}{225,245,238}
\providecolor{ssltealtx}{RGB}{15,110,86}

\providecommand{\sslbadge}[3]{%
      \parbox{\linewidth}{\raggedright
            \colorbox{#1}{\scriptsize\parbox{0.92\linewidth}{\raggedright
                  \textbf{\textcolor{#2}{#3}}}}}}

\providecommand{\kpistrip}[1]{%
      \par\vspace{2pt}\noindent\rule{\linewidth}{0.3pt}%
      \par\noindent\textit{\scriptsize #1}}

\providecommand{\cellhead}[1]{\noindent\textbf{#1}\par\smallskip}
\providecommand{\gatelbl}[1]{\scriptsize\textit{Gate~#1:}}
\providecommand{\subp}[1]{\textit{\scriptsize(#1)}}

\section{Strategic Priority Clusters}
\label{ch:strategic_priority_clusters}

% \subsection{Consolidation Rationale and Structural Logic}
% \label{sec:consolidation_rationale}

% The ten executive priorities identified through the 4-Filter Prioritisation Framework represent
% distinct technology domains, but they share deep structural dependencies and converge on three
% strategic imperatives that define Europe's 15-year cybersecurity trajectory for the Cloud-Edge-IoT
% continuum. Maintaining ten parallel profiles in a dual-audience deliverable risks fragmenting
% the narrative for European Commission technical evaluators while obscuring the systemic
% interdependencies that determine whether individual priorities succeed or fail.

% The consolidation presented in this chapter reorganises the ten priorities into three strategic
% clusters, each with a unifying architectural thesis. To preserve the analytical precision required
% by IEEE Access peer review, each cluster retains its constituent sub-priorities as explicitly
% labelled components within a nested progression matrix, allowing individual technology maturation
% trajectories to be tracked independently while the cluster-level framing communicates strategic
% coherence to policy audiences.

% The mapping is as follows:
% \begin{itemize}
%   \item \textbf{Priority Cluster 1 --- Sovereign Trusted Infrastructure:}
%         P1 (Hardware Roots of Trust), P2 (PQC Agility), P3 (Lightweight
%         Cryptography and Physical-Layer Security), P8 (Confidential
%         Computing), P9 (Resilient-by-Design / Island Mode).
%   \item \textbf{Priority Cluster 2 --- Cognitive Security and Autonomous
%         Cyber Defence:} P6 (Cognitive SecOps / LLMSecOps), P10 (Proactive
%         APT Management), and the AI-as-target dimension drawn from
%         Cross-Layers CL1.2, CL1.3, and CL4.7.
%   \item \textbf{Priority Cluster 3 --- Federated Operational Sovereignty:}
%         P4 (Security Interoperability via Simpl), P5 (Supply Chain
%         Verification), P7 (Pan-European NaaS Federation).
% \end{itemize}

% One critical cross-priority dependency must be made explicit before the cluster profiles are
% presented. Hardware supply chain attestation (P5) and Runtime SBOM verification constitute
% a shared gate condition for both Cluster 1 (TEE and secure element integrity) and Cluster 3
% (federated supply chain governance). Progression in P8 (Confidential Computing) is contingent
% on the availability of trustworthy RISC-V secure elements from P1; P3 (Lightweight PhySec)
% and P2 (PQC) both depend on EU Chips Act fabrication capacity as a common Phase 2 enabler.
% These shared bottlenecks are identified explicitly in the Cross-Priority Bottleneck assessment
% in Section~\ref{sec:bottlenecks} and traced to specific progression matrix gate conditions.

% =========================================================================
% PRIORITY CLUSTER 1: SOVEREIGN TRUSTED INFRASTRUCTURE
% =========================================================================

% \newpage
\section{Priority Cluster 1: Sovereign Trusted Infrastructure}
\label{sec:cluster1}

\noindent\textbf{Constituent Sub-Priorities:} P1 (Hardware Roots of Trust \& EUCC `High'),
P2 (PQC Agility for Edge/IoT), P3 (Lightweight Cryptography \& Physical-Layer Security),
P8 (Confidential Computing \& Privacy-Preserving AI), P9 (Resilient-by-Design \& Island Mode).

\medskip
\noindent\textbf{Cluster Thesis:} Before any federated or AI-driven security capability can be
trusted, the substrate on which it executes must provide cryptographically verifiable integrity,
quantum-resistant confidentiality, and autonomous survivability independent of external
infrastructure. This cluster defines the verifiable hardware and cryptographic foundation that
all higher-order capabilities in Clusters 2 and 3 depend upon as an architectural prerequisite.

\subsection*{7-Dimensional Analysis: Priority Cluster 1}

\begin{enumerate}

% ---- DIMENSION 1 ----
\item \textbf{Baseline \& Target Objective}

\textit{Baseline (2025--2026):}
European critical infrastructure and edge deployments operate under high foreign dependency
across all five constituent sub-priorities. Hardware roots of trust are dominated by non-EU
Trusted Platform Module (TPM) manufacturers --- Infineon SLB9670, NXP SE050, and
STMicroelectronics ST33 series collectively account for the majority of the EU market under
ISO/IEC~11889~\cite{isoiec11889} and TCG TPM~2.0 specifications~\cite{tgcpm20}. The Trusted
Execution Environment (TEE) market is dominated by Intel TDX, AMD SEV-SNP, and ARM TrustZone,
none of which are EU-designed or manufactured. Post-Quantum Cryptography (PQC) standardisation
reached a definitive milestone with NIST's publication of FIPS~203 (ML-KEM)~\cite{fips203},
FIPS~204 (ML-DSA)~\cite{fips204}, and FIPS~205 (SLH-DSA)~\cite{fips205} in August~2024;
however, hardware-optimised implementations for 8/16-bit microcontrollers remain at the
Applied~R\&D stage, as confirmed by NIST IR~8413~\cite{nistir8413}. Physical-layer security
for extreme-edge mMTC devices lacks standardised specifications. Edge nodes universally operate
with hard cloud dependencies and no autonomous fallback architectures meeting the operational
resilience mandates of NIS2~\cite{nis2} and DORA~\cite{dora}.

Current cluster ODP: Applied R\&D (driven by PQC, Lightweight PhySec, and Confidential
Computing sovereign alternatives). Current cluster SSL: High Foreign Dependency.

\textit{Target (2040):}
At least 60\% of EU critical infrastructure edge nodes deploy EU-designed RISC-V open-hardware
secure elements certified at EUCC `High' assurance level~\cite{eucc}. 100\% of new EU critical
IoT deployments support hybrid or full FIPS~203/204/205-compliant PQC. Zero-energy nodes
execute quantum-resistant device authentication within a sub-100\,$\mu$J energy budget per
operation. Trusted Execution Environments built on EU open-hardware architectures (Keystone~\cite{keystone},
OpenTitan~\cite{opentitan}) cover at least 50\% of new EU cloud-edge deployments. All NIS2-essential
and DORA-critical edge nodes demonstrate validated 72-hour autonomous Island Mode survivability
independent of WAN connectivity.

% ---- DIMENSION 2 ----
\item \textbf{Primary Drivers}

\textit{Threat Drivers:}
\begin{itemize}
  \item State-sponsored APT exploitation of JTAG debug interfaces and electromagnetic
        side-channel analysis (EMSCA) on field-deployed edge nodes; the Spectre and Meltdown
        class of microarchitectural side-channel vulnerabilities establishes that even
        production-certified hardware can be compromised at scale, and far-edge nodes
        operating without physical security controls present a materially higher exposure
        surface than hyperscaler data centres.
  \item ``Harvest Now, Decrypt Later'' (HNDL) campaigns: state-level adversaries are
        actively archiving encrypted EU government and industrial communications under
        the expectation of Cryptographically Relevant Quantum Computer (CRQC)
        availability. NIST IR~8413 places the Q-Day horizon in the 10--20 year range
        with a highly uncertain lower bound~\cite{nistir8413}; HNDL renders this a
        present-tense threat against data with multi-decade confidentiality
        requirements.
  \item Physical tampering and environmental attacks on unattended far-edge nodes
        (smart grid substations, O-RAN radio units, industrial gateways) where the
        physical security controls available in hyperscaler facilities are absent.
  \item Cascading WAN disconnection failures propagating through centrally-dependent edge
        architectures, exploited by adversaries to degrade critical infrastructure
        services during hybrid conflict scenarios.
\end{itemize}

\textit{Regulatory Drivers:}
\begin{itemize}
  \item CRA Articles~10 and~11: mandatory vulnerability handling and security update
        obligations for products with digital elements; enforcement begins in 2027 for
        most product categories~\cite{cra}.
  \item NIS2 Article~21(2)(e): cryptographic policies and key management for essential
        and important entities; full enforcement began October 2024~\cite{nis2}.
  \item DORA Article~9: ICT resilience testing including hardware-level resilience
        requirements for financial sector critical infrastructure~\cite{dora}.
  \item AI Act Article~9: risk management system requirements for high-risk AI
        systems; critical infrastructure AI components require hardware integrity
        guarantees that only EUCC `High' certified secure elements can
        provide~\cite{aiact}.
\end{itemize}

% ---- DIMENSION 3 ----
\item \textbf{High Impact Activities}

\begin{enumerate}[label=(\alph*)]
  \item \textbf{[Phase~1]} Develop sovereign European Secure Elements based on
        RISC\-V open-hardware (building on OpenTitan~\cite{opentitan} and
        Keystone~\cite{keystone} frameworks), targeting EUCC `High' assurance level
        evaluation; implement automated hardware-level remote attestation protocols
        interoperable with TPM~2.0~\cite{tgcpm20} and ISO/IEC~11889~\cite{isoiec11889}
        for heterogeneous large-scale edge fleets.

  \item \textbf{[Phase~1--2]} Optimise NIST FIPS~203 (ML-KEM)~\cite{fips203} and
        FIPS~204 (ML-DSA)~\cite{fips204} for 8/16-bit microcontrollers under severe
        energy and memory constraints (target: ML-KEM-512 within 32\,KB RAM on
        Cortex-M0+ class devices); develop crypto-agility OTA update middleware enabling
        algorithm rotation without hardware replacement, compatible with ETSI TS~103~744
        hybrid key exchange profiles~\cite{etsi103744}.

  \item \textbf{[Phase~1--2]} Implement ultra-lightweight PQC and keyless
        Physical-Layer Security (RF fingerprinting exploiting Channel State Information
        and Carrier Frequency Offset, CSI/CFO-based device authentication) for mMTC
        devices operating under zero-energy or near-zero-energy constraints; target
        integration with 6G ISAC (Integrated Sensing and Communication) interfaces per
        SNS~JU~SRIA~\cite{snsju}.

  \item \textbf{[Phase~2]} Deploy multi-party confidential computing pipelines
        using TEE abstraction layers (Keystone enclaves on RISC-V~\cite{keystone})
        combined with Fully Homomorphic Encryption (FHE) schemes for collaborative AI
        training across multi-provider nodes without raw data exposure; align with
        EUCS `High+' certification requirements~\cite{eucs}.

  \item \textbf{[Phase~2--3]} Engineer autonomous Island Mode architecture: soft-state
        edge orchestration using cached credentials with 72-hour validity windows,
        LoRaWAN LPWAN fallback micro-slices, eBPF-based kernel-level policy enforcement
        for autonomous node isolation, and bio-inspired self-healing swarm protocols
        enabling edge nodes to sustain critical service delivery without cloud
        connectivity or human intervention.
\end{enumerate}

% ---- DIMENSION 4 ----
\item \textbf{Catalysts \& Enablers}

\textit{Funding Alignment:}
Horizon Europe Cluster~4 (Digital, Industry and Space) for sovereign secure silicon and
PQC hardware implementation; Horizon Europe Cluster~3 (Civil Security for Society) for
edge resilience and Island Mode architectures; EU Chips Act (Regulation (EU) 2023/1781)
pilot fabrication lines for security-critical silicon via the Chips Joint Undertaking
(KDT~JU); ECCC Strategic Agenda priority areas ``Secure Components and Systems'' and
``Post-Quantum Cryptography transition infrastructure''~\cite{eccc}; Digital Europe
Programme for PQC migration tooling and crypto-agility middleware.

\textit{Standardisation:}
NIST NCCoE ``Migration to Post-Quantum Cryptography'' project (reference architecture
and migration playbooks); CEN/CENELEC JTC~13/WG~8 (hardware security standards aligned
with EUCC Common Criteria); ETSI TC CYBER (ETSI TS~103~744 quantum-safe hybrid key
exchanges~\cite{etsi103744}); ENISA PQC transition guidelines and cryptographic inventory
tooling; 3GPP SA3 for physical-layer security integration with 5G/6G
specifications~\cite{threeGPP33501}.

\textit{Ecosystem Enablers:}
RISC-V International open-hardware community; OpenTitan coalition pilot deployments
for EU sovereign RoT~\cite{opentitan}; SNS~JU pilot networks for 6G-integrated PhySec
validation~\cite{snsju}; federated Cyber Range infrastructure (CL3.5) for Island Mode
and TEE attestation stress-testing.

% ---- DIMENSION 5 ----
\item \textbf{Disruptors \& Systemic Risks}

\begin{itemize}
  \item \textbf{Q-Day cliff:} CRQC becoming operational before EU PQC migration
        completes; NIST IR~8413 places this horizon at 10--20 years with a highly
        uncertain lower bound~\cite{nistir8413}. The EU PQC migration timeline targeting
        critical infrastructure completion by 2030 has limited buffer against the
        pessimistic lower bound, and HNDL campaigns make this a present-tense threat
        against long-lived secrets.
  \item \textbf{EU Chips Act EUV lithography shortfall:} RISC-V secure element
        production at the node geometries required for EUCC `High' certification depends
        on advanced lithography; EU 2030 Chips Act targets (20\% global market share)
        do not specify security-critical silicon capacity independently, and ASML EUV
        tooling remains subject to US--Netherlands export control dynamics that could
        constrain sovereign fabrication timelines beyond 2028.
  \item \textbf{Side-channel analysis breakthroughs against FIPS~203 software
        implementations:} lattice-based algorithms have demonstrated timing and
        cache-timing attack surfaces in software implementations on constrained devices;
        hardware countermeasures are not yet standardised at the time of writing
        [CITATION REQUIRES VERIFICATION --- recent SCA against ML-KEM implementations].
  \item \textbf{TEE architectural monoculture:} if EU open-hardware TEE adoption
        concentrates on a single architecture (e.g., Keystone only~\cite{keystone}), a
        single discovered vulnerability propagates across the entire EU critical
        infrastructure TEE estate without architectural diversity as a containment layer.
  \item \textbf{LPWAN regulatory fragmentation:} Island Mode fallback using LoRaWAN
        operates on unlicensed 868\,MHz spectrum in the EU; spectrum policy
        fragmentation across Member States could prevent coordinated emergency
        micro-slice activation under unified crisis response scenarios.
\end{itemize}

% ---- DIMENSION 6 ----
\item \textbf{Failsafe Mechanisms \& Resilience}

\begin{itemize}
  \item \textbf{Hybrid cryptographic transition:} deploy ML-KEM (FIPS~203)~\cite{fips203}
        combined with ECDH-P384 in a hybrid key encapsulation scheme per
        ETSI TS~103~744~\cite{etsi103744}; if ML-KEM is found vulnerable via cryptanalytic
        breakthrough, ECDH-P384 maintains classical confidentiality while SLH-DSA
        (FIPS~205)~\cite{fips205} provides backup hash-based signature integrity with
        distinct mathematical foundations.
  \item \textbf{TEE implementation diversity:} deploy a minimum of two architecturally
        distinct TEE technologies across critical EU infrastructure (ARM CCA on
        server-class hardware and RISC-V Keystone~\cite{keystone} on edge nodes) to
        prevent single-architecture compromise from cascading across the entire
        confidential computing estate.
  \item \textbf{Autonomous node quarantine via eBPF:} if TEE remote attestation fails,
        eBPF-based network policy enforcement at the Linux kernel level (e.g., via
        Cilium or Falco eBPF rulesets) isolates the affected node from the orchestration
        plane within sub-second latency, blocking lateral movement before human
        intervention is available.
  \item \textbf{Island Mode failsafe:} edge nodes maintain cryptographically signed
        cached credentials with defined 72-hour validity windows; on WAN disconnection,
        a LoRaWAN LPWAN micro-slice activates a minimal viable network carrying
        life-safety telemetry; if LPWAN also fails, nodes revert to local-only operation
        with integrity-protected audit logging for post-incident forensics.
  \item \textbf{RF fingerprinting fallback:} for zero-energy IoT nodes where
        cryptographic processing exceeds the available energy budget, Physical-Layer
        authentication via RF fingerprinting (CSI/CFO-based device identity) provides
        keyless device verification using intrinsic wireless channel characteristics,
        requiring no cryptographic key storage or computation.
\end{itemize}

% ---- DIMENSION 7 ----
\item \textbf{Transversal Meta-Layer}

\textit{Sovereignty \& Autonomy:}
This cluster eliminates approximately 100\% foreign dependency on non-EU TPMs (Infineon,
NXP, STMicroelectronics dominate the EU market under TCG TPM~2.0~\cite{tgcpm20}) and
non-EU TEE architectures (Intel TDX, AMD SEV-SNP, ARM TrustZone, all non-EU designed and
manufactured). Crypto-agility frameworks additionally reduce dependency on US-controlled
NIST algorithm supply chains by enabling algorithm substitution if standardisation processes
are constrained by geopolitical factors.

\textit{Trust \& Ethics:}
GDPR Article~25 (data protection by design and by default) is operationalised through
TEE-enforced data minimisation and in-use encryption in multi-tenant environments,
ensuring that sensitive data processing cannot be accessed by platform operators~\cite{gdpr}.
AI Act Article~9 risk management system requirements for high-risk AI components require
hardware integrity guarantees --- only EUCC `High' certified secure elements provide the
mathematically verifiable trust anchor this obligation demands~\cite{aiact}.

\textit{Sustainable Security (Green IT):}
CRYSTALS-Kyber (ML-KEM) demonstrates approximately 1.5$\times$ computational overhead
relative to ECDH at equivalent security levels on ARM Cortex-M4, per NIST IR~8413~\cite{nistir8413},
representing a manageable energy penalty. Lightweight PQC research targets sub-100\,$\mu$J
per authentication operation for zero-energy IoT nodes. EU-sovereign RISC-V secure
elements designed for edge power constraints enable more energy-efficient cryptographic
operations than general-purpose processor TPM emulation, reducing per-device carbon
footprint at continental deployment scale.

\end{enumerate}



% =========================================================================
% CLUSTER 1: SOVEREIGN TRUSTED INFRASTRUCTURE
% =========================================================================
\newpage

\subsection{Progression Matrix: Priority Cluster 1 --- Sovereign Trusted Infrastructure}

\begin{landscape}
\setlength{\textwidth}{26.5cm}\setlength{\linewidth}{26.5cm}%
\setlength{\hsize}{26.5cm}\setlength{\columnwidth}{26.5cm}
\setlength{\tabcolsep}{5pt}
\renewcommand{\arraystretch}{1.35}

\begin{longtable}{%
|>{\footnotesize\raggedright\arraybackslash}p{2.8cm}%
|>{\footnotesize\raggedright\arraybackslash}p{5.4cm}%
|>{\footnotesize\raggedright\arraybackslash}p{5.4cm}%
|>{\footnotesize\raggedright\arraybackslash}p{5.4cm}%
|>{\footnotesize\raggedright\arraybackslash}p{5.4cm}|}

\caption{\textbf{Cluster~1: Sovereign Trusted Infrastructure} (P1, P2, P3, P8, P9).
Transition matrix across the four methodology dimensions. Sub-priority labels in
parentheses indicate where individual priorities diverge from the cluster-level
descriptor within a cell. ODP scale: Applied R\&D $\to$ Fragmented Piloting
$\to$ Standardised Integration $\to$ Ubiquitous Operation. SSL scale:
High Foreign Dependency $\to$ Emerging EU Ecosystem $\to$ Full Sovereign Autonomy.}
\label{tab:cluster1}\\

\hline
\rowcolor{colheader}
\textbf{Phase / SSL} &
\textbf{R\&I Maturity}{\scriptsize ODP indicator per sub-priority} &
\textbf{Policy, Regulatory \& Standards}{\scriptsize enforcement $\cdot$
  certification $\cdot$ standardisation} &
\textbf{Operational Deployment \& Adoption}{\scriptsize actor types $\cdot$
  scale $\cdot$ sector coverage} &
\textbf{Market Uptake \& Strategic Sovereignty}{\scriptsize SSL indicator
  $\cdot$ dependency reduction $\cdot$ industrial capacity} \\
\hline
\endfirsthead

\hline
\rowcolor{colheader}
\textbf{Phase / SSL} &
\textbf{R\&I Maturity} &
\textbf{Policy, Regulatory \& Standards} &
\textbf{Operational Deployment \& Adoption} &
\textbf{Market Uptake \& Strategic Sovereignty} \\
\hline
\endhead
\hline\multicolumn{5}{r}{\footnotesize\textit{Continued\ldots}}\\ \endfoot
\hline \endlastfoot

% ---- PHASE 1 ----
\rowcolor{phase1row}
\textbf{Phase~1}
\newline\textit{Foundation \&}
\newline\textit{Compliance}
\newline\textit{2025--2030}
\newline\sslbadge{sslredbg}{sslredtx}{SSL: High Foreign Dependency}
&
\cellhead{Cluster ODP: Applied R\&D / Fragmented Piloting}
The cluster is split at baseline. \subp{P1, P9} reach Fragmented
Piloting: RISC-V secure element projects targeting EUCC `High' are
funded under Horizon Europe and EU Chips Act; Island Mode soft-state
protocols are validated in 5G research testbeds.
\subp{P2, P3, P8} remain at Applied R\&D: NIST FIPS~203/204/205
algorithms are being optimised for 8/16-bit microcontrollers under
severe energy and memory constraints; Keystone and OpenTitan
prototypes are under EUCC pre-evaluation; RF fingerprinting
Physical-Layer Security is validated in controlled 5G lab environments;
Fully Homomorphic Encryption approaches practical latency bounds in
research settings only.
\kpistrip{The cluster ODP is Applied R\&D, driven by the majority of
sub-priorities. Deployment claims for P2, P3, P8 must be grounded in
named Horizon Europe or ECCC project outcomes, not technology
previews.}
&
\cellhead{Regulatory environment: enforcing, certification nascent}
NIS2 Article~21(2)(e) cryptographic policies for essential entities
entered full enforcement in October~2024. CRA Articles~10 and~11
(vulnerability handling and security update obligations) begin
enforcement in 2027, creating immediate procurement-level demand for
hardware security certification. AI~Act Articles~9 and~15 establish
hardware integrity requirements for high-risk AI systems and enter
applicability in 2025--2026. ENISA PQC migration guidelines are
published; EUCC evaluation processes are initiated for the first RISC-V
prototypes, but no EU-specific EUCC~`High' certified sovereign secure
element exists yet. CEN/CENELEC JTC~13 standards for hardware security
MIMs are in draft phase.
\kpistrip{The regulatory framework is enforcing existing obligations
but lacks EU-specific hardware certification outputs; the critical gap
is the absence of a ratified EUCC~`High' EU-origin product.}
&
\cellhead{Deployment: research consortia and national agencies only}
RISC-V secure element pilots are operated by 3--5 Horizon Europe
consortia and national cybersecurity agencies in controlled
environments. Intel~TDX, AMD~SEV-SNP, and ARM~TrustZone dominate
all operational TEE deployments across EU cloud and critical
infrastructure. PQC is confined to cloud-side software implementations
in pilot contexts; far-edge IoT devices run classical cryptography
exclusively. Edge nodes operate with hard cloud dependencies and no
validated autonomous fallback architecture meeting DORA or NIS2
business continuity thresholds. SMEs have no access to EU-sovereign
hardware alternatives at viable cost.
\kpistrip{Operational deployment at Phase~1 is institutional and
research-concentrated. The gap between pilot existence and operational
availability for regulated-sector procurement must be explicitly
acknowledged.}
&
\cellhead{SSL: High Foreign Dependency}
Non-EU TPMs (Infineon~SLB9670, NXP~SE050, STMicroelectronics~ST33)
hold $\geq$90\% of EU market share in hardware security modules.
Intel, AMD, and ARM control the entire operational TEE market. No
EU-designed and EU-manufactured secure element is commercially
available. EU Chips Act fabrication investment for security-critical
silicon is initiated but pre-production. US NIST algorithm supply chain
is the sole PQC standardisation reference. Structural vulnerability is
direct: geopolitical restriction of access to any of the three major
TPM suppliers would immediately impair EU critical infrastructure
attestation capability.
\kpistrip{Quantified dependency at Phase~1 baseline: $\geq$90\% foreign
market share, zero EU-certified sovereign hardware alternatives.
Sovereignty progress is measured from this specific figure.}
\\
\hline

% ---- GATE 1->2 ----
\rowcolor{gaterow}
\gatelbl{1 $\rightarrow$ 2} &
\multicolumn{4}{>{\footnotesize\raggedright\arraybackslash}p{22.0cm}|}{%
\cellcolor{gaterow}%
\textbf{R\&I trigger:} At least one EU-designed RISC-V secure element prototype
has been independently validated against EUCC~`High' evaluation criteria and has
entered the formal ENISA/national scheme certification process.
\textbf{Policy trigger:} CRA Articles~10 and~11 are in active enforcement with
documented non-compliance decisions, confirming that market incentives for
certified EU hardware alternatives are financially material.
\textbf{Operational trigger:} Documented operational deployment of hybrid
ML-KEM (FIPS~203) plus ECDH-P384 PQC schemes in at least one EU critical
infrastructure sector (energy or financial) at production scale, not pilot scale.
\textbf{Market trigger:} EU Chips Act security-critical silicon fabrication
pilot line has produced and delivered its first commercially viable batch of
EU-origin secure elements, confirmed by independent yield and cost assessment.}
\\
\hline

% ---- PHASE 2 ----
\rowcolor{phase2row}
\textbf{Phase~2}
\newline\textit{Applied \& Industrial}
\newline\textit{R\&D alongside}
\newline\textit{Federation \&}
\newline\textit{Standardisation}
\newline\textit{2030--2035}
\newline\sslbadge{ssltransbg}{ssltrantx}{SSL: Emerging EU Ecosystem}
&
\cellhead{Cluster ODP: Standardised Integration / continued Applied R\&D}
\subp{P1, P9} reach Standardised Integration: the first EUCC~`High' certified
EU-manufactured RISC-V secure elements are in production and being integrated
into critical infrastructure; automated remote attestation at fleet scale is
validated; 72-hour Island Mode achieves standardised certification under
3GPP/ETSI fallback API profiles.
\subp{P2, P3} transition from Applied R\&D to Fragmented Piloting: hybrid
ML-KEM plus ECDH-P384 schemes are validated at Industrial IoT scale; lightweight
PQC for sub-100~$\mu$J zero-energy nodes is demonstrated under 6G ISAC testbeds.
\subp{P8} reaches Fragmented Piloting: RISC-V Keystone TEEs are operational
in multi-provider Data Spaces; FHE performance approaches the threshold for
selective production use in high-value financial analytics.
\kpistrip{Phase~2 R\&I is not purely consolidation: three of five sub-priorities
are still in active applied and industrial R\&D. Describing this phase as
``standardisation only'' misrepresents the maturation arc.}
&
\cellhead{EUCC certification and PQC standards formalised}
The EUCC~`High' certification of the first EU-manufactured RISC-V secure
element is the definitive Phase~2 policy gate: it transforms applied research
into a market-accessible, legally certifiable product and unlocks the procurement
pathway for EU critical infrastructure operators. NIST FIPS~203/204/205 are
adopted into ETSI~TS profiles for EU deployment contexts. ENISA publishes a
mandatory PQC migration timeline for critical infrastructure operators, referenced
in NIS2 sector-specific implementing acts. EU Chips Act fabrication investments
are active and measurable. CEN/CENELEC hardware security standards for EU-origin
secure elements enter ratification.
\kpistrip{The single Phase~2 policy gate is EUCC~`High' certification of a
EU-origin product. All other policy milestones are enabling conditions, not
gate conditions.}
&
\cellhead{Critical infrastructure and regulated sectors deploying EU alternatives}
EU EUCC~`High' certified secure elements are deployed in energy sector
substations and telecommunications critical infrastructure in at least five
Member States. Hybrid PQC schemes are standard in new Industrial IoT and
financial sector (DORA-compliant) deployments. RISC-V Keystone TEEs operate
in federated multi-provider Data Spaces alongside existing Intel TDX
infrastructure. Validated 72-hour Island Mode is deployed in at least three
EU 5G cross-border corridor pilots under SNS~JU governance. SME-accessible
PQC tooling and SBOM compliance support are available commercially.
\kpistrip{Phase~2 deployment evidence requires named sector and Member State
scope; anonymous ``multi-sector deployment'' claims do not constitute
verifiable operational maturity.}
&
\cellhead{SSL: Emerging EU Ecosystem}
EU-designed secure elements achieve 15--30\% market share in new critical
infrastructure procurement, rising from a near-zero baseline. EU Chips Act
fabrication investments have created a domestic production capacity that is
commercially credible, if not yet dominant. Foreign TPMs remain the standard
for legacy deployments but are no longer the only option for new regulated-sector
procurement. EU open-hardware TEE ecosystem (Keystone, OpenTitan) holds
commercially certifiable products. The conditions for Full Sovereign Autonomy
are structurally in place --- fabrication capacity, certification infrastructure,
standards leadership --- but have not yet reached continental operational scale.
\kpistrip{SSL transition requires industrial substance, not research demonstration.
15--30\% new critical infra market share is the quantified Phase~2 SSL threshold.}
\\
\hline

% ---- GATE 2->3 ----
\rowcolor{gaterow}
\gatelbl{2 $\rightarrow$ 3} &
\multicolumn{4}{>{\footnotesize\raggedright\arraybackslash}p{22.0cm}|}{%
\cellcolor{gaterow}%
\textbf{R\&I trigger:} Production-ready PQC-native RISC-V silicon is available
from EU-domiciled fabrication at commercially viable yields, confirmed by
independent cost-per-unit assessment against equivalent non-EU alternatives.
\textbf{Policy trigger:} EUCC~`High' certification is mandated as a procurement
requirement for all new EU critical infrastructure edge node deployments, with
no waiver provisions for new procurements (legacy replacement programmes are
time-bound).
\textbf{Operational trigger:} Validated 72-hour Island Mode is deployed as a
standard operational baseline --- not a pilot --- in $\geq$50\% of new EU
NIS2-essential entity edge node procurements in the energy and
telecommunications sectors.
\textbf{Market trigger:} EU Chips Act sovereign fabrication capacity achieves
$\geq$50\% self-sufficiency for security-critical silicon in new EU critical
infrastructure procurement, measured by share of new orders fulfilled by
EU-origin fabrication.}
\\
\hline

% ---- PHASE 3 ----
\rowcolor{phase3row}
\textbf{Phase~3}
\newline\textit{Autonomy \&}
\newline\textit{Leadership}
\newline\textit{2035--2040}
\newline\sslbadge{ssltealbg}{ssltealtx}{SSL: Full Sovereign Autonomy}
&
\cellhead{Cluster ODP: Ubiquitous Operation}
All five sub-priorities operate at Ubiquitous Operation. EU-designed
RISC-V secure elements are the hardware default across all new EU ICT
products. Fully Homomorphic Encryption is operationalised for
end-to-end confidential processing from IoT to HPC, not confined to
high-value analytics. Zero-energy PhySec authentication via RF
fingerprinting operates as a standard device identity mechanism for
battery-less mMTC nodes. Bio-inspired swarm protocols enable autonomous
cyber-physical recovery without cloud intervention. Phase~3 R\&I is
oriented toward the next-generation frontier: neuromorphic secure
compute, post-second-generation quantum hardware, and post-CMOS
security architectures --- EU research groups define the agenda rather
than adopting it.
\kpistrip{Ubiquitous Operation requires deployment statistics, not
deployment plans. Phase~3 evidence must cite procurement share data
and operational audit findings, not programme milestone completions.}
&
\cellhead{EU as international hardware security standards reference}
EUCC~`High' certification is mandatory for all new EU ICT products
without waiver provisions. PQC-native silicon is the EU procurement
default. EU hardware security standards are referenced in ISO/IEC
international standardisation and in bilateral digital trade
agreements. Continuous automated EUCC lifecycle maintenance replaces
point-in-time audit models completely. The EU sets the terms of
hardware security compliance for market access, transforming its
internal market standard into a de facto international reference.
\kpistrip{Phase~3 policy maturity requires evidence of international
standards influence: at least one EU hardware security standard adopted
or formally referenced in an ISO/IEC or ITU-T international standard.}
&
\cellhead{Security-by-Design the unconditional hardware and crypto default}
$\geq$60\% of EU critical infrastructure edge nodes deploy EUCC~`High'
certified EU-designed secure elements. FHE is operational end-to-end
from IoT sensor to HPC compute node. Autonomous Island Mode provides
validated 72-hour survivability as a mandatory operational baseline
across all NIS2-essential and DORA-critical edge deployments. No new
EU ICT product enters the market without a hardware root of trust.
The question ``does this deployment have hardware-level integrity?''
has no non-trivial answer for any new EU deployment.
\kpistrip{Phase~3 deployment normalisation: $\geq$60\% market penetration
for EU secure elements in critical infra; 100\% of new EU ICT products
with hardware RoT as a default.}
&
\cellhead{SSL: Full Sovereign Autonomy}
EU-designed and EU-manufactured secure elements hold $\geq$60\% of new
critical infrastructure procurement market share. Foreign legacy
dependency operates under binding replacement programmes with defined
timelines. EU sovereign RISC-V silicon and open-hardware TEE
architectures are exported to and adopted by partner jurisdictions.
EU PQC and hardware TEE standards are the international reference.
Structural exposure to geopolitical disruption of any single non-EU
hardware supplier has been eliminated for new deployments. The EU
has moved from defensive sovereignty (reducing dependency) to offensive
strategic posture (defining the terms of global hardware security
competition).
\kpistrip{Full Sovereign Autonomy is verified by $\geq$60\% market share
in new critical infra procurement plus the existence of binding
replacement programmes for legacy foreign hardware --- not by 100\%
market share.}
\\
\hline

\end{longtable}
\end{landscape}


% ---- PROGRESSION MATRIX: CLUSTER 1 ----

% \noindent
% \resizebox{\textwidth}{!}{%
% \begin{tabular}{|p{4.2cm}|p{5.6cm}|p{5.6cm}|p{5.6cm}|}
% \hline
% \rowcolor{gray!30}
% \multicolumn{4}{|c|}{\textbf{Progression Matrix: Priority Cluster 1 ---
%   Sovereign Trusted Infrastructure (P1, P2, P3, P8, P9)}} \\
% \hline
% \rowcolor{gray!15}
% \textbf{Metric} &
% \textbf{Phase 1: Foundation \& Compliance (2025--2030)} &
% \textbf{Phase 2: Federation \& Standardisation (2030--2035)} &
% \textbf{Phase 3: Autonomy \& Leadership (2035--2040)} \\
% \hline
% \rowcolor{gray!10}
% \multicolumn{4}{|l|}{\textit{\textbf{Cluster-Level Progression Indicators}}} \\
% \hline
% \textbf{Cluster ODP} &
% Applied R\&D &
% Standardised Integration &
% Ubiquitous Operation \\
% \hline
% \textbf{Cluster SSL} &
% High Foreign Dependency &
% Emerging EU Ecosystem &
% Full Sovereign Autonomy \\
% \hline
% \rowcolor{gray!10}
% \multicolumn{4}{|l|}{\textit{\textbf{Sub-Priority Decomposition (Individual ODP / SSL per Phase)}}} \\
% \hline
% \textit{P1: Hardware RoT \newline \& EUCC `High'} &
% ODP: Fragmented Piloting \newline
% SSL: High Foreign Dependency &
% ODP: Standardised Integration \newline
% SSL: Emerging EU Ecosystem &
% ODP: Ubiquitous Operation \newline
% SSL: Full Sovereign Autonomy \\
% \hline
% \textit{P2: PQC Agility \newline for Edge/IoT} &
% ODP: Applied R\&D \newline
% SSL: High Foreign Dependency &
% ODP: Fragmented Piloting \newline
% SSL: Emerging EU Ecosystem &
% ODP: Ubiquitous Operation \newline
% SSL: Full Sovereign Autonomy \\
% \hline
% \textit{P3: Lightweight Crypto \newline \& Physical-Layer Sec.} &
% ODP: Applied R\&D \newline
% SSL: High Foreign Dependency &
% ODP: Standardised Integration \newline
% SSL: Emerging EU Ecosystem &
% ODP: Ubiquitous Operation \newline
% SSL: Full Sovereign Autonomy \\
% \hline
% \textit{P8: Confidential Computing \newline \& Privacy-Pres. AI} &
% ODP: Applied R\&D \newline
% SSL: High Foreign Dependency &
% ODP: Fragmented Piloting \newline
% SSL: Emerging EU Ecosystem &
% ODP: Ubiquitous Operation \newline
% SSL: Full Sovereign Autonomy \\
% \hline
% \textit{P9: Resilient-by-Design \newline \& Island Mode} &
% ODP: Fragmented Piloting \newline
% SSL: High Foreign Dependency &
% ODP: Standardised Integration \newline
% SSL: Contained Local Survivability &
% ODP: Ubiquitous Operation \newline
% SSL: Systemic EU Resilience \\
% \hline
% \rowcolor{blue!7}
% \textbf{Operational State} &
% EU-funded RISC-V secure element pilot projects commence EUCC evaluation;
% NIST FIPS~203/204/205 algorithms are adopted by research consortia and
% optimised for constrained devices in laboratory settings; Intel TDX
% and AMD SEV-SNP dominate all operational TEE deployments while
% EU open-hardware alternatives remain at prototype stage;
% edge nodes operate with hard cloud dependencies and lack
% autonomous fallback architectures validated at operational scale. &
% EU-manufactured EUCC `High' certified secure elements enter
% critical infrastructure deployments across energy and
% telecommunications sectors; hybrid ML-KEM plus ECDH-P384 schemes
% are standardised and deployed in Industrial IoT corridors;
% RISC-V Keystone-based TEEs are operational in federated
% multi-provider Data Spaces; validated 72-hour Island Mode
% survivability demonstrated in designated EU 5G cross-border
% corridor pilots under SNS~JU governance. &
% EU Chips Act sovereign fabrication supplies $\geq$60\% of EU demand
% for security-critical silicon; all new EU ICT products ship with
% EUCC `High' certified hardware roots of trust by default;
% fully homomorphic encryption is operationalised for
% end-to-end confidential processing from IoT to HPC;
% bio-inspired swarm protocols coordinate autonomous
% cyber-physical recovery without any cloud intervention,
% achieving DORA-mandated recovery time objectives
% autonomously at continental scale. \\
% \hline
% \rowcolor{orange!12}
% \textbf{Critical Enabler} &
% CRA Articles~10 and~11 enforcement (2024--2027) creating
% mandatory market demand for hardware security certification,
% establishing the commercial conditions for EU secure element
% pilot investment &
% EUCC `High' certification of the first EU-manufactured
% RISC-V secure element, constituting the gate condition
% for the hardware sovereignty transition and unlocking
% EU Chips Act fabrication investment justification &
% EU Chips Act sovereign EUV-capable fabrication capacity
% achieving $\geq$60\% market self-sufficiency in
% security-critical silicon, removing the last structural
% foreign dependency in the infrastructure trust chain \\
% \hline
% \end{tabular}%
% }

% =========================================================================
% PRIORITY CLUSTER 2: COGNITIVE SECURITY AND AUTONOMOUS CYBER DEFENCE
% =========================================================================

\newpage
\section{Priority Cluster 2: Cognitive Security and Autonomous Cyber Defence}
\label{sec:cluster2}

\noindent\textbf{Constituent Sub-Priorities:} P6 (Cognitive SecOps / LLMSecOps \& Automated
Threat Intelligence), P10 (Proactive APT Management), and the AI-as-target dimension
drawn from CL1.2 (Model Integrity / DataBOMs), CL1.3 (Dynamic Supply Chain Verification /
Runtime SBOMs), and CL4.7 (AI Trustworthiness \& Algorithmic Accountability).

\medskip
\noindent\textbf{Cluster Thesis:} Artificial intelligence is simultaneously the most powerful
instrument available for autonomous cyber defence and the most consequential new attack surface
introduced into the security architecture. A roadmap that deploys AI agents for autonomous
threat containment without hardening those agents against adversarial manipulation creates a
critical second-order vulnerability. This cluster therefore addresses three inseparable
sub-dimensions: AI as security instrument (LLMSecOps and APT hunting), AI as adversary
amplifier (AI-generated APT tooling), and AI as target (model poisoning, adversarial evasion,
and algorithmic accountability for autonomous containment decisions).

\subsection*{7-Dimensional Analysis: Priority Cluster 2}

\begin{enumerate}

% ---- DIMENSION 1 ----
\item \textbf{Baseline \& Target Objective}

\textit{Baseline (2025--2026):}
European Security Operations Centres (SOCs) rely predominantly on rule-based SIEM platforms ---
Splunk (US-domiciled, acquired by Cisco in 2023) and IBM QRadar dominate enterprise
deployments. AI-based tools are deployed as analyst ``copilots'' with human analysts
retaining all containment authority; no EU-sovereign foundation models for security
operations exist at production scale. APT dwell times in European enterprise environments
average 150--200 days, with state-sponsored threat actors demonstrating extended persistence
in EU critical infrastructure. The adversarial robustness of AI security agents --- their
resistance to prompt injection, model evasion, and training data poisoning --- is addressed
only in research contexts without operational certification frameworks. NWDAF integration
per 3GPP TS~23.288~\cite{threeGPP23288} exists in 5G standalone deployments but remains
fragmented across operators, and no EU-mandated framework for the trustworthiness of
autonomous security AI exists prior to AI Act enforcement.

Current cluster ODP: Applied R\&D (AI-as-target and sovereign LLM development); Siloed AI
Copilots (P6); Multi-Layered Alert Correlation (P10). Cluster ODP: Applied R\&D.
Cluster SSL: High Reliance on US Foundation Models / High Dependency on Global Threat Feeds.

\textit{Target (2040):}
\begin{sloppypar}
Swarm-based agentic AI defends the EU Cloud-Edge Continuum autonomously, with APT dwell
time reduced to under 24 hours in certified critical infrastructure deployments. EU-sovereign
edge Small Language Models (SLMs) and Large Context Models (LCMs) for security operations
are certified under AI Act Article~9 and cover at least 50\% of EU critical infrastructure
SOC deployments. All autonomous security agents pass mandatory adversarial robustness
validation before operational deployment. Moving Target Defense reduces APT reconnaissance
success rates below 5\% in validated simulation environments. EU AI security frameworks
are adopted as global reference implementations by partner jurisdictions.
\end{sloppypar}

% ---- DIMENSION 2 ----
\item \textbf{Primary Drivers}

\textit{Threat Drivers:}
\begin{itemize}
  \item Nation-state APT campaigns (attributed by EU agencies to APT28/GRU,
        APT29/SVR, and Sandworm) increasingly exploit LLM hallucination and
        prompt injection vulnerabilities to manipulate autonomous security agents
        into executing incorrect containment actions --- weaponising the defence
        architecture against itself.
  \item AI-generated spear-phishing and poly\-morphic malware
        enabling automated, highly-targeted APT campaigns that adapt
        faster than signature-based
        detection cycles, invalidating the static detection rules that rule-based
        SIEM platforms depend upon.
  \item Supply chain attacks targeting AI model training pipelines: data poisoning
        and backdoor insertion into security AI models deployed at continental scale
        can silently degrade detection capability without triggering observable
        anomalies in standard monitoring.
  \item Systematic abuse of MITRE ATT\&CK-catalogued techniques (T1027 Obfuscated
        Files or Information; T1055 Process Injection) within AI-assisted
        reconnaissance tools, enabling more efficient bypass of detection rules
        derived from those same frameworks.
\end{itemize}

\textit{Regulatory Drivers:}
\begin{itemize}
  \item AI Act Article~9: risk management system requirements for high-risk AI
        systems --- autonomous security containment agents making access revocation
        or network isolation decisions qualify as high-risk systems requiring
        mandatory conformity assessment before operational deployment~\cite{aiact}.
  \item AI Act Article~15: accuracy, robustness, and cybersecurity requirements
        for high-risk AI, including specific resilience requirements against
        adversarial manipulation of inputs and training data~\cite{aiact}.
  \item AI Act Article~13: transparency and provision of information for high-risk
        AI --- autonomous security decisions must be explainable to affected
        entities and competent national authorities~\cite{aiact}.
  \item NIS2 Article~21(2)(b): incident handling procedures, including
        AI-augmented detection and response, for essential entities~\cite{nis2}.
\end{itemize}

% ---- DIMENSION 3 ----
\item \textbf{High Impact Activities}

\begin{enumerate}[label=(\alph*)]
  \item \textbf{[Phase~1]} Deploy NWDAF-integrated (3GPP TS~23.288~\cite{threeGPP23288})
        telemetry pipelines feeding ML-based UEBA (User and Entity Behaviour Analytics)
        for APT indicator correlation; establish behavioural baseline profiles across
        federated enterprise environments targeting reduction of APT dwell time from
        the current 150--200 day average to under 72 hours in pilot deployments via
        ML-assisted Indicator of Compromise (IoC) pattern detection.

  \item \textbf{[Phase~1--2]} Develop EU-sovereign Small Language Models (SLMs,
        under 10B parameters) and Large Context Models (LCMs) for edge-native SOC
        operations (LLMSecOps), trained on MITRE ATT\&CK-aligned threat intelligence
        corpora under GDPR-compliant data governance~\cite{gdpr}; establish mandatory
        AI Act Article~9 conformity assessment pipelines for all autonomous security
        agent deployments prior to production authorisation~\cite{aiact}.

  \item \textbf{[Phase~2]} Deploy edge-native ``Guardian'' nodes executing autonomous
        Zero-Trust reconfigurations per NIST SP~800-207 ZTA principles~\cite{nistsp800207}
        and Moving Target Defense (MTD) --- dynamically shuffling IP addressing, network
        topology, and software diversity to disrupt APT reconnaissance --- governed by
        ETSI GS~ZSM~009 closed-loop automation standards~\cite{etsizsm009}.

  \item \textbf{[Phase~2]} Establish mandatory adversarial robustness validation
        pipelines for all AI security agents prior to operational deployment: automated
        red-team adversarial testing (prompt injection, model evasion, backdoor
        detection via CL1.4 purple-team automation), with continuous certification
        under a CEN/CENELEC-standardised AI security agent assurance framework
        aligned with CL4.7 AI Trustworthiness requirements.

  \item \textbf{[Phase~3]} Deploy swarm-based agentic AI coordination protocols
        enabling decentralised, self-evolving threat intelligence synthesis across
        the EU Cloud-Edge Continuum; integrate with Deception Technology
        infrastructure (CL2.6: honeypots, honeynets, credential canaries, fake API
        endpoints) to generate high-fidelity, low-noise adversarial ground truth for
        continuous autonomous agent retraining.
\end{enumerate}

% ---- DIMENSION 4 ----
\item \textbf{Catalysts \& Enablers}

\textit{Funding Alignment:}
EuroHPC Joint Undertaking AI Factories for training EU-sovereign security foundation
models at the compute scale required; ECCC Strategic Agenda priority areas ``Cybersecurity
of AI'' and ``AI-enabled cybersecurity tools and services''~\cite{eccc}; Horizon Europe
Cluster~3 (Civil Security for Society) calls targeting AI-driven threat intelligence and
autonomous incident response; EU AI Office for conformity assessment infrastructure
supporting high-risk security AI certification under the AI Act.

\textit{Standardisation:}
ETSI GS~ZSM~009 (Zero-Touch Network and Service Management, Closed-Loop
Automation~\cite{etsizsm009}); 3GPP TS~23.288 (NWDAF architecture for AI-native 5G
network analytics~\cite{threeGPP23288}); MITRE ATT\&CK framework (behavioural threat
intelligence taxonomy for agent training corpora); NIST SP~800-207 (Zero Trust
Architecture~\cite{nistsp800207}); CEN/CENELEC JTC~13 working groups on AI security
agent assurance standardisation.

\textit{Ecosystem Enablers:}
ENISA AI cybersecurity threat landscape guidelines; federated Cyber Range infrastructure
(CL3.5) for adversarial agent testing under realistic attack scenarios; CONCORDIA
successor programme for EU-wide SOC federation and cross-border CTI sharing~\cite{concordia};
EU AI Act notified conformity assessment bodies for high-risk AI product certification.

% ---- DIMENSION 5 ----
\item \textbf{Disruptors \& Systemic Risks}

\begin{itemize}
  \item \textbf{Adversarial AI arms race:} threat actors weaponising LLMs for
        automated zero-day discovery and polymorphic exploit generation may outpace
        the defensive AI maturation timeline --- specifically, AI-generated
        vulnerability discovery could exceed human patching bandwidth within the
        Phase~1--2 window, before EU-sovereign autonomous response agents are
        certified for operational deployment.
  \item \textbf{LLM hallucination in containment agents:} false positive isolation
        of critical infrastructure nodes (e.g., a hospital network segment incorrectly
        quarantined during a medical emergency) constitutes an operational safety risk
        that may impose worse harm than the threat the agent was designed to
        contain --- this is the central argument for mandatory HITL escalation in Phase~1.
  \item \textbf{EU sovereign LLM compute deficit:} training security-relevant
        foundation models requires EuroHPC-scale compute; as of 2025--2026, EU frontier
        model training capacity lags US and China by approximately one order of
        magnitude, constraining Phase~1 EU-sovereign models to SLM/LCM scale
        ($<$10B parameters) and limiting capability relative to US foundation models.
  \item \textbf{AI model poisoning via CTI feed corruption:} if adversaries inject
        false indicators into NWDAF analytics pipelines or Cyber Threat Intelligence
        (CTI) sharing platforms (CL2.1), behavioural baselines trained on poisoned
        data will silently misclassify legitimate traffic and fail to detect real APTs
        --- a harder failure mode to detect than outright system unavailability.
  \item \textbf{AI Act certification bottleneck:} mandatory conformity assessment
        for high-risk autonomous security agents may create a 24--36 month regulatory
        gap during which EU organisations operate without certified autonomous
        containment, relying on slower human-in-the-loop processes that cannot match
        AI-assisted adversary speed.
\end{itemize}

% ---- DIMENSION 6 ----
\item \textbf{Failsafe Mechanisms \& Resilience}

\begin{itemize}
  \item \textbf{Shadow mode operation:} all autonomous AI security agents run in
        parallel with human-supervised rule-based SIEM during Phases~1--2;
        containment actions affecting Tier-1 critical infrastructure (energy,
        healthcare, transport per NIS2 Annex~I) require mandatory Human-In-The-Loop
        (HITL) escalation regardless of AI confidence score, with reversible quarantine
        actions bounded by a defined 60-minute rollback window~\cite{nis2}.
  \item \textbf{Model diversity and behavioural circuit breakers:} deploy a minimum
        of two architecturally distinct AI detection models (e.g., transformer-based
        LLM plus graph neural network UEBA) to prevent a single-model adversarial
        exploit from disabling the entire detection capability; behavioural circuit
        breakers automatically suspend any agent exhibiting an anomalous action rate
        or pattern deviating beyond $3\sigma$ from its established operational
        baseline, reverting to rule-based SIEM until re-certification.
  \item \textbf{RLHF continuous correction:} incorporate analyst override decisions
        into continuous Reinforcement Learning from Human Feedback (RLHF) retraining
        to prevent model drift from training distribution; quarterly adversarial
        red-team audits (CL1.4) detect capability degradation before it reaches
        operational significance.
  \item \textbf{Air-gapped CTI validation:} all external CTI feeds ingested by
        NWDAF analytics pipelines pass through an isolated validation environment
        where indicator quality is assessed against independent behavioural ground
        truth before poisoned data can influence operational AI model weights.
  \item \begin{sloppypar}\textbf{Deception-backed calibration:} honeynet and credential canary
        infrastructure (CL2.6) provides high-fidelity, low-noise ground truth against
        which AI behavioural models are continuously calibrated, reducing false
        positive rates and providing an independent signal for detecting model
        poisoning via discrepancy analysis between canary alerts and AI model
        predictions.\end{sloppypar}
\end{itemize}

% ---- DIMENSION 7 ----
\item \textbf{Transversal Meta-Layer}

\textit{Sovereignty \& Autonomy:}
This cluster eliminates reliance on US-domiciled SIEM platforms (Splunk/Cisco, IBM QRadar)
and US foundation models for operational security intelligence. EU-sovereign SLMs trained
under GDPR-compliant data governance~\cite{gdpr} ensure that sensitive threat intelligence
and behavioural telemetry does not transit non-EU cloud infrastructure. EU leadership in
trustworthy AI security tools creates an exportable technology capability aligned with
the AI Act's ambition for global regulatory standard-setting.

\textit{Trust \& Ethics:}
AI Act Article~13 (transparency requirements) is mandated for all autonomous security
decisions --- affected entities must receive meaningful explanations for AI-driven access
revocations or network isolation actions~\cite{aiact}. GDPR Article~22 (restrictions on
automated individual decision-making) applies where AI containment decisions affect
individual data subjects; human review mechanisms must be structurally embedded, not
optional add-ons~\cite{gdpr}. AI Act Article~15 adversarial robustness requirements
ensure security AI systems cannot be weaponised against the infrastructure they are
designed to protect~\cite{aiact}.

\textit{Sustainable Security (Green IT):}
Edge-native SLMs under 10B parameters reduce inference energy consumption by two to
three orders of magnitude compared to centralised GPT-class model inference, enabling
real-time threat classification at the point of detection without data centre energy
overhead. Federated training paradigms reduce repeated data transmission overhead
compared to centralised training architectures and are compatible with the privacy
preservation requirements of GDPR Article~25~\cite{gdpr}.

\end{enumerate}

% =========================================================================
% CLUSTER 2: COGNITIVE SECURITY AND AUTONOMOUS CYBER DEFENCE
% =========================================================================
\newpage
\subsection{Progression Matrix: Priority Cluster 2 --- Cognitive Security and Autonomous
Cyber Defence}

\begin{landscape}
\setlength{\textwidth}{26.5cm}\setlength{\linewidth}{26.5cm}%
\setlength{\hsize}{26.5cm}\setlength{\columnwidth}{26.5cm}
\setlength{\tabcolsep}{5pt}
\renewcommand{\arraystretch}{1.35}

\begin{longtable}{%
|>{\footnotesize\raggedright\arraybackslash}p{2.8cm}%
|>{\footnotesize\raggedright\arraybackslash}p{5.4cm}%
|>{\footnotesize\raggedright\arraybackslash}p{5.4cm}%
|>{\footnotesize\raggedright\arraybackslash}p{5.4cm}%
|>{\footnotesize\raggedright\arraybackslash}p{5.4cm}|}

\caption{\textbf{Cluster~2: Cognitive Security \& Autonomous Cyber Defence}
(P6, P10, AI-as-target dimension from CL1.2, CL1.3, CL4.7).
Transition matrix across the four methodology dimensions. The AI-as-target
dimension (model integrity, adversarial robustness, algorithmic accountability)
is labelled \textit{(AI-T)} where it diverges from P6 or P10 descriptors.}
\label{tab:cluster2}\\

\hline
\rowcolor{colheader}
\textbf{Phase / SSL} &
\textbf{R\&I Maturity}{\scriptsize ODP indicator per sub-priority} &
\textbf{Policy, Regulatory \& Standards}{\scriptsize enforcement $\cdot$
  certification $\cdot$ standardisation} &
\textbf{Operational Deployment \& Adoption}{\scriptsize actor types $\cdot$
  scale $\cdot$ sector coverage} &
\textbf{Market Uptake \& Strategic Sovereignty}{\scriptsize SSL indicator
  $\cdot$ dependency reduction $\cdot$ industrial capacity} \\
\hline
\endfirsthead

\hline
\rowcolor{colheader}
\textbf{Phase / SSL} &
\textbf{R\&I Maturity} &
\textbf{Policy, Regulatory \& Standards} &
\textbf{Operational Deployment \& Adoption} &
\textbf{Market Uptake \& Strategic Sovereignty} \\
\hline
\endhead
\hline\multicolumn{5}{r}{\footnotesize\textit{Continued\ldots}}\\ \endfoot
\hline \endlastfoot

% ---- PHASE 1 ----
\rowcolor{phase1row}
\textbf{Phase~1}
\newline\textit{Foundation \&}
\newline\textit{Compliance}
\newline\textit{2025--2030}
\newline\sslbadge{sslredbg}{sslredtx}{SSL: High Reliance on}
\newline\sslbadge{sslredbg}{sslredtx}{US Foundation Models}
&
\cellhead{Cluster ODP: Applied R\&D / Siloed AI Copilots}
\subp{P6} reaches Siloed AI Copilots: commercial LLMs (GPT-4-class)
assist human SOC analysts in alert triage, threat report generation, and
MITRE ATT\&CK technique mapping; no EU-sovereign LLM for security
operations exists at production scale; NWDAF (3GPP TS~23.288)
telemetry pipelines are integrated in 5G SA pilot deployments for
anomaly detection. \subp{P10} is at Multi-Layered Alert Correlation:
ML-based UEBA is applied to APT IoC detection in research and pilot
contexts; APT dwell time reduction from the 150--200 day average is
demonstrated in controlled settings. \subp{AI-T} is at Applied R\&D:
adversarial robustness testing (prompt injection, model evasion,
DataBOM verification) is an academic research function without
operational certification frameworks.
\kpistrip{Phase~1 R\&I distinguishes three sub-dimensions with
different maturity levels. Conflating them into a single ODP cell
obscures that AI-as-target capability is the least mature and the
most consequential gap.}
&
\cellhead{AI Act enforcing; conformity assessment infrastructure absent}
AI~Act Articles~9 (risk management), 13 (transparency), and 15
(adversarial robustness) enter applicability in 2025--2026, establishing
mandatory requirements for autonomous security AI as high-risk systems.
However, no EU conformity assessment bodies with capacity to certify
autonomous security agents are operational at the start of Phase~1.
ETSI GS~ZSM~009 (closed-loop automation) is published but not mandated.
NIS2 Article~21(2)(b) incident handling procedures are being
operationalised across essential entities. ENISA AI cybersecurity threat
landscape guidelines are published. The critical regulatory gap is the
absence of a standardised adversarial robustness test suite applicable
to security-specific AI agents.
\kpistrip{The AI Act creates obligations before the conformity
assessment infrastructure exists to fulfil them. This 24--36 month gap
is the dominant Phase~1 policy risk for this cluster.}
&
\cellhead{Deployment: SOC copilot use by large enterprises and national CSIRTs}
US-platform LLMs (GPT-4-class, under data processor agreements) operate
as analyst copilots in large enterprise and national CSIRT SOCs.
Rule-based SIEMs (Splunk/Cisco, IBM~QRadar) remain the primary
detection infrastructure across EU regulated sectors. NWDAF-integrated
AI analytics operate in 3--5 5G SA deployments run by major EU
telecom operators. No certified autonomous security containment agent
is in operational deployment anywhere in the EU. APT dwell times of
150--200 days in EU enterprise environments reflect the gap between
AI-assisted detection capability and operational deployment at scale.
SMEs lack access to AI-powered SOC tooling at viable cost.
\kpistrip{Phase~1 deployment is concentrated in the largest institutional
actors (national CSIRTs, Tier-1 operators, major financial entities).
The gap between these deployments and SME access defines the adoption
challenge for Phase~2.}
&
\cellhead{SSL: High Reliance on US Foundation Models}
The EU SOC tooling market is dominated by US-domiciled platforms:
Splunk (Cisco), IBM QRadar, and Microsoft Sentinel collectively control
the EU enterprise SIEM market. The foundation models used for SOC
copilot functions are US-origin (OpenAI, Anthropic, Google) and are
accessed under data processor agreements, creating persistent
extra-territorial data exposure. No EU-sovereign foundation model at
production scale for security operations exists. EuroHPC compute
capacity is insufficient for frontier-class security model training.
The adversarial robustness validation market is pre-commercial.
\kpistrip{Quantified dependency: US platforms hold $\geq$80\% of EU
enterprise SIEM market; US-origin models are the sole source for
LLM-based SOC tooling. This specific dependency structure is what
Phases~2 and~3 must reduce.}
\\
\hline

% ---- GATE 1->2 ----
\rowcolor{gaterow}
\gatelbl{1 $\rightarrow$ 2} &
\multicolumn{4}{>{\footnotesize\raggedright\arraybackslash}p{22.0cm}|}{%
\cellcolor{gaterow}%
\textbf{R\&I trigger:} At least one EU-sovereign Small Language Model
($\leq$10B parameters) trained on GDPR-compliant security corpora has
been validated against an independently benchmarked EU critical infrastructure
SOC use case at production-equivalent scale.
\textbf{Policy trigger:} AI~Act conformity assessment bodies with capacity
to certify autonomous security AI as high-risk systems under Articles~9 and~15
are operational and have issued at least one certificate to a security AI
agent product.
\textbf{Operational trigger:} At least one EU-sovereign LLMSecOps agent is
in certified operational deployment in a named NIS2-essential entity,
demonstrating APT dwell time below 72 hours in a documented production
environment.
\textbf{Market trigger:} A standardised adversarial robustness test suite
for security AI agents, developed under CEN/CENELEC or ETSI TC~CYBER
governance, has been published and is referenced in at least one EU
procurement framework for critical infrastructure AI tools.}
\\
\hline

% ---- PHASE 2 ----
\rowcolor{phase2row}
\textbf{Phase~2}
\newline\textit{Applied \& Industrial}
\newline\textit{R\&D alongside}
\newline\textit{Federation \&}
\newline\textit{Standardisation}
\newline\textit{2030--2035}
\newline\sslbadge{ssltransbg}{ssltrantx}{SSL: Sovereign EU Edge}
\newline\sslbadge{ssltransbg}{ssltrantx}{LLMs Emerging}
&
\cellhead{Cluster ODP: Federated Edge Agents / Automated Triage}
\subp{P6} reaches Federated Edge Agents: EU-sovereign SLMs ($\leq$10B
parameters) trained on MITRE ATT\&CK-aligned corpora are operational
at the edge for real-time SOC analytics without cloud inference
dependency; Guardian nodes execute autonomous Zero-Trust
reconfigurations per NIST~SP~800-207. \subp{P10} reaches Automated
Triage and LLMSecOps Containment: LLMSecOps agents ingest CTI feeds,
track APT lateral movement against behavioural frameworks, and
orchestrate incident triage autonomously under ETSI GS~ZSM~009
governance; APT dwell time $\leq$72 hours in certified pilot deployments.
\subp{AI-T} reaches Standardised Integration: adversarial robustness
validation pipelines (prompt injection, model evasion, backdoor
detection) are standardised and mandatory under AI~Act~Article~9
conformity assessment; continuous RLHF retraining and circuit-breaker
mechanisms are in production.
\kpistrip{Phase~2 R\&I includes active applied and industrial R\&D for
all three sub-dimensions --- swarm AI protocols, post-quantum AI
security, and next-generation adversarial testing frameworks --- not
merely consolidation of Phase~1 results.}
&
\cellhead{AI Act conformity assessment operational: the Phase~2 gate}
The single Phase~2 policy gate is the operationalisation of a
standardised AI~Act conformity assessment framework for autonomous
security agents, enabling certified deployment without per-case
regulatory approval. ETSI GS~ZSM~009 is mandated within EU NaaS
federation governance frameworks. A CEN/CENELEC-standardised adversarial
robustness test suite for security AI is published. MITRE ATT\&CK
behavioural framework is formally referenced in NIS2 sector-specific
incident handling guidance. EuroHPC AI Factory compute contracts for
EU-sovereign security model training are active.
\kpistrip{Without operational AI~Act conformity assessment
infrastructure, autonomous security agents cannot be certified for
regulated-sector deployment regardless of technical readiness.
This is the dominant Phase~2 policy dependency.}
&
\cellhead{Certified EU sovereign AI agents in critical infrastructure SOCs}
EU SLM-based LLMSecOps agents are in certified operational deployment
in energy, financial (DORA), and telecommunications critical
infrastructure SOCs in $\geq$3 Member States. Guardian nodes execute
autonomous Zero-Trust reconfigurations under human-in-the-loop
escalation protocols for Tier-1 critical actions. NWDAF-integrated
analytics are mandatory in 5G SA deployments. Deception infrastructure
(honeynets, credential canaries, fake API endpoints per CL2.6) is
deployed alongside AI detection systems, providing high-fidelity
ground truth for model calibration. APT dwell time $\leq$72 hours
demonstrated in pilot critical infrastructure deployments.
\kpistrip{Phase~2 deployment evidence requires named critical
infrastructure sector, Member State scope, and a documented APT dwell
time metric from an operational --- not simulated --- environment.}
&
\cellhead{SSL: Sovereign EU Edge LLMs Emerging}
EU-sovereign SLMs cover $\geq$30\% of critical infrastructure SOC
deployments in regulated sectors, rising from zero at baseline.
EuroHPC AI Factories supply training infrastructure for EU-origin
security models, reducing dependence on US hyperscaler compute.
US-platform SIEMs retain market share in legacy deployments but
face procurement competition from certified EU alternatives.
Adversarial robustness validation is a commercially available service
with EU-domiciled providers. EU AI security tool exports to partner
jurisdictions have begun, indicating the transition from security
consumer to security producer posture.
\kpistrip{SSL evidence at Phase~2: $\geq$30\% critical infra SOC market
share for EU-sovereign AI tools; at least one documented EU AI security
tool export event to a non-EU partner jurisdiction.}
\\
\hline

% ---- GATE 2->3 ----
\rowcolor{gaterow}
\gatelbl{2 $\rightarrow$ 3} &
\multicolumn{4}{>{\footnotesize\raggedright\arraybackslash}p{22.0cm}|}{%
\cellcolor{gaterow}%
\textbf{R\&I trigger:} EU-certified swarm AI coordination protocols
enabling decentralised autonomous threat intelligence synthesis have passed
independent security evaluation and are deployed in at least one cross-border
EU critical infrastructure corridor under ETSI GS~ZSM~009 governance.
\textbf{Policy trigger:} EU AI security standards (adversarial robustness
profiles, autonomous agent accountability frameworks) are adopted or formally
referenced in at least one ISO/IEC or ITU-T international standard, confirming
EU regulatory leadership beyond domestic jurisdiction.
\textbf{Operational trigger:} APT dwell time $\leq$24 hours is documented in
a named NIS2-essential entity operational deployment (not a simulation),
attributable to autonomous AI-driven detection and containment.
\textbf{Market trigger:} EU-sovereign AI security tools hold $\geq$50\% market
share in new critical infrastructure SOC deployments across at least five
Member States, displacing US-platform SIEMs as the procurement default in
regulated sectors.}
\\
\hline

% ---- PHASE 3 ----
\rowcolor{phase3row}
\textbf{Phase~3}
\newline\textit{Autonomy \&}
\newline\textit{Leadership}
\newline\textit{2035--2040}
\newline\sslbadge{ssltealbg}{ssltealtx}{SSL: Global Leadership in}
\newline\sslbadge{ssltealbg}{ssltealtx}{Trustworthy AI Security}
&
\cellhead{Cluster ODP: Ubiquitous AI Immune System}
All three sub-dimensions reach Ubiquitous Operation. Swarm-based
agentic AI defends the EU Cloud-Edge Continuum via self-evolving,
decentralised threat intelligence without any central coordination
dependency. Moving Target Defense continuously shuffles IP addressing,
network topology, and software diversity; validated APT reconnaissance
success rates fall below 5\% in certified operational environments.
Adversarial robustness validation is continuous and fully automated:
no security AI agent operates in EU critical infrastructure without
an active conformity certification. Phase~3 R\&I addresses
post-swarm frontiers: quantum-native AI security architectures,
brain-inspired neuromorphic threat detection, and AI-to-AI
adversarial dynamics beyond current threat models.
\kpistrip{Phase~3 ODP evidence requires operational statistics from
production deployments: APT dwell time figures, MTD effectiveness
metrics, and autonomous containment accuracy rates from certified
environments --- not from cyber range simulations.}
&
\cellhead{EU AI security governance as international reference}
EU AI Act frameworks for trustworthy autonomous security agents are
adopted internationally and referenced in bilateral digital
partnership agreements. Autonomous agent certification operates at
continental scale without per-case regulatory review. Zero-touch
network management under ETSI GS~ZSM~009 governs all EU NaaS
operations. The EU AI Office and ENISA jointly define the international
standard for security AI accountability. EU AI security governance
has transitioned from reactive adaptation to agenda-setting: other
jurisdictions adopt EU frameworks to access EU market and partnership
programmes.
\kpistrip{Phase~3 policy leadership is verified by international
adoption events --- not by the number of EU instruments in force.
At least one non-EU jurisdiction formally adopting EU security AI
governance standards constitutes the Phase~3 policy evidence threshold.}
&
\cellhead{Autonomous continental-scale AI cyber defence as operational default}
APT dwell time $\leq$24 hours in all certified EU NIS2-essential
critical infrastructure deployments. MTD is deployed as a mandatory
operational baseline. Swarm AI coordination operates without central
cloud dependency: each EU Cloud-Edge node participates in a
self-organising threat intelligence mesh. Deception infrastructure
is fully integrated into the autonomous response loop, providing
continuous model calibration ground truth. Human intervention is
reserved for exceptional legal escalation scenarios; routine security
operations are fully autonomous within the governance framework.
\kpistrip{The Phase~3 deployment test: is autonomous AI defence
the unconditional operational baseline, or is it still a premium
programme? 100\% of new critical infrastructure deployments
incorporating autonomous AI security constitutes the evidence threshold.}
&
\cellhead{SSL: Global Leadership in Trustworthy AI Security}
EU-sovereign AI security tools hold $\geq$50\% of new critical
infrastructure SOC deployments across the EU. US-platform SIEMs
are legacy infrastructure under active replacement programmes with
defined timelines. EU AI security models are exported to and
commercially adopted by partner jurisdictions. EU sets the commercial
terms for security AI market access: certification against EU AI~Act
Article~9/15 profiles is a de facto global procurement requirement.
The EU has transitioned from defensive sovereignty to offensive
strategic posture: it defines the technical terms on which the
global security AI industry must compete.
\kpistrip{Full sovereign leadership in this cluster requires both
domestic dominance ($\geq$50\% critical infra SOC market share)
and international influence (at least two partner jurisdictions
formally adopting EU security AI governance). Both conditions must
be met.}
\\
\hline

\end{longtable}
\end{landscape}

% ---- PROGRESSION MATRIX: CLUSTER 2 ----

% \noindent
% \resizebox{\textwidth}{!}{%
% \begin{tabular}{|p{4.2cm}|p{5.6cm}|p{5.6cm}|p{5.6cm}|}
% \hline
% \rowcolor{gray!30}
% \multicolumn{4}{|c|}{\textbf{Progression Matrix: Priority Cluster 2 ---
%   Cognitive Security and Autonomous Cyber Defence (P6, P10, AI-as-Target)}} \\
% \hline
% \rowcolor{gray!15}
% \textbf{Metric} &
% \textbf{Phase 1: Foundation \& Compliance (2025--2030)} &
% \textbf{Phase 2: Federation \& Standardisation (2030--2035)} &
% \textbf{Phase 3: Autonomy \& Leadership (2035--2040)} \\
% \hline
% \rowcolor{gray!10}
% \multicolumn{4}{|l|}{\textit{\textbf{Cluster-Level Progression Indicators}}} \\
% \hline
% \textbf{Cluster ODP} &
% Applied R\&D / Siloed AI Copilots &
% Federated Edge Agents &
% Ubiquitous AI Immune System \\
% \hline
% \textbf{Cluster SSL} &
% High Reliance on US Foundation Models &
% Sovereign European Edge LLMs / SLMs &
% Global Leadership in Trustworthy AI Security \\
% \hline
% \rowcolor{gray!10}
% \multicolumn{4}{|l|}{\textit{\textbf{Sub-Priority Decomposition (Individual ODP / SSL per Phase)}}} \\
% \hline
% \textit{P6: Cognitive SecOps \newline (LLMSecOps)} &
% ODP: Siloed AI Copilots \newline
% SSL: High Reliance on US Foundation Models &
% ODP: Federated Edge Agents \newline
% SSL: Sovereign EU Edge LLMs Emerging &
% ODP: Ubiquitous AI Immune System \newline
% SSL: Global Leadership in Trustworthy AI \\
% \hline
% \textit{P10: Proactive APT \newline Management} &
% ODP: Multi-Layered Alert \newline Correlation \& Hunting \newline
% SSL: High Dependency on Global Threat Feeds &
% ODP: Automated Triage \& \newline LLMSecOps Containment \newline
% SSL: Sovereign Pan-EU CTI \& Behavioural Analytics &
% ODP: Ubiquitous Moving \newline Target Defense (MTD) \newline
% SSL: Full Autonomous Counter-Adversarial Resilience \\
% \hline
% \textit{AI-as-Target \newline (CL1.2, CL1.3, CL4.7)} &
% ODP: Applied R\&D \newline
% SSL: High Foreign Dependency &
% ODP: Standardised Integration \newline
% SSL: Emerging EU Ecosystem &
% ODP: Ubiquitous Operation \newline
% SSL: Full Sovereign Autonomy \\
% \hline
% \rowcolor{blue!7}
% \textbf{Operational State} &
% LLMs assist human SOC analysts in alert triage and threat report
% synthesis; NWDAF telemetry is routed to centralised AI for anomaly
% detection in 5G standalone deployments; APT dwell times are
% reduced via ML-assisted IoC correlation in pilot enterprise environments;
% adversarial robustness validation of AI security agents is conducted
% in academic and research settings without operational certification
% frameworks in force. &
% Edge-native EU-sovereign SLMs execute autonomous Zero-Trust
% reconfigurations and traffic throttling in certified critical
% infrastructure deployments; LLMSecOps agents ingest CTI and
% track lateral movement against MITRE ATT\&CK behavioural
% frameworks in real time; adversarial robustness validation
% is standardised under AI Act Article~9 conformity assessment;
% APT dwell time reduced to under 72 hours in EU critical
% infrastructure pilot deployments. &
% Swarm-based agentic AI defends the EU Cloud-Edge Continuum
% via self-evolving, decentralised threat intelligence without
% central coordination dependency; Moving Target Defense
% continuously shuffles network topology, IP addressing, and
% software diversity, achieving validated APT reconnaissance
% success rates below 5\% in operational simulation environments;
% EU AI security models are formally adopted as reference
% implementations by partner jurisdictions. \\
% \hline
% \rowcolor{orange!12}
% \textbf{Critical Enabler} &
% AI Act Articles~9 and~15 enforcement (2025--2026) establishing
% mandatory risk management and adversarial robustness
% requirements for autonomous security AI, creating the
% regulatory framework for certified Phase~2 deployment &
% Standardised AI Act conformity assessment framework for
% autonomous security agents, enabling certified operational
% deployment without per-case regulatory approval and
% unlocking Phase~3 swarm-level autonomy &
% EU-certified self-evolving AI swarm protocols achieving
% full operational deployment under ETSI GS~ZSM~009
% zero-touch automation governance with Pan-EU
% CTI-sharing infrastructure \\
% \hline
% \end{tabular}%
% }

% =========================================================================
% PRIORITY CLUSTER 3: FEDERATED OPERATIONAL SOVEREIGNTY
% =========================================================================

\newpage
\section{Priority Cluster 3: Federated Operational Sovereignty}
\label{sec:cluster3}

\noindent\textbf{Constituent Sub-Priorities:} P4 (Security Interoperability via Simpl
Middleware), P5 (Supply Chain Verification and Secured Infrastructure), P7 (Pan-European
NaaS Federation --- ``Schengen for Cloud-Edge/AI'').

\medskip
\noindent\textbf{Cluster Thesis:} Technical excellence in Cluster~1 and autonomous
defence capability in Cluster~2 deliver no strategic value if they cannot operate
verifiably across borders, across providers, and across supply chains. This cluster
defines the governance fabric, interoperability standards, and supply chain integrity
mechanisms that transform individually sovereign components into a collectively sovereign
European computing continuum. Critically, it addresses the legal and semantic as well as
the technical interoperability problem --- without a ``Schengen for Cloud'' legal framework,
cross-border compute federation remains commercially fragile regardless of technical
standardisation.

\subsection*{7-Dimensional Analysis: Priority Cluster 3}

\begin{enumerate}

% ---- DIMENSION 1 ----
\item \textbf{Baseline \& Target Objective}

\textit{Baseline (2025--2026):}
European cloud and edge infrastructure remains highly fragmented across proprietary provider
silos with incompatible security policies, authentication mechanisms, and data exchange
protocols. Common European Data Spaces (Manufacturing, Health, Energy, Agriculture) are in
early piloting stages, with security interoperability relying on bilateral agreements rather
than standardised automated mechanisms aligned with the Simpl open-source middleware
stack~\cite{cloud_edge_roadmap}. Supply chain verification is largely static: CRA-mandated
Software Bills of Materials (SBOMs)~\cite{cra} are point-in-time declarations without runtime
verification; Digital Product Passport (DPP) frameworks remain in pilot stage under the EU
Ecodesign Regulation. NaaS federation across EU telecom operators does not exist at
operational scale --- cross-border edge compute sharing requires bespoke commercial agreements.
CAMARA Network APIs are being standardised by GSMA but remain fragmentedly implemented
across European operators.

Current cluster ODP: Fragmented Piloting (P4, P7); early supply chain disclosure phase
(P5). Cluster SSL: Emerging EU Ecosystem (P4); Hyperscaler Platform Dominance (P7);
High reliance on fragmented non-EU supply chain auditing (P5).

\textit{Target (2040):}
A unified European NaaS platform enables seamless cross-border compute roaming under a
``Schengen for Cloud'' legal framework, legally operable as a single jurisdiction for
data workloads. All Common European Data Spaces operate on Simpl-based security
interoperability with automated MIM compliance --- provider-agnostic, cryptographically
verifiable, and free of vendor lock-in. Runtime SBOMs with Zero-Knowledge Proof (ZKP)
verification eliminate unverified components from EU critical infrastructure supply chains.
EU achieves full sovereign control over supply chain transparency with pan-European Cyber
Threat Intelligence sharing for supply chain compromise indicators.

% ---- DIMENSION 2 ----
\item \textbf{Primary Drivers}

\textit{Threat Drivers:}
\begin{itemize}
  \item Supply chain attacks exploiting the disaggregation complexity of Open RAN
        deployments, where a single compromised firmware component sourced from one
        of 50-plus hardware and software vendors per macro cell deployment can affect
        thousands of base stations across multiple operators before detection.
  \item ``Weakest link'' cascading failures in federated NaaS environments: a
        single provider with a compromised Minimal Interoperability Mechanism (MIM)
        implementation can propagate malicious policies across the entire federated
        fabric before semantic validation detects the enforcement inconsistency.
  \item US CLOUD Act extra-territorial reach over non-EU cloud providers processing
        EU data creates a structural sovereignty risk that persists regardless of
        technical security measures if the legal framework enabling EU-sovereign
        federated compute is not established.
  \item Silent policy translation failures: automated security policy translation
        between heterogeneous provider MIM implementations may produce semantically
        equivalent but enforcement-different policies, creating exploitable gaps
        invisible to standard compliance auditing procedures.
\end{itemize}

\textit{Regulatory Drivers:}
\begin{itemize}
  \item Data Act interoperability obligations (Articles~23--31): mandatory
        requirements for cloud service providers covering security interface
        standardisation and data portability~\cite{nis2}.
        [CITATION REQUIRES VERIFICATION --- confirm final Data Act article numbering
        for interoperability provisions.]
  \item NIS2 Article~21(2)(h): supply chain security measures mandatory for essential
        entities, requiring documented risk assessment of all ICT suppliers~\cite{nis2}.
  \item CRA Annex~I Part~II: supply chain vulnerability handling requirements for
        products with digital elements, including SBOM disclosure
        obligations~\cite{cra}.
  \item DORA Article~28: ICT third-party risk management requiring continuous
        monitoring and documented exit strategies for critical ICT
        suppliers~\cite{dora}.
\end{itemize}

% ---- DIMENSION 3 ----
\item \textbf{High Impact Activities}

\begin{enumerate}[label=(\alph*)]
  \item \textbf{[Phase~1]} Define and standardise security-specific Minimal
        Interoperability Mechanisms (MIMs) within the Simpl open-source middleware
        stack~\cite{cloud_edge_roadmap}; deploy pilots in three Common European Data
        Spaces (Manufacturing, Health, Energy) with automated cross-provider policy
        enforcement, federated identity attestation, and cryptographically verifiable
        compliance reporting.

  \item \textbf{[Phase~1--2]} Implement Runtime SBOMs with cryptographic proof of
        execution for all containerised microservices and AI components migrating across
        the multi-provider continuum; integrate Digital Product Passports (DPP) for
        hardware component traceability from silicon fabrication to operational
        deployment, aligned with CRA SBOM obligations~\cite{cra}.

  \item \textbf{[Phase~1--2]} Deploy harmonised CAMARA Network APIs across a minimum
        of five EU telecom operators establishing standardised interfaces for cross-border
        edge compute roaming; test Automated Compliance Negotiators in designated EU
        5G cross-border corridor deployments under IPCEI-CIS governance.

  \item \textbf{[Phase~2]} Establish ``Data Visa'' mechanisms enabling dynamic
        cross-border AI model and data migration without manual legal friction, governed
        by a common semantic security ontology (CL3.4) aligned with STIX~2.1/TAXII
        threat intelligence standards and automated EUCS compliance verification~\cite{eucs}.

  \item \textbf{[Phase~2--3]} Deploy Zero-Knowledge Proofs (ZKPs) for IP-protected
        supply chain component verification: hardware vendors prove FIPS~203/204/205
        algorithm compliance and absence of known vulnerabilities without revealing
        proprietary design details~\cite{fips203,fips204,fips205}; operationalise
        pan-European supply chain compromise indicator sharing via standardised
        CTI lifecycle platforms (CL2.5).
\end{enumerate}

% ---- DIMENSION 4 ----
\item \textbf{Catalysts \& Enablers}

\textit{Funding Alignment:}
IPCEI-CIS (Important Project of Common European Interest on Next Generation Cloud
Infrastructure and Services) as the primary funding vehicle for cross-border NaaS federation
pilots; Digital Europe Programme for Simpl development, Common Data Space pilots, and
MIM standardisation infrastructure; ECCC Strategic Agenda priority areas ``Supply Chain
Security'' and ``Secure Interoperability of Systems and Networks''~\cite{eccc}; European
Alliance for Industrial Data, Edge and Cloud for governance and pilot ecosystem
coordination~\cite{cloud_edge_roadmap}.

\textit{Standardisation:}
CEN/CENELEC JTC~13 for security MIM standardisation within the Simpl ecosystem;
GSMA CAMARA API Alliance for harmonised Network API standardisation across operators;
ETSI TC CYBER for supply chain security standards; OASIS STIX~2.1/TAXII for threat
intelligence interoperability; W3C Verifiable Credentials Data Model for SSI-based
supply chain attestation chains.

\textit{Ecosystem Enablers:}
Gaia-X Federation Services trust framework alignment with Simpl MIM security policies;
Digital Product Passport pilot programmes under EU Ecodesign Regulation~2024;
ENISA supply chain risk management guidelines and threat landscape reports for ICT
sectors; dedicated SME compliance support programmes for CRA-mandated SBOM tooling,
recognising that sub-50 person software vendors represent the most vulnerable link in
the supply chain verification architecture~\cite{cra}.

% ---- DIMENSION 5 ----
\item \textbf{Disruptors \& Systemic Risks}

\begin{itemize}
  \item \textbf{Hyperscaler resistance to Simpl MIM adoption:} AWS, Azure, and
        GCP collectively control approximately 65\% of the EU cloud market as of
        2024 and have strong commercial incentives to maintain proprietary APIs;
        without mandatory EU Cloud Rulebook provisions requiring MIM compliance,
        voluntary adoption rates are unlikely to achieve genuine interoperability
        within Phase~1 timelines.
  \item \textbf{IPCEI-CIS implementation delays:} complex EU state aid notification
        procedures have historically delayed IPCEI projects by 18--36 months
        (IPCEI-H2 delayed by state aid clearance processes); P7 NaaS federation
        Phase~1 milestones carry material schedule risk if IPCEI-CIS governance
        does not streamline implementation mechanisms.
  \item \textbf{Semantic interoperability gap:} automated security policy translation
        between heterogeneous MIM implementations is architecturally harder than
        syntactic interoperability --- subtle semantic differences in ``deny'' versus
        ``drop'' versus ``quarantine'' semantics across provider implementations
        can create silent enforcement gaps undetectable through standard compliance
        auditing.
  \item \textbf{ZKP computational overhead:} current ZKP schemes for supply chain
        verification carry 10--100$\times$ computational overhead relative to
        standard cryptographic verification, limiting practical deployment to
        batch verification rather than real-time per-transaction attestation
        in Phases~1--2.
  \item \textbf{EU semiconductor supply chain concentration:} advanced node
        fabrication (sub-7\,nm) required for cutting-edge supply chain attestation
        hardware remains concentrated in Taiwan (TSMC) and South Korea (Samsung);
        a Taiwan Strait contingency would immediately constrain EU supply chain
        verification capability at the silicon level, exposing a gap the EU Chips Act
        does not address within the 2030 timeframe.
\end{itemize}

% ---- DIMENSION 6 ----
\item \textbf{Failsafe Mechanisms \& Resilience}

\begin{itemize}
  \item \textbf{Autonomous unfederation:} if a provider fails Runtime SBOM
        verification or MIM compliance checks, eBPF-based zero-trust policy
        enforcement severs that provider from the federated NaaS fabric and redirects
        workloads to verified compliant nodes within a sub-60-second SLA; provider
        re-admission requires fresh ZKP-verified attestation, not manual
        recertification.
  \item \textbf{Graceful national fallback:} cross-border SLA failure or ``Data
        Visa'' mechanism collapse triggers workload redirect to national infrastructure
        maintained at sufficient sovereign capacity for essential services, fulfilling
        NIS2 Article~21(2)(c) business continuity obligations and DORA Article~9 ICT
        resilience requirements~\cite{nis2,dora}.
  \item \textbf{Cryptographic attestation of MIM enforcement:} all Simpl MIM security
        controls include verifiable proofs that policies are enforced as declared --- not
        merely configured --- using TEE-backed attestation integrated with Cluster~1
        secure element infrastructure~\cite{cloud_edge_roadmap}, eliminating
        compliance-theatre scenarios where declared and actual enforcement diverge.
  \item \textbf{Supply chain quarantine:} unverified or high-risk components are
        isolated into quarantined network slices (VLAN/VxLAN-isolated segments with
        eBPF-enforced east-west traffic blocking) pending ZKP verification;
        quarantined components cannot interact with production workloads or access
        sensitive data paths.
  \item \textbf{Multi-jurisdictional residency fallback:} if cross-border ``Data
        Visa'' mechanisms fail legally or technically, automated GDPR Article~44
        compliance enforcement pins workloads to the originating jurisdiction,
        preventing unauthorised cross-border data transfers under federation failure
        conditions~\cite{gdpr}.
\end{itemize}

% ---- DIMENSION 7 ----
\item \textbf{Transversal Meta-Layer}

\textit{Sovereignty \& Autonomy:}
This cluster breaks non-EU hyperscaler API lock-in (AWS, Azure, GCP controlling
approximately 65\% of EU cloud market as of 2024) by enabling European providers to
compete on standardised Simpl-based infrastructure. A ``Schengen for Cloud'' legal
framework eliminates US CLOUD Act extra-territorial jurisdiction exposure for EU data
processed on federated EU-sovereign NaaS infrastructure. Runtime SBOM verification and
ZKP-based supply chain attestation reduce EU dependency on non-EU vendor self-certification
declarations --- the dominant model under current CRA-baseline SBOM obligations.

\textit{Trust \& Ethics:}
GDPR Article~25 (privacy by design and by default) is operationalised through
MIM-enforced data minimisation policies that are cryptographically verifiable, not
declarative~\cite{gdpr}. Data Act interoperability obligations are fulfilled through
Simpl MIM compliance, providing auditable trust chains for GDPR Article~5(2) data
processing accountability. Transparent Runtime SBOM verification provides the
supply chain integrity that GDPR Article~25 implicitly requires for processors
handling EU personal data at scale~\cite{gdpr}.

\textit{Sustainable Security (Green IT):}
Federated NaaS resource pooling across EU operators eliminates redundant edge compute
over-provisioning by individual operators, reducing total EU edge infrastructure carbon
footprint through shared utilisation efficiency. DPP-based supply chain traceability
aligns with EU Ecodesign Regulation circular economy requirements for ICT hardware,
enabling component reuse lifecycle tracking that reduces electronic waste from
premature infrastructure replacement driven by opacity in component provenance.

\end{enumerate}


% =========================================================================
% CLUSTER 3: FEDERATED OPERATIONAL SOVEREIGNTY
% =========================================================================
\newpage
\subsection{Progression Matrix: Priority Cluster 3 --- Federated Operational Sovereignty}

\begin{landscape}
\setlength{\textwidth}{26.5cm}\setlength{\linewidth}{26.5cm}%
\setlength{\hsize}{26.5cm}\setlength{\columnwidth}{26.5cm}
\setlength{\tabcolsep}{5pt}
\renewcommand{\arraystretch}{1.35}

\begin{longtable}{%
|>{\footnotesize\raggedright\arraybackslash}p{2.8cm}%
|>{\footnotesize\raggedright\arraybackslash}p{5.4cm}%
|>{\footnotesize\raggedright\arraybackslash}p{5.4cm}%
|>{\footnotesize\raggedright\arraybackslash}p{5.4cm}%
|>{\footnotesize\raggedright\arraybackslash}p{5.4cm}|}

\caption{\textbf{Cluster~3: Federated Operational Sovereignty} (P4, P5, P7).
Transition matrix across the four methodology dimensions. P4 (Simpl
Interoperability), P5 (Supply Chain Verification), and P7 (NaaS Federation)
are labelled where their maturity trajectories diverge within a cell.}
\label{tab:cluster3}\\

\hline
\rowcolor{colheader}
\textbf{Phase / SSL} &
\textbf{R\&I Maturity}{\scriptsize ODP indicator per sub-priority} &
\textbf{Policy, Regulatory \& Standards}{\scriptsize enforcement $\cdot$
  certification $\cdot$ standardisation} &
\textbf{Operational Deployment \& Adoption}{\scriptsize actor types $\cdot$
  scale $\cdot$ sector coverage} &
\textbf{Market Uptake \& Strategic Sovereignty}{\scriptsize SSL indicator
  $\cdot$ dependency reduction $\cdot$ industrial capacity} \\
\hline
\endfirsthead

\hline
\rowcolor{colheader}
\textbf{Phase / SSL} &
\textbf{R\&I Maturity} &
\textbf{Policy, Regulatory \& Standards} &
\textbf{Operational Deployment \& Adoption} &
\textbf{Market Uptake \& Strategic Sovereignty} \\
\hline
\endhead
\hline\multicolumn{5}{r}{\footnotesize\textit{Continued\ldots}}\\ \endfoot
\hline \endlastfoot

% ---- PHASE 1 ----
\rowcolor{phase1row}
\textbf{Phase~1}
\newline\textit{Foundation \&}
\newline\textit{Compliance}
\newline\textit{2025--2030}
\newline\sslbadge{sslredbg}{sslredtx}{SSL: Fragmented; Hyperscaler}
\newline\sslbadge{sslredbg}{sslredtx}{Dominance (P7); Emerging}
\newline\sslbadge{sslredbg}{sslredtx}{EU Ecosystem (P4)}
&
\cellhead{Cluster ODP: Fragmented Piloting (P4, P7) / Disclosure Phase (P5)}
\subp{P4} is at Fragmented Piloting: Simpl security MIM pilots are
launched in three Common European Data Spaces (Manufacturing, Health,
Energy) with manual policy coordination and proof-of-concept
implementations. \subp{P7} is at Bilateral Interconnection: harmonised
CAMARA Network APIs are deployed in 2--3 bilateral 5G corridor agreements
between individual EU telecom operators; cross-border edge compute
roaming does not exist at operational scale. \subp{P5} is at Supply
Chain Disclosure: CRA-mandated static SBOM declarations are being
submitted as point-in-time compliance artefacts without runtime
verification; Digital Product Passport pilot programmes are initiated
under the EU Ecodesign Regulation; Zero-Knowledge Proof verification
for supply chain attestation is at Applied R\&D stage.
\kpistrip{The three sub-priorities have meaningfully different Phase~1
maturity levels. P5 is the least mature and the most consequential
shared dependency: Runtime SBOM verification is a gate condition for
both P4 (MIM compliance trust) and Cluster~1 (TEE integrity claims).}
&
\cellhead{Multiple legislative obligations activating; standards nascent}
CRA~SBOM obligations (Annex~I Part~II) enter enforcement in 2027,
mandating static SBOM disclosure for all products with digital elements
entering the EU market. NIS2~Article~21(2)(h) makes supply chain risk
assessment mandatory for essential entities. DORA~Article~28 requires
continuous ICT third-party risk monitoring for financial sector entities.
Data Act interoperability obligations activate, creating procurement-level
pressure for cross-provider security interface standardisation. The
Simpl-Open official launch with initial MIM definitions constitutes the
Phase~1 technical policy trigger for P4. IPCEI-CIS enters its
deployment phase. GSMA CAMARA API Alliance standardisation work is
active but not yet mandated in EU regulatory instruments.
\kpistrip{Phase~1 policy activity is high-volume but fragmented:
multiple legislative instruments activate simultaneously without
a unified interoperability mandate that would create coherent market
pull for Simpl-based standardisation.}
&
\cellhead{Deployment: bilateral pilots and Data Space proof-of-concepts}
Three Common European Data Space security pilots operate Simpl-based
MIM proof-of-concepts with manual policy coordination; cross-provider
security policy enforcement is not automated. CAMARA API deployments
are confined to 2--3 bilateral 5G corridor agreements between single
operator pairs. Static SBOM declarations are submitted as CRA
compliance artefacts but are not verified at runtime; supply chain
risk assessment is vendor-declared rather than independently attested.
Hyperscalers (AWS, Azure, GCP) retain $\approx$65\% of EU cloud market
through proprietary API ecosystems. Cross-border AI model migration
requires bespoke bilateral legal review per transfer event. SME-accessible
SBOM tooling and MIM compliance support are not yet standardised or
commercially affordable.
\kpistrip{Phase~1 deployment is bilateral and pilot-scale. The gap
between proof-of-concept MIM implementations and automated enforcement
across the majority of Data Spaces defines the core operational
challenge for Phase~2.}
&
\cellhead{SSL: Mixed --- Emerging EU Ecosystem (P4) / Hyperscaler Dominance (P7)}
\subp{P4} has reached Emerging EU Ecosystem: the Simpl open-source
stack exists and is being piloted with EU institutional backing.
\subp{P7} remains at Hyperscaler Platform Dominance: AWS, Azure, and GCP
collectively control $\approx$65\% of EU cloud market; no EU NaaS
federation operates at commercially relevant scale; cross-border edge
compute is fragmented into national silos. \subp{P5} is at High
Reliance on Fragmented Non-EU Auditing: supply chain verification
is vendor-declared under US-dominated certification ecosystems (Common
Criteria national schemes, SOC2). The cluster's dominant SSL profile
is foreign dependency, driven by the weight of P7 market structure.
\kpistrip{P4 and P7 require separate SSL tracking in Phase~1: conflating
them into a single cluster SSL figure obscures that Simpl has achieved
Emerging Ecosystem while NaaS federation has not.}
\\
\hline

% ---- GATE 1->2 ----
\rowcolor{gaterow}
\gatelbl{1 $\rightarrow$ 2} &
\multicolumn{4}{>{\footnotesize\raggedright\arraybackslash}p{22.0cm}|}{%
\cellcolor{gaterow}%
\textbf{R\&I trigger:} At least one Common European Data Space has deployed
automated (not manually coordinated) Simpl MIM-based cross-provider security
policy enforcement, with documented evidence that enforcement is cryptographically
verifiable rather than declaratively asserted.
\textbf{Policy trigger:} Simpl security MIM definitions are published in
candidate CEN/CENELEC standard form and referenced as a compliance pathway in
at least one EU sector-specific NIS2 or Data Act implementing act.
\textbf{Operational trigger:} CAMARA-based cross-border edge compute roaming is
operational between at least three distinct EU telecom operator pairs (not a single
bilateral agreement), with documented SLA performance and compliance evidence.
\textbf{Market trigger:} Runtime SBOM tooling with cryptographic execution proofs
is available as a commercially supported product or service from at least two
EU-domiciled vendors, establishing the market infrastructure for Phase~2
mandatory adoption.}
\\
\hline

% ---- PHASE 2 ----
\rowcolor{phase2row}
\textbf{Phase~2}
\newline\textit{Applied \& Industrial}
\newline\textit{R\&D alongside}
\newline\textit{Federation \&}
\newline\textit{Standardisation}
\newline\textit{2030--2035}
\newline\sslbadge{ssltransbg}{ssltrantx}{SSL: Growing Pan-EU}
\newline\sslbadge{ssltransbg}{ssltrantx}{Sovereignty}
&
\cellhead{Cluster ODP: Standardised Integration (P4) / Federated Edge Roaming (P7) / Dynamic Runtime Verification (P5)}
\subp{P4} reaches Standardised Integration: automated policy translation
via standardised Simpl MIMs is deployed across the majority of active
Common European Data Spaces; federated identity attestation and
cross-provider compliance reporting are automated. \subp{P7} reaches
Federated Edge Roaming: CAMARA-based cross-border roaming operates
across $\geq$5 EU operator pairs; ``Data Visa'' mechanisms enable
dynamic AI model migration without manual legal review in certified
corridors. \subp{P5} transitions from disclosure to dynamic runtime
verification: Runtime SBOMs with cryptographic proof-of-execution
verify containerised components as they migrate across multi-provider
nodes; ZKP batch verification of supply chain attestation is operationally
validated; Digital Product Passports are mandatory under the EU
Ecodesign Regulation.
\kpistrip{Phase~2 R\&I activity includes industrial R\&D for ZKP
real-time verification (currently batch-only due to 10--100$\times$
overhead), semantic security ontology standardisation, and legal
engineering of the ``Schengen for Cloud'' framework --- not solely
standardisation of existing prototypes.}
&
\cellhead{Security MIMs as CEN/CENELEC standards: the Phase~2 gate}
The single Phase~2 policy gate is the adoption of Simpl security MIMs
as ratified CEN/CENELEC standards, referenced as mandatory compliance
requirements in the EU Cloud Rulebook. This transforms MIM adoption
from voluntary best-practice to market-access obligation for providers
serving EU-regulated sectors. DPP becomes mandatory under the EU
Ecodesign Regulation~2024 for ICT hardware categories. Data Act full
enforcement creates binding procurement-level interoperability
obligations. DORA continuous ICT third-party monitoring is fully
operationalised in the financial sector. IPCEI-CIS cross-border NaaS
governance framework is operational.
\kpistrip{The Phase~2 policy gate is the CEN/CENELEC ratification of
security MIMs. Voluntary Simpl adoption without this standards event
does not constitute Phase~2 policy maturity.}
&
\cellhead{Cross-sector rollout; SME participation; Runtime SBOM in production}
Simpl MIMs are deployed across the majority of active Common European
Data Spaces with automated cross-provider enforcement. Runtime SBOMs
with cryptographic execution proofs verify supply chain components in
critical infrastructure deployments in at least five Member States.
CAMARA cross-border roaming operates across $\geq$5 EU operator pairs
with documented SLA performance. ``Data Visa'' mechanisms are
operational in certified cross-border AI model migration corridors.
SME-accessible Simpl MIM compliance tooling and Runtime SBOM services
are commercially available from EU-domiciled providers. Unverified
supply chain components are automatically quarantined in isolated
network slices.
\kpistrip{Phase~2 deployment must show geographic breadth ($\geq$5 Member
States), cross-sector presence ($\geq$3 Data Spaces under automated MIM
enforcement), and SME participation in at least one deployment context.
Single-sector rollout does not constitute Phase~2 operational maturity.}
&
\cellhead{SSL: Growing Pan-EU Sovereignty}
EU Data Space infrastructure operating under Simpl achieves $\geq$30\%
non-hyperscaler share in new regulated-sector cloud procurement. Runtime
SBOM ecosystem is commercially mature with multiple competing EU-domiciled
vendors. EU NaaS consortium (under IPCEI-CIS) operates cross-border
compute at pilot commercial scale, demonstrating that EU providers
can federate to compete with hyperscaler geographic coverage. DPP
traceability reduces procurement risk for EU critical infrastructure
buyers by making supply chain provenance verifiable. Foreign hyperscaler
market share in EU regulated sectors declines for the first time from
its $\approx$65\% Phase~1 baseline.
\kpistrip{Phase~2 SSL evidence: $\geq$30\% non-hyperscaler share in new
regulated procurement; first documented market share decline for
non-EU hyperscalers in EU regulated cloud procurement.}
\\
\hline

% ---- GATE 2->3 ----
\rowcolor{gaterow}
\gatelbl{2 $\rightarrow$ 3} &
\multicolumn{4}{>{\footnotesize\raggedright\arraybackslash}p{22.0cm}|}{%
\cellcolor{gaterow}%
\textbf{R\&I trigger:} ZKP verification for supply chain attestation has
achieved real-time per-transaction performance (eliminating the current
batch-only constraint) in at least one production critical infrastructure
deployment, confirmed by independent performance benchmarking.
\textbf{Policy trigger:} A ``Schengen for Cloud'' legal framework --- or
equivalent EU legal instrument resolving US CLOUD Act extra-territorial
jurisdiction conflict for federated EU-sovereign NaaS workloads --- has
been enacted, providing the jurisdictional basis for single-market operation
of federated compute.
\textbf{Operational trigger:} The EU NaaS federation (under IPCEI-CIS or its
successor) operates at commercial scale across $\geq$10 EU telecom operator
members, offering cross-border edge compute services as a standard commercial
product with documented customer deployments outside its founding consortium.
\textbf{Market trigger:} EU supply chain governance frameworks (Runtime SBOM
plus ZKP attestation) are adopted or formally referenced by at least two
non-EU partner jurisdictions as part of digital trade or partnership agreements,
confirming that EU governance leadership is internationally exportable.}
\\
\hline

% ---- PHASE 3 ----
\rowcolor{phase3row}
\textbf{Phase~3}
\newline\textit{Autonomy \&}
\newline\textit{Leadership}
\newline\textit{2035--2040}
\newline\sslbadge{ssltealbg}{ssltealtx}{SSL: Full Federated}
\newline\sslbadge{ssltealbg}{ssltealtx}{Sovereign Autonomy}
&
\cellhead{Cluster ODP: Ubiquitous Operation across all sub-priorities}
All three sub-priorities reach Ubiquitous Operation. Simpl-based
security is the universal standard for cross-provider Data Space
operation in the EU; no alternative interoperability model holds
market relevance in regulated sectors. ZKP-verified supply chain
attestation operates at real-time transaction scale for all critical
infrastructure component verification. The EU NaaS federation operates
as a single commercial entity legally equivalent to a unified
jurisdiction for data workloads. Phase~3 R\&I pursues next-generation
frontiers: post-ZKP cryptographic verification, semantic AI-to-AI
policy negotiation, and legal engineering of the next-generation
``Digital Single Market'' framework beyond cloud-edge compute.
\kpistrip{Phase~3 ODP evidence: Simpl MIM compliance rate in new regulated
EU cloud procurement; ZKP verification throughput in production critical
infrastructure deployments; EU NaaS federation revenue relative to
hyperscaler EU revenue in regulated sectors.}
&
\cellhead{EU federated governance as international model}
The ``Schengen for Cloud'' legal framework is enacted and operational,
resolving US CLOUD Act jurisdiction conflicts for federated EU-sovereign
workloads and enabling single-market legal operation of cross-border
compute. EU Cloud Rulebook is mandatory for all new cloud deployments
in EU regulated sectors without exception. EU supply chain governance
standards (Runtime SBOM plus ZKP) are referenced in international
supply chain security standards (ISO/IEC) and in digital trade
agreements with partner jurisdictions. The EU has transitioned from
rule-taker to rule-setter in international cloud and supply chain
governance.
\kpistrip{Phase~3 governance leadership is verified by two events:
enactment of the ``Schengen for Cloud'' framework (domestic) and
formal adoption of EU supply chain governance standards by $\geq$2
non-EU partner jurisdictions (international).}
&
\cellhead{Provider-agnostic EU trust fabric as continental operational default}
A single provider-agnostic EU trust fabric operates continuously across
all active Common European Data Spaces with real-time automated security
orchestration and zero manual policy coordination. ZKP-verified supply
chain attestation eliminates all unverified components from EU critical
infrastructure deployments in regulated sectors. The EU NaaS federation
is a standard commercial service operating at hyperscaler-comparable
scale, with EU operators collectively competitive on cost and coverage.
Security-by-Design is not an add-on: every new cloud-edge deployment
in the EU incorporates Simpl MIM compliance, Runtime SBOM verification,
and cross-border sovereignty assurance as architectural defaults.
\kpistrip{Phase~3 deployment normalisation: Simpl MIM compliance in
$\geq$90\% of new regulated EU cloud deployments; ZKP supply chain
verification in 100\% of new critical infrastructure component
procurement; EU NaaS federation operating as a single commercial entity
with documented non-founding-member customers.}
&
\cellhead{SSL: Full Federated Sovereign Autonomy}
Non-EU hyperscaler market share in EU regulated cloud procurement
falls below 40\% for new deployments (from $\approx$65\% at baseline),
with a binding trajectory toward continued reduction. EU NaaS
federation achieves hyperscaler-comparable geographic coverage and
commercial viability across the EU. Supply chain governance
internationally adopted: at least two partner jurisdictions have
formally adopted EU Runtime SBOM plus ZKP attestation standards as
domestic procurement requirements. The EU has removed the structural
dependency that allowed US CLOUD Act jurisdiction to apply to
EU-regulated data: federated sovereignty is legally guaranteed, not
only technically implemented.
\kpistrip{Full Federated Sovereign Autonomy requires three
concurrent conditions: $<$40\% hyperscaler new-procurement share;
``Schengen for Cloud'' legally enacted; $\geq$2 partner jurisdictions
formally adopting EU supply chain governance standards. Achieving
any two without the third does not constitute this SSL level.}
\\
\hline

\end{longtable}
\end{landscape}

% ---- PROGRESSION MATRIX: CLUSTER 3 ----


% \noindent
% \resizebox{\textwidth}{!}{%
% \begin{tabular}{|p{4.2cm}|p{5.6cm}|p{5.6cm}|p{5.6cm}|}
% \hline
% \rowcolor{gray!30}
% \multicolumn{4}{|c|}{\textbf{Progression Matrix: Priority Cluster 3 ---
%   Federated Operational Sovereignty (P4, P5, P7)}} \\
% \hline
% \rowcolor{gray!15}
% \textbf{Metric} &
% \textbf{Phase 1: Foundation \& Compliance (2025--2030)} &
% \textbf{Phase 2: Federation \& Standardisation (2030--2035)} &
% \textbf{Phase 3: Autonomy \& Leadership (2035--2040)} \\
% \hline
% \rowcolor{gray!10}
% \multicolumn{4}{|l|}{\textit{\textbf{Cluster-Level Progression Indicators}}} \\
% \hline
% \textbf{Cluster ODP} &
% Fragmented Piloting &
% Standardised Integration &
% Ubiquitous Operation \\
% \hline
% \textbf{Cluster SSL} &
% Emerging EU Ecosystem &
% Growing Pan-EU Sovereignty &
% Full Federated Sovereign Autonomy \\
% \hline
% \rowcolor{gray!10}
% \multicolumn{4}{|l|}{\textit{\textbf{Sub-Priority Decomposition (Individual ODP / SSL per Phase)}}} \\
% \hline
% \textit{P4: Security Interoperability \newline via Simpl Middleware} &
% ODP: Fragmented Piloting \newline
% SSL: Emerging EU Ecosystem &
% ODP: Standardised Integration \newline
% SSL: Full Sovereign Autonomy &
% ODP: Ubiquitous Operation \newline
% SSL: Full Sovereign Autonomy \\
% \hline
% \textit{P5: Supply Chain \newline Verification \& Infra.} &
% ODP: Full Supply Chain Disclosure \newline
% SSL: High Reliance on \newline Fragmented Non-EU Auditing &
% ODP: Partial Restriction \& \newline Dynamic Runtime Verification \newline
% SSL: Emerging EU Supply Chain \newline Tracking Frameworks &
% ODP: Full Restriction Process \& \newline ZKP-Verified Attestation \newline
% SSL: Full Sovereign Supply \newline Chain Control \& Transparency \\
% \hline
% \textit{P7: Pan-European NaaS \newline Federation} &
% ODP: Bilateral Interconnection \newline
% SSL: Hyperscaler Platform Dominance &
% ODP: Federated Edge Roaming \newline
% SSL: Pan-EU Consortium Competes &
% ODP: True Digital Single Market \newline
% SSL: European Scale Parity Achieved \\
% \hline
% \rowcolor{blue!7}
% \textbf{Operational State} &
% Initial Simpl security pilots are deployed in three Common European
% Data Spaces with manual policy coordination and proof-of-concept
% MIM implementations; CRA-baseline SBOM declarations are submitted
% by hardware and software vendors as static point-in-time documents
% without runtime verification; harmonised CAMARA APIs are tested in
% bilateral 5G corridor agreements between individual EU telecom
% operators; hyperscalers retain dominant market share through
% proprietary API ecosystems. &
% Automated policy translation via standardised Simpl MIMs is
% deployed across the majority of active Common European Data Spaces;
% Runtime SBOMs with cryptographic execution proofs verify
% containerised components dynamically as they migrate across
% multi-provider continuum nodes; harmonised CAMARA-based cross-border
% edge roaming operates across five or more EU operator pairs;
% ``Data Visa'' mechanisms enable AI model migration without
% manual legal review in certified corridor deployments. &
% A single provider-agnostic European trust fabric operates
% continuously across all EU Common Data Spaces with real-time
% automated security orchestration; ZKP-verified supply chain
% attestation eliminates unverified components from all
% EU critical infrastructure deployments; a unified NaaS
% platform legally operable as a single data jurisdiction
% enables EU operators to achieve hyperscaler-comparable
% scale; EU supply chain governance frameworks are adopted
% as reference implementations internationally. \\
% \hline
% \rowcolor{orange!12}
% \textbf{Critical Enabler} &
% Simpl-Open official launch with initial security MIM definitions
% published as open standards, establishing the technical
% baseline for cross-provider security policy automation &
% Security MIMs adopted as CEN/CENELEC standards with
% mandatory reference in EU Cloud Rulebook, creating the
% regulatory basis for non-voluntary hyperscaler compliance &
% ``Schengen for Cloud'' legal framework enacted, resolving
% US CLOUD Act extra-territorial jurisdiction conflict and
% enabling single-jurisdiction data workload operation at
% continental scale \\
% \hline
% \end{tabular}%
% }

% =========================================================================
% CROSS-PRIORITY BOTTLENECK ASSESSMENT
% =========================================================================

\newpage
\section{Cross-Priority Bottleneck Assessment}
\label{sec:bottlenecks}

Three shared structural dependencies determine whether individual sub-priority
progressions can advance on their stated timelines, regardless of cluster-level investment.
These bottlenecks are identified here because they constitute cross-cutting gate conditions
that cannot be resolved by any single cluster in isolation.

\textbf{Bottleneck 1 --- EU Chips Act Sovereign Fabrication Capacity.}
The Phase~2 gate conditions for P1 (EUCC `High' certified RISC-V secure elements),
P2 (PQC-native RISC-V silicon), P3 (zero-energy secure hardware), and P8 (RISC-V
Keystone TEE scaling) all converge on the same physical dependency: sovereign EU
fabrication capacity for security-critical silicon at commercially viable yields.
The Phase~2 critical enabler for P1 (``EUCC `High' certification of first
EU-manufactured RISC-V secure element'') and the Phase~3 enabler (``EU Chips Act
sovereign EUV-capable fabrication achieving $\geq$60\% self-sufficiency'') constitute
a single two-stage gate that controls the SSL trajectory of four of the five Cluster~1
sub-priorities simultaneously.

\textbf{Bottleneck 2 --- AI Act Conformity Assessment Infrastructure.}
The Phase~2 gate condition for Cluster~2 (``standardised AI Act conformity assessment
framework for autonomous security agents'') is a prerequisite not only for P6 and P10
but also for any AI-driven capability within Cluster~1 (autonomous Island Mode swarm
protocols) and Cluster~3 (Automated Compliance Negotiators in NaaS federation) that
incorporates autonomous decision-making. The absence of EU AI Act conformity assessment
bodies with sufficient capacity to certify high-risk security AI at scale represents a
systemic bottleneck that, if unresolved by 2028--2030, will delay the Phase~2 SSL
transition across all three clusters.

\textbf{Bottleneck 3 --- Runtime SBOM and Supply Chain Attestation Infrastructure.}
Runtime SBOM verification (P5) is a prerequisite for both the Cluster~1 TEE integrity
claim (P8 confidential computing pipelines cannot be trusted without verified component
provenance) and the Cluster~3 MIM compliance model (Simpl-based interoperability
cannot provide cryptographic enforcement guarantees if component integrity is
unverifiable). The Phase~1 critical enabler for Cluster~3 (Simpl-Open MIM definitions)
and the Phase~1 ODP for P5 (full supply chain disclosure) must therefore proceed in
coordination, not sequentially, with the Phase~1 applied research activities in
Cluster~1. This interdependency must be explicitly reflected in Horizon Europe project
call design to ensure that SBOM infrastructure investment is not siloed within supply
chain security calls but is structurally integrated into hardware security and
interoperability funding instruments.

% =========================================================================
% NEW BIBLIOGRAPHY ENTRIES
% =========================================================================

% Add the following entries to the master \begin{thebibliography}{99} block:
 in master document
%
% Required packages (add to master preamble if not present):
%   \usepackage[table]{xcolor}
%   \usepackage{multirow}
%   \usepackage{array}
%
% New bibliography keys required (entries provided at end of file):
%   fips203, fips204, fips205, nistir8413, nistsp800207,
%   etsi103744, etsizsmoo9, tgcpm20, isoiec11889,
%   opentitan, keystone, threeGPP23288, threeGPP33501, dora, gdpr
% =========================================================================

% ---- Shared colors/macro helpers when included from a master document ----
\providecolor{gaterow}{RGB}{240,240,238}
\providecolor{phase1row}{RGB}{237,245,253}
\providecolor{phase2row}{RGB}{237,249,244}
\providecolor{phase3row}{RGB}{245,244,254}
\providecolor{colheader}{RGB}{245,245,243}
\providecolor{clusterheader}{RGB}{220,232,248}

\providecolor{sslredbg}{RGB}{252,235,235}
\providecolor{sslredtx}{RGB}{163,45,45}
\providecolor{sslamberbg}{RGB}{250,238,218}
\providecolor{sslambtx}{RGB}{133,79,11}
\providecolor{ssltransbg}{RGB}{230,241,251}
\providecolor{ssltrantx}{RGB}{24,95,165}
\providecolor{ssltealbg}{RGB}{225,245,238}
\providecolor{ssltealtx}{RGB}{15,110,86}

\providecommand{\sslbadge}[3]{%
      \parbox{\linewidth}{\raggedright
            \colorbox{#1}{\scriptsize\parbox{0.92\linewidth}{\raggedright
                  \textbf{\textcolor{#2}{#3}}}}}}

\providecommand{\kpistrip}[1]{%
      \par\vspace{2pt}\noindent\rule{\linewidth}{0.3pt}%
      \par\noindent\textit{\scriptsize #1}}

\providecommand{\cellhead}[1]{\noindent\textbf{#1}\par\smallskip}
\providecommand{\gatelbl}[1]{\scriptsize\textit{Gate~#1:}}
\providecommand{\subp}[1]{\textit{\scriptsize(#1)}}

\section{Strategic Priority Clusters}
\label{ch:strategic_priority_clusters}

% \subsection{Consolidation Rationale and Structural Logic}
% \label{sec:consolidation_rationale}

% The ten executive priorities identified through the 4-Filter Prioritisation Framework represent
% distinct technology domains, but they share deep structural dependencies and converge on three
% strategic imperatives that define Europe's 15-year cybersecurity trajectory for the Cloud-Edge-IoT
% continuum. Maintaining ten parallel profiles in a dual-audience deliverable risks fragmenting
% the narrative for European Commission technical evaluators while obscuring the systemic
% interdependencies that determine whether individual priorities succeed or fail.

% The consolidation presented in this chapter reorganises the ten priorities into three strategic
% clusters, each with a unifying architectural thesis. To preserve the analytical precision required
% by IEEE Access peer review, each cluster retains its constituent sub-priorities as explicitly
% labelled components within a nested progression matrix, allowing individual technology maturation
% trajectories to be tracked independently while the cluster-level framing communicates strategic
% coherence to policy audiences.

% The mapping is as follows:
% \begin{itemize}
%   \item \textbf{Priority Cluster 1 --- Sovereign Trusted Infrastructure:}
%         P1 (Hardware Roots of Trust), P2 (PQC Agility), P3 (Lightweight
%         Cryptography and Physical-Layer Security), P8 (Confidential
%         Computing), P9 (Resilient-by-Design / Island Mode).
%   \item \textbf{Priority Cluster 2 --- Cognitive Security and Autonomous
%         Cyber Defence:} P6 (Cognitive SecOps / LLMSecOps), P10 (Proactive
%         APT Management), and the AI-as-target dimension drawn from
%         Cross-Layers CL1.2, CL1.3, and CL4.7.
%   \item \textbf{Priority Cluster 3 --- Federated Operational Sovereignty:}
%         P4 (Security Interoperability via Simpl), P5 (Supply Chain
%         Verification), P7 (Pan-European NaaS Federation).
% \end{itemize}

% One critical cross-priority dependency must be made explicit before the cluster profiles are
% presented. Hardware supply chain attestation (P5) and Runtime SBOM verification constitute
% a shared gate condition for both Cluster 1 (TEE and secure element integrity) and Cluster 3
% (federated supply chain governance). Progression in P8 (Confidential Computing) is contingent
% on the availability of trustworthy RISC-V secure elements from P1; P3 (Lightweight PhySec)
% and P2 (PQC) both depend on EU Chips Act fabrication capacity as a common Phase 2 enabler.
% These shared bottlenecks are identified explicitly in the Cross-Priority Bottleneck assessment
% in Section~\ref{sec:bottlenecks} and traced to specific progression matrix gate conditions.

% =========================================================================
% PRIORITY CLUSTER 1: SOVEREIGN TRUSTED INFRASTRUCTURE
% =========================================================================

% \newpage
\section{Priority Cluster 1: Sovereign Trusted Infrastructure}
\label{sec:cluster1}

\noindent\textbf{Constituent Sub-Priorities:} P1 (Hardware Roots of Trust \& EUCC `High'),
P2 (PQC Agility for Edge/IoT), P3 (Lightweight Cryptography \& Physical-Layer Security),
P8 (Confidential Computing \& Privacy-Preserving AI), P9 (Resilient-by-Design \& Island Mode).

\medskip
\noindent\textbf{Cluster Thesis:} Before any federated or AI-driven security capability can be
trusted, the substrate on which it executes must provide cryptographically verifiable integrity,
quantum-resistant confidentiality, and autonomous survivability independent of external
infrastructure. This cluster defines the verifiable hardware and cryptographic foundation that
all higher-order capabilities in Clusters 2 and 3 depend upon as an architectural prerequisite.

\subsection*{7-Dimensional Analysis: Priority Cluster 1}

\begin{enumerate}

% ---- DIMENSION 1 ----
\item \textbf{Baseline \& Target Objective}

\textit{Baseline (2025--2026):}
European critical infrastructure and edge deployments operate under high foreign dependency
across all five constituent sub-priorities. Hardware roots of trust are dominated by non-EU
Trusted Platform Module (TPM) manufacturers --- Infineon SLB9670, NXP SE050, and
STMicroelectronics ST33 series collectively account for the majority of the EU market under
ISO/IEC~11889~\cite{isoiec11889} and TCG TPM~2.0 specifications~\cite{tgcpm20}. The Trusted
Execution Environment (TEE) market is dominated by Intel TDX, AMD SEV-SNP, and ARM TrustZone,
none of which are EU-designed or manufactured. Post-Quantum Cryptography (PQC) standardisation
reached a definitive milestone with NIST's publication of FIPS~203 (ML-KEM)~\cite{fips203},
FIPS~204 (ML-DSA)~\cite{fips204}, and FIPS~205 (SLH-DSA)~\cite{fips205} in August~2024;
however, hardware-optimised implementations for 8/16-bit microcontrollers remain at the
Applied~R\&D stage, as confirmed by NIST IR~8413~\cite{nistir8413}. Physical-layer security
for extreme-edge mMTC devices lacks standardised specifications. Edge nodes universally operate
with hard cloud dependencies and no autonomous fallback architectures meeting the operational
resilience mandates of NIS2~\cite{nis2} and DORA~\cite{dora}.

Current cluster ODP: Applied R\&D (driven by PQC, Lightweight PhySec, and Confidential
Computing sovereign alternatives). Current cluster SSL: High Foreign Dependency.

\textit{Target (2040):}
At least 60\% of EU critical infrastructure edge nodes deploy EU-designed RISC-V open-hardware
secure elements certified at EUCC `High' assurance level~\cite{eucc}. 100\% of new EU critical
IoT deployments support hybrid or full FIPS~203/204/205-compliant PQC. Zero-energy nodes
execute quantum-resistant device authentication within a sub-100\,$\mu$J energy budget per
operation. Trusted Execution Environments built on EU open-hardware architectures (Keystone~\cite{keystone},
OpenTitan~\cite{opentitan}) cover at least 50\% of new EU cloud-edge deployments. All NIS2-essential
and DORA-critical edge nodes demonstrate validated 72-hour autonomous Island Mode survivability
independent of WAN connectivity.

% ---- DIMENSION 2 ----
\item \textbf{Primary Drivers}

\textit{Threat Drivers:}
\begin{itemize}
  \item State-sponsored APT exploitation of JTAG debug interfaces and electromagnetic
        side-channel analysis (EMSCA) on field-deployed edge nodes; the Spectre and Meltdown
        class of microarchitectural side-channel vulnerabilities establishes that even
        production-certified hardware can be compromised at scale, and far-edge nodes
        operating without physical security controls present a materially higher exposure
        surface than hyperscaler data centres.
  \item ``Harvest Now, Decrypt Later'' (HNDL) campaigns: state-level adversaries are
        actively archiving encrypted EU government and industrial communications under
        the expectation of Cryptographically Relevant Quantum Computer (CRQC)
        availability. NIST IR~8413 places the Q-Day horizon in the 10--20 year range
        with a highly uncertain lower bound~\cite{nistir8413}; HNDL renders this a
        present-tense threat against data with multi-decade confidentiality
        requirements.
  \item Physical tampering and environmental attacks on unattended far-edge nodes
        (smart grid substations, O-RAN radio units, industrial gateways) where the
        physical security controls available in hyperscaler facilities are absent.
  \item Cascading WAN disconnection failures propagating through centrally-dependent edge
        architectures, exploited by adversaries to degrade critical infrastructure
        services during hybrid conflict scenarios.
\end{itemize}

\textit{Regulatory Drivers:}
\begin{itemize}
  \item CRA Articles~10 and~11: mandatory vulnerability handling and security update
        obligations for products with digital elements; enforcement begins in 2027 for
        most product categories~\cite{cra}.
  \item NIS2 Article~21(2)(e): cryptographic policies and key management for essential
        and important entities; full enforcement began October 2024~\cite{nis2}.
  \item DORA Article~9: ICT resilience testing including hardware-level resilience
        requirements for financial sector critical infrastructure~\cite{dora}.
  \item AI Act Article~9: risk management system requirements for high-risk AI
        systems; critical infrastructure AI components require hardware integrity
        guarantees that only EUCC `High' certified secure elements can
        provide~\cite{aiact}.
\end{itemize}

% ---- DIMENSION 3 ----
\item \textbf{High Impact Activities}

\begin{enumerate}[label=(\alph*)]
  \item \textbf{[Phase~1]} Develop sovereign European Secure Elements based on
        RISC\-V open-hardware (building on OpenTitan~\cite{opentitan} and
        Keystone~\cite{keystone} frameworks), targeting EUCC `High' assurance level
        evaluation; implement automated hardware-level remote attestation protocols
        interoperable with TPM~2.0~\cite{tgcpm20} and ISO/IEC~11889~\cite{isoiec11889}
        for heterogeneous large-scale edge fleets.

  \item \textbf{[Phase~1--2]} Optimise NIST FIPS~203 (ML-KEM)~\cite{fips203} and
        FIPS~204 (ML-DSA)~\cite{fips204} for 8/16-bit microcontrollers under severe
        energy and memory constraints (target: ML-KEM-512 within 32\,KB RAM on
        Cortex-M0+ class devices); develop crypto-agility OTA update middleware enabling
        algorithm rotation without hardware replacement, compatible with ETSI TS~103~744
        hybrid key exchange profiles~\cite{etsi103744}.

  \item \textbf{[Phase~1--2]} Implement ultra-lightweight PQC and keyless
        Physical-Layer Security (RF fingerprinting exploiting Channel State Information
        and Carrier Frequency Offset, CSI/CFO-based device authentication) for mMTC
        devices operating under zero-energy or near-zero-energy constraints; target
        integration with 6G ISAC (Integrated Sensing and Communication) interfaces per
        SNS~JU~SRIA~\cite{snsju}.

  \item \textbf{[Phase~2]} Deploy multi-party confidential computing pipelines
        using TEE abstraction layers (Keystone enclaves on RISC-V~\cite{keystone})
        combined with Fully Homomorphic Encryption (FHE) schemes for collaborative AI
        training across multi-provider nodes without raw data exposure; align with
        EUCS `High+' certification requirements~\cite{eucs}.

  \item \textbf{[Phase~2--3]} Engineer autonomous Island Mode architecture: soft-state
        edge orchestration using cached credentials with 72-hour validity windows,
        LoRaWAN LPWAN fallback micro-slices, eBPF-based kernel-level policy enforcement
        for autonomous node isolation, and bio-inspired self-healing swarm protocols
        enabling edge nodes to sustain critical service delivery without cloud
        connectivity or human intervention.
\end{enumerate}

% ---- DIMENSION 4 ----
\item \textbf{Catalysts \& Enablers}

\textit{Funding Alignment:}
Horizon Europe Cluster~4 (Digital, Industry and Space) for sovereign secure silicon and
PQC hardware implementation; Horizon Europe Cluster~3 (Civil Security for Society) for
edge resilience and Island Mode architectures; EU Chips Act (Regulation (EU) 2023/1781)
pilot fabrication lines for security-critical silicon via the Chips Joint Undertaking
(KDT~JU); ECCC Strategic Agenda priority areas ``Secure Components and Systems'' and
``Post-Quantum Cryptography transition infrastructure''~\cite{eccc}; Digital Europe
Programme for PQC migration tooling and crypto-agility middleware.

\textit{Standardisation:}
NIST NCCoE ``Migration to Post-Quantum Cryptography'' project (reference architecture
and migration playbooks); CEN/CENELEC JTC~13/WG~8 (hardware security standards aligned
with EUCC Common Criteria); ETSI TC CYBER (ETSI TS~103~744 quantum-safe hybrid key
exchanges~\cite{etsi103744}); ENISA PQC transition guidelines and cryptographic inventory
tooling; 3GPP SA3 for physical-layer security integration with 5G/6G
specifications~\cite{threeGPP33501}.

\textit{Ecosystem Enablers:}
RISC-V International open-hardware community; OpenTitan coalition pilot deployments
for EU sovereign RoT~\cite{opentitan}; SNS~JU pilot networks for 6G-integrated PhySec
validation~\cite{snsju}; federated Cyber Range infrastructure (CL3.5) for Island Mode
and TEE attestation stress-testing.

% ---- DIMENSION 5 ----
\item \textbf{Disruptors \& Systemic Risks}

\begin{itemize}
  \item \textbf{Q-Day cliff:} CRQC becoming operational before EU PQC migration
        completes; NIST IR~8413 places this horizon at 10--20 years with a highly
        uncertain lower bound~\cite{nistir8413}. The EU PQC migration timeline targeting
        critical infrastructure completion by 2030 has limited buffer against the
        pessimistic lower bound, and HNDL campaigns make this a present-tense threat
        against long-lived secrets.
  \item \textbf{EU Chips Act EUV lithography shortfall:} RISC-V secure element
        production at the node geometries required for EUCC `High' certification depends
        on advanced lithography; EU 2030 Chips Act targets (20\% global market share)
        do not specify security-critical silicon capacity independently, and ASML EUV
        tooling remains subject to US--Netherlands export control dynamics that could
        constrain sovereign fabrication timelines beyond 2028.
  \item \textbf{Side-channel analysis breakthroughs against FIPS~203 software
        implementations:} lattice-based algorithms have demonstrated timing and
        cache-timing attack surfaces in software implementations on constrained devices;
        hardware countermeasures are not yet standardised at the time of writing
        [CITATION REQUIRES VERIFICATION --- recent SCA against ML-KEM implementations].
  \item \textbf{TEE architectural monoculture:} if EU open-hardware TEE adoption
        concentrates on a single architecture (e.g., Keystone only~\cite{keystone}), a
        single discovered vulnerability propagates across the entire EU critical
        infrastructure TEE estate without architectural diversity as a containment layer.
  \item \textbf{LPWAN regulatory fragmentation:} Island Mode fallback using LoRaWAN
        operates on unlicensed 868\,MHz spectrum in the EU; spectrum policy
        fragmentation across Member States could prevent coordinated emergency
        micro-slice activation under unified crisis response scenarios.
\end{itemize}

% ---- DIMENSION 6 ----
\item \textbf{Failsafe Mechanisms \& Resilience}

\begin{itemize}
  \item \textbf{Hybrid cryptographic transition:} deploy ML-KEM (FIPS~203)~\cite{fips203}
        combined with ECDH-P384 in a hybrid key encapsulation scheme per
        ETSI TS~103~744~\cite{etsi103744}; if ML-KEM is found vulnerable via cryptanalytic
        breakthrough, ECDH-P384 maintains classical confidentiality while SLH-DSA
        (FIPS~205)~\cite{fips205} provides backup hash-based signature integrity with
        distinct mathematical foundations.
  \item \textbf{TEE implementation diversity:} deploy a minimum of two architecturally
        distinct TEE technologies across critical EU infrastructure (ARM CCA on
        server-class hardware and RISC-V Keystone~\cite{keystone} on edge nodes) to
        prevent single-architecture compromise from cascading across the entire
        confidential computing estate.
  \item \textbf{Autonomous node quarantine via eBPF:} if TEE remote attestation fails,
        eBPF-based network policy enforcement at the Linux kernel level (e.g., via
        Cilium or Falco eBPF rulesets) isolates the affected node from the orchestration
        plane within sub-second latency, blocking lateral movement before human
        intervention is available.
  \item \textbf{Island Mode failsafe:} edge nodes maintain cryptographically signed
        cached credentials with defined 72-hour validity windows; on WAN disconnection,
        a LoRaWAN LPWAN micro-slice activates a minimal viable network carrying
        life-safety telemetry; if LPWAN also fails, nodes revert to local-only operation
        with integrity-protected audit logging for post-incident forensics.
  \item \textbf{RF fingerprinting fallback:} for zero-energy IoT nodes where
        cryptographic processing exceeds the available energy budget, Physical-Layer
        authentication via RF fingerprinting (CSI/CFO-based device identity) provides
        keyless device verification using intrinsic wireless channel characteristics,
        requiring no cryptographic key storage or computation.
\end{itemize}

% ---- DIMENSION 7 ----
\item \textbf{Transversal Meta-Layer}

\textit{Sovereignty \& Autonomy:}
This cluster eliminates approximately 100\% foreign dependency on non-EU TPMs (Infineon,
NXP, STMicroelectronics dominate the EU market under TCG TPM~2.0~\cite{tgcpm20}) and
non-EU TEE architectures (Intel TDX, AMD SEV-SNP, ARM TrustZone, all non-EU designed and
manufactured). Crypto-agility frameworks additionally reduce dependency on US-controlled
NIST algorithm supply chains by enabling algorithm substitution if standardisation processes
are constrained by geopolitical factors.

\textit{Trust \& Ethics:}
GDPR Article~25 (data protection by design and by default) is operationalised through
TEE-enforced data minimisation and in-use encryption in multi-tenant environments,
ensuring that sensitive data processing cannot be accessed by platform operators~\cite{gdpr}.
AI Act Article~9 risk management system requirements for high-risk AI components require
hardware integrity guarantees --- only EUCC `High' certified secure elements provide the
mathematically verifiable trust anchor this obligation demands~\cite{aiact}.

\textit{Sustainable Security (Green IT):}
CRYSTALS-Kyber (ML-KEM) demonstrates approximately 1.5$\times$ computational overhead
relative to ECDH at equivalent security levels on ARM Cortex-M4, per NIST IR~8413~\cite{nistir8413},
representing a manageable energy penalty. Lightweight PQC research targets sub-100\,$\mu$J
per authentication operation for zero-energy IoT nodes. EU-sovereign RISC-V secure
elements designed for edge power constraints enable more energy-efficient cryptographic
operations than general-purpose processor TPM emulation, reducing per-device carbon
footprint at continental deployment scale.

\end{enumerate}



% =========================================================================
% CLUSTER 1: SOVEREIGN TRUSTED INFRASTRUCTURE
% =========================================================================
\newpage

\subsection{Progression Matrix: Priority Cluster 1 --- Sovereign Trusted Infrastructure}

\begin{landscape}
\setlength{\textwidth}{26.5cm}\setlength{\linewidth}{26.5cm}%
\setlength{\hsize}{26.5cm}\setlength{\columnwidth}{26.5cm}
\setlength{\tabcolsep}{5pt}
\renewcommand{\arraystretch}{1.35}

\begin{longtable}{%
|>{\footnotesize\raggedright\arraybackslash}p{2.8cm}%
|>{\footnotesize\raggedright\arraybackslash}p{5.4cm}%
|>{\footnotesize\raggedright\arraybackslash}p{5.4cm}%
|>{\footnotesize\raggedright\arraybackslash}p{5.4cm}%
|>{\footnotesize\raggedright\arraybackslash}p{5.4cm}|}

\caption{\textbf{Cluster~1: Sovereign Trusted Infrastructure} (P1, P2, P3, P8, P9).
Transition matrix across the four methodology dimensions. Sub-priority labels in
parentheses indicate where individual priorities diverge from the cluster-level
descriptor within a cell. ODP scale: Applied R\&D $\to$ Fragmented Piloting
$\to$ Standardised Integration $\to$ Ubiquitous Operation. SSL scale:
High Foreign Dependency $\to$ Emerging EU Ecosystem $\to$ Full Sovereign Autonomy.}
\label{tab:cluster1}\\

\hline
\rowcolor{colheader}
\textbf{Phase / SSL} &
\textbf{R\&I Maturity}{\scriptsize ODP indicator per sub-priority} &
\textbf{Policy, Regulatory \& Standards}{\scriptsize enforcement $\cdot$
  certification $\cdot$ standardisation} &
\textbf{Operational Deployment \& Adoption}{\scriptsize actor types $\cdot$
  scale $\cdot$ sector coverage} &
\textbf{Market Uptake \& Strategic Sovereignty}{\scriptsize SSL indicator
  $\cdot$ dependency reduction $\cdot$ industrial capacity} \\
\hline
\endfirsthead

\hline
\rowcolor{colheader}
\textbf{Phase / SSL} &
\textbf{R\&I Maturity} &
\textbf{Policy, Regulatory \& Standards} &
\textbf{Operational Deployment \& Adoption} &
\textbf{Market Uptake \& Strategic Sovereignty} \\
\hline
\endhead
\hline\multicolumn{5}{r}{\footnotesize\textit{Continued\ldots}}\\ \endfoot
\hline \endlastfoot

% ---- PHASE 1 ----
\rowcolor{phase1row}
\textbf{Phase~1}
\newline\textit{Foundation \&}
\newline\textit{Compliance}
\newline\textit{2025--2030}
\newline\sslbadge{sslredbg}{sslredtx}{SSL: High Foreign Dependency}
&
\cellhead{Cluster ODP: Applied R\&D / Fragmented Piloting}
The cluster is split at baseline. \subp{P1, P9} reach Fragmented
Piloting: RISC-V secure element projects targeting EUCC `High' are
funded under Horizon Europe and EU Chips Act; Island Mode soft-state
protocols are validated in 5G research testbeds.
\subp{P2, P3, P8} remain at Applied R\&D: NIST FIPS~203/204/205
algorithms are being optimised for 8/16-bit microcontrollers under
severe energy and memory constraints; Keystone and OpenTitan
prototypes are under EUCC pre-evaluation; RF fingerprinting
Physical-Layer Security is validated in controlled 5G lab environments;
Fully Homomorphic Encryption approaches practical latency bounds in
research settings only.
\kpistrip{The cluster ODP is Applied R\&D, driven by the majority of
sub-priorities. Deployment claims for P2, P3, P8 must be grounded in
named Horizon Europe or ECCC project outcomes, not technology
previews.}
&
\cellhead{Regulatory environment: enforcing, certification nascent}
NIS2 Article~21(2)(e) cryptographic policies for essential entities
entered full enforcement in October~2024. CRA Articles~10 and~11
(vulnerability handling and security update obligations) begin
enforcement in 2027, creating immediate procurement-level demand for
hardware security certification. AI~Act Articles~9 and~15 establish
hardware integrity requirements for high-risk AI systems and enter
applicability in 2025--2026. ENISA PQC migration guidelines are
published; EUCC evaluation processes are initiated for the first RISC-V
prototypes, but no EU-specific EUCC~`High' certified sovereign secure
element exists yet. CEN/CENELEC JTC~13 standards for hardware security
MIMs are in draft phase.
\kpistrip{The regulatory framework is enforcing existing obligations
but lacks EU-specific hardware certification outputs; the critical gap
is the absence of a ratified EUCC~`High' EU-origin product.}
&
\cellhead{Deployment: research consortia and national agencies only}
RISC-V secure element pilots are operated by 3--5 Horizon Europe
consortia and national cybersecurity agencies in controlled
environments. Intel~TDX, AMD~SEV-SNP, and ARM~TrustZone dominate
all operational TEE deployments across EU cloud and critical
infrastructure. PQC is confined to cloud-side software implementations
in pilot contexts; far-edge IoT devices run classical cryptography
exclusively. Edge nodes operate with hard cloud dependencies and no
validated autonomous fallback architecture meeting DORA or NIS2
business continuity thresholds. SMEs have no access to EU-sovereign
hardware alternatives at viable cost.
\kpistrip{Operational deployment at Phase~1 is institutional and
research-concentrated. The gap between pilot existence and operational
availability for regulated-sector procurement must be explicitly
acknowledged.}
&
\cellhead{SSL: High Foreign Dependency}
Non-EU TPMs (Infineon~SLB9670, NXP~SE050, STMicroelectronics~ST33)
hold $\geq$90\% of EU market share in hardware security modules.
Intel, AMD, and ARM control the entire operational TEE market. No
EU-designed and EU-manufactured secure element is commercially
available. EU Chips Act fabrication investment for security-critical
silicon is initiated but pre-production. US NIST algorithm supply chain
is the sole PQC standardisation reference. Structural vulnerability is
direct: geopolitical restriction of access to any of the three major
TPM suppliers would immediately impair EU critical infrastructure
attestation capability.
\kpistrip{Quantified dependency at Phase~1 baseline: $\geq$90\% foreign
market share, zero EU-certified sovereign hardware alternatives.
Sovereignty progress is measured from this specific figure.}
\\
\hline

% ---- GATE 1->2 ----
\rowcolor{gaterow}
\gatelbl{1 $\rightarrow$ 2} &
\multicolumn{4}{>{\footnotesize\raggedright\arraybackslash}p{22.0cm}|}{%
\cellcolor{gaterow}%
\textbf{R\&I trigger:} At least one EU-designed RISC-V secure element prototype
has been independently validated against EUCC~`High' evaluation criteria and has
entered the formal ENISA/national scheme certification process.
\textbf{Policy trigger:} CRA Articles~10 and~11 are in active enforcement with
documented non-compliance decisions, confirming that market incentives for
certified EU hardware alternatives are financially material.
\textbf{Operational trigger:} Documented operational deployment of hybrid
ML-KEM (FIPS~203) plus ECDH-P384 PQC schemes in at least one EU critical
infrastructure sector (energy or financial) at production scale, not pilot scale.
\textbf{Market trigger:} EU Chips Act security-critical silicon fabrication
pilot line has produced and delivered its first commercially viable batch of
EU-origin secure elements, confirmed by independent yield and cost assessment.}
\\
\hline

% ---- PHASE 2 ----
\rowcolor{phase2row}
\textbf{Phase~2}
\newline\textit{Applied \& Industrial}
\newline\textit{R\&D alongside}
\newline\textit{Federation \&}
\newline\textit{Standardisation}
\newline\textit{2030--2035}
\newline\sslbadge{ssltransbg}{ssltrantx}{SSL: Emerging EU Ecosystem}
&
\cellhead{Cluster ODP: Standardised Integration / continued Applied R\&D}
\subp{P1, P9} reach Standardised Integration: the first EUCC~`High' certified
EU-manufactured RISC-V secure elements are in production and being integrated
into critical infrastructure; automated remote attestation at fleet scale is
validated; 72-hour Island Mode achieves standardised certification under
3GPP/ETSI fallback API profiles.
\subp{P2, P3} transition from Applied R\&D to Fragmented Piloting: hybrid
ML-KEM plus ECDH-P384 schemes are validated at Industrial IoT scale; lightweight
PQC for sub-100~$\mu$J zero-energy nodes is demonstrated under 6G ISAC testbeds.
\subp{P8} reaches Fragmented Piloting: RISC-V Keystone TEEs are operational
in multi-provider Data Spaces; FHE performance approaches the threshold for
selective production use in high-value financial analytics.
\kpistrip{Phase~2 R\&I is not purely consolidation: three of five sub-priorities
are still in active applied and industrial R\&D. Describing this phase as
``standardisation only'' misrepresents the maturation arc.}
&
\cellhead{EUCC certification and PQC standards formalised}
The EUCC~`High' certification of the first EU-manufactured RISC-V secure
element is the definitive Phase~2 policy gate: it transforms applied research
into a market-accessible, legally certifiable product and unlocks the procurement
pathway for EU critical infrastructure operators. NIST FIPS~203/204/205 are
adopted into ETSI~TS profiles for EU deployment contexts. ENISA publishes a
mandatory PQC migration timeline for critical infrastructure operators, referenced
in NIS2 sector-specific implementing acts. EU Chips Act fabrication investments
are active and measurable. CEN/CENELEC hardware security standards for EU-origin
secure elements enter ratification.
\kpistrip{The single Phase~2 policy gate is EUCC~`High' certification of a
EU-origin product. All other policy milestones are enabling conditions, not
gate conditions.}
&
\cellhead{Critical infrastructure and regulated sectors deploying EU alternatives}
EU EUCC~`High' certified secure elements are deployed in energy sector
substations and telecommunications critical infrastructure in at least five
Member States. Hybrid PQC schemes are standard in new Industrial IoT and
financial sector (DORA-compliant) deployments. RISC-V Keystone TEEs operate
in federated multi-provider Data Spaces alongside existing Intel TDX
infrastructure. Validated 72-hour Island Mode is deployed in at least three
EU 5G cross-border corridor pilots under SNS~JU governance. SME-accessible
PQC tooling and SBOM compliance support are available commercially.
\kpistrip{Phase~2 deployment evidence requires named sector and Member State
scope; anonymous ``multi-sector deployment'' claims do not constitute
verifiable operational maturity.}
&
\cellhead{SSL: Emerging EU Ecosystem}
EU-designed secure elements achieve 15--30\% market share in new critical
infrastructure procurement, rising from a near-zero baseline. EU Chips Act
fabrication investments have created a domestic production capacity that is
commercially credible, if not yet dominant. Foreign TPMs remain the standard
for legacy deployments but are no longer the only option for new regulated-sector
procurement. EU open-hardware TEE ecosystem (Keystone, OpenTitan) holds
commercially certifiable products. The conditions for Full Sovereign Autonomy
are structurally in place --- fabrication capacity, certification infrastructure,
standards leadership --- but have not yet reached continental operational scale.
\kpistrip{SSL transition requires industrial substance, not research demonstration.
15--30\% new critical infra market share is the quantified Phase~2 SSL threshold.}
\\
\hline

% ---- GATE 2->3 ----
\rowcolor{gaterow}
\gatelbl{2 $\rightarrow$ 3} &
\multicolumn{4}{>{\footnotesize\raggedright\arraybackslash}p{22.0cm}|}{%
\cellcolor{gaterow}%
\textbf{R\&I trigger:} Production-ready PQC-native RISC-V silicon is available
from EU-domiciled fabrication at commercially viable yields, confirmed by
independent cost-per-unit assessment against equivalent non-EU alternatives.
\textbf{Policy trigger:} EUCC~`High' certification is mandated as a procurement
requirement for all new EU critical infrastructure edge node deployments, with
no waiver provisions for new procurements (legacy replacement programmes are
time-bound).
\textbf{Operational trigger:} Validated 72-hour Island Mode is deployed as a
standard operational baseline --- not a pilot --- in $\geq$50\% of new EU
NIS2-essential entity edge node procurements in the energy and
telecommunications sectors.
\textbf{Market trigger:} EU Chips Act sovereign fabrication capacity achieves
$\geq$50\% self-sufficiency for security-critical silicon in new EU critical
infrastructure procurement, measured by share of new orders fulfilled by
EU-origin fabrication.}
\\
\hline

% ---- PHASE 3 ----
\rowcolor{phase3row}
\textbf{Phase~3}
\newline\textit{Autonomy \&}
\newline\textit{Leadership}
\newline\textit{2035--2040}
\newline\sslbadge{ssltealbg}{ssltealtx}{SSL: Full Sovereign Autonomy}
&
\cellhead{Cluster ODP: Ubiquitous Operation}
All five sub-priorities operate at Ubiquitous Operation. EU-designed
RISC-V secure elements are the hardware default across all new EU ICT
products. Fully Homomorphic Encryption is operationalised for
end-to-end confidential processing from IoT to HPC, not confined to
high-value analytics. Zero-energy PhySec authentication via RF
fingerprinting operates as a standard device identity mechanism for
battery-less mMTC nodes. Bio-inspired swarm protocols enable autonomous
cyber-physical recovery without cloud intervention. Phase~3 R\&I is
oriented toward the next-generation frontier: neuromorphic secure
compute, post-second-generation quantum hardware, and post-CMOS
security architectures --- EU research groups define the agenda rather
than adopting it.
\kpistrip{Ubiquitous Operation requires deployment statistics, not
deployment plans. Phase~3 evidence must cite procurement share data
and operational audit findings, not programme milestone completions.}
&
\cellhead{EU as international hardware security standards reference}
EUCC~`High' certification is mandatory for all new EU ICT products
without waiver provisions. PQC-native silicon is the EU procurement
default. EU hardware security standards are referenced in ISO/IEC
international standardisation and in bilateral digital trade
agreements. Continuous automated EUCC lifecycle maintenance replaces
point-in-time audit models completely. The EU sets the terms of
hardware security compliance for market access, transforming its
internal market standard into a de facto international reference.
\kpistrip{Phase~3 policy maturity requires evidence of international
standards influence: at least one EU hardware security standard adopted
or formally referenced in an ISO/IEC or ITU-T international standard.}
&
\cellhead{Security-by-Design the unconditional hardware and crypto default}
$\geq$60\% of EU critical infrastructure edge nodes deploy EUCC~`High'
certified EU-designed secure elements. FHE is operational end-to-end
from IoT sensor to HPC compute node. Autonomous Island Mode provides
validated 72-hour survivability as a mandatory operational baseline
across all NIS2-essential and DORA-critical edge deployments. No new
EU ICT product enters the market without a hardware root of trust.
The question ``does this deployment have hardware-level integrity?''
has no non-trivial answer for any new EU deployment.
\kpistrip{Phase~3 deployment normalisation: $\geq$60\% market penetration
for EU secure elements in critical infra; 100\% of new EU ICT products
with hardware RoT as a default.}
&
\cellhead{SSL: Full Sovereign Autonomy}
EU-designed and EU-manufactured secure elements hold $\geq$60\% of new
critical infrastructure procurement market share. Foreign legacy
dependency operates under binding replacement programmes with defined
timelines. EU sovereign RISC-V silicon and open-hardware TEE
architectures are exported to and adopted by partner jurisdictions.
EU PQC and hardware TEE standards are the international reference.
Structural exposure to geopolitical disruption of any single non-EU
hardware supplier has been eliminated for new deployments. The EU
has moved from defensive sovereignty (reducing dependency) to offensive
strategic posture (defining the terms of global hardware security
competition).
\kpistrip{Full Sovereign Autonomy is verified by $\geq$60\% market share
in new critical infra procurement plus the existence of binding
replacement programmes for legacy foreign hardware --- not by 100\%
market share.}
\\
\hline

\end{longtable}
\end{landscape}


% ---- PROGRESSION MATRIX: CLUSTER 1 ----

% \noindent
% \resizebox{\textwidth}{!}{%
% \begin{tabular}{|p{4.2cm}|p{5.6cm}|p{5.6cm}|p{5.6cm}|}
% \hline
% \rowcolor{gray!30}
% \multicolumn{4}{|c|}{\textbf{Progression Matrix: Priority Cluster 1 ---
%   Sovereign Trusted Infrastructure (P1, P2, P3, P8, P9)}} \\
% \hline
% \rowcolor{gray!15}
% \textbf{Metric} &
% \textbf{Phase 1: Foundation \& Compliance (2025--2030)} &
% \textbf{Phase 2: Federation \& Standardisation (2030--2035)} &
% \textbf{Phase 3: Autonomy \& Leadership (2035--2040)} \\
% \hline
% \rowcolor{gray!10}
% \multicolumn{4}{|l|}{\textit{\textbf{Cluster-Level Progression Indicators}}} \\
% \hline
% \textbf{Cluster ODP} &
% Applied R\&D &
% Standardised Integration &
% Ubiquitous Operation \\
% \hline
% \textbf{Cluster SSL} &
% High Foreign Dependency &
% Emerging EU Ecosystem &
% Full Sovereign Autonomy \\
% \hline
% \rowcolor{gray!10}
% \multicolumn{4}{|l|}{\textit{\textbf{Sub-Priority Decomposition (Individual ODP / SSL per Phase)}}} \\
% \hline
% \textit{P1: Hardware RoT \newline \& EUCC `High'} &
% ODP: Fragmented Piloting \newline
% SSL: High Foreign Dependency &
% ODP: Standardised Integration \newline
% SSL: Emerging EU Ecosystem &
% ODP: Ubiquitous Operation \newline
% SSL: Full Sovereign Autonomy \\
% \hline
% \textit{P2: PQC Agility \newline for Edge/IoT} &
% ODP: Applied R\&D \newline
% SSL: High Foreign Dependency &
% ODP: Fragmented Piloting \newline
% SSL: Emerging EU Ecosystem &
% ODP: Ubiquitous Operation \newline
% SSL: Full Sovereign Autonomy \\
% \hline
% \textit{P3: Lightweight Crypto \newline \& Physical-Layer Sec.} &
% ODP: Applied R\&D \newline
% SSL: High Foreign Dependency &
% ODP: Standardised Integration \newline
% SSL: Emerging EU Ecosystem &
% ODP: Ubiquitous Operation \newline
% SSL: Full Sovereign Autonomy \\
% \hline
% \textit{P8: Confidential Computing \newline \& Privacy-Pres. AI} &
% ODP: Applied R\&D \newline
% SSL: High Foreign Dependency &
% ODP: Fragmented Piloting \newline
% SSL: Emerging EU Ecosystem &
% ODP: Ubiquitous Operation \newline
% SSL: Full Sovereign Autonomy \\
% \hline
% \textit{P9: Resilient-by-Design \newline \& Island Mode} &
% ODP: Fragmented Piloting \newline
% SSL: High Foreign Dependency &
% ODP: Standardised Integration \newline
% SSL: Contained Local Survivability &
% ODP: Ubiquitous Operation \newline
% SSL: Systemic EU Resilience \\
% \hline
% \rowcolor{blue!7}
% \textbf{Operational State} &
% EU-funded RISC-V secure element pilot projects commence EUCC evaluation;
% NIST FIPS~203/204/205 algorithms are adopted by research consortia and
% optimised for constrained devices in laboratory settings; Intel TDX
% and AMD SEV-SNP dominate all operational TEE deployments while
% EU open-hardware alternatives remain at prototype stage;
% edge nodes operate with hard cloud dependencies and lack
% autonomous fallback architectures validated at operational scale. &
% EU-manufactured EUCC `High' certified secure elements enter
% critical infrastructure deployments across energy and
% telecommunications sectors; hybrid ML-KEM plus ECDH-P384 schemes
% are standardised and deployed in Industrial IoT corridors;
% RISC-V Keystone-based TEEs are operational in federated
% multi-provider Data Spaces; validated 72-hour Island Mode
% survivability demonstrated in designated EU 5G cross-border
% corridor pilots under SNS~JU governance. &
% EU Chips Act sovereign fabrication supplies $\geq$60\% of EU demand
% for security-critical silicon; all new EU ICT products ship with
% EUCC `High' certified hardware roots of trust by default;
% fully homomorphic encryption is operationalised for
% end-to-end confidential processing from IoT to HPC;
% bio-inspired swarm protocols coordinate autonomous
% cyber-physical recovery without any cloud intervention,
% achieving DORA-mandated recovery time objectives
% autonomously at continental scale. \\
% \hline
% \rowcolor{orange!12}
% \textbf{Critical Enabler} &
% CRA Articles~10 and~11 enforcement (2024--2027) creating
% mandatory market demand for hardware security certification,
% establishing the commercial conditions for EU secure element
% pilot investment &
% EUCC `High' certification of the first EU-manufactured
% RISC-V secure element, constituting the gate condition
% for the hardware sovereignty transition and unlocking
% EU Chips Act fabrication investment justification &
% EU Chips Act sovereign EUV-capable fabrication capacity
% achieving $\geq$60\% market self-sufficiency in
% security-critical silicon, removing the last structural
% foreign dependency in the infrastructure trust chain \\
% \hline
% \end{tabular}%
% }

% =========================================================================
% PRIORITY CLUSTER 2: COGNITIVE SECURITY AND AUTONOMOUS CYBER DEFENCE
% =========================================================================

\newpage
\section{Priority Cluster 2: Cognitive Security and Autonomous Cyber Defence}
\label{sec:cluster2}

\noindent\textbf{Constituent Sub-Priorities:} P6 (Cognitive SecOps / LLMSecOps \& Automated
Threat Intelligence), P10 (Proactive APT Management), and the AI-as-target dimension
drawn from CL1.2 (Model Integrity / DataBOMs), CL1.3 (Dynamic Supply Chain Verification /
Runtime SBOMs), and CL4.7 (AI Trustworthiness \& Algorithmic Accountability).

\medskip
\noindent\textbf{Cluster Thesis:} Artificial intelligence is simultaneously the most powerful
instrument available for autonomous cyber defence and the most consequential new attack surface
introduced into the security architecture. A roadmap that deploys AI agents for autonomous
threat containment without hardening those agents against adversarial manipulation creates a
critical second-order vulnerability. This cluster therefore addresses three inseparable
sub-dimensions: AI as security instrument (LLMSecOps and APT hunting), AI as adversary
amplifier (AI-generated APT tooling), and AI as target (model poisoning, adversarial evasion,
and algorithmic accountability for autonomous containment decisions).

\subsection*{7-Dimensional Analysis: Priority Cluster 2}

\begin{enumerate}

% ---- DIMENSION 1 ----
\item \textbf{Baseline \& Target Objective}

\textit{Baseline (2025--2026):}
European Security Operations Centres (SOCs) rely predominantly on rule-based SIEM platforms ---
Splunk (US-domiciled, acquired by Cisco in 2023) and IBM QRadar dominate enterprise
deployments. AI-based tools are deployed as analyst ``copilots'' with human analysts
retaining all containment authority; no EU-sovereign foundation models for security
operations exist at production scale. APT dwell times in European enterprise environments
average 150--200 days, with state-sponsored threat actors demonstrating extended persistence
in EU critical infrastructure. The adversarial robustness of AI security agents --- their
resistance to prompt injection, model evasion, and training data poisoning --- is addressed
only in research contexts without operational certification frameworks. NWDAF integration
per 3GPP TS~23.288~\cite{threeGPP23288} exists in 5G standalone deployments but remains
fragmented across operators, and no EU-mandated framework for the trustworthiness of
autonomous security AI exists prior to AI Act enforcement.

Current cluster ODP: Applied R\&D (AI-as-target and sovereign LLM development); Siloed AI
Copilots (P6); Multi-Layered Alert Correlation (P10). Cluster ODP: Applied R\&D.
Cluster SSL: High Reliance on US Foundation Models / High Dependency on Global Threat Feeds.

\textit{Target (2040):}
\begin{sloppypar}
Swarm-based agentic AI defends the EU Cloud-Edge Continuum autonomously, with APT dwell
time reduced to under 24 hours in certified critical infrastructure deployments. EU-sovereign
edge Small Language Models (SLMs) and Large Context Models (LCMs) for security operations
are certified under AI Act Article~9 and cover at least 50\% of EU critical infrastructure
SOC deployments. All autonomous security agents pass mandatory adversarial robustness
validation before operational deployment. Moving Target Defense reduces APT reconnaissance
success rates below 5\% in validated simulation environments. EU AI security frameworks
are adopted as global reference implementations by partner jurisdictions.
\end{sloppypar}

% ---- DIMENSION 2 ----
\item \textbf{Primary Drivers}

\textit{Threat Drivers:}
\begin{itemize}
  \item Nation-state APT campaigns (attributed by EU agencies to APT28/GRU,
        APT29/SVR, and Sandworm) increasingly exploit LLM hallucination and
        prompt injection vulnerabilities to manipulate autonomous security agents
        into executing incorrect containment actions --- weaponising the defence
        architecture against itself.
  \item AI-generated spear-phishing and poly\-morphic malware
        enabling automated, highly-targeted APT campaigns that adapt
        faster than signature-based
        detection cycles, invalidating the static detection rules that rule-based
        SIEM platforms depend upon.
  \item Supply chain attacks targeting AI model training pipelines: data poisoning
        and backdoor insertion into security AI models deployed at continental scale
        can silently degrade detection capability without triggering observable
        anomalies in standard monitoring.
  \item Systematic abuse of MITRE ATT\&CK-catalogued techniques (T1027 Obfuscated
        Files or Information; T1055 Process Injection) within AI-assisted
        reconnaissance tools, enabling more efficient bypass of detection rules
        derived from those same frameworks.
\end{itemize}

\textit{Regulatory Drivers:}
\begin{itemize}
  \item AI Act Article~9: risk management system requirements for high-risk AI
        systems --- autonomous security containment agents making access revocation
        or network isolation decisions qualify as high-risk systems requiring
        mandatory conformity assessment before operational deployment~\cite{aiact}.
  \item AI Act Article~15: accuracy, robustness, and cybersecurity requirements
        for high-risk AI, including specific resilience requirements against
        adversarial manipulation of inputs and training data~\cite{aiact}.
  \item AI Act Article~13: transparency and provision of information for high-risk
        AI --- autonomous security decisions must be explainable to affected
        entities and competent national authorities~\cite{aiact}.
  \item NIS2 Article~21(2)(b): incident handling procedures, including
        AI-augmented detection and response, for essential entities~\cite{nis2}.
\end{itemize}

% ---- DIMENSION 3 ----
\item \textbf{High Impact Activities}

\begin{enumerate}[label=(\alph*)]
  \item \textbf{[Phase~1]} Deploy NWDAF-integrated (3GPP TS~23.288~\cite{threeGPP23288})
        telemetry pipelines feeding ML-based UEBA (User and Entity Behaviour Analytics)
        for APT indicator correlation; establish behavioural baseline profiles across
        federated enterprise environments targeting reduction of APT dwell time from
        the current 150--200 day average to under 72 hours in pilot deployments via
        ML-assisted Indicator of Compromise (IoC) pattern detection.

  \item \textbf{[Phase~1--2]} Develop EU-sovereign Small Language Models (SLMs,
        under 10B parameters) and Large Context Models (LCMs) for edge-native SOC
        operations (LLMSecOps), trained on MITRE ATT\&CK-aligned threat intelligence
        corpora under GDPR-compliant data governance~\cite{gdpr}; establish mandatory
        AI Act Article~9 conformity assessment pipelines for all autonomous security
        agent deployments prior to production authorisation~\cite{aiact}.

  \item \textbf{[Phase~2]} Deploy edge-native ``Guardian'' nodes executing autonomous
        Zero-Trust reconfigurations per NIST SP~800-207 ZTA principles~\cite{nistsp800207}
        and Moving Target Defense (MTD) --- dynamically shuffling IP addressing, network
        topology, and software diversity to disrupt APT reconnaissance --- governed by
        ETSI GS~ZSM~009 closed-loop automation standards~\cite{etsizsm009}.

  \item \textbf{[Phase~2]} Establish mandatory adversarial robustness validation
        pipelines for all AI security agents prior to operational deployment: automated
        red-team adversarial testing (prompt injection, model evasion, backdoor
        detection via CL1.4 purple-team automation), with continuous certification
        under a CEN/CENELEC-standardised AI security agent assurance framework
        aligned with CL4.7 AI Trustworthiness requirements.

  \item \textbf{[Phase~3]} Deploy swarm-based agentic AI coordination protocols
        enabling decentralised, self-evolving threat intelligence synthesis across
        the EU Cloud-Edge Continuum; integrate with Deception Technology
        infrastructure (CL2.6: honeypots, honeynets, credential canaries, fake API
        endpoints) to generate high-fidelity, low-noise adversarial ground truth for
        continuous autonomous agent retraining.
\end{enumerate}

% ---- DIMENSION 4 ----
\item \textbf{Catalysts \& Enablers}

\textit{Funding Alignment:}
EuroHPC Joint Undertaking AI Factories for training EU-sovereign security foundation
models at the compute scale required; ECCC Strategic Agenda priority areas ``Cybersecurity
of AI'' and ``AI-enabled cybersecurity tools and services''~\cite{eccc}; Horizon Europe
Cluster~3 (Civil Security for Society) calls targeting AI-driven threat intelligence and
autonomous incident response; EU AI Office for conformity assessment infrastructure
supporting high-risk security AI certification under the AI Act.

\textit{Standardisation:}
ETSI GS~ZSM~009 (Zero-Touch Network and Service Management, Closed-Loop
Automation~\cite{etsizsm009}); 3GPP TS~23.288 (NWDAF architecture for AI-native 5G
network analytics~\cite{threeGPP23288}); MITRE ATT\&CK framework (behavioural threat
intelligence taxonomy for agent training corpora); NIST SP~800-207 (Zero Trust
Architecture~\cite{nistsp800207}); CEN/CENELEC JTC~13 working groups on AI security
agent assurance standardisation.

\textit{Ecosystem Enablers:}
ENISA AI cybersecurity threat landscape guidelines; federated Cyber Range infrastructure
(CL3.5) for adversarial agent testing under realistic attack scenarios; CONCORDIA
successor programme for EU-wide SOC federation and cross-border CTI sharing~\cite{concordia};
EU AI Act notified conformity assessment bodies for high-risk AI product certification.

% ---- DIMENSION 5 ----
\item \textbf{Disruptors \& Systemic Risks}

\begin{itemize}
  \item \textbf{Adversarial AI arms race:} threat actors weaponising LLMs for
        automated zero-day discovery and polymorphic exploit generation may outpace
        the defensive AI maturation timeline --- specifically, AI-generated
        vulnerability discovery could exceed human patching bandwidth within the
        Phase~1--2 window, before EU-sovereign autonomous response agents are
        certified for operational deployment.
  \item \textbf{LLM hallucination in containment agents:} false positive isolation
        of critical infrastructure nodes (e.g., a hospital network segment incorrectly
        quarantined during a medical emergency) constitutes an operational safety risk
        that may impose worse harm than the threat the agent was designed to
        contain --- this is the central argument for mandatory HITL escalation in Phase~1.
  \item \textbf{EU sovereign LLM compute deficit:} training security-relevant
        foundation models requires EuroHPC-scale compute; as of 2025--2026, EU frontier
        model training capacity lags US and China by approximately one order of
        magnitude, constraining Phase~1 EU-sovereign models to SLM/LCM scale
        ($<$10B parameters) and limiting capability relative to US foundation models.
  \item \textbf{AI model poisoning via CTI feed corruption:} if adversaries inject
        false indicators into NWDAF analytics pipelines or Cyber Threat Intelligence
        (CTI) sharing platforms (CL2.1), behavioural baselines trained on poisoned
        data will silently misclassify legitimate traffic and fail to detect real APTs
        --- a harder failure mode to detect than outright system unavailability.
  \item \textbf{AI Act certification bottleneck:} mandatory conformity assessment
        for high-risk autonomous security agents may create a 24--36 month regulatory
        gap during which EU organisations operate without certified autonomous
        containment, relying on slower human-in-the-loop processes that cannot match
        AI-assisted adversary speed.
\end{itemize}

% ---- DIMENSION 6 ----
\item \textbf{Failsafe Mechanisms \& Resilience}

\begin{itemize}
  \item \textbf{Shadow mode operation:} all autonomous AI security agents run in
        parallel with human-supervised rule-based SIEM during Phases~1--2;
        containment actions affecting Tier-1 critical infrastructure (energy,
        healthcare, transport per NIS2 Annex~I) require mandatory Human-In-The-Loop
        (HITL) escalation regardless of AI confidence score, with reversible quarantine
        actions bounded by a defined 60-minute rollback window~\cite{nis2}.
  \item \textbf{Model diversity and behavioural circuit breakers:} deploy a minimum
        of two architecturally distinct AI detection models (e.g., transformer-based
        LLM plus graph neural network UEBA) to prevent a single-model adversarial
        exploit from disabling the entire detection capability; behavioural circuit
        breakers automatically suspend any agent exhibiting an anomalous action rate
        or pattern deviating beyond $3\sigma$ from its established operational
        baseline, reverting to rule-based SIEM until re-certification.
  \item \textbf{RLHF continuous correction:} incorporate analyst override decisions
        into continuous Reinforcement Learning from Human Feedback (RLHF) retraining
        to prevent model drift from training distribution; quarterly adversarial
        red-team audits (CL1.4) detect capability degradation before it reaches
        operational significance.
  \item \textbf{Air-gapped CTI validation:} all external CTI feeds ingested by
        NWDAF analytics pipelines pass through an isolated validation environment
        where indicator quality is assessed against independent behavioural ground
        truth before poisoned data can influence operational AI model weights.
  \item \begin{sloppypar}\textbf{Deception-backed calibration:} honeynet and credential canary
        infrastructure (CL2.6) provides high-fidelity, low-noise ground truth against
        which AI behavioural models are continuously calibrated, reducing false
        positive rates and providing an independent signal for detecting model
        poisoning via discrepancy analysis between canary alerts and AI model
        predictions.\end{sloppypar}
\end{itemize}

% ---- DIMENSION 7 ----
\item \textbf{Transversal Meta-Layer}

\textit{Sovereignty \& Autonomy:}
This cluster eliminates reliance on US-domiciled SIEM platforms (Splunk/Cisco, IBM QRadar)
and US foundation models for operational security intelligence. EU-sovereign SLMs trained
under GDPR-compliant data governance~\cite{gdpr} ensure that sensitive threat intelligence
and behavioural telemetry does not transit non-EU cloud infrastructure. EU leadership in
trustworthy AI security tools creates an exportable technology capability aligned with
the AI Act's ambition for global regulatory standard-setting.

\textit{Trust \& Ethics:}
AI Act Article~13 (transparency requirements) is mandated for all autonomous security
decisions --- affected entities must receive meaningful explanations for AI-driven access
revocations or network isolation actions~\cite{aiact}. GDPR Article~22 (restrictions on
automated individual decision-making) applies where AI containment decisions affect
individual data subjects; human review mechanisms must be structurally embedded, not
optional add-ons~\cite{gdpr}. AI Act Article~15 adversarial robustness requirements
ensure security AI systems cannot be weaponised against the infrastructure they are
designed to protect~\cite{aiact}.

\textit{Sustainable Security (Green IT):}
Edge-native SLMs under 10B parameters reduce inference energy consumption by two to
three orders of magnitude compared to centralised GPT-class model inference, enabling
real-time threat classification at the point of detection without data centre energy
overhead. Federated training paradigms reduce repeated data transmission overhead
compared to centralised training architectures and are compatible with the privacy
preservation requirements of GDPR Article~25~\cite{gdpr}.

\end{enumerate}

% =========================================================================
% CLUSTER 2: COGNITIVE SECURITY AND AUTONOMOUS CYBER DEFENCE
% =========================================================================
\newpage
\subsection{Progression Matrix: Priority Cluster 2 --- Cognitive Security and Autonomous
Cyber Defence}

\begin{landscape}
\setlength{\textwidth}{26.5cm}\setlength{\linewidth}{26.5cm}%
\setlength{\hsize}{26.5cm}\setlength{\columnwidth}{26.5cm}
\setlength{\tabcolsep}{5pt}
\renewcommand{\arraystretch}{1.35}

\begin{longtable}{%
|>{\footnotesize\raggedright\arraybackslash}p{2.8cm}%
|>{\footnotesize\raggedright\arraybackslash}p{5.4cm}%
|>{\footnotesize\raggedright\arraybackslash}p{5.4cm}%
|>{\footnotesize\raggedright\arraybackslash}p{5.4cm}%
|>{\footnotesize\raggedright\arraybackslash}p{5.4cm}|}

\caption{\textbf{Cluster~2: Cognitive Security \& Autonomous Cyber Defence}
(P6, P10, AI-as-target dimension from CL1.2, CL1.3, CL4.7).
Transition matrix across the four methodology dimensions. The AI-as-target
dimension (model integrity, adversarial robustness, algorithmic accountability)
is labelled \textit{(AI-T)} where it diverges from P6 or P10 descriptors.}
\label{tab:cluster2}\\

\hline
\rowcolor{colheader}
\textbf{Phase / SSL} &
\textbf{R\&I Maturity}{\scriptsize ODP indicator per sub-priority} &
\textbf{Policy, Regulatory \& Standards}{\scriptsize enforcement $\cdot$
  certification $\cdot$ standardisation} &
\textbf{Operational Deployment \& Adoption}{\scriptsize actor types $\cdot$
  scale $\cdot$ sector coverage} &
\textbf{Market Uptake \& Strategic Sovereignty}{\scriptsize SSL indicator
  $\cdot$ dependency reduction $\cdot$ industrial capacity} \\
\hline
\endfirsthead

\hline
\rowcolor{colheader}
\textbf{Phase / SSL} &
\textbf{R\&I Maturity} &
\textbf{Policy, Regulatory \& Standards} &
\textbf{Operational Deployment \& Adoption} &
\textbf{Market Uptake \& Strategic Sovereignty} \\
\hline
\endhead
\hline\multicolumn{5}{r}{\footnotesize\textit{Continued\ldots}}\\ \endfoot
\hline \endlastfoot

% ---- PHASE 1 ----
\rowcolor{phase1row}
\textbf{Phase~1}
\newline\textit{Foundation \&}
\newline\textit{Compliance}
\newline\textit{2025--2030}
\newline\sslbadge{sslredbg}{sslredtx}{SSL: High Reliance on}
\newline\sslbadge{sslredbg}{sslredtx}{US Foundation Models}
&
\cellhead{Cluster ODP: Applied R\&D / Siloed AI Copilots}
\subp{P6} reaches Siloed AI Copilots: commercial LLMs (GPT-4-class)
assist human SOC analysts in alert triage, threat report generation, and
MITRE ATT\&CK technique mapping; no EU-sovereign LLM for security
operations exists at production scale; NWDAF (3GPP TS~23.288)
telemetry pipelines are integrated in 5G SA pilot deployments for
anomaly detection. \subp{P10} is at Multi-Layered Alert Correlation:
ML-based UEBA is applied to APT IoC detection in research and pilot
contexts; APT dwell time reduction from the 150--200 day average is
demonstrated in controlled settings. \subp{AI-T} is at Applied R\&D:
adversarial robustness testing (prompt injection, model evasion,
DataBOM verification) is an academic research function without
operational certification frameworks.
\kpistrip{Phase~1 R\&I distinguishes three sub-dimensions with
different maturity levels. Conflating them into a single ODP cell
obscures that AI-as-target capability is the least mature and the
most consequential gap.}
&
\cellhead{AI Act enforcing; conformity assessment infrastructure absent}
AI~Act Articles~9 (risk management), 13 (transparency), and 15
(adversarial robustness) enter applicability in 2025--2026, establishing
mandatory requirements for autonomous security AI as high-risk systems.
However, no EU conformity assessment bodies with capacity to certify
autonomous security agents are operational at the start of Phase~1.
ETSI GS~ZSM~009 (closed-loop automation) is published but not mandated.
NIS2 Article~21(2)(b) incident handling procedures are being
operationalised across essential entities. ENISA AI cybersecurity threat
landscape guidelines are published. The critical regulatory gap is the
absence of a standardised adversarial robustness test suite applicable
to security-specific AI agents.
\kpistrip{The AI Act creates obligations before the conformity
assessment infrastructure exists to fulfil them. This 24--36 month gap
is the dominant Phase~1 policy risk for this cluster.}
&
\cellhead{Deployment: SOC copilot use by large enterprises and national CSIRTs}
US-platform LLMs (GPT-4-class, under data processor agreements) operate
as analyst copilots in large enterprise and national CSIRT SOCs.
Rule-based SIEMs (Splunk/Cisco, IBM~QRadar) remain the primary
detection infrastructure across EU regulated sectors. NWDAF-integrated
AI analytics operate in 3--5 5G SA deployments run by major EU
telecom operators. No certified autonomous security containment agent
is in operational deployment anywhere in the EU. APT dwell times of
150--200 days in EU enterprise environments reflect the gap between
AI-assisted detection capability and operational deployment at scale.
SMEs lack access to AI-powered SOC tooling at viable cost.
\kpistrip{Phase~1 deployment is concentrated in the largest institutional
actors (national CSIRTs, Tier-1 operators, major financial entities).
The gap between these deployments and SME access defines the adoption
challenge for Phase~2.}
&
\cellhead{SSL: High Reliance on US Foundation Models}
The EU SOC tooling market is dominated by US-domiciled platforms:
Splunk (Cisco), IBM QRadar, and Microsoft Sentinel collectively control
the EU enterprise SIEM market. The foundation models used for SOC
copilot functions are US-origin (OpenAI, Anthropic, Google) and are
accessed under data processor agreements, creating persistent
extra-territorial data exposure. No EU-sovereign foundation model at
production scale for security operations exists. EuroHPC compute
capacity is insufficient for frontier-class security model training.
The adversarial robustness validation market is pre-commercial.
\kpistrip{Quantified dependency: US platforms hold $\geq$80\% of EU
enterprise SIEM market; US-origin models are the sole source for
LLM-based SOC tooling. This specific dependency structure is what
Phases~2 and~3 must reduce.}
\\
\hline

% ---- GATE 1->2 ----
\rowcolor{gaterow}
\gatelbl{1 $\rightarrow$ 2} &
\multicolumn{4}{>{\footnotesize\raggedright\arraybackslash}p{22.0cm}|}{%
\cellcolor{gaterow}%
\textbf{R\&I trigger:} At least one EU-sovereign Small Language Model
($\leq$10B parameters) trained on GDPR-compliant security corpora has
been validated against an independently benchmarked EU critical infrastructure
SOC use case at production-equivalent scale.
\textbf{Policy trigger:} AI~Act conformity assessment bodies with capacity
to certify autonomous security AI as high-risk systems under Articles~9 and~15
are operational and have issued at least one certificate to a security AI
agent product.
\textbf{Operational trigger:} At least one EU-sovereign LLMSecOps agent is
in certified operational deployment in a named NIS2-essential entity,
demonstrating APT dwell time below 72 hours in a documented production
environment.
\textbf{Market trigger:} A standardised adversarial robustness test suite
for security AI agents, developed under CEN/CENELEC or ETSI TC~CYBER
governance, has been published and is referenced in at least one EU
procurement framework for critical infrastructure AI tools.}
\\
\hline

% ---- PHASE 2 ----
\rowcolor{phase2row}
\textbf{Phase~2}
\newline\textit{Applied \& Industrial}
\newline\textit{R\&D alongside}
\newline\textit{Federation \&}
\newline\textit{Standardisation}
\newline\textit{2030--2035}
\newline\sslbadge{ssltransbg}{ssltrantx}{SSL: Sovereign EU Edge}
\newline\sslbadge{ssltransbg}{ssltrantx}{LLMs Emerging}
&
\cellhead{Cluster ODP: Federated Edge Agents / Automated Triage}
\subp{P6} reaches Federated Edge Agents: EU-sovereign SLMs ($\leq$10B
parameters) trained on MITRE ATT\&CK-aligned corpora are operational
at the edge for real-time SOC analytics without cloud inference
dependency; Guardian nodes execute autonomous Zero-Trust
reconfigurations per NIST~SP~800-207. \subp{P10} reaches Automated
Triage and LLMSecOps Containment: LLMSecOps agents ingest CTI feeds,
track APT lateral movement against behavioural frameworks, and
orchestrate incident triage autonomously under ETSI GS~ZSM~009
governance; APT dwell time $\leq$72 hours in certified pilot deployments.
\subp{AI-T} reaches Standardised Integration: adversarial robustness
validation pipelines (prompt injection, model evasion, backdoor
detection) are standardised and mandatory under AI~Act~Article~9
conformity assessment; continuous RLHF retraining and circuit-breaker
mechanisms are in production.
\kpistrip{Phase~2 R\&I includes active applied and industrial R\&D for
all three sub-dimensions --- swarm AI protocols, post-quantum AI
security, and next-generation adversarial testing frameworks --- not
merely consolidation of Phase~1 results.}
&
\cellhead{AI Act conformity assessment operational: the Phase~2 gate}
The single Phase~2 policy gate is the operationalisation of a
standardised AI~Act conformity assessment framework for autonomous
security agents, enabling certified deployment without per-case
regulatory approval. ETSI GS~ZSM~009 is mandated within EU NaaS
federation governance frameworks. A CEN/CENELEC-standardised adversarial
robustness test suite for security AI is published. MITRE ATT\&CK
behavioural framework is formally referenced in NIS2 sector-specific
incident handling guidance. EuroHPC AI Factory compute contracts for
EU-sovereign security model training are active.
\kpistrip{Without operational AI~Act conformity assessment
infrastructure, autonomous security agents cannot be certified for
regulated-sector deployment regardless of technical readiness.
This is the dominant Phase~2 policy dependency.}
&
\cellhead{Certified EU sovereign AI agents in critical infrastructure SOCs}
EU SLM-based LLMSecOps agents are in certified operational deployment
in energy, financial (DORA), and telecommunications critical
infrastructure SOCs in $\geq$3 Member States. Guardian nodes execute
autonomous Zero-Trust reconfigurations under human-in-the-loop
escalation protocols for Tier-1 critical actions. NWDAF-integrated
analytics are mandatory in 5G SA deployments. Deception infrastructure
(honeynets, credential canaries, fake API endpoints per CL2.6) is
deployed alongside AI detection systems, providing high-fidelity
ground truth for model calibration. APT dwell time $\leq$72 hours
demonstrated in pilot critical infrastructure deployments.
\kpistrip{Phase~2 deployment evidence requires named critical
infrastructure sector, Member State scope, and a documented APT dwell
time metric from an operational --- not simulated --- environment.}
&
\cellhead{SSL: Sovereign EU Edge LLMs Emerging}
EU-sovereign SLMs cover $\geq$30\% of critical infrastructure SOC
deployments in regulated sectors, rising from zero at baseline.
EuroHPC AI Factories supply training infrastructure for EU-origin
security models, reducing dependence on US hyperscaler compute.
US-platform SIEMs retain market share in legacy deployments but
face procurement competition from certified EU alternatives.
Adversarial robustness validation is a commercially available service
with EU-domiciled providers. EU AI security tool exports to partner
jurisdictions have begun, indicating the transition from security
consumer to security producer posture.
\kpistrip{SSL evidence at Phase~2: $\geq$30\% critical infra SOC market
share for EU-sovereign AI tools; at least one documented EU AI security
tool export event to a non-EU partner jurisdiction.}
\\
\hline

% ---- GATE 2->3 ----
\rowcolor{gaterow}
\gatelbl{2 $\rightarrow$ 3} &
\multicolumn{4}{>{\footnotesize\raggedright\arraybackslash}p{22.0cm}|}{%
\cellcolor{gaterow}%
\textbf{R\&I trigger:} EU-certified swarm AI coordination protocols
enabling decentralised autonomous threat intelligence synthesis have passed
independent security evaluation and are deployed in at least one cross-border
EU critical infrastructure corridor under ETSI GS~ZSM~009 governance.
\textbf{Policy trigger:} EU AI security standards (adversarial robustness
profiles, autonomous agent accountability frameworks) are adopted or formally
referenced in at least one ISO/IEC or ITU-T international standard, confirming
EU regulatory leadership beyond domestic jurisdiction.
\textbf{Operational trigger:} APT dwell time $\leq$24 hours is documented in
a named NIS2-essential entity operational deployment (not a simulation),
attributable to autonomous AI-driven detection and containment.
\textbf{Market trigger:} EU-sovereign AI security tools hold $\geq$50\% market
share in new critical infrastructure SOC deployments across at least five
Member States, displacing US-platform SIEMs as the procurement default in
regulated sectors.}
\\
\hline

% ---- PHASE 3 ----
\rowcolor{phase3row}
\textbf{Phase~3}
\newline\textit{Autonomy \&}
\newline\textit{Leadership}
\newline\textit{2035--2040}
\newline\sslbadge{ssltealbg}{ssltealtx}{SSL: Global Leadership in}
\newline\sslbadge{ssltealbg}{ssltealtx}{Trustworthy AI Security}
&
\cellhead{Cluster ODP: Ubiquitous AI Immune System}
All three sub-dimensions reach Ubiquitous Operation. Swarm-based
agentic AI defends the EU Cloud-Edge Continuum via self-evolving,
decentralised threat intelligence without any central coordination
dependency. Moving Target Defense continuously shuffles IP addressing,
network topology, and software diversity; validated APT reconnaissance
success rates fall below 5\% in certified operational environments.
Adversarial robustness validation is continuous and fully automated:
no security AI agent operates in EU critical infrastructure without
an active conformity certification. Phase~3 R\&I addresses
post-swarm frontiers: quantum-native AI security architectures,
brain-inspired neuromorphic threat detection, and AI-to-AI
adversarial dynamics beyond current threat models.
\kpistrip{Phase~3 ODP evidence requires operational statistics from
production deployments: APT dwell time figures, MTD effectiveness
metrics, and autonomous containment accuracy rates from certified
environments --- not from cyber range simulations.}
&
\cellhead{EU AI security governance as international reference}
EU AI Act frameworks for trustworthy autonomous security agents are
adopted internationally and referenced in bilateral digital
partnership agreements. Autonomous agent certification operates at
continental scale without per-case regulatory review. Zero-touch
network management under ETSI GS~ZSM~009 governs all EU NaaS
operations. The EU AI Office and ENISA jointly define the international
standard for security AI accountability. EU AI security governance
has transitioned from reactive adaptation to agenda-setting: other
jurisdictions adopt EU frameworks to access EU market and partnership
programmes.
\kpistrip{Phase~3 policy leadership is verified by international
adoption events --- not by the number of EU instruments in force.
At least one non-EU jurisdiction formally adopting EU security AI
governance standards constitutes the Phase~3 policy evidence threshold.}
&
\cellhead{Autonomous continental-scale AI cyber defence as operational default}
APT dwell time $\leq$24 hours in all certified EU NIS2-essential
critical infrastructure deployments. MTD is deployed as a mandatory
operational baseline. Swarm AI coordination operates without central
cloud dependency: each EU Cloud-Edge node participates in a
self-organising threat intelligence mesh. Deception infrastructure
is fully integrated into the autonomous response loop, providing
continuous model calibration ground truth. Human intervention is
reserved for exceptional legal escalation scenarios; routine security
operations are fully autonomous within the governance framework.
\kpistrip{The Phase~3 deployment test: is autonomous AI defence
the unconditional operational baseline, or is it still a premium
programme? 100\% of new critical infrastructure deployments
incorporating autonomous AI security constitutes the evidence threshold.}
&
\cellhead{SSL: Global Leadership in Trustworthy AI Security}
EU-sovereign AI security tools hold $\geq$50\% of new critical
infrastructure SOC deployments across the EU. US-platform SIEMs
are legacy infrastructure under active replacement programmes with
defined timelines. EU AI security models are exported to and
commercially adopted by partner jurisdictions. EU sets the commercial
terms for security AI market access: certification against EU AI~Act
Article~9/15 profiles is a de facto global procurement requirement.
The EU has transitioned from defensive sovereignty to offensive
strategic posture: it defines the technical terms on which the
global security AI industry must compete.
\kpistrip{Full sovereign leadership in this cluster requires both
domestic dominance ($\geq$50\% critical infra SOC market share)
and international influence (at least two partner jurisdictions
formally adopting EU security AI governance). Both conditions must
be met.}
\\
\hline

\end{longtable}
\end{landscape}

% ---- PROGRESSION MATRIX: CLUSTER 2 ----

% \noindent
% \resizebox{\textwidth}{!}{%
% \begin{tabular}{|p{4.2cm}|p{5.6cm}|p{5.6cm}|p{5.6cm}|}
% \hline
% \rowcolor{gray!30}
% \multicolumn{4}{|c|}{\textbf{Progression Matrix: Priority Cluster 2 ---
%   Cognitive Security and Autonomous Cyber Defence (P6, P10, AI-as-Target)}} \\
% \hline
% \rowcolor{gray!15}
% \textbf{Metric} &
% \textbf{Phase 1: Foundation \& Compliance (2025--2030)} &
% \textbf{Phase 2: Federation \& Standardisation (2030--2035)} &
% \textbf{Phase 3: Autonomy \& Leadership (2035--2040)} \\
% \hline
% \rowcolor{gray!10}
% \multicolumn{4}{|l|}{\textit{\textbf{Cluster-Level Progression Indicators}}} \\
% \hline
% \textbf{Cluster ODP} &
% Applied R\&D / Siloed AI Copilots &
% Federated Edge Agents &
% Ubiquitous AI Immune System \\
% \hline
% \textbf{Cluster SSL} &
% High Reliance on US Foundation Models &
% Sovereign European Edge LLMs / SLMs &
% Global Leadership in Trustworthy AI Security \\
% \hline
% \rowcolor{gray!10}
% \multicolumn{4}{|l|}{\textit{\textbf{Sub-Priority Decomposition (Individual ODP / SSL per Phase)}}} \\
% \hline
% \textit{P6: Cognitive SecOps \newline (LLMSecOps)} &
% ODP: Siloed AI Copilots \newline
% SSL: High Reliance on US Foundation Models &
% ODP: Federated Edge Agents \newline
% SSL: Sovereign EU Edge LLMs Emerging &
% ODP: Ubiquitous AI Immune System \newline
% SSL: Global Leadership in Trustworthy AI \\
% \hline
% \textit{P10: Proactive APT \newline Management} &
% ODP: Multi-Layered Alert \newline Correlation \& Hunting \newline
% SSL: High Dependency on Global Threat Feeds &
% ODP: Automated Triage \& \newline LLMSecOps Containment \newline
% SSL: Sovereign Pan-EU CTI \& Behavioural Analytics &
% ODP: Ubiquitous Moving \newline Target Defense (MTD) \newline
% SSL: Full Autonomous Counter-Adversarial Resilience \\
% \hline
% \textit{AI-as-Target \newline (CL1.2, CL1.3, CL4.7)} &
% ODP: Applied R\&D \newline
% SSL: High Foreign Dependency &
% ODP: Standardised Integration \newline
% SSL: Emerging EU Ecosystem &
% ODP: Ubiquitous Operation \newline
% SSL: Full Sovereign Autonomy \\
% \hline
% \rowcolor{blue!7}
% \textbf{Operational State} &
% LLMs assist human SOC analysts in alert triage and threat report
% synthesis; NWDAF telemetry is routed to centralised AI for anomaly
% detection in 5G standalone deployments; APT dwell times are
% reduced via ML-assisted IoC correlation in pilot enterprise environments;
% adversarial robustness validation of AI security agents is conducted
% in academic and research settings without operational certification
% frameworks in force. &
% Edge-native EU-sovereign SLMs execute autonomous Zero-Trust
% reconfigurations and traffic throttling in certified critical
% infrastructure deployments; LLMSecOps agents ingest CTI and
% track lateral movement against MITRE ATT\&CK behavioural
% frameworks in real time; adversarial robustness validation
% is standardised under AI Act Article~9 conformity assessment;
% APT dwell time reduced to under 72 hours in EU critical
% infrastructure pilot deployments. &
% Swarm-based agentic AI defends the EU Cloud-Edge Continuum
% via self-evolving, decentralised threat intelligence without
% central coordination dependency; Moving Target Defense
% continuously shuffles network topology, IP addressing, and
% software diversity, achieving validated APT reconnaissance
% success rates below 5\% in operational simulation environments;
% EU AI security models are formally adopted as reference
% implementations by partner jurisdictions. \\
% \hline
% \rowcolor{orange!12}
% \textbf{Critical Enabler} &
% AI Act Articles~9 and~15 enforcement (2025--2026) establishing
% mandatory risk management and adversarial robustness
% requirements for autonomous security AI, creating the
% regulatory framework for certified Phase~2 deployment &
% Standardised AI Act conformity assessment framework for
% autonomous security agents, enabling certified operational
% deployment without per-case regulatory approval and
% unlocking Phase~3 swarm-level autonomy &
% EU-certified self-evolving AI swarm protocols achieving
% full operational deployment under ETSI GS~ZSM~009
% zero-touch automation governance with Pan-EU
% CTI-sharing infrastructure \\
% \hline
% \end{tabular}%
% }

% =========================================================================
% PRIORITY CLUSTER 3: FEDERATED OPERATIONAL SOVEREIGNTY
% =========================================================================

\newpage
\section{Priority Cluster 3: Federated Operational Sovereignty}
\label{sec:cluster3}

\noindent\textbf{Constituent Sub-Priorities:} P4 (Security Interoperability via Simpl
Middleware), P5 (Supply Chain Verification and Secured Infrastructure), P7 (Pan-European
NaaS Federation --- ``Schengen for Cloud-Edge/AI'').

\medskip
\noindent\textbf{Cluster Thesis:} Technical excellence in Cluster~1 and autonomous
defence capability in Cluster~2 deliver no strategic value if they cannot operate
verifiably across borders, across providers, and across supply chains. This cluster
defines the governance fabric, interoperability standards, and supply chain integrity
mechanisms that transform individually sovereign components into a collectively sovereign
European computing continuum. Critically, it addresses the legal and semantic as well as
the technical interoperability problem --- without a ``Schengen for Cloud'' legal framework,
cross-border compute federation remains commercially fragile regardless of technical
standardisation.

\subsection*{7-Dimensional Analysis: Priority Cluster 3}

\begin{enumerate}

% ---- DIMENSION 1 ----
\item \textbf{Baseline \& Target Objective}

\textit{Baseline (2025--2026):}
European cloud and edge infrastructure remains highly fragmented across proprietary provider
silos with incompatible security policies, authentication mechanisms, and data exchange
protocols. Common European Data Spaces (Manufacturing, Health, Energy, Agriculture) are in
early piloting stages, with security interoperability relying on bilateral agreements rather
than standardised automated mechanisms aligned with the Simpl open-source middleware
stack~\cite{cloud_edge_roadmap}. Supply chain verification is largely static: CRA-mandated
Software Bills of Materials (SBOMs)~\cite{cra} are point-in-time declarations without runtime
verification; Digital Product Passport (DPP) frameworks remain in pilot stage under the EU
Ecodesign Regulation. NaaS federation across EU telecom operators does not exist at
operational scale --- cross-border edge compute sharing requires bespoke commercial agreements.
CAMARA Network APIs are being standardised by GSMA but remain fragmentedly implemented
across European operators.

Current cluster ODP: Fragmented Piloting (P4, P7); early supply chain disclosure phase
(P5). Cluster SSL: Emerging EU Ecosystem (P4); Hyperscaler Platform Dominance (P7);
High reliance on fragmented non-EU supply chain auditing (P5).

\textit{Target (2040):}
A unified European NaaS platform enables seamless cross-border compute roaming under a
``Schengen for Cloud'' legal framework, legally operable as a single jurisdiction for
data workloads. All Common European Data Spaces operate on Simpl-based security
interoperability with automated MIM compliance --- provider-agnostic, cryptographically
verifiable, and free of vendor lock-in. Runtime SBOMs with Zero-Knowledge Proof (ZKP)
verification eliminate unverified components from EU critical infrastructure supply chains.
EU achieves full sovereign control over supply chain transparency with pan-European Cyber
Threat Intelligence sharing for supply chain compromise indicators.

% ---- DIMENSION 2 ----
\item \textbf{Primary Drivers}

\textit{Threat Drivers:}
\begin{itemize}
  \item Supply chain attacks exploiting the disaggregation complexity of Open RAN
        deployments, where a single compromised firmware component sourced from one
        of 50-plus hardware and software vendors per macro cell deployment can affect
        thousands of base stations across multiple operators before detection.
  \item ``Weakest link'' cascading failures in federated NaaS environments: a
        single provider with a compromised Minimal Interoperability Mechanism (MIM)
        implementation can propagate malicious policies across the entire federated
        fabric before semantic validation detects the enforcement inconsistency.
  \item US CLOUD Act extra-territorial reach over non-EU cloud providers processing
        EU data creates a structural sovereignty risk that persists regardless of
        technical security measures if the legal framework enabling EU-sovereign
        federated compute is not established.
  \item Silent policy translation failures: automated security policy translation
        between heterogeneous provider MIM implementations may produce semantically
        equivalent but enforcement-different policies, creating exploitable gaps
        invisible to standard compliance auditing procedures.
\end{itemize}

\textit{Regulatory Drivers:}
\begin{itemize}
  \item Data Act interoperability obligations (Articles~23--31): mandatory
        requirements for cloud service providers covering security interface
        standardisation and data portability~\cite{nis2}.
        [CITATION REQUIRES VERIFICATION --- confirm final Data Act article numbering
        for interoperability provisions.]
  \item NIS2 Article~21(2)(h): supply chain security measures mandatory for essential
        entities, requiring documented risk assessment of all ICT suppliers~\cite{nis2}.
  \item CRA Annex~I Part~II: supply chain vulnerability handling requirements for
        products with digital elements, including SBOM disclosure
        obligations~\cite{cra}.
  \item DORA Article~28: ICT third-party risk management requiring continuous
        monitoring and documented exit strategies for critical ICT
        suppliers~\cite{dora}.
\end{itemize}

% ---- DIMENSION 3 ----
\item \textbf{High Impact Activities}

\begin{enumerate}[label=(\alph*)]
  \item \textbf{[Phase~1]} Define and standardise security-specific Minimal
        Interoperability Mechanisms (MIMs) within the Simpl open-source middleware
        stack~\cite{cloud_edge_roadmap}; deploy pilots in three Common European Data
        Spaces (Manufacturing, Health, Energy) with automated cross-provider policy
        enforcement, federated identity attestation, and cryptographically verifiable
        compliance reporting.

  \item \textbf{[Phase~1--2]} Implement Runtime SBOMs with cryptographic proof of
        execution for all containerised microservices and AI components migrating across
        the multi-provider continuum; integrate Digital Product Passports (DPP) for
        hardware component traceability from silicon fabrication to operational
        deployment, aligned with CRA SBOM obligations~\cite{cra}.

  \item \textbf{[Phase~1--2]} Deploy harmonised CAMARA Network APIs across a minimum
        of five EU telecom operators establishing standardised interfaces for cross-border
        edge compute roaming; test Automated Compliance Negotiators in designated EU
        5G cross-border corridor deployments under IPCEI-CIS governance.

  \item \textbf{[Phase~2]} Establish ``Data Visa'' mechanisms enabling dynamic
        cross-border AI model and data migration without manual legal friction, governed
        by a common semantic security ontology (CL3.4) aligned with STIX~2.1/TAXII
        threat intelligence standards and automated EUCS compliance verification~\cite{eucs}.

  \item \textbf{[Phase~2--3]} Deploy Zero-Knowledge Proofs (ZKPs) for IP-protected
        supply chain component verification: hardware vendors prove FIPS~203/204/205
        algorithm compliance and absence of known vulnerabilities without revealing
        proprietary design details~\cite{fips203,fips204,fips205}; operationalise
        pan-European supply chain compromise indicator sharing via standardised
        CTI lifecycle platforms (CL2.5).
\end{enumerate}

% ---- DIMENSION 4 ----
\item \textbf{Catalysts \& Enablers}

\textit{Funding Alignment:}
IPCEI-CIS (Important Project of Common European Interest on Next Generation Cloud
Infrastructure and Services) as the primary funding vehicle for cross-border NaaS federation
pilots; Digital Europe Programme for Simpl development, Common Data Space pilots, and
MIM standardisation infrastructure; ECCC Strategic Agenda priority areas ``Supply Chain
Security'' and ``Secure Interoperability of Systems and Networks''~\cite{eccc}; European
Alliance for Industrial Data, Edge and Cloud for governance and pilot ecosystem
coordination~\cite{cloud_edge_roadmap}.

\textit{Standardisation:}
CEN/CENELEC JTC~13 for security MIM standardisation within the Simpl ecosystem;
GSMA CAMARA API Alliance for harmonised Network API standardisation across operators;
ETSI TC CYBER for supply chain security standards; OASIS STIX~2.1/TAXII for threat
intelligence interoperability; W3C Verifiable Credentials Data Model for SSI-based
supply chain attestation chains.

\textit{Ecosystem Enablers:}
Gaia-X Federation Services trust framework alignment with Simpl MIM security policies;
Digital Product Passport pilot programmes under EU Ecodesign Regulation~2024;
ENISA supply chain risk management guidelines and threat landscape reports for ICT
sectors; dedicated SME compliance support programmes for CRA-mandated SBOM tooling,
recognising that sub-50 person software vendors represent the most vulnerable link in
the supply chain verification architecture~\cite{cra}.

% ---- DIMENSION 5 ----
\item \textbf{Disruptors \& Systemic Risks}

\begin{itemize}
  \item \textbf{Hyperscaler resistance to Simpl MIM adoption:} AWS, Azure, and
        GCP collectively control approximately 65\% of the EU cloud market as of
        2024 and have strong commercial incentives to maintain proprietary APIs;
        without mandatory EU Cloud Rulebook provisions requiring MIM compliance,
        voluntary adoption rates are unlikely to achieve genuine interoperability
        within Phase~1 timelines.
  \item \textbf{IPCEI-CIS implementation delays:} complex EU state aid notification
        procedures have historically delayed IPCEI projects by 18--36 months
        (IPCEI-H2 delayed by state aid clearance processes); P7 NaaS federation
        Phase~1 milestones carry material schedule risk if IPCEI-CIS governance
        does not streamline implementation mechanisms.
  \item \textbf{Semantic interoperability gap:} automated security policy translation
        between heterogeneous MIM implementations is architecturally harder than
        syntactic interoperability --- subtle semantic differences in ``deny'' versus
        ``drop'' versus ``quarantine'' semantics across provider implementations
        can create silent enforcement gaps undetectable through standard compliance
        auditing.
  \item \textbf{ZKP computational overhead:} current ZKP schemes for supply chain
        verification carry 10--100$\times$ computational overhead relative to
        standard cryptographic verification, limiting practical deployment to
        batch verification rather than real-time per-transaction attestation
        in Phases~1--2.
  \item \textbf{EU semiconductor supply chain concentration:} advanced node
        fabrication (sub-7\,nm) required for cutting-edge supply chain attestation
        hardware remains concentrated in Taiwan (TSMC) and South Korea (Samsung);
        a Taiwan Strait contingency would immediately constrain EU supply chain
        verification capability at the silicon level, exposing a gap the EU Chips Act
        does not address within the 2030 timeframe.
\end{itemize}

% ---- DIMENSION 6 ----
\item \textbf{Failsafe Mechanisms \& Resilience}

\begin{itemize}
  \item \textbf{Autonomous unfederation:} if a provider fails Runtime SBOM
        verification or MIM compliance checks, eBPF-based zero-trust policy
        enforcement severs that provider from the federated NaaS fabric and redirects
        workloads to verified compliant nodes within a sub-60-second SLA; provider
        re-admission requires fresh ZKP-verified attestation, not manual
        recertification.
  \item \textbf{Graceful national fallback:} cross-border SLA failure or ``Data
        Visa'' mechanism collapse triggers workload redirect to national infrastructure
        maintained at sufficient sovereign capacity for essential services, fulfilling
        NIS2 Article~21(2)(c) business continuity obligations and DORA Article~9 ICT
        resilience requirements~\cite{nis2,dora}.
  \item \textbf{Cryptographic attestation of MIM enforcement:} all Simpl MIM security
        controls include verifiable proofs that policies are enforced as declared --- not
        merely configured --- using TEE-backed attestation integrated with Cluster~1
        secure element infrastructure~\cite{cloud_edge_roadmap}, eliminating
        compliance-theatre scenarios where declared and actual enforcement diverge.
  \item \textbf{Supply chain quarantine:} unverified or high-risk components are
        isolated into quarantined network slices (VLAN/VxLAN-isolated segments with
        eBPF-enforced east-west traffic blocking) pending ZKP verification;
        quarantined components cannot interact with production workloads or access
        sensitive data paths.
  \item \textbf{Multi-jurisdictional residency fallback:} if cross-border ``Data
        Visa'' mechanisms fail legally or technically, automated GDPR Article~44
        compliance enforcement pins workloads to the originating jurisdiction,
        preventing unauthorised cross-border data transfers under federation failure
        conditions~\cite{gdpr}.
\end{itemize}

% ---- DIMENSION 7 ----
\item \textbf{Transversal Meta-Layer}

\textit{Sovereignty \& Autonomy:}
This cluster breaks non-EU hyperscaler API lock-in (AWS, Azure, GCP controlling
approximately 65\% of EU cloud market as of 2024) by enabling European providers to
compete on standardised Simpl-based infrastructure. A ``Schengen for Cloud'' legal
framework eliminates US CLOUD Act extra-territorial jurisdiction exposure for EU data
processed on federated EU-sovereign NaaS infrastructure. Runtime SBOM verification and
ZKP-based supply chain attestation reduce EU dependency on non-EU vendor self-certification
declarations --- the dominant model under current CRA-baseline SBOM obligations.

\textit{Trust \& Ethics:}
GDPR Article~25 (privacy by design and by default) is operationalised through
MIM-enforced data minimisation policies that are cryptographically verifiable, not
declarative~\cite{gdpr}. Data Act interoperability obligations are fulfilled through
Simpl MIM compliance, providing auditable trust chains for GDPR Article~5(2) data
processing accountability. Transparent Runtime SBOM verification provides the
supply chain integrity that GDPR Article~25 implicitly requires for processors
handling EU personal data at scale~\cite{gdpr}.

\textit{Sustainable Security (Green IT):}
Federated NaaS resource pooling across EU operators eliminates redundant edge compute
over-provisioning by individual operators, reducing total EU edge infrastructure carbon
footprint through shared utilisation efficiency. DPP-based supply chain traceability
aligns with EU Ecodesign Regulation circular economy requirements for ICT hardware,
enabling component reuse lifecycle tracking that reduces electronic waste from
premature infrastructure replacement driven by opacity in component provenance.

\end{enumerate}


% =========================================================================
% CLUSTER 3: FEDERATED OPERATIONAL SOVEREIGNTY
% =========================================================================
\newpage
\subsection{Progression Matrix: Priority Cluster 3 --- Federated Operational Sovereignty}

\begin{landscape}
\setlength{\textwidth}{26.5cm}\setlength{\linewidth}{26.5cm}%
\setlength{\hsize}{26.5cm}\setlength{\columnwidth}{26.5cm}
\setlength{\tabcolsep}{5pt}
\renewcommand{\arraystretch}{1.35}

\begin{longtable}{%
|>{\footnotesize\raggedright\arraybackslash}p{2.8cm}%
|>{\footnotesize\raggedright\arraybackslash}p{5.4cm}%
|>{\footnotesize\raggedright\arraybackslash}p{5.4cm}%
|>{\footnotesize\raggedright\arraybackslash}p{5.4cm}%
|>{\footnotesize\raggedright\arraybackslash}p{5.4cm}|}

\caption{\textbf{Cluster~3: Federated Operational Sovereignty} (P4, P5, P7).
Transition matrix across the four methodology dimensions. P4 (Simpl
Interoperability), P5 (Supply Chain Verification), and P7 (NaaS Federation)
are labelled where their maturity trajectories diverge within a cell.}
\label{tab:cluster3}\\

\hline
\rowcolor{colheader}
\textbf{Phase / SSL} &
\textbf{R\&I Maturity}{\scriptsize ODP indicator per sub-priority} &
\textbf{Policy, Regulatory \& Standards}{\scriptsize enforcement $\cdot$
  certification $\cdot$ standardisation} &
\textbf{Operational Deployment \& Adoption}{\scriptsize actor types $\cdot$
  scale $\cdot$ sector coverage} &
\textbf{Market Uptake \& Strategic Sovereignty}{\scriptsize SSL indicator
  $\cdot$ dependency reduction $\cdot$ industrial capacity} \\
\hline
\endfirsthead

\hline
\rowcolor{colheader}
\textbf{Phase / SSL} &
\textbf{R\&I Maturity} &
\textbf{Policy, Regulatory \& Standards} &
\textbf{Operational Deployment \& Adoption} &
\textbf{Market Uptake \& Strategic Sovereignty} \\
\hline
\endhead
\hline\multicolumn{5}{r}{\footnotesize\textit{Continued\ldots}}\\ \endfoot
\hline \endlastfoot

% ---- PHASE 1 ----
\rowcolor{phase1row}
\textbf{Phase~1}
\newline\textit{Foundation \&}
\newline\textit{Compliance}
\newline\textit{2025--2030}
\newline\sslbadge{sslredbg}{sslredtx}{SSL: Fragmented; Hyperscaler}
\newline\sslbadge{sslredbg}{sslredtx}{Dominance (P7); Emerging}
\newline\sslbadge{sslredbg}{sslredtx}{EU Ecosystem (P4)}
&
\cellhead{Cluster ODP: Fragmented Piloting (P4, P7) / Disclosure Phase (P5)}
\subp{P4} is at Fragmented Piloting: Simpl security MIM pilots are
launched in three Common European Data Spaces (Manufacturing, Health,
Energy) with manual policy coordination and proof-of-concept
implementations. \subp{P7} is at Bilateral Interconnection: harmonised
CAMARA Network APIs are deployed in 2--3 bilateral 5G corridor agreements
between individual EU telecom operators; cross-border edge compute
roaming does not exist at operational scale. \subp{P5} is at Supply
Chain Disclosure: CRA-mandated static SBOM declarations are being
submitted as point-in-time compliance artefacts without runtime
verification; Digital Product Passport pilot programmes are initiated
under the EU Ecodesign Regulation; Zero-Knowledge Proof verification
for supply chain attestation is at Applied R\&D stage.
\kpistrip{The three sub-priorities have meaningfully different Phase~1
maturity levels. P5 is the least mature and the most consequential
shared dependency: Runtime SBOM verification is a gate condition for
both P4 (MIM compliance trust) and Cluster~1 (TEE integrity claims).}
&
\cellhead{Multiple legislative obligations activating; standards nascent}
CRA~SBOM obligations (Annex~I Part~II) enter enforcement in 2027,
mandating static SBOM disclosure for all products with digital elements
entering the EU market. NIS2~Article~21(2)(h) makes supply chain risk
assessment mandatory for essential entities. DORA~Article~28 requires
continuous ICT third-party risk monitoring for financial sector entities.
Data Act interoperability obligations activate, creating procurement-level
pressure for cross-provider security interface standardisation. The
Simpl-Open official launch with initial MIM definitions constitutes the
Phase~1 technical policy trigger for P4. IPCEI-CIS enters its
deployment phase. GSMA CAMARA API Alliance standardisation work is
active but not yet mandated in EU regulatory instruments.
\kpistrip{Phase~1 policy activity is high-volume but fragmented:
multiple legislative instruments activate simultaneously without
a unified interoperability mandate that would create coherent market
pull for Simpl-based standardisation.}
&
\cellhead{Deployment: bilateral pilots and Data Space proof-of-concepts}
Three Common European Data Space security pilots operate Simpl-based
MIM proof-of-concepts with manual policy coordination; cross-provider
security policy enforcement is not automated. CAMARA API deployments
are confined to 2--3 bilateral 5G corridor agreements between single
operator pairs. Static SBOM declarations are submitted as CRA
compliance artefacts but are not verified at runtime; supply chain
risk assessment is vendor-declared rather than independently attested.
Hyperscalers (AWS, Azure, GCP) retain $\approx$65\% of EU cloud market
through proprietary API ecosystems. Cross-border AI model migration
requires bespoke bilateral legal review per transfer event. SME-accessible
SBOM tooling and MIM compliance support are not yet standardised or
commercially affordable.
\kpistrip{Phase~1 deployment is bilateral and pilot-scale. The gap
between proof-of-concept MIM implementations and automated enforcement
across the majority of Data Spaces defines the core operational
challenge for Phase~2.}
&
\cellhead{SSL: Mixed --- Emerging EU Ecosystem (P4) / Hyperscaler Dominance (P7)}
\subp{P4} has reached Emerging EU Ecosystem: the Simpl open-source
stack exists and is being piloted with EU institutional backing.
\subp{P7} remains at Hyperscaler Platform Dominance: AWS, Azure, and GCP
collectively control $\approx$65\% of EU cloud market; no EU NaaS
federation operates at commercially relevant scale; cross-border edge
compute is fragmented into national silos. \subp{P5} is at High
Reliance on Fragmented Non-EU Auditing: supply chain verification
is vendor-declared under US-dominated certification ecosystems (Common
Criteria national schemes, SOC2). The cluster's dominant SSL profile
is foreign dependency, driven by the weight of P7 market structure.
\kpistrip{P4 and P7 require separate SSL tracking in Phase~1: conflating
them into a single cluster SSL figure obscures that Simpl has achieved
Emerging Ecosystem while NaaS federation has not.}
\\
\hline

% ---- GATE 1->2 ----
\rowcolor{gaterow}
\gatelbl{1 $\rightarrow$ 2} &
\multicolumn{4}{>{\footnotesize\raggedright\arraybackslash}p{22.0cm}|}{%
\cellcolor{gaterow}%
\textbf{R\&I trigger:} At least one Common European Data Space has deployed
automated (not manually coordinated) Simpl MIM-based cross-provider security
policy enforcement, with documented evidence that enforcement is cryptographically
verifiable rather than declaratively asserted.
\textbf{Policy trigger:} Simpl security MIM definitions are published in
candidate CEN/CENELEC standard form and referenced as a compliance pathway in
at least one EU sector-specific NIS2 or Data Act implementing act.
\textbf{Operational trigger:} CAMARA-based cross-border edge compute roaming is
operational between at least three distinct EU telecom operator pairs (not a single
bilateral agreement), with documented SLA performance and compliance evidence.
\textbf{Market trigger:} Runtime SBOM tooling with cryptographic execution proofs
is available as a commercially supported product or service from at least two
EU-domiciled vendors, establishing the market infrastructure for Phase~2
mandatory adoption.}
\\
\hline

% ---- PHASE 2 ----
\rowcolor{phase2row}
\textbf{Phase~2}
\newline\textit{Applied \& Industrial}
\newline\textit{R\&D alongside}
\newline\textit{Federation \&}
\newline\textit{Standardisation}
\newline\textit{2030--2035}
\newline\sslbadge{ssltransbg}{ssltrantx}{SSL: Growing Pan-EU}
\newline\sslbadge{ssltransbg}{ssltrantx}{Sovereignty}
&
\cellhead{Cluster ODP: Standardised Integration (P4) / Federated Edge Roaming (P7) / Dynamic Runtime Verification (P5)}
\subp{P4} reaches Standardised Integration: automated policy translation
via standardised Simpl MIMs is deployed across the majority of active
Common European Data Spaces; federated identity attestation and
cross-provider compliance reporting are automated. \subp{P7} reaches
Federated Edge Roaming: CAMARA-based cross-border roaming operates
across $\geq$5 EU operator pairs; ``Data Visa'' mechanisms enable
dynamic AI model migration without manual legal review in certified
corridors. \subp{P5} transitions from disclosure to dynamic runtime
verification: Runtime SBOMs with cryptographic proof-of-execution
verify containerised components as they migrate across multi-provider
nodes; ZKP batch verification of supply chain attestation is operationally
validated; Digital Product Passports are mandatory under the EU
Ecodesign Regulation.
\kpistrip{Phase~2 R\&I activity includes industrial R\&D for ZKP
real-time verification (currently batch-only due to 10--100$\times$
overhead), semantic security ontology standardisation, and legal
engineering of the ``Schengen for Cloud'' framework --- not solely
standardisation of existing prototypes.}
&
\cellhead{Security MIMs as CEN/CENELEC standards: the Phase~2 gate}
The single Phase~2 policy gate is the adoption of Simpl security MIMs
as ratified CEN/CENELEC standards, referenced as mandatory compliance
requirements in the EU Cloud Rulebook. This transforms MIM adoption
from voluntary best-practice to market-access obligation for providers
serving EU-regulated sectors. DPP becomes mandatory under the EU
Ecodesign Regulation~2024 for ICT hardware categories. Data Act full
enforcement creates binding procurement-level interoperability
obligations. DORA continuous ICT third-party monitoring is fully
operationalised in the financial sector. IPCEI-CIS cross-border NaaS
governance framework is operational.
\kpistrip{The Phase~2 policy gate is the CEN/CENELEC ratification of
security MIMs. Voluntary Simpl adoption without this standards event
does not constitute Phase~2 policy maturity.}
&
\cellhead{Cross-sector rollout; SME participation; Runtime SBOM in production}
Simpl MIMs are deployed across the majority of active Common European
Data Spaces with automated cross-provider enforcement. Runtime SBOMs
with cryptographic execution proofs verify supply chain components in
critical infrastructure deployments in at least five Member States.
CAMARA cross-border roaming operates across $\geq$5 EU operator pairs
with documented SLA performance. ``Data Visa'' mechanisms are
operational in certified cross-border AI model migration corridors.
SME-accessible Simpl MIM compliance tooling and Runtime SBOM services
are commercially available from EU-domiciled providers. Unverified
supply chain components are automatically quarantined in isolated
network slices.
\kpistrip{Phase~2 deployment must show geographic breadth ($\geq$5 Member
States), cross-sector presence ($\geq$3 Data Spaces under automated MIM
enforcement), and SME participation in at least one deployment context.
Single-sector rollout does not constitute Phase~2 operational maturity.}
&
\cellhead{SSL: Growing Pan-EU Sovereignty}
EU Data Space infrastructure operating under Simpl achieves $\geq$30\%
non-hyperscaler share in new regulated-sector cloud procurement. Runtime
SBOM ecosystem is commercially mature with multiple competing EU-domiciled
vendors. EU NaaS consortium (under IPCEI-CIS) operates cross-border
compute at pilot commercial scale, demonstrating that EU providers
can federate to compete with hyperscaler geographic coverage. DPP
traceability reduces procurement risk for EU critical infrastructure
buyers by making supply chain provenance verifiable. Foreign hyperscaler
market share in EU regulated sectors declines for the first time from
its $\approx$65\% Phase~1 baseline.
\kpistrip{Phase~2 SSL evidence: $\geq$30\% non-hyperscaler share in new
regulated procurement; first documented market share decline for
non-EU hyperscalers in EU regulated cloud procurement.}
\\
\hline

% ---- GATE 2->3 ----
\rowcolor{gaterow}
\gatelbl{2 $\rightarrow$ 3} &
\multicolumn{4}{>{\footnotesize\raggedright\arraybackslash}p{22.0cm}|}{%
\cellcolor{gaterow}%
\textbf{R\&I trigger:} ZKP verification for supply chain attestation has
achieved real-time per-transaction performance (eliminating the current
batch-only constraint) in at least one production critical infrastructure
deployment, confirmed by independent performance benchmarking.
\textbf{Policy trigger:} A ``Schengen for Cloud'' legal framework --- or
equivalent EU legal instrument resolving US CLOUD Act extra-territorial
jurisdiction conflict for federated EU-sovereign NaaS workloads --- has
been enacted, providing the jurisdictional basis for single-market operation
of federated compute.
\textbf{Operational trigger:} The EU NaaS federation (under IPCEI-CIS or its
successor) operates at commercial scale across $\geq$10 EU telecom operator
members, offering cross-border edge compute services as a standard commercial
product with documented customer deployments outside its founding consortium.
\textbf{Market trigger:} EU supply chain governance frameworks (Runtime SBOM
plus ZKP attestation) are adopted or formally referenced by at least two
non-EU partner jurisdictions as part of digital trade or partnership agreements,
confirming that EU governance leadership is internationally exportable.}
\\
\hline

% ---- PHASE 3 ----
\rowcolor{phase3row}
\textbf{Phase~3}
\newline\textit{Autonomy \&}
\newline\textit{Leadership}
\newline\textit{2035--2040}
\newline\sslbadge{ssltealbg}{ssltealtx}{SSL: Full Federated}
\newline\sslbadge{ssltealbg}{ssltealtx}{Sovereign Autonomy}
&
\cellhead{Cluster ODP: Ubiquitous Operation across all sub-priorities}
All three sub-priorities reach Ubiquitous Operation. Simpl-based
security is the universal standard for cross-provider Data Space
operation in the EU; no alternative interoperability model holds
market relevance in regulated sectors. ZKP-verified supply chain
attestation operates at real-time transaction scale for all critical
infrastructure component verification. The EU NaaS federation operates
as a single commercial entity legally equivalent to a unified
jurisdiction for data workloads. Phase~3 R\&I pursues next-generation
frontiers: post-ZKP cryptographic verification, semantic AI-to-AI
policy negotiation, and legal engineering of the next-generation
``Digital Single Market'' framework beyond cloud-edge compute.
\kpistrip{Phase~3 ODP evidence: Simpl MIM compliance rate in new regulated
EU cloud procurement; ZKP verification throughput in production critical
infrastructure deployments; EU NaaS federation revenue relative to
hyperscaler EU revenue in regulated sectors.}
&
\cellhead{EU federated governance as international model}
The ``Schengen for Cloud'' legal framework is enacted and operational,
resolving US CLOUD Act jurisdiction conflicts for federated EU-sovereign
workloads and enabling single-market legal operation of cross-border
compute. EU Cloud Rulebook is mandatory for all new cloud deployments
in EU regulated sectors without exception. EU supply chain governance
standards (Runtime SBOM plus ZKP) are referenced in international
supply chain security standards (ISO/IEC) and in digital trade
agreements with partner jurisdictions. The EU has transitioned from
rule-taker to rule-setter in international cloud and supply chain
governance.
\kpistrip{Phase~3 governance leadership is verified by two events:
enactment of the ``Schengen for Cloud'' framework (domestic) and
formal adoption of EU supply chain governance standards by $\geq$2
non-EU partner jurisdictions (international).}
&
\cellhead{Provider-agnostic EU trust fabric as continental operational default}
A single provider-agnostic EU trust fabric operates continuously across
all active Common European Data Spaces with real-time automated security
orchestration and zero manual policy coordination. ZKP-verified supply
chain attestation eliminates all unverified components from EU critical
infrastructure deployments in regulated sectors. The EU NaaS federation
is a standard commercial service operating at hyperscaler-comparable
scale, with EU operators collectively competitive on cost and coverage.
Security-by-Design is not an add-on: every new cloud-edge deployment
in the EU incorporates Simpl MIM compliance, Runtime SBOM verification,
and cross-border sovereignty assurance as architectural defaults.
\kpistrip{Phase~3 deployment normalisation: Simpl MIM compliance in
$\geq$90\% of new regulated EU cloud deployments; ZKP supply chain
verification in 100\% of new critical infrastructure component
procurement; EU NaaS federation operating as a single commercial entity
with documented non-founding-member customers.}
&
\cellhead{SSL: Full Federated Sovereign Autonomy}
Non-EU hyperscaler market share in EU regulated cloud procurement
falls below 40\% for new deployments (from $\approx$65\% at baseline),
with a binding trajectory toward continued reduction. EU NaaS
federation achieves hyperscaler-comparable geographic coverage and
commercial viability across the EU. Supply chain governance
internationally adopted: at least two partner jurisdictions have
formally adopted EU Runtime SBOM plus ZKP attestation standards as
domestic procurement requirements. The EU has removed the structural
dependency that allowed US CLOUD Act jurisdiction to apply to
EU-regulated data: federated sovereignty is legally guaranteed, not
only technically implemented.
\kpistrip{Full Federated Sovereign Autonomy requires three
concurrent conditions: $<$40\% hyperscaler new-procurement share;
``Schengen for Cloud'' legally enacted; $\geq$2 partner jurisdictions
formally adopting EU supply chain governance standards. Achieving
any two without the third does not constitute this SSL level.}
\\
\hline

\end{longtable}
\end{landscape}

% ---- PROGRESSION MATRIX: CLUSTER 3 ----


% \noindent
% \resizebox{\textwidth}{!}{%
% \begin{tabular}{|p{4.2cm}|p{5.6cm}|p{5.6cm}|p{5.6cm}|}
% \hline
% \rowcolor{gray!30}
% \multicolumn{4}{|c|}{\textbf{Progression Matrix: Priority Cluster 3 ---
%   Federated Operational Sovereignty (P4, P5, P7)}} \\
% \hline
% \rowcolor{gray!15}
% \textbf{Metric} &
% \textbf{Phase 1: Foundation \& Compliance (2025--2030)} &
% \textbf{Phase 2: Federation \& Standardisation (2030--2035)} &
% \textbf{Phase 3: Autonomy \& Leadership (2035--2040)} \\
% \hline
% \rowcolor{gray!10}
% \multicolumn{4}{|l|}{\textit{\textbf{Cluster-Level Progression Indicators}}} \\
% \hline
% \textbf{Cluster ODP} &
% Fragmented Piloting &
% Standardised Integration &
% Ubiquitous Operation \\
% \hline
% \textbf{Cluster SSL} &
% Emerging EU Ecosystem &
% Growing Pan-EU Sovereignty &
% Full Federated Sovereign Autonomy \\
% \hline
% \rowcolor{gray!10}
% \multicolumn{4}{|l|}{\textit{\textbf{Sub-Priority Decomposition (Individual ODP / SSL per Phase)}}} \\
% \hline
% \textit{P4: Security Interoperability \newline via Simpl Middleware} &
% ODP: Fragmented Piloting \newline
% SSL: Emerging EU Ecosystem &
% ODP: Standardised Integration \newline
% SSL: Full Sovereign Autonomy &
% ODP: Ubiquitous Operation \newline
% SSL: Full Sovereign Autonomy \\
% \hline
% \textit{P5: Supply Chain \newline Verification \& Infra.} &
% ODP: Full Supply Chain Disclosure \newline
% SSL: High Reliance on \newline Fragmented Non-EU Auditing &
% ODP: Partial Restriction \& \newline Dynamic Runtime Verification \newline
% SSL: Emerging EU Supply Chain \newline Tracking Frameworks &
% ODP: Full Restriction Process \& \newline ZKP-Verified Attestation \newline
% SSL: Full Sovereign Supply \newline Chain Control \& Transparency \\
% \hline
% \textit{P7: Pan-European NaaS \newline Federation} &
% ODP: Bilateral Interconnection \newline
% SSL: Hyperscaler Platform Dominance &
% ODP: Federated Edge Roaming \newline
% SSL: Pan-EU Consortium Competes &
% ODP: True Digital Single Market \newline
% SSL: European Scale Parity Achieved \\
% \hline
% \rowcolor{blue!7}
% \textbf{Operational State} &
% Initial Simpl security pilots are deployed in three Common European
% Data Spaces with manual policy coordination and proof-of-concept
% MIM implementations; CRA-baseline SBOM declarations are submitted
% by hardware and software vendors as static point-in-time documents
% without runtime verification; harmonised CAMARA APIs are tested in
% bilateral 5G corridor agreements between individual EU telecom
% operators; hyperscalers retain dominant market share through
% proprietary API ecosystems. &
% Automated policy translation via standardised Simpl MIMs is
% deployed across the majority of active Common European Data Spaces;
% Runtime SBOMs with cryptographic execution proofs verify
% containerised components dynamically as they migrate across
% multi-provider continuum nodes; harmonised CAMARA-based cross-border
% edge roaming operates across five or more EU operator pairs;
% ``Data Visa'' mechanisms enable AI model migration without
% manual legal review in certified corridor deployments. &
% A single provider-agnostic European trust fabric operates
% continuously across all EU Common Data Spaces with real-time
% automated security orchestration; ZKP-verified supply chain
% attestation eliminates unverified components from all
% EU critical infrastructure deployments; a unified NaaS
% platform legally operable as a single data jurisdiction
% enables EU operators to achieve hyperscaler-comparable
% scale; EU supply chain governance frameworks are adopted
% as reference implementations internationally. \\
% \hline
% \rowcolor{orange!12}
% \textbf{Critical Enabler} &
% Simpl-Open official launch with initial security MIM definitions
% published as open standards, establishing the technical
% baseline for cross-provider security policy automation &
% Security MIMs adopted as CEN/CENELEC standards with
% mandatory reference in EU Cloud Rulebook, creating the
% regulatory basis for non-voluntary hyperscaler compliance &
% ``Schengen for Cloud'' legal framework enacted, resolving
% US CLOUD Act extra-territorial jurisdiction conflict and
% enabling single-jurisdiction data workload operation at
% continental scale \\
% \hline
% \end{tabular}%
% }

% =========================================================================
% CROSS-PRIORITY BOTTLENECK ASSESSMENT
% =========================================================================

\newpage
\section{Cross-Priority Bottleneck Assessment}
\label{sec:bottlenecks}

Three shared structural dependencies determine whether individual sub-priority
progressions can advance on their stated timelines, regardless of cluster-level investment.
These bottlenecks are identified here because they constitute cross-cutting gate conditions
that cannot be resolved by any single cluster in isolation.

\textbf{Bottleneck 1 --- EU Chips Act Sovereign Fabrication Capacity.}
The Phase~2 gate conditions for P1 (EUCC `High' certified RISC-V secure elements),
P2 (PQC-native RISC-V silicon), P3 (zero-energy secure hardware), and P8 (RISC-V
Keystone TEE scaling) all converge on the same physical dependency: sovereign EU
fabrication capacity for security-critical silicon at commercially viable yields.
The Phase~2 critical enabler for P1 (``EUCC `High' certification of first
EU-manufactured RISC-V secure element'') and the Phase~3 enabler (``EU Chips Act
sovereign EUV-capable fabrication achieving $\geq$60\% self-sufficiency'') constitute
a single two-stage gate that controls the SSL trajectory of four of the five Cluster~1
sub-priorities simultaneously.

\textbf{Bottleneck 2 --- AI Act Conformity Assessment Infrastructure.}
The Phase~2 gate condition for Cluster~2 (``standardised AI Act conformity assessment
framework for autonomous security agents'') is a prerequisite not only for P6 and P10
but also for any AI-driven capability within Cluster~1 (autonomous Island Mode swarm
protocols) and Cluster~3 (Automated Compliance Negotiators in NaaS federation) that
incorporates autonomous decision-making. The absence of EU AI Act conformity assessment
bodies with sufficient capacity to certify high-risk security AI at scale represents a
systemic bottleneck that, if unresolved by 2028--2030, will delay the Phase~2 SSL
transition across all three clusters.

\textbf{Bottleneck 3 --- Runtime SBOM and Supply Chain Attestation Infrastructure.}
Runtime SBOM verification (P5) is a prerequisite for both the Cluster~1 TEE integrity
claim (P8 confidential computing pipelines cannot be trusted without verified component
provenance) and the Cluster~3 MIM compliance model (Simpl-based interoperability
cannot provide cryptographic enforcement guarantees if component integrity is
unverifiable). The Phase~1 critical enabler for Cluster~3 (Simpl-Open MIM definitions)
and the Phase~1 ODP for P5 (full supply chain disclosure) must therefore proceed in
coordination, not sequentially, with the Phase~1 applied research activities in
Cluster~1. This interdependency must be explicitly reflected in Horizon Europe project
call design to ensure that SBOM infrastructure investment is not siloed within supply
chain security calls but is structurally integrated into hardware security and
interoperability funding instruments.

% =========================================================================
% NEW BIBLIOGRAPHY ENTRIES
% =========================================================================

% Add the following entries to the master \begin{thebibliography}{99} block:
